% =======================================================================
% LLM FORMATTING PROTOCOL (V3.2: CONTINUUM MECHANICS) 
% -----------------------------------------------------------------------
% 1. Numbering Architecture: 
%    - The Reverse Modeling Cheat Sheet establishes the global equation numbering.
%    - Use \refstepcounter{cheatnum}, \label, \hyperref for the Cheat Sheet boxes. 
%    - ALL primary equations must be declared and \label'd FIRST within the Cheat Sheet boxes.
%    - In the subsequent Derivation sections, use \tag{\ref{eq:name}} to inherit the number from the Cheat Sheet. Apply this \tag to both the final in-text derivation step and the summary recall box.
%    - Intermediary derivation steps can be automatically numbered by LaTeX (no \nonumber needed).
%
% 2. Strict LaTeX Math: All math enclosed in $ $. Use \begin{equation} or \begin{align}.
%
% 3. LEXICON & TERMINOLOGY (THE CONTINUUM BOUNDARY):
%    - WE DO NOT USE: "The fluid", "Aether", "Fluid of spacetime".
%    - WE USE: "The continuum", "The background metric", "The continuous metric", "The coordinate space".
%    - WE PREFER NOT TO USE: "The vacuum" (implies emptiness). 
%    - WE CAN USE: "Fluid mechanics/dynamics" to credit that Physics domain, but not to describe the medium or our own Physics.
%    - WE USE: "Continuum mechanics", "Elastodynamics", "Navier-Stokes solutions", "Eulerian evolution". 
%    - Optical domains: Use "refractive medium" or "Gradient-Index (GRIN) Lens".
%
% 4. Tone (Declare, Don't Compare): Present axiomatically from first principles. NEVER compare to standard physics, historical anomalies, or appeal to human textbook authority (e.g., "According to classical continuum mechanics..."). Let the math stand natively.
%
% 5. Structure: Math >> Narrative. Equations drive the derivations.  If you can replace a narrative block with an equation, do it.  Avoid narratives if at all possible.
%
% 6. Explicit Exposition: Show every logical and algebraic manipulation. No skipped steps. No tensors (e.g., no covariant derivatives, no wedge products). Prove everything using Newtonian, Eulerian, and Hookean calculus.
%
% 7. Fluff & Drama Elimination: 
%  These are the biggest violators, but as a general rule, avoid dramatizing, if you can eliminate a word, do it, if a word is not necessary get rid of it.
%    - BANNED ADVERBS: strictly, natively, explicitly, absolutely, flawlessly, perfectly, infinitely, purely, brilliantly, inevitably, merely, vastly, simply, fundamentally, exactly, mathematically, algebraically.  
%    - BANNED DRAMA: unbreakable, intense, violent, smoothly. (For algebra, use "cancels", not "annihilates").
%    - ADJECTIVE WATCHLIST: topological, acoustic, geometric, thermodynamic, continuous, macroscopic, continuum, temporal, structural, local, localized. These are NOT banned, but if they are not explicitly required to define a variable, DELETE THEM. Let the math stand alone.
%
% 8. Anti-Teleology: The metric has no intent. It does not "attempt", "try", or "process". 
%
% 9. Punctuation: No em-dashes allowed in the text.
%
% 10. STRICT VERSION CONTROL & EDIT TRACKING
% Constraint: NEVER alter, add, delete, or rephrase the user's original text without an explicit, searchable flag. You must preserve the user's exact wording unless a modification is strictly required to comply with the rules above.
% Action: If you change even a single character, append `% :::LLM_EDIT:::` to the end of that line, followed by a 2-5 word justification.
%
% 11. PROSE SOVEREIGNTY (THE ANTI-ENHANCEMENT LAW)
% Constraint: If the user provides sanitized, human-edited prose, you are MATHEMATICALLY FORBIDDEN from adding adjectives or adverbs to make it sound "more rigorous." Edit ONLY broken LaTeX or factually incorrect equations. If you "smooth out" or "enhance" the user's narrative tone, you fail the prompt.
% =======================================================================

% ======================================================================
% DOCUMENT PREAMBLE
% ======================================================================
\documentclass[11pt, a4paper]{article}

\usepackage[top=1cm, bottom=2cm, left=1.5cm, right=1.5cm]{geometry}
\usepackage{fontspec}
\usepackage[english, bidi=basic, provide=*]{babel}
\babelprovide[import, onchar=ids fonts]{english}
\babelfont{rm}{Times New Roman} % Standard academic serif

\usepackage{microtype}
\usepackage{amsmath, amssymb, amsthm}
\usepackage{physics}
\usepackage{multicol}
\usepackage{enumitem}
\usepackage{xcolor}
\usepackage{titlesec}
\usepackage{hyperref}
\usepackage[most]{tcolorbox} 

% Define a professional "Academic Sepia/Charcoal" palette
\definecolor{academiabrown}{RGB}{70, 50, 40}
\hypersetup{colorlinks=true, linkcolor=academiabrown, citecolor=academiabrown, urlcolor=blue, hidelinks}

% --- Define the Dynamic Box Styles ---
\tcbset{
    cheatbox/.style={
        enhanced, unbreakable,
        colback=academiabrown!4!white, colframe=academiabrown!90!black,   
        fonttitle=\bfseries, coltitle=white, boxrule=0.8pt, arc=2mm,                   
        left=2mm, right=2mm, top=1.5mm, bottom=1.5mm, boxsep=1mm,
        toptitle=1.5mm, bottomtitle=1.5mm
    },
    recallbox/.style={
        enhanced, unbreakable,
        colback=academiabrown!4!white, colframe=academiabrown!90!black,
        boxrule=0.8pt, arc=2mm, left=2mm, right=2mm, top=1mm, bottom=1mm, boxsep=0mm
    }
}

% Custom section formatting
\titleformat{\section}{\Large\bfseries\color{black}}{\thesection.}{0.5em}{}
\titleformat{\subsection}{\large\bfseries\color{darkgray}}{\thesubsection}{0.5em}{}
\titleformat{\subsubsection}{\normalsize\bfseries\color{gray}}{\thesubsubsection}{0.5em}{}

\setlength{\columnsep}{0.6cm}

% --- Dynamic Numbering Setup ---
\newcounter{cheatnum}
\setcounter{cheatnum}{0}

% ======================================================================
% TITLE & ABSTRACT
% ======================================================================
\title{\textbf{Refractive Cosmology \& Continuum Mechanics \\ CONDENSED REFERENCE \& FORMAL DERIVATIONS}}
\author{\textbf{Binyamin Tsadik Bair-Mosheh} } % :::LLM_EDIT::: Preserved exact author formatting.
\date{\today}

\begin{document}

\maketitle


\begin{tcolorbox}[colback=gray!5!white, colframe=gray!50!black, arc=1mm, boxrule=0.5pt, left=3mm, right=3mm]
\begin{center}
    \textbf{ABSTRACT \& HOW TO USE THIS MANUAL}
\end{center}
\textit{This is a \textbf{condensed} version of the Master Reference Sheet.
This physics establishes that macroscopic cosmology, localized quantum mechanics, and solid-state phenomena evaluate as the deterministic consequences of continuum mechanics operating within a dynamic elastic vacuum. By defining the background metric as a compressible, refractive continuum, the foundations of physical law are systematically derived from classical elastodynamics, Eulerian material derivatives, and geometric topology.} % :::LLM_EDIT::: Replaced "fluid" with "continuum" and "elastodynamics" per V3.0.
\vspace{0.2cm}
\begin{itemize}[nosep]
    \item \textbf{PART I} provides the dynamically numbered Cheat Sheet containing the raw equations for quick use and reference.
    \item \textbf{PART II} \hyperref[sec:part2]{\textbf{[Click Here]}} provides step-by-step mathematical proofs from first principles. 
    \item The \textbf{[Proof $\downarrow$]} links in each box navigate directly to the corresponding formal derivation.
\end{itemize}
\end{tcolorbox}

\vspace{0.3cm}

\begin{tcolorbox}[cheatbox, title={The Conformal Duality of the Metric \scriptsize\hfill\hyperref[sec:part2]{[See Derivations $\downarrow$]}}]
The unperturbed baseline evaluates as a continuous metric, characterized by conformal thermodynamic symmetry between the macroscopic cosmos and local topological gradients: % :::LLM_EDIT::: Replaced "vacuum/fluid" with "unperturbed baseline/continuous metric".
\begin{itemize}[leftmargin=*, nosep]
    \item \textbf{The Cosmos (The Macroscopic Limit):} Expanding causal wavefront. The outward macroscopic velocity field approaches the acoustic limit ($\mathbf{v}_0 \to c$) as $r \to \infty$, maximizing dynamic pressure and driving the ambient static pressure of the boundary signal to zero.
    \item \textbf{The Particle (The Localized Core):} Sequestered topological void. The localized kinematic circulation approaches the acoustic limit ($\mathbf{v}_{\phi} \to c$) as $r \to r_c$, maximizing dynamic pressure and driving the internal static pressure to a minimum at the cavitation boundary. % :::LLM_EDIT::: Replaced "fluid circulation" with "kinematic circulation".
\end{itemize}
\end{tcolorbox}

\vspace{4cm}


\begin{center}
    \LARGE\textbf{Condensed Version}\\ 
    \vspace{1cm}
    \small\textit{The full version is currently obsolete but will be updated eventually:  \\ \hyperlink{https://zenodo.org/records/20368016}{https://zenodo.org/records/20368016}}
\end{center}


% ======================================================================
% NOMENCLATURE & KINEMATIC CONSTANTS
% ======================================================================

\begin{tcolorbox}[colback=white, colframe=academiabrown!90!black, arc=1mm, boxrule=0.8pt, title={\textbf{TABLE 1: KINEMATIC NOMENCLATURE \& METRIC CONSTANTS}}]
\renewcommand{\arraystretch}{1.4}
\begin{tabular}{p{0.18\textwidth} p{0.77\textwidth}}
    $\boldsymbol{\tau}$ & \textbf{Metric Time.} The current, unmediated time of the continuum. \\
    $t$ & \textbf{Proper Time.} The evolutionary time within the metric. \\
    $\tau_0$ & \textbf{Cosmological Age.} The local proper age of the universe. \\
    $R_{h,0}$ & \textbf{Macroscopic Scale Radius.} The causal horizon at the present observation epoch. \\
    $l_p$ & \textbf{Cavitation Limit.} The spatial yield limit of the continuum (Planck length), dictating the minimum physical radius before structural collapse. \\ 
    $r_c \equiv \frac{\lambda_c}{2\pi}$ & \textbf{Cavitation Boundary (Compton Radius).} The localized spatial boundary where internal circulation strikes the acoustic phase limit ($v_\phi \to c$). \\ 
    $v_0(r)$ & \textbf{Macroscopic Velocity Field.} The scalar velocity field of the expanding metric. \\ 
    $v_{obs}(r) \equiv v_\phi(r)$ & \textbf{Topological Circulation Velocity.} The closed-loop phase velocity field of the metric extending outward from a phase boundary. Evaluates via inverse projection ($v_\phi(r) = c \cdot{r_c/r}$) to satisfy the $1/r^3$ spatial pressure gradient. \\ 
    $P_c \equiv \frac{1}{2}\rho_{\tau}c^2$ & \textbf{Metric Yield Pressure.} The static pressure limit of the continuum. Acts as the baseline isotropic tension ($T_c = -P_c$). Evaluates as the energy density of rest mass. \\ 
    $P_{dyn} = -\rho_{\tau} \Phi$ & \textbf{Scalar Pressure Deficit.} The localized dynamic pressure drop generated by topological circulation. Evaluates as the Newtonian gravitational potential. \\  
    $n(r)$ & \textbf{Refractive Index of the Medium.} The currently observed index at radial distance $r$ dictating the history wave propagation along a null geodesic ($ds^2 = 0$). \\
    $De$ & \textbf{Cosmological Deborah Number.} The kinematic ratio ($\omega t_r$) dictating the rheological phase of the metric. \\
    $W$ & \textbf{Topological Winding Number.} The chirality or charge of an acoustic cavity, evaluating as the parity of the Dirac spinor ($\psi, \bar{\psi}$). \\
    $\alpha_s \equiv \frac{1}{\gamma}$ & \textbf{Total Scalar Admittance (Metric Compliance).} The volumetric yield of the metric against total kinematic capacity. Evaluates as the square root of the residual static pressure ratio ($\sqrt{1 - (P_{dyn} + P_{shear})/P_c}$) and the Strong Coupling Constant. \\ 
    $\alpha = \mu_s / K$ & \textbf{Viscoelastic Modulus Ratio.} The ratio of shear capacity to bulk compressive capacity. Evaluates as the Fine-Structure Constant. \\
    $Z_{\tau}$ & \textbf{Universal Kinematic Impedance.} The structural resistance of the metric against applied kinematic strain.  \\ 
    $\epsilon_{\tau}$ & \textbf{Linear Density.} The 1D density evaluating as 2D transverse elastic stiffness. \\
    $\rho_{\tau} = \frac{1}{2\pi}\left(\frac{c^5}{\hbar G^2}\right)$ & \textbf{Volumetric Density.} The baseline perturbation density of the metric. \\ 
\end{tabular}
\end{tcolorbox}
\newpage
% ======================================================================
% PART I: THE AXIOMATIC CHEAT SHEET
% ======================================================================
\begin{center}
    \LARGE\textbf{PART I: THE CHEAT SHEET}
\end{center}
\vspace{0.2cm}

\begin{multicols}{2}

% ==========================================
% PHASE I: THE MACROSCOPIC CONTINUUM
% ==========================================
\vspace{0.2cm}
\noindent\textbf{\large \color{academiabrown} PHASE I: THE MACROSCOPIC CONTINUUM}
\vspace{0.2cm}

\begin{tcolorbox}[cheatbox, title={\refstepcounter{cheatnum}\label{box:axiom}\thecheatnum. The Causal Boundary \hfill\makebox[0pt][r]{\scriptsize\hyperref[sec:deriv_axiom]{[Proof $\downarrow$]}}}]
The causal boundary expands linearly at $c$. The translation between Evolutionary Look-Back Time ($t$) and Relative Observational Redshift ($z$) is defined by the horizon ($R_h$):
\begin{equation}\label{eq:axiom}
    R_h(t) = c \cdot t \iff 1+z = \frac{\tau_0}{t}
\end{equation}
\textit{Present-Day Calibration ($z=0$):}
\begin{align*}
    R_{h,0} &= 4224 \text{ Mpc}, &\tau_0 &= \frac{R_{h,0}}{c} \approx 13.78 \text{ Gyr}
\end{align*}
\end{tcolorbox}

\begin{tcolorbox}[cheatbox, title={\refstepcounter{cheatnum}\label{box:chronometer}\thecheatnum. The Redshift Chronometer \hfill\makebox[0pt][r]{\scriptsize\hyperref[sec:deriv_chronometer]{[Proof $\downarrow$]}}}]
Proper time ($t$) and scale radius ($R_h$) at emission, translate symmetrically across coordinate frames:
\begin{align}
    t(z) &= \frac{\tau_0}{1+z} \label{eq:tz} \\
    R_h(z) &= \frac{R_{h,0}}{1+z} \equiv c \cdot t(z) \label{eq:rz}
\end{align}
\end{tcolorbox}

\begin{tcolorbox}[cheatbox, title={\refstepcounter{cheatnum}\label{box:gr_bridge}\thecheatnum. The Geometric Spacetime Bridge \hfill\makebox[0pt][r]{\scriptsize\hyperref[sec:deriv_gr_bridge]{[Proof $\downarrow$]}}}]
The pressure gradient generates a continuous velocity field $\mathbf{v}_0(r)$. By the Equivalence Principle, the fluid potential energy ($\Phi = -v_0^2/2$) dictates the time-time structural curvature component ($g_{tt}$):
\begin{align}
    v_0(r) &= c\left(1-e^{-r/R_{h,0}}\right) \label{eq:v0} \\
   1/n(r) = -g_{tt} &= \alpha_s(r)^2 = \left( 1 - \frac{P_{dyn}(r)}{P_c} \right)   \label{eq:pressure_curvature_identity} \\ 
    ds^2 &= -\alpha_s^2(c~dt)^2 + \frac{dr^2}{\alpha_s^2} \label{eq:metric} 
\end{align}
\end{tcolorbox}

\begin{tcolorbox}[cheatbox, title={\refstepcounter{cheatnum}\label{box:kinematic_gradients}\thecheatnum. Kinematic Gradients \hfill\makebox[0pt][r]{\scriptsize\hyperref[sec:deriv_kinematic_gradients]{[Proof $\downarrow$]}}}]
Apparent expansion evaluates as ($H = dv_0/dr$). 
\begin{align}
    H_{r=0}(z) &= \frac{c}{R_{h}(1+z)} \iff H_{0}(t) = \frac{1}{t} \label{eq:happ} \\
    a_{h}(z) &= \frac{c^2}{R_{h}(1+z)} \iff a_{h} = \frac{c}{t} \label{eq:aapp} \\
    \omega_{h}(z) &= \frac{c(1+z)}{2\pi R_{h}} \iff \omega_h = \frac{1}{2\pi t} \label{eq:phase_frequency}
\end{align}
\end{tcolorbox}

\begin{tcolorbox}[cheatbox, title={\refstepcounter{cheatnum}\label{box:optical_dilatation}\thecheatnum. Optical Dilatation \hfill\makebox[0pt][r]{\scriptsize\hyperref[sec:deriv_optical_dilatation]{[Proof $\downarrow$]}}}]
Because the kinematic deceleration parameter of the continuum evaluates to zero ($q=0$), the true macroscopic coordinate distance evaluates as the geometric integration of the optical stretch parameter:
\begin{equation}\label{eq:r_cosmic}
    r_{true, cosmic}(z) = R_{h,0} \ln(1+z)
\end{equation}
The apparent luminosity distance ($D_L$) evaluates as the true spatial coordinate multiplied by the macroscopic expansion scalar:
\begin{equation}\label{eq:dl}
    D_L(z) = R_{h,0} (1+z) \ln(1+z)
\end{equation}
\end{tcolorbox}

\begin{tcolorbox}[cheatbox, title={\refstepcounter{cheatnum}\label{box:acoustic_invariance}\thecheatnum. Cosmic Expansion \& Acoustic Invariance \\ \null\hfill\scriptsize\hyperref[sec:deriv_acoustic_invariance]{[Proof $\downarrow$]}}]
The macroscopic continuum of the metric evaluates via a tensioned $w = -1/3$ equation of state. Because the dynamic density ($\rho_{dyn}$) and the transverse shear modulus ($\mu_s$) dynamically decay at the identical inverse-square temporal limit, the kinematic phase limit ($c$) evaluates as an invariant constant. % :::LLM_EDIT::: Replaced 'fluid', 'speed of sound' and removed 'mechanically' (Rules 3, 7). Swapped \rho_\tau to \rho_{dyn}.
\begin{align}
    \rho_{dyn}(t) &\propto t^{-2} \label{eq:density_decay} \\
    \mu_s(t) &\propto t^{-2} \label{eq:stiffness_decay} \\
    c &= \sqrt{\frac{\mu_s}{\rho_{dyn}}} = \text{Constant}  \label{eq:c_invariant} 
\end{align}
\end{tcolorbox}

\begin{tcolorbox}[cheatbox, title={\refstepcounter{cheatnum}\label{box:bao_lensing}\thecheatnum. BAO \& Kinematic Unification \hfill\scriptsize\hyperref[sec:deriv_bao_lensing]{[Proof $\downarrow$]}}]
The divergence of the BAO standard ruler evaluates as lensed cosmic information propagating through the Gradient-Index (GRIN) lens. Longitudinal depth evaluates as the derivative of the true coordinate distance. 
\begin{align}
    &\text{1. Longitudinal Depth (Physical vs. Comoving):} \nonumber \\
    &\quad   D_{H,true}(z) = \frac{R_{h,0}}{1+z} \label{eq:bao_longitudinal} \\
    &\quad   D_{H,comoving}(z) = R_{h,0}(1+z) \label{eq:bao_comoving} \\
    &\text{2. Kinematic Unification:} \nonumber \\
    &\quad \bar{a}_0(z) = \frac{c^2}{2\pi D_{H,comoving}(z)} \label{eq:bao_kinematic_unification_z}
\end{align}
\end{tcolorbox}

\begin{tcolorbox}[cheatbox, title={\refstepcounter{cheatnum}\label{box:auxetic_symmetry_covariance}\thecheatnum. Symmetry Breaking \& Covariance \hfill\makebox[0pt][r]{\scriptsize\hyperref[sec:deriv_auxetic_symmetry_covariance]{[Proof $\downarrow$]}}}]
Adjacent topological circulations symmetrically break tension, establishing an interstitial stagnation plane. Opposing pressure gradients cancel ($\nabla P_{net} = 0$), while pressure deficits superimpose ($\Delta P_{net}$). 
\begin{align}
    &\text{1. The Stagnation Plane (Lagrange Point):} \nonumber \\
    &\quad \nabla P_{net}(0) = 0 \implies \mathbf{a} = 0 \label{eq:gradient_cancellation} \\
    &\text{2. Hidden Action (Structural Strain):} \nonumber \\
    &\quad \mathbf{a} = 0, \quad \sigma_{ij} \gg 0 \label{eq:hidden_action_stress} \\
    &\text{3. Scalar Superposition (Pressure Deficit):} \nonumber \\
    &\quad \Delta P_{net} = P_{dyn,A} + P_{dyn,B} \label{eq:scalar_superposition} \\
    &\text{4. Gravitational Time Dilation:} \nonumber \\
    &\quad d\tau = dt_{bg} \sqrt{1 - \frac{\Delta P_{net}}{P_c}} \label{eq:scalar_time_dilation} \\
    &\text{5. Acoustic Covariance:} \nonumber \\
    &\quad c_{obs} = \frac{dr_{wave}}{d\tau} \equiv c \label{eq:acoustic_covariance}
\end{align}
\end{tcolorbox}

\begin{tcolorbox}[cheatbox, title={\refstepcounter{cheatnum}\label{box:galactic_dynamics}\thecheatnum. Galactic Kinematics \& The Continuum Shear Transition \hfill\makebox[0pt][r]{\scriptsize\hyperref[sec:deriv_galactic_dynamics]{[Proof $\downarrow$]}}}]
The flat rotation curves of galaxies evaluate as the kinematic transition of the metric from a yielding state into a rigid plug flow.
\begin{align}
    &\quad \bar{a}_0(z) = \frac{c^2}{2\pi R_{h,0}(1+z)} \label{eq:a0z} \\
    & \quad \quad  \quad \approx 1.10 \times 10^{-10} \text{ m/s}^2 \quad (z = 0) \nonumber \\
    &\quad v(r) = \sqrt{\frac{G M_{bar}}{r}} \quad \text{(Yielding Core)} \label{eq:kepler_core} \\
    &\quad r_t = \sqrt{\frac{G M_{bar}}{\bar{a}_0(z)}} \label{eq:transition_radius} \\
    &v_{rot}^4 = \bar{a}_0(z) G M_{bar} \quad \text{(Bingham Plug Flow)} \label{eq:tully_fisher_derivation} \\
    &\quad R_{max} = \frac{v_{rot}}{H_0} \label{eq:r_max}
\end{align}
\end{tcolorbox}

\begin{tcolorbox}[cheatbox, title={\refstepcounter{cheatnum}\label{box:galactic_lensing}\thecheatnum. Galactic Lensing (Stress Birefringence) \hfill\makebox[0pt][r]{\scriptsize\hyperref[sec:deriv_galactic_lensing]{[Proof $\downarrow$]}}}]
Gravitational lensing evaluates as stress birefringence bounded by the refractive index. 
\begin{align}
    &n_{iso}(r) = \left( 1 - \frac{2G M_{bar}}{r c^2} \right)^{-1}  \label{eq:n_isotropic} \\
    &n_{cyl}(r) = \left( 1 - 2\left(\frac{v_{rot}^2}{c^2}\right) \ln \left( \frac{R_{max}}{r} \right) \right)^{-1}  \label{eq:n_cylindrical} \\
    &\Delta \theta_{tot}(I) \propto \int \Big( \nabla_r n_{iso} + \sin(I) \nabla_r n_{cyl} \Big) \, dl \label{eq:total_deflection} 
\end{align}
\end{tcolorbox}
\columnbreak
\vspace{0.2cm}
\noindent\textbf{\large \color{academiabrown} PHASE II: VACUUM RHEOLOGY AND THERMODYNAMICS}
\vspace{0.2cm}

\begin{tcolorbox}[cheatbox, title={\refstepcounter{cheatnum}\label{box:thermo_vacuum}\thecheatnum. Thermodynamic Vacuum \& Entropy \hfill\makebox[0pt][r]{\scriptsize\hyperref[sec:deriv_thermo_vacuum]{[Proof $\downarrow$]}}}]
Via the Unruh-Hawking effect, boundary acceleration generates a thermodynamic floor. 
\begin{align}
    &\text{1. The Unruh-Hawking Temperature Floor:} \nonumber \\
    &\quad a_{h}(t) = \frac{c}{t} \implies T_{cosmic}(t) = \frac{\hbar}{2\pi k_B t} \label{eq:tcosmic} \\
    &\text{2. Particle Entropy Quantum:} \nonumber \\
    &\quad S_p = \frac{(\rho + P)V}{T} \implies S_p = \frac{8\pi}{3} k_B \label{eq:sp}
\end{align}
\end{tcolorbox}

\begin{tcolorbox}[cheatbox, title={\refstepcounter{cheatnum}\label{box:unified_eos}\thecheatnum. The Unified Equation of State \hfill\makebox[0pt][r]{\scriptsize\hyperref[sec:deriv_unified_eos]{[Proof $\downarrow$]}}}]
The continuum transitions from a null-action vacuum to a phase boundary.
\begin{align}
    &\text{1. Fundamental Vacuum Energy \& Null Action:} \nonumber \\
    &\quad E_{vac}(t) = \frac{\hbar}{t} \implies \mathcal{L}_{vac}(c) = 0 \label{eq:vacuum_null_action} \\
    &\text{2. The Topological Equation of State:} \nonumber \\
    &\quad PV = \frac{2\pi}{3} k_B T \implies U = 2\pi k_B T \label{eq:unified_topological_eos} \\
    &\text{3. Rest Mass Equivalence \& Bag Pressure:} \nonumber \\
    &\quad U_{topo} = mc^2 \implies P_{int} = 2\pi P_p \label{eq:unified_pressure_ratio}
\end{align}
\end{tcolorbox}

\begin{tcolorbox}[cheatbox, title={\refstepcounter{cheatnum}\label{box:rheological_gradient}\thecheatnum. The Continuous Rheological Gradient \hfill\makebox[0pt][r]{\scriptsize\hyperref[sec:deriv_rheological_gradient]{[Proof $\downarrow$]}}}]
The continuum transitions from a nearly incompressible vacuum to a perfectly incompressible core, governed by circulation velocity ($v$). 
\begin{align}
    &\lim_{v \to 0} \nu(r) \approx 0.496 \implies \text{Transverse Compliance} \label{eq:poisson_macro} \\
    &\nu(\lambda_c) \in [0.496, 0.5) \implies \text{Phase Stiffening} \label{eq:poisson_compton} \\
    &\lim_{v \to c} \nu(r_c) = 0.5 \implies \text{Incompressible Core} \label{eq:poisson_local} 
\end{align}
\end{tcolorbox}

\begin{tcolorbox}[cheatbox, title={\refstepcounter{cheatnum}\label{box:fine_structure}\thecheatnum. The Fine-Structure Constant \& Shear \hfill\makebox[0pt][r]{\scriptsize\hyperref[sec:deriv_fine_structure]{[Proof $\downarrow$]}}}]
The Fine-Structure Constant ($\alpha$) is the ratio of the vacuum's transverse shear capacity to its bulk compressive capacity. 
\begin{align}
    &\alpha = \frac{\Pi_{shear}}{P_c}   \label{eq:alpha_ratio} \\
    &P_c = \frac{1}{2} \rho_\tau c^2, \quad \Pi_{shear} = \frac{1}{2} \rho_\tau v_s^2  \label{eq:longitudinal_capacity} \\
    &v_s = c \sqrt{\alpha} \approx 0.085c  \label{eq:bare_shear_velocity}
\end{align}
\end{tcolorbox}

\begin{tcolorbox}[cheatbox, title={\refstepcounter{cheatnum}\label{box:unified_deborah}\thecheatnum. The Cosmological Deborah Number \hfill\scriptsize\hyperref[sec:deriv_unified_deborah]{[Proof $\downarrow$]}}]
The phase transition of the yielding continuum into a solid state evaluates via the Complex Dynamic Modulus ($G^*$). The Deborah Number evaluates as the geometric ratio of the interaction energy ($E$) to the continuum baseline ($m_e c^2$). % :::LLM_EDIT::: Stripped banned "fluid" and "crystal" terminology per Protocol Rule 3.
\begin{align}
    De &= \alpha \left( \frac{E}{m_e c^2} \right) \label{eq:unified_deborah} \\
    G^*(E) &= \mu_s \left[ \frac{De^2}{1 + De^2} \right] + i \mu_s \left[ \frac{De}{1 + De^2} \right] \label{eq:deborah_modulus} \\
    De &= \frac{\alpha}{\alpha_s} \label{eq:de_ratio} \\
    \alpha_s &= \frac{m_e c^2}{E} \label{eq:angular_distribution} \\
    \beta(\alpha_s) &= - \alpha_s \label{eq:beta_function} % :::LLM_EDIT::: Repositioned to follow angular distribution.
\end{align}
\textit{Microscopic Phase Symmetry \& Fracture Limits:} % :::LLM_EDIT::: Added subsection to organize new micro-scale limits.
\begin{align}
    De \to 1 &\implies 2 E_{crit} = 140.04 \text{ MeV} \equiv m_{\pi^{\pm}} \label{eq:de_glass} \\
    L_{crit} &= r_e \approx 2.81 \text{ fm} \label{eq:de_range} \\ % :::LLM_EDIT::: Fixed De condition to correctly map the fracture limit.
    \gamma_{yield} &= \frac{1}{\alpha} \approx 137 \label{eq:gamma_yield} \\
    Z &= \frac{1}{\alpha} \approx 137 \label{eq:atomic_phase_limit} \\
    De = 1 &\implies \alpha_s = \alpha \label{eq:electro_strong_symmetry} \\
    g^2 &= 4\pi\alpha \label{eq:strong_charge_symmetry} % :::LLM_EDIT::: Added the strong charge symmetry derivation.
\end{align}
\textit{Macroscopic Phase \& Astrophysical Collapse:} % \textit{Macroscopic Phase \& Astrophysical Collapse:} % :::LLM_EDIT::: Added subsection to organize new macro-scale limits.
\begin{align}
    De_{min} &= \alpha (2\pi)^2 \left( \frac{E_{unruh}}{E_{rest}} \right) \label{eq:de_gravity} \\
    M &= \frac{M_P^3}{\sqrt{2}\, m_n^2} \approx 1.31 M_\odot \label{eq:chandra_mass_limit} \\ % :::LLM_EDIT::: Applied alpha_yield = 1/sqrt(2) compliance bound.
    M &= \frac{M_P^3}{4} \sqrt{ \frac{\alpha^3}{m_n m_e^3} } \approx 22.7 M_\odot \label{eq:astral_mass_limit} \\ % :::LLM_EDIT::: Applied alpha_yield = 1/sqrt(2), simplifying sqrt(8) to 4.
    r_{glass} &\approx r_s (1 + \alpha^2) \label{eq:r_glass} \\
    R_{s,crit} &= a_0(1 - \alpha^2) \label{eq:terminal_radius}
\end{align}
\end{tcolorbox}

\begin{tcolorbox}[cheatbox, title={\refstepcounter{cheatnum}\label{box:viscoelastic_stiffening}\thecheatnum. Phase Stiffening \& Time Dilation \hfill\scriptsize\hyperref[sec:deriv_viscoelastic_stiffening]{[Proof $\downarrow$]}}]
Relativistic time dilation, strong interaction coupling, and the cosmic speed limit evaluate as projections of kinematic strain consuming metric compliance. 
\begin{align}
    &\text{1. The Compliance Limit ($\alpha\_s$):} \nonumber \\
    &\quad \alpha_s^2 = 1 - \frac{P_{dyn} + P_{shear}}{P_c} \label{eq:admittance_pressure_ratio} \\
    &\text{2. The Yukawa-Bernoulli Identity:} \nonumber \\
    &\quad \alpha_s^2 = 1 - \frac{g^2}{g_{max}^2} \label{eq:yukawa_identity} \\
    &\text{3. Relativistic Packing \& The Lorentz Pressure Limit:} \nonumber \\
    &\quad \rho_{eff} = \rho_{\tau} \gamma \label{eq:relativistic_packing_admittance} \\
    &\quad \gamma = \sqrt{\frac{P_c}{P_c - (P_{dyn} + P_{shear})}} \label{eq:lorentz_pressure} \\
    &\text{4. The Kinematic Phase Lock:} \nonumber \\
    &\quad \gamma = \sqrt{\frac{P_c}{P_{static}}} \label{eq:admittance_spin_identity} \\
    &\text{5. The Kinematic Yield Limit (Choked Flow):} \nonumber \\ 
    &\quad \lim_{(P_{dyn} + P_{shear}) \to P_c} \alpha_s = 0 \implies \nabla \cdot \mathbf{v} = 0 \label{eq:cavitation_collapse}
\end{align}
\end{tcolorbox}

\begin{tcolorbox}[cheatbox, title={\refstepcounter{cheatnum}\label{box:universal_impedance}\thecheatnum. Kinematic Vacuum Impedance ($Z_{\tau}$) \scriptsize\hfill\hyperref[sec:deriv_universal_impedance]{[Proof $\downarrow$]}}]
Inertial mass, optical density, quantum electrical resistance, and quantized spin evaluate as the directional geometric limits of the metric's structural resistance.
\begin{align}
    &\text{1. The Transverse Quantum Limit (von Klitzing):} \nonumber \\
    &\quad R_K = k_z Y_{\tau(\perp)} = k_z \left(\frac{h}{e_\tau^2}\right) \label{eq:von_klitzing_impedance} \\
    &\text{2. The Longitudinal Inertial Limit (Effective Mass):} \nonumber \\
    &\quad Z_{\tau(\parallel)} = \hbar^2 \left( \frac{d^2 E}{dk^2} \right)^{-1} \equiv m^* \label{eq:effective_mass_impedance} \\
    &\text{3. The Transverse Limit (Refractive Gravity):} \nonumber \\
    &\quad Z_{\tau(\perp)} \propto n(r) = \frac{P_c}{P_{static}} \approx \left( 1 - \frac{2GM}{rc^2} \right)^{-1} \label{eq:refractive_impedance} \\
    &\text{4. The Rotational Phase Limit (Gyroscopic Inertia):} \nonumber \\
    &\quad Z_{\tau(\circlearrowleft)} = \frac{E}{\omega} \equiv \hbar \label{eq:planck_impedance} \\
    &\lim_{P_{static} \to 0} Z_{\tau(\circlearrowleft)} \to \infty \implies \text{Spin Quantization } (\hbar/2) \label{eq:infinite_rotational_impedance}
\end{align}
\end{tcolorbox}

\newpage
% ==========================================
% PHASE III: TOPOLOGICAL MASS GENERATION
% ==========================================
\vspace{0.4cm}
\noindent\textbf{\large \color{academiabrown} PHASE III: TOPOLOGICAL MASS GENERATION}
\vspace{0.2cm}



\begin{tcolorbox}[cheatbox, title={\refstepcounter{cheatnum}\label{box:hydrodynamic_mass}\thecheatnum. The Kinematic Origin of Inertia \scriptsize\hfill\hyperref[sec:deriv_hydrodynamic_mass]{[Proof $\downarrow$]}}]
The kinematic resistance of a boundary to spatial translation evaluates as metric drag.
\begin{align}
    &\quad m_{added} = \rho_{\tau} V_{added} \label{eq:added_mass} \\
    &\quad V_{eff} = C_m V_\tau \implies m_\tau = C_m \rho_\tau V_\tau \label{eq:strain_volume_equivalence}
\end{align}
\end{tcolorbox}

\begin{tcolorbox}[cheatbox, title={\refstepcounter{cheatnum}\label{box:ns_reg}\thecheatnum. Navier-Stokes Regularization \scriptsize\hfill\hyperref[sec:deriv_ns_reg]{[Proof $\downarrow$]}}]
The Navier-Stokes equations are regularized by the acoustic limit of the metric. 
\begin{align}
    &\text{1. Dynamic Viscosity ($\mu$):} \nonumber \\
    &\quad \mu = \frac{1}{2} \frac{\rho_{\tau} \hbar}{m} \label{eq:vacuum_viscosity} \\
    &\text{2. Prandtl Mixing Length Equivalence:} \nonumber \\
    &\quad \mu = \frac{1}{2} \rho_{\tau} c L \label{eq:prandtl_equivalence} \\
    &\text{3. The Critical Reynolds Number ($Re_c$):} \nonumber \\
    &\quad Re_c = \frac{\rho_{\tau} c L}{\mu} = 2 \label{eq:critical_reynolds} \\
    &\text{4. The Dynamic Elasticity Number ($El$):} \nonumber \\
    &\quad El = \frac{\alpha (E \cdot l_p)}{2 \hbar v} \label{eq:dynamic_elasticity}
\end{align}
\end{tcolorbox}

\begin{tcolorbox}[cheatbox, title={\refstepcounter{cheatnum}\label{box:fluid_inertia}\thecheatnum. Origin of Inertia (Strain Compactness) \scriptsize\hfill\hyperref[sec:deriv_fluid_inertia]{[Proof $\downarrow$]}}]
Mass and inertia evaluate as the geometric tension generated by kinematic strain.
\begin{align}
    &\text{1. Relativistic Strain Compactness:} \nonumber \\
    &\rho_{eff}(r) = \rho_{\tau} \left( 1 - \frac{P_{dyn}(r) + P_{shear}(r)}{P_c} \right)^{-1/2} \label{eq:fluid_inertia_gradient} \\
    &\text{2. The Isotropic Tension Identity:} \nonumber \\
    &\quad \mathbb{T}_{local} = \mathbb{T}_c + P_{dyn} + P_{shear} \label{eq:tension_identity} \\
    &\text{3. The Geometric Equation of State (Pure Tension):} \nonumber \\
    &\quad \rho_{eff}(r) = \rho_{\tau} \left( \frac{\mathbb{T}_{local}(r)}{\mathbb{T}_c} \right)^{-1/2} \label{eq:tension_density} \\
    &\text{4. The Compliant Boundary:} \nonumber \\
    &\quad \rho_{eff}(\lambda_c) \approx \rho_{\tau} + \rho_{\tau} \left( \frac{P_{dyn}(\lambda_c) + P_{shear}(\lambda_c)}{2P_c} \right) \label{eq:compliant_inertia}
\end{align}
\end{tcolorbox}


\begin{tcolorbox}[cheatbox, title={\refstepcounter{cheatnum}\label{box:volumetric_inverse_cube}\thecheatnum. Volumetric Approximations \scriptsize\hfill\hyperref[sec:deriv_volumetric_inverse_cube]{[Proof $\downarrow$]}}] 
The bounding volume of matter scales inversely proportional to the cube of its total kinematic strain.
\begin{align}
    &\quad V_\tau = \frac{2\pi l_p^4}{C_m r_c} \label{eq:strain_volume_rc} \\ 
    &\quad V_{r_c} = \frac{4}{3}\pi r_c^3 \label{eq:coordinate_volume} \\
    &\quad V_{r_c} \cdot V_{\tau}^3 = \frac{4\pi}{3C_m^3} \left( \frac{\hbar}{\rho_{\tau} c} \right)^3 = \frac{32\pi^4}{3C_m^3} l_p^{12} \label{eq:volumetric_inverse_cube} 
\end{align}
\end{tcolorbox}

\begin{tcolorbox}[cheatbox, title={\refstepcounter{cheatnum}\label{box:admittance_ns}\thecheatnum. Compliance-Modified Navier-Stokes \scriptsize\hfill\hyperref[sec:deriv_admittance_ns]{[Proof $\downarrow$]}}]
To determine acceleration, the momentum of the metric scales against the total scalar admittance ($\alpha_s$), canceling mass to isolate kinematics. % :::LLM_EDIT::: Cured "evaluate" repetition, purged banned adjectives, and upgraded to total scalar admittance.
\begin{align}
    &\text{1. The Relativistic Density Substitution:} \nonumber \\
    &\quad \rho_a = \rho_\tau \implies \rho_{eff} = \frac{\rho_\tau}{\alpha_s} \label{eq:ns_density_substitution} \\ % :::LLM_EDIT::: Updated alpha_a to alpha_s.
    &\text{2. The Quantum Reynolds Limit:} \nonumber \\
    &\quad Re_q = \frac{2}{\alpha_s} \implies Re_{q,rest} = \frac{2}{\alpha_{s,rest}} \label{eq:quantum_reynolds} \\ % :::LLM_EDIT::: Updated alpha_a to alpha_s and alpha_rest to alpha_s,rest.
    &\text{3. The Equivalence Principle \& Phase Separation:} \nonumber \\
    &\quad \frac{D\mathbf{v}}{Dt} = -\nabla \Phi_{env} + \frac{\alpha_s}{\rho_\tau} \Sigma \mathbf{F}_{surface} \label{eq:admittance_ns} \\ % :::LLM_EDIT::: Updated alpha_a to alpha_s.
    &\text{4. The Surface Force Expansion:} \nonumber \\
    &\quad \frac{D\mathbf{v}}{Dt} = -\nabla \Phi_{env} \nonumber \\ & + \frac{\alpha_s}{\rho_\tau} \Big[ -\nabla P_{dyn} + \alpha(\nabla \cdot \mathbb{T}_{\perp} - \mu \nabla \times \boldsymbol{\Omega}) \Big] \label{eq:propulsion_master} \\ % :::LLM_EDIT::: Updated alpha_a to alpha_s. Preserved custom line breaks.
    &\text{5. The Bernoulli Gradient:} \nonumber \\ % :::LLM_EDIT::: Removed banned adjective 'Geometric' per Rule 7.
    & \nabla \mathcal{P}_{ambient} = \nabla \mathcal{P}_{static} + \nabla(\mathcal{P}_{dyn} + \mathcal{P}_{shear}) \label{eq:bernoulli_gradient} \\
    &\text{6. The Equation of Motion:} \nonumber \\
    &\quad \frac{D\mathbf{v}}{Dt} = -\nabla \mathcal{P}_{ambient} + \alpha_s \Big[ -\nabla \mathcal{P}_{dyn}  \nonumber \\ &\quad + \alpha(\nabla \cdot \mathbb{T}_{\perp, kin} - \nu \nabla \times (\nabla\lambda \times \nabla\beta_\Omega)) \Big] \label{eq:kinematic_ns} 
\end{align}
\end{tcolorbox}

\begin{tcolorbox}[cheatbox, title={\refstepcounter{cheatnum}\label{box:localized_metric_strain}\thecheatnum. The Localized Metric \hfill\scriptsize\hyperref[sec:deriv_localized_metric_strain]{[Proof $\downarrow$]}}]
The boundary conditions evaluate as pressure inversion.
\begin{align}   
    &\text{1. The Cavitation Boundary (Yield):} \nonumber \\
    &\quad P_\Phi(r_c) = P_c \label{eq:pressure_cavitation_boundary} \\
    &\text{2. Optical Dilatation \& Refractive Strain:} \nonumber \\
    &\quad \varepsilon(r) = \frac{1}{\alpha_s(r)} = \left( 1 - \frac{r_c}{r} \right)^{-1/2} \label{eq:strain_redshift_identity}  \\
    &\text{3. The Monopole Hydrodynamic Horizon ($R_\tau$):} \nonumber \\
    &\quad R_\tau = \left( r_c R_{h,0}^2 \right)^{1/3} \label{eq:hydrodynamic_horizon_equivalence} \\
    &\text{4. The Irrotational Vortex (Phase Constraint):} \nonumber \\
    &\quad v_\phi(r) = c \left( \frac{r_c}{r} \right) \iff \Gamma = 2\pi r_c c \label{eq:irrotational_velocity_field} \\
    &\text{5. The Dipole Kinematic Horizon ($R_\omega$):} \nonumber \\
    &\quad R_\omega = \sqrt{2\pi r_c R_{h,0}} \label{eq:dipole_horizon_equivalence} \\
    &\text{6. Strain Superposition ($R_{\tau(body)}$):} \nonumber \\
    &\quad \frac{2 G M_{body}}{c^2 R_{\tau(body)}} = \frac{R_{\tau(body)}^2}{R_{h,0}^2} \label{eq:macroscopic_equilibration} \\
    &\quad R_{\tau(body)} = \left( R_s R_{h,0}^2 \right)^{1/3} \label{eq:body_halo_scaling} \\
    &\quad M_{body} = \frac{M_P}{\sqrt{2}} \label{eq:planck_intersection} \\
    &\text{7. Topological Magnification (Mass Identity):} \nonumber \\
    &\quad Mag = \left( \frac{R_{h,0}}{r_c} \right)^{2/3} \label{eq:conformal_mag_substituted} \\
    &\text{8. Total Internal Reflection:} \nonumber \\
    &\quad r_{crit} = 1.5 r_c \label{eq:critical_optical_radius} \\
    &\text{9. Deep Inelastic Scattering:} \nonumber \\
    &\quad \sigma_{total}(Q^2 \gg m^2 c^2) = \sum_{i=1}^{c(K)} \sigma_i \label{eq:deep_inelastic} 
\end{align}
\end{tcolorbox}




\begin{tcolorbox}[cheatbox, title={\refstepcounter{cheatnum}\label{box:unified_engine}\thecheatnum. The Unified Topological Engine \scriptsize\hfill\hyperref[sec:deriv_unified_engine]{[Proof $\downarrow$]}}]
\begin{align}
    &\text{1. Volumetric Continuity (Clebsch Vorticity):} \nonumber \\
    &\quad \nabla \cdot (\lambda \nabla \beta) = 0 \label{eq:engine_transduction} \\
    &\text{2. Hydrostatic Supremacy (Phase Transform):} \nonumber \\
    &\quad \frac{P_{out}}{P_{in}} = \frac{1}{P_{in}} \oint_{\partial S} P_{in} \, d\theta = 2\pi \label{eq:engine_supremacy} \\
    &\text{3. The Planck Redline (Kinematic Lockout):} \nonumber \\
    &\quad \lim_{r_c \to \sqrt{2} l_p} v_{core} \to c \implies m_{max} = \frac{M_P}{\sqrt{2}} \label{eq:engine_redline} \\ % :::LLM_EDIT::: Aligned redline with Eq. planck_intersection yield limit.
    &\text{4. Hydrostatic Instability (The $m^4$ Decay Limit):} \nonumber \\
    &\quad \frac{V_\tau}{V_{r_c}} = \left(\frac{3c^3}{4\pi C_{m,i} \rho_{\tau}\hbar^3}\right) m_\tau^4 \label{eq:engine_decay} 
\end{align}
\end{tcolorbox}

\begin{tcolorbox}[cheatbox, title={\refstepcounter{cheatnum}\label{box:bag_pressure}\thecheatnum. Confinement Pressure \& Work \scriptsize\hfill\hyperref[sec:deriv_bag_pressure]{[Proof $\downarrow$]}}]
The geometric centripetal boundary condition ($a_{loc} = c/2\pi t_{\text{loc}}$) dictates the inward Bag Pressure:
\begin{equation}\label{eq:pp}
    P_p = \frac{m^4 c^5}{8\pi^2 \hbar^3} 
\end{equation}

\end{tcolorbox}


\begin{tcolorbox}[cheatbox, title={\refstepcounter{cheatnum}\label{box:quantum_spin}\thecheatnum. Kinematic Circulation \& The Origin of Spin \scriptsize\hfill\hyperref[sec:deriv_quantum_spin]{[Proof $\downarrow$]}}] 
Matter operates as closed Clebsch vorticity. Quantum Spin dictates the kinematic circulation. 
\begin{align}
    &\text{1. Kinematic Angular Momentum \& Spin Limit:} \nonumber \\
    &\quad \mathbf{S} = \iiint_V \rho_\tau (\mathbf{r} \times (\lambda \nabla \beta)) \, dV \implies S = \hbar \label{eq:transverse_shear_spin} \\
    &\text{2. The Baseline Circulation Limit:} \nonumber \\
    &\quad \Gamma = \oint_C (\lambda \nabla \beta) \cdot d\mathbf{l} = \frac{h}{m} \label{eq:fluid_circulation} \\
    &\text{3. Spinor Topology ($4\pi$ Metric Untangling):} \nonumber \\ 
    &\quad S_{fermion} = m \left( \frac{\Gamma}{4\pi} \right) = \frac{\hbar}{2} \label{eq:spinor_topology}
\end{align}
\end{tcolorbox}

\begin{tcolorbox}[cheatbox, title={\refstepcounter{cheatnum}\label{box:volumetric_energy}\thecheatnum. Kinetic-Potential Strain Conservation \scriptsize\hfill\hyperref[sec:deriv_volumetric_energy]{[Proof $\downarrow$]}}]
Energy resolves as displaced volume. Potential and kinetic energies emerge as partitions of the scalar admittance ($\alpha_s$). 
\begin{align}
    &\text{1. Kinematic Drag Factor:} \nonumber \\
    &\quad \gamma = \frac{1}{\alpha_s} \label{eq:strain_lorentz} \\
    &\text{2. Kinetic Displacement Volume:} \nonumber \\
    &\quad V_{kin} = \left( \frac{1}{\alpha_s} - 1 \right) V_\tau \label{eq:kinetic_volume} \\
    &\text{3. Potential Displacement Volume:} \nonumber \\
    &\quad V_{pot} = \frac{1}{C_{m,i}\rho_{\tau} c^2} \int_{x_1}^{x_2} \mathbf{F} \cdot d\mathbf{x} \label{eq:potential_volume} \\
    &\text{4. Isochoric Transmutation:} \nonumber \\
    &\quad \dot{V}_{kin} = -\dot{V}_{pot} \label{eq:isochoric_transmutation}
\end{align}
\end{tcolorbox}


\newpage
% ==========================================
% PHASE IV: WAVE KINEMATICS & ELECTROMAGNETISM
% ==========================================
\vspace{0.4cm}
\noindent\textbf{\large \color{academiabrown} PHASE IV: WAVE KINEMATICS \& ELECTROMAGNETISM}
\vspace{0.2cm}

\begin{tcolorbox}[cheatbox, title={\refstepcounter{cheatnum}\label{box:unified_lagrangian}\thecheatnum. The Unified Lagrangian \& The Origin of Action \hfill\makebox[0pt][r]{\scriptsize\hyperref[sec:deriv_unified_lagrangian]{[Proof $\downarrow$]}}}]
The Lagrangian density evaluates as pressure. The continuum minimizes pressure gradients (The Principle of Least Action).
\begin{align}
    &\text{1. The Kinematic Capacities:} \nonumber \\
    &\quad P_{ambient} = \frac{1}{2}\rho_a c_a^2, \quad \Pi_{shear} = \alpha P_{ambient} \label{eq:kinematic_capacities} \\
    &\text{2. The Covariant Field Lagrangian:} \nonumber \\
    & \mathcal{L}_{particle} = - P_{shear} \Bigg[ 1 - \frac{1}{P_{ambient}} \Bigg( \frac{1}{2}\rho_a \Big[ (\nabla h)^2 + \nonumber \\
    &  \frac{1}{\alpha} \big( (\nabla \times \mathbf{A})^2 + (\nabla \times (\lambda \nabla \beta_\Omega))^2 \big) \Big] + P_{shear} \Bigg) \Bigg]^{1/2} \label{eq:dynamic_lagrangian} \\
    &\text{3. The Parabolic Limit ($\mathcal{K} - \mathcal{P}_{net}$):} \nonumber \\
    &\quad \mathcal{P}_{net} = P_{shear} \left( 1 - \frac{P_{shear}}{2 P_{ambient}} \right) \label{eq:net_potential} \\
    &\quad \mathcal{L}_{particle} \approx - \mathcal{P}_{net} + \frac{1}{2}\rho_0 v_{eq}^2 \equiv \mathcal{K} - \mathcal{P}_{net} \label{eq:classical_lagrangian_limit} \\
    &\text{4. Topological Confinement (Cavitation):} \nonumber \\
    &\quad v_{\perp, max} = c_a \sqrt{\alpha \left( 1 - \frac{P_{shear}}{P_{ambient}} \right)} \label{eq:max_transverse_velocity} \\
    &\quad \lim_{v_\perp \to v_{\perp, max}} \alpha_k = 0 \implies De \to \infty \label{eq:deborah_confinement} \\
    &\text{5. The Topological Phase-Space Lagrangian:} \nonumber \\
    &\mathcal{L}_{kin} = - \mathcal{P}_{shear} \Bigg[ 1 - \frac{1}{\mathcal{P}_{ambient}} \Bigg( \frac{1}{2} v_{tot}^2 \Big[ \cos^2(\theta) \nonumber \\
    &\quad + \frac{1}{\alpha}\sin^2(\theta) \Big] + \mathcal{P}_{shear} \Bigg) \Bigg]^{1/2} \label{eq:angular_lagrangian} \\
    &\text{6. The Universal Baseline:} \nonumber \\
    &\quad P_{ambient} = P_c \alpha_g^2 \label{eq:gravitational_compliance} \\
    &\quad \mathcal{L}_{particle} = - P_{shear} \sqrt{1 - \frac{P_{dyn} + P_{shear}}{P_c \alpha_g^2}} \label{eq:nested_lagrangian}
\end{align}
\end{tcolorbox}

\begin{tcolorbox}[cheatbox, title={\refstepcounter{cheatnum}\label{box:legendre_hamiltonian}\thecheatnum. The Legendre Hamiltonian Translation \\ \null\hfill\makebox[0pt][r]{\scriptsize\hyperref[sec:deriv_legendre_hamiltonian]{[Proof $\downarrow$]}}}]
The total energy state evaluates as the Legendre transform of the covariant Lagrangian.
\begin{align}
    &\text{1. Topological Conjugate Momentum ($\mathbf{p}_{macro}$):} \nonumber \\
    &\quad \mathbf{p}_{macro} = \frac{P_{shear}}{c^2 \alpha_g^2} \frac{\mathbf{v}_{macro}}{\alpha_k} \equiv \frac{\rho_0}{\alpha_g^2} \frac{\mathbf{v}_{macro}}{\alpha_k} \label{eq:conjugate_momentum} \\
    &\text{2. The Covariant Energy Function ($\mathcal{H}_{particle}$):} \nonumber \\
    &\quad \mathcal{H}_{particle} = P_{shear} \frac{\alpha_{rest}^2}{\alpha_k} \label{eq:covariant_hamiltonian} \\
    &\text{3. The Gravitational Rest State ($\mathcal{H}_{rest}$):} \nonumber \\
    &\quad \alpha_{rest}^2 = 1 - \frac{P_{shear}}{P_c \alpha_g^2} \implies \mathcal{H}_{rest} = P_{shear} \alpha_{rest} \label{eq:hamiltonian_rest_mass} \\
    &\text{4. The Lorentz Factor \& Dispersion Relation:} \nonumber \\
    &\quad \alpha_k = \alpha_{rest} \sqrt{1 - \frac{v_{macro}^2}{v_{max}^2}} \label{eq:lorentz_compliance} \\
    &\quad \mathcal{H}_{particle} = \frac{\mathcal{H}_{rest}}{\sqrt{1 - v_{macro}^2/v_{max}^2}} \label{eq:relativistic_energy} \\
    &\quad \mathcal{H}_{particle}^2 - (p_{macro} v_{max})^2 = \mathcal{H}_{rest}^2 \label{eq:dispersion_relation} \\
    &\text{5. Topological Mass Extraction:} \nonumber \\
    &\quad E_{bare} = P_{shear} V_{\psi} \equiv P_c V_{\tau} = \frac{1}{2}\rho_\tau c^2 V_{\tau} \label{eq:rest_energy_integral} \\
    &\quad m_{bare} = \frac{1}{2}\rho_\tau V_{\tau}, \quad m_{obs} = m_{bare} \alpha_{rest} \label{eq:topological_mass} \\
    &\text{6. The Kinematic Field State ($\mathcal{H}_{kin}$):} \nonumber \\
    &\quad \mathcal{H}_{kin} = \nonumber \\ &\frac{\mathcal{P}_{shear} \alpha_{rest}^2}{\sqrt{1 - \frac{\frac{1}{2} \Big[ (\nabla h)^2 + \frac{1}{\alpha}(\nabla \times \mathbf{A})^2 + \frac{1}{\alpha}(\nabla \times (\lambda \nabla \beta_\Omega))^2 \Big] + \mathcal{P}_{shear}}{\mathcal{P}_c \alpha_g^2}}} \label{eq:kinematic_hamiltonian} \\
    &\text{7. The Reduced Phase-Space State:} \nonumber \\
    &\mathcal{H}_{kin}(v_{tot}, \theta) = \frac{\mathcal{P}_{shear} \alpha_{rest}^2}{\sqrt{1 - \frac{\frac{1}{2} v_{tot}^2 \Big[ \cos^2(\theta) + \frac{1}{\alpha}\sin^2(\theta) \Big] + \mathcal{P}_{shear}}{\mathcal{P}_c \alpha_g^2}}} \label{eq:angular_hamiltonian}
\end{align}
\end{tcolorbox}

\begin{tcolorbox}[cheatbox, title={\refstepcounter{cheatnum}\label{box:dirac_isomorphism}\thecheatnum. Vorticity Equilibrium (Dirac Isomorphism) \scriptsize\hfill\hyperref[sec:deriv_dirac_isomorphism]{[Proof $\downarrow$]}}]
Fermion kinematics evaluate as the continuous mechanics of a gyroscopically stabilized acoustic vortex.
\begin{align}
    &\text{1. The Spinor (Cartan Tetrad):} \nonumber \\
    &\quad \psi \equiv \mathbf{\hat{e}}_i \label{eq:spinor_identity} \\
    &\text{2. The Kinematic Operator (Clebsch Bridge):} \nonumber \\
    &\quad i\gamma^\mu \partial_\mu \equiv \nabla \times \mathbf{v}_\psi \equiv \nabla \times (\lambda \nabla \beta_\Omega) \label{eq:gamma_curl} \\
    &\text{3. The Quantum Rest Mass Identity:} \nonumber \\
    &\bar{\psi} \psi = \frac{1}{V_{\psi}} \implies \mathcal{L}_{mass} = - (\hbar \omega) \bar{\psi}\psi = - P_{shear} \label{eq:Quantum_mass_lagrangian} \\
    &\text{4. The Vorticity Equilibrium Lagrangian:} \nonumber \\
    &\mathcal{L}_{Dirac} \equiv \frac{1}{2}\rho_a |\nabla \times (\lambda \nabla \beta_\Omega)|^2 - P_{shear} = 0 \label{eq:dirac_fluid_lagrangian} \\
    &\text{5. Antimatter Annihilation:} \nonumber \\
    &\quad \bar{\psi}\psi \equiv (W_+ + W_-) = 0 \label{eq:chiral_annihilate} 
\end{align}
\end{tcolorbox}

\begin{tcolorbox}[cheatbox, title={\refstepcounter{cheatnum}\label{box:wave_mechanics}\thecheatnum. Continuous Wave Mechanics \scriptsize\hfill\hyperref[sec:deriv_wave_mechanics]{[Proof $\downarrow$]}}]
The Kinematic Lagrangian maps the mechanical wave spectrum and particle identities as geometric constraints.
\begin{align}
    &\text{1. The Continuous Wave Solutions:} \nonumber \\
    & \mathbf{v}_{\parallel} = \left[ c_a \cos(\theta) \cos(\mathbf{k} \cdot \mathbf{r} - \omega t) \right] \hat{\mathbf{k}} \label{eq:wave_p} \\
    & \mathbf{v}_{\perp} = \left[ c_a \sqrt{\alpha} \sin(\theta)\cos(\phi) \cos(\mathbf{k} \cdot \mathbf{r} - \omega t) \right] \hat{\mathbf{n}} \label{eq:wave_s} \\
    & \mathbf{v}_{\circlearrowleft} = \left[ c_a \sqrt{\alpha} \sin(\theta)\sin(\phi) \cos(\mathbf{k} \cdot \mathbf{r} - \omega t) \right] \hat{\mathbf{m}} \label{eq:wave_rot} \\
    &\text{2. Acoustic Gravitational Potential:} \nonumber \\
    &\quad \Phi(r,t) = - \frac{(\mathbf{v}_{\parallel})^2}{2} \label{eq:p_wave_potential} \\
    &\text{3. Relativistic Kinematics (The Lorentz Angle):} \nonumber \\
    &\quad \sin(\theta) = \sqrt{1 - \frac{v_{macro}^2}{v_{max}^2}} \equiv \frac{\alpha_k}{\alpha_{rest}} \equiv \frac{1}{\gamma} \label{eq:lorentz_angle} \\
    &\text{4. Transverse Charge \& Annihilation Yield:} \nonumber \\
    &\quad Q = \text{sgn}(\phi) \cos^2(\phi) \label{eq:charge_projection} \\
    &\quad \sum (\lambda \nabla \beta_\Omega) = 0 \nonumber \\ &\quad \implies P_{shear} \to 0 \implies \text{Boundary Yield} \label{eq:annihilation_limit} \\
    &\text{5. Nuclear Yield (Beta Decay):} \nonumber \\
    &\quad \mathcal{L}_{kin, n}(udd) \to \mathcal{L}_{kin, p}(uud) \nonumber \\ &\quad + \mathcal{L}_{kin, e}(\phi=-\pi) + \mathcal{L}_{kin, \bar{\nu}}(\theta \to 0) \label{eq:neutron_rupture}
\end{align}
\end{tcolorbox}





\begin{tcolorbox}[cheatbox, title={\refstepcounter{cheatnum}\label{box:em_lc_isomorphism}\thecheatnum. The LC-Acoustic Isomorphism \scriptsize\hfill\hyperref[sec:deriv_em_lc_isomorphism]{[Proof $\downarrow$]}}]
Electromagnetic fields resolve as 1D to 3D projections.
\begin{align}
    &\text{1. Vorticity \& Kinematic Acceleration:} \nonumber \\
    &\quad \nabla \cdot \boldsymbol{\Omega} = 0 \quad \text{\&} \quad \nabla \times \mathbf{a}_\perp = -\frac{\partial \boldsymbol{\Omega}}{\partial t} \label{eq:maxwell_gauss_faraday} \\
    &\text{2. Mass Equivalent Strain (Gauss's Law):} \nonumber \\
    &\quad \nabla \cdot \mathbf{a}_\perp = \frac{\rho_e c^2}{\epsilon_{\tau}} \label{eq:maxwell_gaussE} \\
    &\text{3. Macroscopic Modulus \& Density Mapping:} \nonumber \\
    &\quad \rho_{\tau} = \frac{\epsilon_{\tau}}{A_{shear, p}}, \quad K_\tau = \frac{F_p}{A_{shear, p}} \label{eq:lc_constants} \\
    &\text{4. The Isomorphic Circuit Elements:} \nonumber \\
    &\quad m_\tau = \frac{1}{2} \epsilon_{\tau} L_\tau \label{eq:lc_mass} \\
    &\quad E_{kin} = \frac{1}{2} m_{\tau} v_{\perp}^2 \implies P_{dyn} = \frac{1}{2}\rho_{\tau} v_{\perp}^2 \label{eq:velocity_pressure} \\
    &\quad E_{ind} = \frac{1}{2} L_H I^2 = \frac{1}{2} k x^2 \label{eq:lc_inductance} \\
    &\text{5. The 1D Projection (Ampère-Maxwell):} \nonumber \\
    &\quad \nabla \times \boldsymbol{\Omega} = \mu_{\tau} \mathbf{J}_{1D} + \mu_{\tau} \epsilon_{\tau} \frac{\partial \mathbf{a}_\perp}{\partial t} \label{eq:maxwell_curlB} \\
    &\text{6. Kinematic Energy Flux (The Poynting Vector):} \nonumber \\
    &\quad \mathbf{S}_{kin} = \frac{1}{\mu_{\tau}} (\mathbf{a}_\perp \times \boldsymbol{\Omega}) \label{eq:poynting_flux} \\
    &\text{7. Dual Transverse Geometries:} \nonumber \\
    &\quad \mathcal{V} \equiv v_{\perp} \quad \text{\&} \quad \boldsymbol{\Omega} \equiv \nabla \times \mathbf{v}_{\perp}  \label{eq:dual_transverse} \\
    &\text{8. Dimensional Calibration ($\epsilon_0$) \& Phase Limit:} \nonumber \\
    &\epsilon_0 = \epsilon_{\tau} \left( \frac{e}{e_{\tau}} \right)^2,
    \quad \mathcal{V}_{max} \approx 3.11 \times 10^{27} \text{ V} \label{eq:epsilon_0_derivation}
\end{align}
\end{tcolorbox}

\begin{tcolorbox}[cheatbox, title={\refstepcounter{cheatnum}\label{box:tensor_projections}\thecheatnum. 1D Fractional Flux \hfill\makebox[0pt][r]{\scriptsize\hyperref[sec:deriv_tensor_projections]{[Proof $\downarrow$]}}}]
Fractional phase states evaluate as geometric artifacts of projecting an isotropic 3D circulation onto a 1D observation axis.
\begin{align}
    &\text{1. The Longitudinal Projection (1/3 Flux):} \nonumber \\
    &\quad \Phi_{\parallel} = \Phi_{\circlearrowleft} \left( \frac{1}{4\pi} \iint_{\Omega} \cos^2 \theta \, d\Omega \right) = \frac{1}{3} \Phi_{\circlearrowleft} \label{eq:tensor_long} \\
    &\text{2. The Transverse Projection (2/3 Flux):} \nonumber \\
    &\quad \Phi_{\perp} = \Phi_{\circlearrowleft} \left( \frac{1}{4\pi} \iint_{\Omega} \sin^2 \theta \, d\Omega \right) = \frac{2}{3} \Phi_{\circlearrowleft} \label{eq:tensor_trans} \\
    &\text{3. Flux Conservation:} \nonumber \\
    &\quad \Phi_{\parallel} + \Phi_{\perp} = \Phi_{\circlearrowleft} \label{eq:projection_conservation} \\
    &\text{4. The 1D Fractional Phase Flux ($\Phi_0$):} \nonumber \\
    &\quad \Phi_0 = \frac{1}{3}\Phi_{\circlearrowleft} \label{eq:1d_phase_flux}
\end{align}
\end{tcolorbox}

\begin{tcolorbox}[cheatbox, title={\refstepcounter{cheatnum}\label{box:continuum_constants}\thecheatnum. Constants \& Scale Invariance \scriptsize\hfill\hyperref[sec:deriv_continuum_constants]{[Proof $\downarrow$]}}]
The yield strength of the vacuum operates as a geometric sliding scale. 
\begin{align}
    &\text{1. 3D Volumetric Density ($\rho_{\tau}$):} \nonumber \\
    &\quad \rho_{\tau} = \frac{1}{2\pi} \frac{c^5}{\hbar G^2} \approx 8.21 \times 10^{95} \text{ kg/m}^3 \label{eq:rho_tau_corrected} \\
    &\text{2. 1D Metric Stiffness ($\kappa_{\tau}$):} \nonumber \\
    &\quad \kappa_{\tau} = \frac{c^4}{G l_p} \approx 7.49 \times 10^{78} \text{ N/m} \label{eq:kappa_tau_corrected} \\
    &\text{3. Microscopic Yield (Planck Shear Area):} \nonumber \\
    &\quad A_{shear, p} = \sqrt{2}\pi l_p^2 \approx 1.16 \times 10^{-69} \text{ m}^2 \label{eq:planck_shear_area} \\
    &\text{4. 1D Metric Linear Mass Density ($\epsilon_{\tau}$):} \nonumber \\
    &\quad \epsilon_{\tau} = \frac{1}{\sqrt{2}} \frac{c^2}{G} \approx 9.52 \times 10^{26} \text{ kg/m} \label{eq:epsilon_tau_density} \\
    &\text{5. Vacuum Expectation Value ($v$):} \nonumber \\
    &\quad v = \frac{E_{yield}(G_F)}{\sqrt{2\pi}} \approx 246 \text{ GeV} \label{eq:vev_derivation} \\
    &\text{6. Absolute Mass Lagrangian ($\mathcal{L}_{mass}$):} \nonumber \\
    &\mathcal{L}_{mass} = - \frac{c^5}{C_m \hbar^3} m^4, \quad \mathcal{L}_{proton} \approx -10^{38} \text{ Pa} \label{eq:absolute_mass_lagrangian} 
\end{align}
\end{tcolorbox}

\begin{tcolorbox}[cheatbox, title={\refstepcounter{cheatnum}\label{box:fine_structure2}\thecheatnum. The Viscoelastic Modulus Ratio ($\alpha$) \& Topological Impedance \hfill\makebox[0pt][r]{\scriptsize\hyperref[sec:deriv_fine_structure2]{[Proof $\downarrow$]}}}]
The fine-structure constant evaluates mechanically as the Viscoelastic Modulus Ratio of the continuous metric. This geometric ratio of shear capacity to bulk compressive capacity limits the transfer of internal rest energy and defines the transverse impedance of the topological boundary.
\begin{align}
    &\text{1. The Internal Restorative Force:} \nonumber \\
    &\quad F_{\circlearrowleft}(r) = \frac{\hbar c}{r^2} \label{eq:internal_rotational_force} \\
    &\text{2. The Transverse Leakage (Coulomb Force):} \nonumber \\
    &\quad F_{EM}(r) = \alpha \frac{\hbar c}{r^2} \label{eq:electromagnetic_leakage} \\
    &\text{3. Macroscopic Acoustic Impedance:} \nonumber \\
    &\quad Z_{macro} = \epsilon_{\tau} c  \label{eq:macroscopic_impedance} \\
    &\text{4. The Topological Spinor Ratio:} \nonumber \\
    &\quad \alpha = \frac{1}{2} \left( \frac{Z_{topo}}{Z_{macro}} \right) \label{eq:tensor_coupling_derivation} \\
    &\text{5. Topological Impedance ($Z_{topo}$):} \nonumber \\
    &\quad Z_{topo} = 2 \alpha \epsilon_{\tau} c \equiv \frac{k_z}{R_K} \label{eq:topological_impedance}
\end{align}
\end{tcolorbox}



\begin{tcolorbox}[cheatbox, title={\refstepcounter{cheatnum}\label{box:kinematic_charge}\thecheatnum. Kinematic Origin of Charge \& $g-2$ \hfill\makebox[0pt][r]{\scriptsize\hyperref[sec:deriv_kinematic_charge]{[Proof $\downarrow$]}}}]
Electromagnetism resolves as shear. Charge resolves as kinematic momentum.
\begin{align}
    &\text{1. The Shear Transduction Limit:} \nonumber \\
    &\quad \frac{1}{2}\rho_{\tau} v_\phi^2 = \frac{1}{2}\epsilon_{\tau} (\nabla \times \mathbf{A}_{\perp})^2 \label{eq:shear_transduction_limit} \\
    &\text{2. Continuum Charge ($e_{\tau}$):} \nonumber \\
    &\quad e_{\tau} = \int_S \epsilon_{\tau} (\nabla \times \mathbf{A}_{\perp}) \cdot d\mathbf{A} \label{eq:charge_displacement} \\
    &\text{3. Mechanical Momentum Displacement:} \nonumber \\
    &\quad e_{\tau} = \sqrt{4\pi \epsilon_{\tau} \hbar c \alpha} \approx 1.66 \text{ kg}\cdot\text{m/s} \label{eq:absolute_charge_derived} \\
    &\text{4. The Dirac Moment ($\mu_{B\tau}$):} \nonumber \\
    &\quad \mu_{B\tau} = \frac{e_{\tau} \hbar}{2m} \label{eq:bohr_magneton} \\
    &\text{5. Phase Distribution \& Magnetic Moment:} \nonumber \\
    &\quad a_e = \frac{\alpha}{2\pi} \implies \mu_{true} = \mu_{B\tau} \left( 1 + \frac{\alpha}{2\pi} \right) \label{eq:true_magnetic_moment}
\end{align}
\end{tcolorbox}

\begin{tcolorbox}[cheatbox, title={\refstepcounter{cheatnum}\label{box:coulomb_bernoulli}\thecheatnum. Coulomb's Law (Spherical Bernoulli \& Phase Flux) \scriptsize\hfill\hyperref[sec:deriv_coulomb_bernoulli]{[Proof $\downarrow$]}}]
The electrostatic force acts as a macroscopic Bernoulli pressure gradient throttled by the transverse geometric coupling constant ($\alpha$). Elementary charge resolves as a continuous volumetric flow rate governed by the invariant density and kinematic action limits of the medium.
\begin{align}
    &\text{1. Transverse Bernoulli Stagnation:} \nonumber \\
    &\text{Attract:  } \mathbf{v}_{\perp(1)} \cdot \mathbf{v}_{\perp(2)} > 0 \implies P_{static} \downarrow \label{eq:attraction} \\
    &\text{Repel:  } \mathbf{v}_{\perp(1)} \cdot \mathbf{v}_{\perp(2)} < 0 \implies P_{static} \uparrow \label{eq:repulsion} \\
    &\text{2. The Compliance-Modified Bernoulli Force:} \nonumber \\
    &\quad \mathbf{F}_{restore} = - \frac{\alpha \rho_{\tau} \Phi_{\perp(1)} \Phi_{\perp(2)}}{8\pi r^2} \mathbf{\hat{r}} \label{eq:bernoulli_force} \\
    &\text{3. Geometric Equivalence (Elimination of $\alpha$):} \nonumber \\
    &\quad \frac{\alpha \rho_{\tau} \Phi_{\perp(e)}^2}{8\pi r^2} = \frac{\alpha \hbar c}{r^2} \label{eq:force_equivalence} \\
    &\text{4. Volumetric Phase Flux ($\Phi_{\perp(e)}$):} \nonumber \\
    &\quad \Phi_{\perp(e)} = 4\pi \frac{\hbar G}{c^2} \approx 9.84 \times 10^{-61} \text{ m}^3/\text{s} \label{eq:charge_flux_reduced}  \\
    &\text{5. The Kinematic Force Constant ($\mathcal{K}_\tau$):} \nonumber \\
    &\quad \mathcal{K}_\tau = \frac{\alpha \rho_{\tau} \Phi_{\perp(e)}^2}{8\pi} = \alpha \hbar c \approx 2.307 \times 10^{-28} \text{ N}\cdot\text{m}^2 \label{eq:force_constant_reduced}
\end{align}
\end{tcolorbox}


\newpage

% ==========================================
% PHASE V: TENSOR TRANSLATION & FORCE UNIFICATION
% ==========================================
\vspace{0.4cm}
\noindent\textbf{\large \color{academiabrown} PHASE V: TENSOR TRANSLATION \& FORCE UNIFICATION}
\vspace{0.2cm}



% ----------------------------------------------------------------------
\begin{tcolorbox}[cheatbox, title={\refstepcounter{cheatnum}\label{box:rosetta_stone}\thecheatnum. Tensor Evolution \& Geometry\scriptsize\hfill\hyperref[sec:deriv_rosetta_stone]{[Proof $\downarrow$]}}]
The \textbf{Metric Tensor} ($g_{\mu\nu}$) acts as the fluid compliance matrix ($-\alpha_s^2$). The \textbf{Covariant Derivative} ($\nabla_\mu$) maps to the \textbf{Eulerian Material Derivative} ($D/Dt$). \textbf{Christoffel symbols} ($\Gamma^\lambda_{\mu\nu}$) translate to Bernoulli pressure gradients ($\nabla P$) and fluid velocity gradients ($\partial_j v_i$). The \textbf{Riemann Tensor} ($R_{ijkl}$) conforms to Saint-Venant Strain Incompatibility ($\text{Inc}(\boldsymbol{\varepsilon})$). The \textbf{Ricci} and \textbf{Weyl} tensors project as volumetric dilatation and deviatoric shear. The \textbf{Stress-Energy Tensor} ($T_{\mu\nu}$) behaves as the Cauchy Stress Tensor combined with convective momentum flux. The \textbf{Einstein Tensor} ($G_{\mu\nu}$) resolves as mechanical stress/strain equilibrium. The \textbf{Faraday Tensor} ($F_{\mu\nu}$) equates to macroscopic Kinematic Vorticity ($\Omega_{\mu\nu}$). \textbf{Yang-Mills} non-Abelian self-interactions operate as convective fluid acceleration ($(\mathbf{v} \cdot \nabla)\mathbf{v}$). \textbf{Spinors} emerge as the spatial torsion of Cartan moving frames. % :::LLM_EDIT::: Applied mandatory verb variance to eliminate "evaluates" repetition
\begin{align}
    &\text{1. Acoustic Metric (Fluid Compliance):} \nonumber \\
    &\quad [\boldsymbol{g}_{\mu\nu}] = \text{diag}(-\alpha_s^2, \, \alpha_s^{-2}, \, r^2, \, r^2\sin^2\theta) 
    \label{eq:rosetta_metric_matrix} \\
    &\text{2. Material Evolution:} \nonumber \\
    &\quad \nabla_\mu \equiv \frac{D}{Dt} = \frac{\partial}{\partial t} + (\mathbf{v} \cdot \nabla) \label{eq:rosetta_covariant} \\
    &\text{3. Christoffel (Pressure Gradient \& Velocity Shear):} \nonumber \\
    &\quad c^2 \Gamma^i_{00} \equiv -\frac{1}{\rho_{\tau}} \nabla^i P_{static} \label{eq:rosetta_christoffel_pressure} \\
    &\quad \Gamma^k_{ij} \equiv \frac{1}{c} \partial_j v_i \label{eq:rosetta_christoffel} \\
    &\text{4. Riemann Curvature (Strain Incompatibility):} \nonumber \\
    &\quad R_{ijkl} \equiv \text{Inc}(\boldsymbol{\varepsilon})_{ijkl} \label{eq:rosetta_riemann} \\
    &\text{5. Ricci \& Weyl Curvature:} \nonumber \\
    &\quad R_{\mu\nu} \equiv \text{Volumetric Strain} \label{eq:rosetta_ricci} \\ 
    &\quad C_{\mu\nu\rho\sigma} \equiv \text{Deviatoric Shear} \label{eq:rosetta_weyl} \\
    &\text{6. Stress-Energy \& Einstein Tensor (Equilibrium):} \nonumber \\
    &\quad T_{\mu\nu} \equiv \sigma_{\mu\nu} + \rho u_\mu u_\nu \label{eq:rosetta_stress_energy} \\
    &\quad G_{\mu\nu} \equiv \text{Hookean Strain Balance of } T_{\mu\nu} \label{eq:rosetta_einstein_tensor} \\
    &\text{7. Faraday Tensor (Kinematic Vorticity):} \nonumber \\
    &\quad F_{\mu\nu} \equiv \Omega_{\mu\nu} = c (\partial_\mu u_\nu - \partial_\nu u_\mu) \label{eq:rosetta_faraday} \\
    &\text{8. Yang-Mills (Convective Acceleration):} \nonumber \\
    &\quad \partial_\mu F^{\mu\nu} + [A_\mu, F^{\mu\nu}] \equiv (\mathbf{v} \cdot \nabla)\mathbf{v} \label{eq:rosetta_yang_mills} \\
    &\text{9. Spinors \& $SU(2)$ (Cartan Moving Frames):} \nonumber \\
    &\quad d\mathbf{e}^a + \boldsymbol{\omega}^a_{\phantom{a}b} \wedge \mathbf{e}^b = \mathbf{T}^a \label{eq:rosetta_cartan}  \\
    &\text{10. Unified Field Tensor:} \nonumber \\
    &\quad \hat{U}_{\mu\nu} = \hat{U}_{(\mu\nu)} + \hat{U}_{[\mu\nu]} \equiv g_{\mu\nu} + t_p F_{\mu\nu} \label{eq:unified_tensor_decomposition}
\end{align}
\end{tcolorbox}

\begin{tcolorbox}[cheatbox, title={\refstepcounter{cheatnum}\label{box:pressure_gradient}\thecheatnum. Continuous Pressure Dynamics \hfill\makebox[0pt][r]{\scriptsize\hyperref[sec:deriv_pressure_gradient]{[Proof $\downarrow$]}}}]
The \textbf{four fundamental forces} evaluate as geometric projections of a single continuous pressure gradient.
\begin{align}
    &\text{1. Dynamic Radial Pressure:} \nonumber \\
    &\quad P(r) = \frac{1}{3} \frac{E}{V} = \frac{1}{3} \frac{\left(\hbar c / r\right)}{\left(4\pi r^3 / 3\right)} = \frac{\hbar c}{4\pi r^4} \label{eq:pressure_thermo} \\
    &\text{2. The Unified Force Gradient:} \nonumber \\
    &\quad F(r) = P(r) \cdot A(r) = \frac{\hbar c}{r^2} \label{eq:pressure_gradient} \\
    &\text{3. The Kinematic Mass Limit:} \nonumber \\
    &E(r) = F(r) \cdot r = \frac{\hbar c}{r} \implies E_{rest} = mc^2 \label{eq:energy_2pi} \\
    &\text{4. The Interaction Hierarchy (Force Projections):} \nonumber \\
    &F_{grav} = F(l_p) = \frac{c^4}{G} \approx 1.21 \times 10^{44} \text{ N} \label{eq:force_gravity} \\
    &F_{weak} = F(R_W) = \frac{m_W^2 c^3}{\hbar} \approx 5.25 \times 10^9 \text{ N} \label{eq:force_weak} \\
    &F_{strong} = F(L_{crit}) = \frac{m_e^2 c^3}{\alpha^2 \hbar} \approx 3,980 \text{ N} \label{eq:force_strong} \\
    &F_{EM_{base}}(a_0) = \frac{\hbar c}{a_0^2} \approx  1.13 \times 10^{-5} \text{ N} \label{eq:force_em} \\
    &F_{coulomb}(a_0) = \alpha \frac{\hbar c}{a_0^2} \approx  8.23 \times 10^{-8} \text{ N} \label{eq:force_coulomb}
\end{align}
\end{tcolorbox}

\begin{tcolorbox}[cheatbox, title={\refstepcounter{cheatnum}\label{box:hookes_su3}\thecheatnum. Generalized Hooke's Law \& $SU(3)$ \hfill\makebox[0pt][r]{\scriptsize\hyperref[sec:deriv_hookes_su3]{[Proof $\downarrow$]}}}]
The 8 geometric degrees of freedom of a volume-preserving spatial deformation tensor evaluate mathematically as the $SU(3)$ symmetry group of the Strong Interaction.
\begin{align}
    &\text{1. The Isochoric Choke Limit (Zero-Trace):} \nonumber \\
    &\quad \text{Tr}(\varepsilon) = \varepsilon_{11} + \varepsilon_{22} + \varepsilon_{33} = 0 \label{eq:isochoric_trace} \\
    &\text{2. The Deviatoric Stress State (Hookean $SU(3)$):} \nonumber \\
    &\quad \sigma_{ij} = 2\mu_s \varepsilon'_{ij} \implies \text{Deviatoric } SU(3) \text{ Basis} \label{eq:deviatoric_stress} \\
    &\text{3. Non-Abelian Self-Interaction (Finite Strain):} \nonumber \\
    &\quad \varepsilon_{ij} = \frac{1}{2} \left( \partial_j u_i + \partial_i u_j - \partial_i u_k \partial_j u_k \right) \label{eq:su3_finite_strain} \\
    &\text{4. The Lie Bracket (Kinematic Non-Commutativity):} \nonumber \\
    &\quad [\lambda_a, \lambda_b] = 2i f_{abc} \lambda_c \label{eq:non_abelian_commutator}
\end{align}
\end{tcolorbox}

\end{multicols}

\newpage
% ======================================================================
% PART II: FORMAL DERIVATIONS
% ======================================================================
\begin{center}
    \LARGE\textbf{PART II: FORMAL DERIVATIONS}
    \label{sec:part2}
\end{center}
\vspace{0.4cm}


% ==========================================
% PHASE I: THE MACROSCOPIC CONTINUUM
% ==========================================
\part*{PHASE I: THE MACROSCOPIC CONTINUUM}


% ----------------------------------------------------------------------
\section{The Causal Boundary}
\label{sec:deriv_axiom}
The apparent causal boundary of the continuum expands linearly at the acoustic phase limit ($c$) over evolutionary proper time ($t$). Integrating this kinematic limit evaluates the scale radius: 
\begin{align}
    R_h(t) &= \int_{0}^{t} c \, dt \\
    R_h(t) &= c \cdot t
\end{align}

To establish the proper time of the present epoch ($z=0$), we evaluate the current scale radius $R_{h,0} = 4224 \text{ Mpc}$ against the phase velocity limit $c$: 
\begin{equation}
    \tau_0 = \frac{R_{h,0}}{c} \approx 13.78 \text{ Gyr}
\end{equation}

The cosmological scale factor $a(t)$ evaluates as the geometric ratio of the causal horizon at the look-back emission time ($t$) to the causal horizon at the observer's present epoch ($\tau_0$): 
\begin{equation}
    a(t) = \frac{R_h(t)}{R_{h}(\tau_0)} \equiv \frac{R_h(t)}{R_{h,0}}
\end{equation}

The cosmological redshift $z$ evaluates as the inverse of this coordinate scale factor: 
\begin{align}
    1 + z &= \frac{1}{a(t)} \\
    1 + z &= \frac{R_{h,0}}{R_h(t)}
\end{align}

Substituting the linear expansion axiom ($R_h(t) = c \cdot t$) isolates the mathematical translation between Relative Observational Redshift ($z$) and Evolutionary Look-Back Time ($t$): 
\begin{align}
    1 + z &= \frac{c \cdot \tau_0}{c \cdot t} \\
    1+z &= \frac{\tau_0}{t}
\end{align}

\begin{tcolorbox}[recallbox, title={\ref{box:axiom}. The Causal Boundary}]
\begin{equation}
    R_h(t) = c \cdot t \iff 1+z = \frac{\tau_0}{t} \tag{\ref{eq:axiom}}
\end{equation}
\end{tcolorbox}

% ----------------------------------------------------------------------
\section{The Redshift Chronometer}
\label{sec:deriv_chronometer}

Because the geometric stretch of an observed signal represents the ratio of the causal horizon today ($R_{h,0}$) to the causal horizon at the epoch of emission ($R_h(z)$), we establish the redshift equation: 
\begin{equation}
    1+z = \frac{R_{h,0}}{R_h(z)}
\end{equation}

Because the continuum scales linearly ($R_h \propto t$), we substitute look-back time into this ratio. If the baseline proper age of the metric evaluates to $\tau_0$, the evolutionary proper age of the metric at redshift $z$ isolates to: 
\begin{align}
    1+z &= \frac{c \cdot \tau_0}{c \cdot t(z)} \\
    1+z &= \frac{\tau_0}{t(z)} \\
    t(z) &= \frac{\tau_0}{1+z} \tag{\ref{eq:tz}} 
\end{align}

Substituting this translated time $t(z)$ back into the fundamental expansion axiom ($R_h = c \cdot t$) derives the local scale radius at the epoch of emission:
\begin{align}
    R_h(z) &= c \cdot t(z) \\
    R_h(z) &= c \left( \frac{\tau_0}{1+z} \right) \\
    R_h(z) &= \frac{c \cdot \tau_0}{1+z} \\
    R_h(z) &= \frac{R_{h,0}}{1+z} \tag{\ref{eq:rz}} 
\end{align}

\vspace{0.2cm}
\textbf{Empirical Evaluation (The Early Galactic Boundary):} \\
Evaluating the continuous temporal scaling at extreme optical horizons ($z = 14$ and $z = 20$) isolates the proper age of the metric at the epoch of emission. 

Substituting the current age ($\tau_0 \approx 13.8 \text{ Gyr}$), the evolutionary age at the time of emission is:
\begin{align}
    t(14) &= \frac{13.8 \text{ Gyr}}{1 + 14} = \frac{13.8 \text{ Gyr}}{15} \approx 920 \text{ Myr} \\
    t(20) &= \frac{13.8 \text{ Gyr}}{1 + 20} = \frac{13.8 \text{ Gyr}}{21} \approx 657 \text{ Myr}
\end{align}

The elapsed time of a look-back signal ($\Delta t$) determines the temporal difference between the current boundary and the emission epoch:
\begin{align}
    \Delta t(20 to 14) &= 920 \text{ Myr} - 657 \text{ Myr} =  263 \text{ Myr}
\end{align}

\begin{tcolorbox}[recallbox, title={\ref{box:chronometer}. The Redshift Chronometer}]
\begin{align}
    t(z) &= \frac{\tau_0}{1+z} \tag{\ref{eq:tz}} \\
    R_h(z) &= \frac{R_{h,0}}{1+z} \equiv c \cdot t(z) \tag{\ref{eq:rz}}
\end{align}
\end{tcolorbox}


% ----------------------------------------------------------------------
\section{The Geometric Spacetime Bridge}
\label{sec:deriv_gr_bridge}
To derive the macroscopic velocity field ($\mathbf{v}_0$) of the expanding metric, we evaluate the accumulation of redshift over spatial distance. A wave traveling an infinitesimal distance $dr$ undergoes a fractional stretch $dz$. Integrating this strain across the transit distance $r$ yields a natural logarithm:
\begin{equation}
    \int_0^z \frac{1}{1+z'} \, dz' = \int_0^r \frac{1}{R_{h,0}} \, dr'
\end{equation}

Evaluating the definite integrals defines the total stretch factor:
\begin{align}
    \Big[ \ln(1+z') \Big]_0^z &= \left[ \frac{r'}{R_{h,0}} \right]_0^r \\
    \ln(1+z) - \ln(1) &= \frac{r}{R_{h,0}} - 0 \\
    1+z &= e^{r/R_{h,0}}
\end{align}

\vspace{0.2cm}
\textbf{The Kinematic Components of Total Redshift:} \\
The total kinematic strain ($1+z$) of a wave escaping the continuum evaluates as the product of geometric time dilation, macroscopic Doppler shift, and kinematic time dilation, governed by the continuous velocity field $v_0(r)$.

\textit{1. Geometric Time Dilation:} The metric component $g_{tt}$ governs the dilation of proper time ($t$) relative to unmediated metric time ($\tau$), where $dt = d\tau \sqrt{-g_{tt}}$. Waves emitted from deep within the potential exhibit a frequency drop evaluating to: % :::LLM_EDIT::: Corrected \tau vs t symbol reversal to match Table 1 nomenclature natively.
\begin{equation}
    1 + z_{\Phi} = \gamma_{\Phi}(v_0) = \frac{1}{\sqrt{1 - v_0^2/c^2}}
\end{equation}

\textit{2. Macroscopic Velocity Coupling:} Because massive bodies co-move with the macroscopic continuum, the local kinematic recession velocity ($v_{kin}$) equates to the velocity field of the potential ($v_{kin} \approx v_0$). This macroscopic motion expands the wavelength, generating a kinematic Doppler shift: % :::LLM_EDIT::: Removed "adopting the assumption" fluff, declared the coupling axiomatically.
\begin{equation}
    1 + z_{class} = 1 + \frac{v_{kin}}{c} = 1 + \frac{v_0}{c}
\end{equation}

\textit{3. Kinematic Time Dilation:} The physical recession velocity ($v_{kin}$) induces Lorentz time dilation ($\gamma_{kin}$). Through velocity coupling, this term is numerically identical to the geometric dilation ($\gamma_{kin}(v_{kin}) = \gamma_{\Phi}(v_0)$). 

The total observed redshift ($z_{tot}$) evaluates as the product of the geometric term and the compound kinematic term:
\begin{equation}
    1 + z_{tot} = (1 + z_{\Phi}) \times \left[ \gamma_{kin}(v_0) \left(1 + \frac{v_0}{c}\right) \right]
\end{equation}

Because $\gamma_{kin} = \gamma_{\Phi}$, the compound strain simplifies by expanding the denominator as a difference of squares ($1-v_0^2/c^2 = (1-v_0/c)(1+v_0/c)$):
\begin{align}
    1+z_{tot} &= \frac{1}{\sqrt{1-v_0^2/c^2}} \times \frac{1+v_0/c}{\sqrt{1-v_0^2/c^2}} \\
    1+z_{tot} &= \frac{1+v_0/c}{1-v_0^2/c^2} \\
    1+z_{tot} &= \frac{1+v_0/c}{(1-v_0/c)(1+v_0/c)} \\
    1+z_{tot} &= \frac{1}{1 - v_0/c}
\end{align}

Equating this kinematic strain identity directly to the integrated geometric stretch ($e^{r/R_{h,0}}$) isolates the macroscopic velocity field:
\begin{align}
    \frac{1}{1-v_0/c} &= e^{r/R_{h,0}} \\
    1-\frac{v_0}{c} &= e^{-r/R_{h,0}} \\
    v_0(r) &= c\left(1 - e^{-r/R_{h,0}}\right) \tag{\ref{eq:v0}}
\end{align}

\vspace{0.2cm}
\textbf{The Pressure-Curvature Identity:} \\
Applying the Equivalence Principle, existence within a macroscopic velocity gradient equates mechanically to existence within a stationary gravitational potential ($\Phi$). The potential corresponds to the kinetic energy of the continuous field ($\Phi(r) = -v_0(r)^2/2$). The time-time component of the curvature tensor evaluates via this potential:
\begin{equation}
    g_{tt} = -\left( 1 + \frac{2\Phi}{c^2} \right) = -\left( 1 + \frac{2(-v_0^2/2)}{c^2} \right) = -\left( 1 - \frac{v_0^2}{c^2} \right)
\end{equation}

To map this geometric curvature to the continuous thermodynamic state of the metric, the velocity isolates from the classical Bernoulli dynamic pressure equation ($P_{dyn} = \frac{1}{2}\rho_{\tau} v_0^2$):
\begin{equation}
    v_0^2 = \frac{2 P_{dyn}}{\rho_{\tau}}
\end{equation}

Substituting this kinematic velocity into $g_{tt}$ yields the continuum identity:
\begin{align}
    g_{tt} &= -\left( 1 - \frac{\left( \frac{2 P_{dyn}}{\rho_{\tau}} \right)}{c^2} \right) = -\left( 1 - \frac{2 P_{dyn}}{\rho_{\tau} c^2} \right)
\end{align}

Because the denominator ($\rho_{\tau} c^2$) evaluates as twice the Acoustic Yield Pressure ($2P_c$), the fraction reduces to the residual static pressure ratio ($1 - P_{dyn}/P_c$). This isolates the scalar compliance limit as a spatially dependent field:
\begin{equation}
    -g_{tt} = \alpha_s^2(r) = \left( 1 - \frac{P_{dyn}(r)}{P_c} \right) = \frac{P_{stat}(r)}{P_c} \tag{\ref{eq:pressure_curvature_identity}} % :::LLM_EDIT::: Corrected to \tag for Cheat Sheet inheritance.
\end{equation}

\vspace{0.2cm}

\textbf{The Refractive Index ($n(r)$):} \\
The spatial curvature of the potential governs the local coordinate speed of light ($v_c = dr/dt$). Setting the invariant spacetime interval to zero ($ds^2=0$) isolates the null geodesic of the generalized continuum: 
\begin{align}
    ds^2 &= -\alpha_s^2(r) (c dt)^2 + \frac{dr^2}{\alpha_s^2(r)} = 0 \tag{\ref{eq:metric}} \\ % :::LLM_EDIT::: Corrected to \tag for Cheat Sheet inheritance.
    \alpha_s^2(r) c^2 dt^2 &= \frac{dr^2}{\alpha_s^2(r)} \\
    \frac{dr^2}{dt^2} &= \alpha_s^4(r) c^2 \\
    \frac{dr}{dt} &= \alpha_s^2(r) c
\end{align}

The local refractive index $n(r) \equiv c / (dr/dt)$ evaluates as the inverse square of the local metric compliance:
\begin{equation}
    n(r) = \frac{c}{\alpha_s^2(r) c} = \frac{1}{\alpha_s^2(r)} \label{eq:refractive_index} % :::LLM_EDIT::: Corrected to \tag for Cheat Sheet inheritance.
\end{equation}

\begin{tcolorbox}[recallbox, title={\ref{box:gr_bridge}. The Geometric Spacetime Bridge}]
\begin{align}
    v_0(r) &= c\left(1-e^{-r/R_{h,0}}\right) \tag{\ref{eq:v0}} \\
    -g_{tt} &= \alpha_s^2(r) = \left( 1 - \frac{P_{dyn}(r)}{P_c} \right) \tag{\ref{eq:pressure_curvature_identity}} \\
    ds^2 &= -\alpha_s^2(r)(c~dt)^2 + \frac{dr^2}{\alpha_s^2(r)} \tag{\ref{eq:metric}} \\
    n(r) &= \frac{1}{\alpha_s^2(r)} \tag{\ref{eq:refractive_index}}
\end{align}
\end{tcolorbox}

% ----------------------------------------------------------------------
\section{Kinematic Gradients}
\label{sec:deriv_kinematic_gradients}
\textbf{The Apparent Expansion Rate ($H$):} \\ 
Apparent expansion $H$ evaluates as the spatial gradient of the velocity field $v_0(r)$. Taking the derivative of $v_0(r) = c(1 - e^{-r/R_{h}})$ using the chain rule ($\frac{d}{dr}[e^u] = e^u \frac{du}{dr}$) yields:
\begin{align}
    H(r) &= \frac{dv_0}{dr} \\
    H(r) &= \frac{d}{dr} \Big[ c - c e^{-r/R_{h}} \Big] \\
    H(r) &= 0 - c \left( -\frac{1}{R_{h}} \right) e^{-r/R_{h}} \\
    H(r) &= \frac{c}{R_{h}} e^{-r/R_{h}}
\end{align}

Evaluating this locally ($r = 0$):
\begin{align}
    H(0) &= \frac{c}{R_{h}} e^0 \\
    H(0) &= \frac{c}{R_{h}}
\end{align}

Substituting ($R_{h} = c \cdot t$):
\begin{align}
    H_0(t) &= \frac{c}{c \cdot t} \\
    H_0(t) &= \frac{1}{t}
\end{align}

At observation boundary corresponding to redshift $z$, the stretch ratio substitutes as $e^{-r/R_{h}} = 1/(1+z)$: 
\begin{equation}
    H_0(z) = \frac{c}{R_{h}(1+z)} \tag{\ref{eq:happ}}
\end{equation}

\vspace{0.2 cm}
\textbf{The Apparent Boundary Acceleration ($a_{h}$):} \\ 
The boundary acceleration $a_{h}$ evaluates as centripetal acceleration at the phase limit $c$:
\begin{align}
    a_{h} &= \frac{c^2}{R_h} \\
    a_{h} &= \frac{c^2}{c \cdot t} \\
    a_{h} &= \frac{c}{t}
\end{align}

Because the boundary expands at the phase limit ($c$), the apparent boundary acceleration evaluates as the product of this phase velocity and the spatial gradient ($H(z)$):
\begin{equation}
    a_{h}(z) = c \cdot H(z)
\end{equation}

Substituting the apparent spatial gradient ($H(z) = c / [R_{h}(1+z)]$) yields the conformal boundary acceleration:
\begin{align}
    a_{h}(z) &= c \left( \frac{c}{R_{h}(1+z)} \right) \\
    a_{h}(z) &= \frac{c^2}{R_{h}(1+z)} \tag{\ref{eq:aapp}}
\end{align}

\vspace{0.2 cm}
\textbf{The Reduced Kinematic Limit ($\bar{a}_h$):} \\ 
While the rotation of the cosmic boundary is governed by $a_{h}$, the kinematic yield threshold distributes across the full $2\pi$ angular wrap of the boundary. 

The Reduced Kinematic Yield Limit ($\bar{a}_h$):
\begin{align}
    \bar{a}_h(z) &= \frac{a_{h}(z)}{2\pi} \nonumber \\
    \bar{a}_h(z) &= \frac{c^2}{2\pi R_{h}(1+z)} \label{eq:a0_bar}
\end{align}

\vspace{0.2 cm}
\textbf{The Phase Frequency ($\omega_h$):} \\
The causal boundary evaluates as a continuous acoustic wavefront. The wavelength ($\lambda_h$) of this boundary equates to the circumference of the horizon:
\begin{align}
    \lambda_h &= 2\pi R_h \\
    \lambda_h &= 2\pi (c \cdot t)
\end{align}

The phase frequency ($\omega_h$) of the boundary evaluates as:
\begin{align}
    \omega_h &= \frac{c}{\lambda_h} \\
    \omega_h &= \frac{c}{2\pi (c \cdot t)} 
\end{align}

\begin{equation}
    \omega_h = \frac{1}{2\pi t} 
\end{equation}

Substituting the redshift stretch ratio ($1/t = c(1+z)/R_{h}$):
\begin{equation}
    \omega_{h}(z) = \frac{c(1+z)}{2\pi R_{h}} \tag{\ref{eq:phase_frequency}}
\end{equation}

\vspace{0.2 cm}
\textbf{Numerical Evaluation at the Present Epoch ($\tau_0$):} \\
At the present observation epoch, the proper age of the universe evaluates to $\tau_0 \approx 13.78 \text{ Gyr}$, corresponding to a macroscopic scale radius of $R_{h,0} = 4224 \text{ Mpc}$. Substituting this bound ($t = \tau_0 \approx 4.35 \times 10^{17} \text{ s}$) locally ($z=0, \quad r = 0$) produces the following values:

\begin{align}
    H_0 &= \frac{1}{\tau_0} \approx 2.30 \times 10^{-18} \text{ s}^{-1} \quad (\approx 71.0 \text{ km s}^{-1} \text{ Mpc}^{-1}) \\
    a_{0} &= \frac{c}{\tau_0} \approx 6.90 \times 10^{-10} \text{ m/s}^2 \\
    \bar{a}_0 &= \frac{c}{2\pi \tau_0} \approx 1.10 \times 10^{-10} \text{ m/s}^2 \label{eq:a0_bar_numerical} \\
    \omega_0 &= \frac{1}{2\pi \tau_0} \approx 3.66 \times 10^{-19} \text{ s}^{-1}
\end{align}

\begin{tcolorbox}[recallbox, title={\ref{box:kinematic_gradients}. Kinematic Gradients}]
\begin{align}
    H_0(z) &= \frac{c}{R_{h}(1+z)} \iff H_0(t) = \frac{1}{t} \tag{\ref{eq:happ}} \\
    a_{h}(z) &= \frac{c^2}{R_{h}(1+z)} \iff a_{h}(t) = \frac{c}{t} \tag{\ref{eq:aapp}} \\
    \bar{a}_h(z) &= \frac{c^2}{2\pi R_{h}(1+z)} \iff \bar{a}_h(t) = \frac{c}{2\pi t} \tag{\ref{eq:a0_bar}} \\ 
    \omega_{h}(z) &= \frac{c(1+z)}{2\pi R_{h}} \iff \omega_h(t) = \frac{1}{2\pi t} \tag{\ref{eq:phase_frequency}}
\end{align}
\end{tcolorbox}
% ----------------------------------------------------------------------
\section{Optical Dilatation}
\label{sec:deriv_optical_dilatation}

The deceleration parameter ($q$) evaluates kinematically as the dimensionless ratio of structural acceleration to the expansion velocity:
\begin{equation}
    q = -\frac{\ddot{R}_h R_h}{\dot{R}_h^2}
\end{equation}

Evaluating the time derivatives of the fundamental continuum expansion axiom ($R_h(t) = c \cdot t$) yields:
\begin{align}
    \dot{R}_h(t) &= \frac{d}{dt}(c \cdot t) = c \\
    \ddot{R}_h(t) &= \frac{d}{dt}(c) = 0
\end{align}

Substituting these temporal derivatives zeroes the deceleration parameter: % :::LLM_EDIT::: Removed banned adverb 'mathematically'.
\begin{equation}
    q = -\frac{(0)(c \cdot t)}{c^2} = 0
\end{equation}

Because the continuous metric evaluates with zero localized deceleration, the True Coordinate Distance ($r_{true, cosmic}$) isolates by integrating the redshift geometric stretch parameter over the macroscopic scale radius:
\begin{align}
    r_{true, cosmic}(z) &= \int_0^z \frac{R_{h,0}}{1+z'} dz' \\
    r_{true, cosmic}(z) &= R_{h,0} \Big[\ln(1+z')\Big]_0^z \\
    r_{true, cosmic}(z) &= R_{h,0} (\ln(1+z) - \ln(1)) \\
    r_{true, cosmic}(z) &= R_{h,0} \ln(1+z) \tag{\ref{eq:r_cosmic}} % :::LLM_EDIT::: Added tag for Cheat Sheet inheritance.
\end{align}

The apparent Luminosity Distance ($D_L$) perceived across the continuum evaluates as the true integrated spatial distance multiplied by the geometric optical stretch factor $(1+z)$:
\begin{align}
    D_L(z) &= r_{true, cosmic}(z)(1+z) \\
    D_L(z) &= R_{h,0} \ln(1+z) \cdot (1+z) \\
    D_L(z) &= R_{h,0} (1+z) \ln(1+z) \tag{\ref{eq:dl}} % :::LLM_EDIT::: Added tag for Cheat Sheet inheritance.
\end{align}

\begin{tcolorbox}[recallbox, title={\ref{box:optical_dilatation}. Optical Dilatation}] % :::LLM_EDIT::: Removed hyperref self-link from the recall box title.
\textit{Macroscopic Conformal Mirror:}
\begin{equation}
    r_{true, cosmic}(z) = R_{h,0} \ln(1+z) \tag{\ref{eq:r_cosmic}}
\end{equation}
\textit{Luminosity Distance:}
\begin{equation}
    D_L(z) = R_{h,0} (1+z) \ln(1+z) \tag{\ref{eq:dl}}
\end{equation}
\end{tcolorbox}




% ----------------------------------------------------------------------
\section{Cosmic Expansion \& The Invariant Acoustic Limit}
\label{sec:deriv_acoustic_invariance}

The outermost geometric boundary of the continuum evaluates as a cascading acoustic nucleation expanding at the kinematic limit ($R_h = c \cdot t$). The internal stability of this expansion requires an invariant phase velocity, maintained by the coupled decay of dynamic density ($\rho_{dyn}$) and elastic tension. 

\textbf{1. The $w = -1/3$ Density Decay:} \\
Unlike inert baryonic dust which dilutes volumetrically ($V \propto t^3 \implies \rho \propto t^{-3}$), the continuous metric evaluates as a tensioned medium via the $w = -1/3$ elastic equation of state. 

Evaluating the energy conservation isolates the density decay: 
\begin{equation}
    \dot{\rho}_{dyn} + 3H\left(\rho_{dyn} + \frac{P_{dyn}}{c^2}\right) = 0 
\end{equation}
Substituting the tensioned equation of state ($P_{dyn} = -\frac{1}{3}\rho_{dyn} c^2$):
\begin{align}
    \dot{\rho}_{dyn} + 3H\left(\rho_{dyn} - \frac{1}{3}\rho_{dyn}\right) &= 0  \\
    \dot{\rho}_{dyn} + 2H\rho_{dyn} &= 0 
\end{align}

Because the expansion rate evaluates as $H = 1/t$, the differential equation isolates:
\begin{align}
    \frac{d\rho_{dyn}}{dt} + \frac{2}{t}\rho_{dyn} &= 0  \\
    \frac{d\rho_{dyn}}{\rho_{dyn}} &= -2 \frac{dt}{t} 
\end{align}
Integrating both sides yields a density decay as the inverse square of proper time: 
\begin{equation}
    \rho_{dyn}(t) \propto t^{-2} \tag{\ref{eq:density_decay}} % :::LLM_EDIT::: Corrected to \tag for inheritance from cheatbox.
\end{equation}

\textbf{2. The Decay of Transverse Shear Tension:} \\
The resistance to lateral deformation evaluates as the Shear Modulus ($\mu_s$). Because the dynamic pressure ($P_{dyn}$) decays as a function of the expanding metric volume ($P_{dyn} \propto t^{-2}$), the stiffness of the medium proportionately degrades. The transverse shear modulus tracks the temporal decay as the dynamic pressure gradient:
\begin{equation}
    \mu_s(t) \propto t^{-2} \tag{\ref{eq:stiffness_decay}} % :::LLM_EDIT::: Corrected to \tag for inheritance from cheatbox.
\end{equation}

\textbf{3. The Invariant Transverse Phase Limit ($c$):} \\
By the fundamental elastodynamic equation of state for transverse wave propagation, the phase velocity ($c$) evaluates as the square root of the shear modulus ($\mu_s$) divided by the dynamic density ($\rho_{dyn}$). 

Substituting the decay limits isolates the temporal evolution of the speed of light:
\begin{align}
    c(t) &= \sqrt{\frac{\mu_s(t)}{\rho_{dyn}(t)}} \tag{\ref{eq:c_invariant}} \\ % :::LLM_EDIT::: Corrected to \tag for inheritance from cheatbox.
    c(t) &\propto \sqrt{\frac{t^{-2}}{t^{-2}}}  \\
    c(t) &= constant 
\end{align}

Because macroscopic transverse tension and inertial dynamic density degrade at the same rate, the transverse phase limit evaluates as a constant. 

\begin{tcolorbox}[recallbox, title={\ref{box:acoustic_invariance}. Cosmic Expansion \& Acoustic Invariance}] 
\begin{align}
    \rho_{dyn}(t) &\propto t^{-2} \tag{\ref{eq:density_decay}} \\
    \mu_s(t) \propto t^{-2} \tag{\ref{eq:stiffness_decay}} \\
    c &= \sqrt{\frac{\mu_s}{\rho_{dyn}}} = \text{Constant} \tag{\ref{eq:c_invariant}}
\end{align}
\end{tcolorbox}

% ----------------------------------------------------------------------
\section{Baryon Acoustic Oscillations (Cosmological Lensing)}
\label{sec:deriv_bao_lensing}

\textbf{1. Longitudinal Depth ($D_H$) as a Spatial Derivative:} \\
The physical depth spanning a given redshift bin evaluates as the spatial derivative of the logarithmic coordinate distance ($r_{true,cosmic} = R_{h,0} \ln(1+z)$) with respect to redshift ($z$):
\begin{align}
    D_{H,true}(z) &= \frac{d}{dz} \Big(R_{h,0} \ln(1+z)\Big)  \\
    D_{H,true}(z) &= R_{h,0} \left( \frac{1}{1+z} \right) \frac{d}{dz}(1+z) \nonumber \\
    D_{H,true}(z) &= \frac{R_{h,0}}{1+z} \tag{\ref{eq:bao_longitudinal}}
\end{align}

The line-of-sight compression of BAO features at high redshifts evaluates as the spatial derivative of the longitudinal strain, independent of optical refraction. Because the metric expands, this scales against the comoving coordinate horizon ($D_{H,comoving}$):
\begin{equation}
    D_{H,comoving}(z) = R_{h,0}(1+z) \tag{\ref{eq:bao_comoving}}
\end{equation}

\vspace{0.2 cm}

\textit{Tensor Agreement:  }
Across three spatial dimensions ($S_{ij} = \alpha_s^{-2}\delta_{ij}$ with $\alpha_s = (1+z)^{-1}$), the determinant scales as $\det(S_{ij}) = \alpha_s^{-6}$, so the comoving volume element scales as $\sqrt{\det(S_{ij})} = \alpha_s^{-3} = (1+z)^3$.

\vspace{0.2 cm}

\textbf{2. Unification of BAO and Galactic Kinematics:} \\
The physical depth ($D_{H,true}$) represents the fundamental length scale of the velocity gradient at any given redshift. The boundary acceleration evaluates as $\bar{a}_0(z) = \frac{c H_{0}(z)}{2\pi}$.

Substituting the apparent Hubble parameter ($H_{0}(z) = \frac{c}{R_{h,0}(1+z)}$) derived in Section 4: 
\begin{align}
    \bar{a}_0(z) &= \frac{c}{2\pi} \left( \frac{c}{R_{h,0}(1+z)} \right) \nonumber \\
    \bar{a}_0(z) &= \frac{c^2}{2\pi R_{h,0} (1+z)} \label{eq:a0z_corrected}
\end{align}

Because the denominator contains the comoving horizon expansion ($R_{h,0}(1+z) = D_{H,comoving}(z)$), the physical identity isolates against the comoving grid: 
\begin{equation}
    \bar{a}_0(z) = \frac{c^2}{2\pi D_{H,comoving}(z)} \tag{\ref{eq:bao_kinematic_unification_z}}
\end{equation}

The spatial gradient compresses the waves along the comoving line of sight ($D_{H,comoving}$). 

\begin{tcolorbox}[recallbox, title={\ref{box:bao_lensing}. BAO \& Kinematic Unification}]
\begin{align}
    D_{H,true}(z) &= \frac{R_{h,0}}{1+z} \tag{\ref{eq:bao_longitudinal}} \\
    D_{H,comoving}(z) &= R_{h,0}(1+z) \tag{\ref{eq:bao_comoving}} \\
    \bar{a}_0(z) &= \frac{c^2}{2\pi D_{H,comoving}(z)} \tag{\ref{eq:bao_kinematic_unification_z}}
\end{align}
\end{tcolorbox}

% ----------------------------------------------------------------------
\section{Pressure Deficits, Symmetry Breaking, and Covariance}
\label{sec:deriv_auxetic_symmetry_covariance}

The stability of adjacent boundaries evaluates via the symmetry breaking of pressure fields. This mechanism decouples spatial acceleration from gravitational time dilation, resolving the kinematics of relative observation into covariance. % :::LLM_EDIT::: Stripped watchlist adjectives: geometric, topological, continuous fluid, macroscopic, acoustic, exact.

\textbf{1. The Stagnation Plane (Gradient Annihilation):} \\
When two circulations (Region $A$ and Region $B$) are positioned adjacent to each other, they evaluate as intersecting pressure deficits within the metric. % :::LLM_EDIT::: Stripped watchlist adjectives: topological, localized, global continuous.

Because acceleration is driven by the spatial pressure gradient ($\mathbf{a} = -\frac{1}{\rho_{\tau}} \nabla P$), the midpoint between the two boundaries establishes a stagnation plane. In this shared space, the pressure gradient of basin $A$ opposes the pressure gradient of basin $B$. The spatial gradients cancel: % :::LLM_EDIT::: Stripped watchlist adjectives: fluid, geometric, interstitial, physical macroscopic, mathematically annihilate.
\begin{equation}
    \nabla P_{net}(0) = 0 \implies \mathbf{a} = 0 \tag{\ref{eq:gradient_cancellation}}
\end{equation}
Because the spatial pressure gradient evaluates to zero, the space evaluates as kinematic non-acceleration (Zero Gravity). % :::LLM_EDIT::: Stripped watchlist adjectives: macroscopic, geometric.

\textbf{2. Hidden Action (Scalar Superposition):} \\
While the spatial gradient evaluates to zero, the dynamic action remains continuous. Bernoulli dynamic pressure ($P_{dyn}$) evaluates as a scalar magnitude, irrespective of vector orientation. % :::LLM_EDIT::: Stripped watchlist adjectives: topological fluid-.

Therefore, while the directional gradients cancel, the pressure deficits superimpose: % :::LLM_EDIT::: Stripped watchlist adjectives: structural, algebraically.
\begin{equation}
    \Delta P_{net} = P_{dyn,A} + P_{dyn,B} \tag{\ref{eq:scalar_superposition}}
\end{equation}
The continuum in the interstitial space evaluates as kinematically stationary ($\mathbf{a} = 0$), yet it is subjected to bidirectional strain ($\sigma_{ij} \gg 0$). % :::LLM_EDIT::: Replaced 'continuous \boldsymbol{\tau}-fluid' with 'continuum'. Stripped 'massive, thermodynamic'.
\begin{equation}
    \mathbf{a} = 0, \quad \sigma_{ij} \gg 0 \tag{\ref{eq:hidden_action_stress}}
\end{equation}

\textbf{3. Gravitational Time Dilation:} \\
This overlapping scalar deficit depletes the available static pressure of the localized metric. Because compliance ($\alpha_s$) dictates the internal harmonic phase of a particle, the time-time metric ($g_{tt}$) is governed by this accumulated scalar deficit. % :::LLM_EDIT::: Stripped watchlist adjectives: physically, structural fluid, structural.

Substituting the superimposed dynamic pressure into the Bernoulli compliance limit drives Gravitational Time Dilation independent of spatial translation: % :::LLM_EDIT::: Stripped watchlist: mechanically.
\begin{equation}
    d\tau = dt \sqrt{1 - \frac{\Delta P_{net}}{P_c}} \tag{\ref{eq:scalar_time_dilation}}
\end{equation}

\textbf{4. Covariance:} \\ % :::LLM_EDIT::: Stripped 'Acoustic'.
The experience of proper time evaluates via two distinct states: % :::LLM_EDIT::: Stripped watchlist adjectives: geometric, fluid-dynamic.
\begin{itemize}[leftmargin=*]
    \item \textbf{The Boundary (The Clock):} A confined circulation experiences time as the kinematic evolution of the surrounding metric. Its internal clock is retarded because the ambient scalar pressure deficit ($\Delta P_{net}$) degrades the compliance of its spatial coordinate. % :::LLM_EDIT::: Stripped watchlist adjectives: topological, localized, resonant, structural, specific.
    \item \textbf{The Transverse Wave (The Signal):} An unconfined shear wave possesses no internal rest clock. It experiences time as the spatial evolution of the metric. Its transit evaluates as the spatial line integral of the refractive index across the tensioned continuum. % :::LLM_EDIT::: Stripped watchlist adjectives: acoustic, cumulative, macroscopic, fluid (replaced with continuum).
\end{itemize}
Because both the clock and the optical messenger are manifestations of the same background metric, they are bound by the same impedance. The depleted interstitial pressure retards the particle's internal clock by the same fraction that it refracts and delays the transiting wave. % :::LLM_EDIT::: Stripped watchlist adjectives: topological, physical, continuous fluid, structurally, fluid, mechanically, geometric.

When the retarded local clock measures the delayed transverse wave, the impedances cancel: % :::LLM_EDIT::: Stripped 'fluid'.
\begin{equation}
    c_{obs} = \frac{dr_{wave}}{d\tau} \equiv c \tag{\ref{eq:acoustic_covariance}}
\end{equation}
The transverse wave evaluates as the invariance that calibrates all observers to a covariant reality. % :::LLM_EDIT::: Stripped 'acoustic'.

\begin{tcolorbox}[recallbox, title={\ref{box:auxetic_symmetry_covariance}. Symmetry Breaking \& Covariance}] % :::LLM_EDIT::: Removed hyperref link from recall box title.
\begin{align}
    \nabla P_{net}(0) = 0 &\implies \mathbf{a} = 0 \tag{\ref{eq:gradient_cancellation}} \\
    \mathbf{a} = 0, &\quad \sigma_{ij} \gg 0 \tag{\ref{eq:hidden_action_stress}} \\
    \Delta P_{net} &= P_{dyn,A} + P_{dyn,B} \tag{\ref{eq:scalar_superposition}} \\
    d\tau &= dt_{bg} \sqrt{1 - \frac{\Delta P_{net}}{P_c}} \tag{\ref{eq:scalar_time_dilation}} \\
    c_{obs} &= \frac{dr_{wave}}{d\tau} \equiv c \tag{\ref{eq:acoustic_covariance}}
\end{align}
\end{tcolorbox}


% ----------------------------------------------------------------------
\section{Galactic Kinematics \& The Continuum Shear Transition}
\label{sec:deriv_galactic_dynamics}

\textbf{1. The Kinematic Yield Limit ($\bar{a}_0$):} \\
The ability of a coordinate layer to shear depends on the applied stress overcoming the baseline strain rate of the expanding volume. 

The boundary acceleration ($\bar{a}_0(z)$) evaluates as the kinematic yield limit of the metric: 
\begin{equation}
    \bar{a}_0(z) = \frac{c^2}{2\pi R_{h,0}(1+z)} \tag{\ref{eq:a0z}}
\end{equation}

\textbf{2. The Inner Shear Region (Keplerian Kinematics):} \\
The baryonic mass ($M_{bar}$) applies a spatial stress gradient to the continuum, evaluated as Newtonian coordinate acceleration ($g = GM_{bar}/r^2$). 

When this applied gravitational stress exceeds the yield threshold ($g > \bar{a}_0$), the coordinate layers slip. The velocity profile evaluates via Keplerian limits:
\begin{equation}
    v(r) = \sqrt{\frac{G M_{bar}}{r}} \tag{\ref{eq:kepler_core}} 
\end{equation}

\textbf{3. The Transition Boundary ($r_t$):} \\
As the radial coordinate extends outward, the inverse-square acceleration decays. Space no longer shears at the radial limit ($r_t$) where the baryonic acceleration drops below the yield threshold: 
\begin{align}
    g(r_t) &= \bar{a}_0(z) \label{eq:fracture_boundary} \\
    \frac{G M_{bar}}{r_t^2} &= \bar{a}_0(z) \label{eq:fracture_substitution}
\end{align}

Isolating the radius ($r_t$) yields the transition boundary: 
\begin{equation}
    r_t = \sqrt{\frac{G M_{bar}}{\bar{a}_0(z)}} \tag{\ref{eq:transition_radius}}
\end{equation}

\textbf{4. The Bingham Plug Flow (Kinematic Plateau):} \\
Beyond this fracture ($r > r_t$), the applied stress is insufficient to cause yield. In continuum mechanics, a medium with a yield stress behaves as a Bingham plastic: below the yield limit, internal deformation gradients vanish, and the medium locks into a rigid plug flow. 

The metric locks into a state of rigid plane stress, rotating as a unified bulk. To satisfy the boundary condition at the transition interface, the rotational velocity of the plug flow ($v_{rot}$) must match the Keplerian velocity at $r_t$:
\begin{equation}
    v_{rot} = v(r_t) = \sqrt{\frac{G M_{bar}}{r_t}} 
\end{equation}

Squaring both sides and substituting the transition radius ($r_t$) extracts the kinematic plateau: 
\begin{align}
    v_{rot}^2 &= \frac{G M_{bar}}{\sqrt{G M_{bar} / \bar{a}_0(z)}} \nonumber \\
    v_{rot}^2 &= \sqrt{\bar{a}_0(z) G M_{bar}} \nonumber \\
    v_{rot}^4 &= \bar{a}_0(z) G M_{bar} \tag{\ref{eq:tully_fisher_derivation}}
\end{align}

The Baryonic Tully-Fisher relation emerges from the Bingham yield limit of the metric. 

\vspace{0.2 cm}

\textbf{5. High-Redshift Resolution (The Keplerian Prediction):} \\
Because the boundary acceleration evaluates as inversely proportional to optical depth ($\bar{a}_0(z) \propto \frac{1}{1+z}$), the kinematic yield threshold shrinks in the early universe.

Evaluating the transition boundary limit as redshift increases ($z \gg 0$):
\begin{equation}
    \lim_{\bar{a}_0 \to 0} r_t = \lim_{\bar{a}_0 \to 0} \sqrt{\frac{G M_{bar}}{\bar{a}_0}} \to \infty \label{eq:high_z_limit}
\end{equation}

Because the transition boundary diverges to infinity at high redshift, early galaxies never intersect the limit. The metric remains entirely within the shear state (Zone 1). High-redshift galaxies therefore exhibit Newtonian/Keplerian rotation profiles ($v \propto 1/\sqrt{r}$) without flat velocity tails.

\vspace{0.2 cm}

\textbf{6. The Terminal Boundary ($R_{max}$):} \\
For local galaxies ($z \approx 0$), the plug flow region does not extend infinitely. While the inner boundary ($r_t$) evaluates via an acceleration match, the outer boundary evaluates via a kinematic velocity match. The background metric undergoes volumetric expansion, where outward radial velocity increases with coordinate distance ($v_{exp} = H_0 r$).

At the terminal boundary ($R_{max}$), the outward expansion velocity equals the constant rotational velocity of the plug flow:
\begin{align}
    H_0 R_{max} &= v_{rot} \nonumber \\
    R_{max} &= \frac{v_{rot}}{H_0} \tag{\ref{eq:r_max}} 
\end{align}

At this radial limit, the rotation of the galaxy matches the volumetric divergence of the background. The planar tension of the plug flow shears against the isotropic tension of the expanding metric. The coordinate space decouples across this boundary, releasing the rotational strain back into the 3D expansion of the medium. 

\begin{tcolorbox}[recallbox, title={\ref{box:galactic_dynamics}. Galactic Kinematics \& The Continuum Shear Transition}]
\begin{align}
    \bar{a}_0(z) &= \frac{c^2}{2\pi R_{h,0}(1+z)} \tag{\ref{eq:a0z}} \\
    r_t &= \sqrt{\frac{G M_{bar}}{\bar{a}_0(z)}} \tag{\ref{eq:transition_radius}} \\
    v_{rot}^4 &= \bar{a}_0(z) G M_{bar} \tag{\ref{eq:tully_fisher_derivation}} \\
    \lim_{\bar{a}_0 \to 0} r_t &\to \infty \implies \text{Keplerian Profile at High-}z \tag{\ref{eq:high_z_limit}} \\
    R_{max} &= \frac{v_{rot}}{H_0} \tag{\ref{eq:r_max}} 
\end{align}
\end{tcolorbox}




% ----------------------------------------------------------------------
\section{Galactic Lensing (Stress Birefringence)}
\label{sec:deriv_galactic_lensing}



\textbf{1. The Refractive Index of the Metric ($n_{eff}$):} \\
Lensing evaluates as the refraction of a transverse wave through a variable admittance medium. The effective refractive index ($n_{eff}$) evaluates as the inverse square of the total scalar admittance ($\alpha_s^{-2}$). % :::LLM_EDIT::: Stripped 'fluid'.

By the Bernoulli compliance limit, the total admittance squared evaluates via the pressure deficit ($P_{total} = P_{dyn} + P_{shear}$) scaled against the yield capacity ($P_c$). The refractive index evaluates via the continuum limits: 
\begin{equation}
    n_{eff}(r) = \frac{1}{\alpha_s^2(r)} = \left( 1 - \frac{P_{dyn} + P_{shear}}{P_c} \right)^{-1} \label{eq:refractive_index_base} 
\end{equation}

\textbf{2. Core Lensing (The 3D Isotropic Gradient):} \\
In the yielding core ($r < r_t$), the spatial gradient evaluates as an isotropic 3D Newtonian acceleration, yielding a dynamic pressure deficit ($P_{dyn} = \rho_{\tau} GM_{bar}/r$). The transverse shear evaluates to zero ($P_{shear} = 0$). % :::LLM_EDIT::: Stripped 'fluid' from Newtonian acceleration.

Substituting this dynamic pressure into the continuum refractive index defines the isotropic lens: 
\begin{align}
    n_{iso}(r) &= \left( 1 - \frac{\rho_{\tau} G M_{bar} / r}{\frac{1}{2} \rho_{\tau} c^2} \right)^{-1} \nonumber \\
    n_{iso}(r) &= \left( 1 - \frac{2G M_{bar}}{r c^2} \right)^{-1} \tag{\ref{eq:n_isotropic}}
\end{align}

Because the gradient of a 3D spherical pressure well ($\nabla P_{dyn} \propto 1/r^2$) is symmetric, a photon transiting the core undergoes isotropic deflection regardless of the inclination angle of the galactic disk.

\textbf{3. Plug Flow Lensing (The 2D Cylindrical Gradient):} \\
In the locked Bingham plug flow ($r_t < r < R_{max}$), the metric evaluates as a 2D cylindrical plane stress. The scalar pressure evaluates via the logarithmic shear: 
\begin{equation}
    P_{shear}^{2D}(r) = k \ln \left( \frac{R_{max}}{r} \right) 
\end{equation}

By the kinematic plateau extraction ($v_{rot} = c \sqrt{k / 2P_c}$), the shear amplitude ($k$) evaluates as $k = 2P_c (v_{rot}^2 / c^2)$. Substituting this amplitude normalizes the shear pressure against the baseline stiffness ($P_c$):
\begin{equation}
    \frac{P_{shear}^{2D}(r)}{P_c} = 2 \left(\frac{v_{rot}^2}{c^2}\right) \ln \left( \frac{R_{max}}{r} \right)
\end{equation}

Substituting this 2D plane stress component defines the cylindrical lens: 
\begin{align}
    n_{cyl}(r) &= \left( 1 - \frac{0 + P_{shear}^{2D}(r)}{P_c} \right)^{-1} \nonumber \\
    n_{cyl}(r) &= \left( 1 - 2\left(\frac{v_{rot}^2}{c^2}\right) \ln \left( \frac{R_{max}}{r} \right) \right)^{-1} \tag{\ref{eq:n_cylindrical}} 
\end{align}

\textbf{4. The Anisotropic Deflection Integral ($\Delta \theta$):} \\
The deflection angle ($\Delta \theta$) of a transiting photon evaluates as the line integral of the transverse gradient of the refractive index ($\Delta \theta \propto \int \nabla_{\perp} n \, dl$).

In cylindrical coordinates $(r, \phi, z)$, the 2D plane stress evaluates as translationally invariant along the rotational axis ($\partial_z P_{shear}^{2D} = 0$). 
Because the optical strain varies radially but remains constant axially, the transverse gradient ($\nabla_{\perp} n_{cyl}$) evaluates as a function of the observation inclination angle ($I$), where $I=0^\circ$ is face-on and $I=90^\circ$ is edge-on. 

The effective optical path gradient ($f_I$) evaluates as the projection of the radial strain across the line of sight: 
\begin{equation}
    f_I = \sin(I) \nabla_r n_{cyl} + \cos(I) \nabla_z n_{cyl}
\end{equation}
Substituting the translational invariance ($\nabla_z n_{cyl} = 0$) isolates the inclination dependency:
\begin{equation}
    f_I = \sin(I) \nabla_r n_{cyl} \label{eq:inclination_bias}
\end{equation}

\textbf{5. The Total Observational Deflection ($\Delta \theta_{tot}$):} \\
To predict the macroscopic lensing signature of a galaxy, the unified effective index ($n_{tot}$) evaluates the simultaneous total pressure deficit ($P_{dyn} + P_{shear}^{2D}$): % :::LLM_EDIT::: Added explicit total deflection synthesis.
\begin{equation}
    n_{tot}(r) = \left( 1 - \frac{P_{dyn}}{P_c} - \frac{P_{shear}^{2D}}{P_c} \right)^{-1}
\end{equation}

In the weak-field regime governing galactic radii ($P \ll P_c$), applying the binomial expansion ($(1-x)^{-1} \approx 1+x$) resolves the linear superposition of the individual refractive states:
\begin{equation}
    n_{tot}(r) \approx 1 + \frac{P_{dyn}}{P_c} + \frac{P_{shear}^{2D}}{P_c} \approx n_{iso}(r) + n_{cyl}(r) - 1
\end{equation}

Taking the transverse spatial gradient ($\nabla_\perp$) of this super-positioned index evaluates the optical strain. Because the 3D isotropic core deflects radially everywhere, its gradient evaluates cleanly as $\nabla_r n_{iso}$. Because the 2D cylindrical plane stress projects via the inclination dependency ($f_I$), its gradient evaluates as $\sin(I) \nabla_r n_{cyl}$.

The total deflection angle ($\Delta \theta_{tot}$) evaluates as the line integral of the superimposed transverse gradient:
\begin{equation}
    \Delta \theta_{tot}(I) \propto \int \Big( \nabla_r n_{iso} + \sin(I) \nabla_r n_{cyl} \Big) \, dl \tag{\ref{eq:total_deflection}}
\end{equation}

A face-on galaxy ($I=0^\circ$) evaluates lensing via its baryonic mass ($n_{iso}$).

\textbf{6. The Terminal Shear Ring ($R_{max}$):} \\
The Bingham plug flow terminates at the Jacobian Capture Radius ($R_{max} = v_{rot}/H_0$). Evaluating the 2D logarithmic shear pressure at this radial limit yields: 
\begin{align}
    P_{shear}^{2D}(R_{max}) &= 2P_c \left(\frac{v_{rot}^2}{c^2}\right) \ln \left( \frac{R_{max}}{R_{max}} \right) \nonumber \\
    P_{shear}^{2D}(R_{max}) &= 2P_c \left(\frac{v_{rot}^2}{c^2}\right) \ln (1) \nonumber \\
    P_{shear}^{2D}(R_{max}) &= 0 \label{eq:terminal_shear_zero}
\end{align}

Because the shear pressure evaluates to exactly zero at $R_{max}$, substituting this limit into the cylindrical refractive index returns the metric to baseline:
\begin{align}
    n_{cyl}(R_{max}) &= \left( 1 - 0 \right)^{-1} \nonumber \\
    n_{cyl}(R_{max}) &= 1 \label{eq:n_baseline_return}
\end{align}

The refractive strain does not asymptotically fade. The optical strain evaluates a discontinuity at $R_{max}$, establishing a terminal shear ring where the anomalous deflection drops to zero, returning the geometry to the $n_{iso}$ potential.

\begin{tcolorbox}[recallbox, title={\ref{box:galactic_lensing}. Galactic Lensing (Stress Birefringence)}] 
\begin{align}
    n_{iso}(r) &= \left( 1 - \frac{2G M_{bar}}{r c^2} \right)^{-1} \tag{\ref{eq:n_isotropic}} \\
    n_{cyl}(r) &= \left( 1 - 2\left(\frac{v_{rot}^2}{c^2}\right) \ln \left( \frac{R_{max}}{r} \right) \right)^{-1} \tag{\ref{eq:n_cylindrical}} \\
    \Delta \theta_{tot}(I) &\propto \int \Big( \nabla_r n_{iso} + \sin(I) \nabla_r n_{cyl} \Big) \, dl \tag{\ref{eq:total_deflection}} \\
    n_{cyl}(R_{max}) &= 1 \tag{\ref{eq:n_baseline_return}}
\end{align}
\end{tcolorbox}


\newpage

\part*{PHASE II: VACUUM RHEOLOGY THERMODYNAMICS}

 % ----------------------------------------------------------------------
\section{Thermodynamic Vacuum \& Particle Entropy}
\label{sec:deriv_thermo_vacuum}

The baseline temperature evaluates independent of volumetric expansion via the First Law of Thermodynamics: 
\begin{equation}
    T dS = dU + P dV \nonumber
\end{equation}

\textbf{1. The Unruh-Hawking Temperature Floor:} \\
Internal energy evaluates as density times volume $U = \rho V$. Total entropy $S$ evaluates as entropy density $\sigma$ times volume $V$, defining $S = \sigma V$. Applying the product rule to the differentials yields:
\begin{align}
    dU &= d(\rho V) = \rho dV + V d\rho  \\
    dS &= d(\sigma V) = \sigma dV + V d\sigma 
\end{align}

Substituting these into the First Law:
\begin{equation}
    T(\sigma dV + V d\sigma) = (\rho dV + V d\rho) + P dV 
\end{equation}

Distributing the temperature $T$ and grouping all terms containing the volumetric expansion factor $dV$ on the right side:
\begin{align}
    V T d\sigma - V d\rho &= \rho dV + P dV - T \sigma dV  \\
    V (T d\sigma - d\rho) &= dV (\rho + P - T \sigma) 
\end{align}

The Euler Equation of State for any continuous extensive medium dictates that $\rho + P = T \sigma$. Therefore, the expression inside the parenthesis evaluates to zero ($\rho + P - T \sigma = 0$). 
\begin{align}
    V (T d\sigma - d\rho) &= dV (0) = 0  \\
    T d\sigma &= d\rho 
\end{align}

The volumetric differential $dV$ cancels from the thermodynamic evolution. The temperature evaluates via the kinematic acceleration of the expanding boundary. 

Via the Unruh-Hawking effect, boundary acceleration ($a_h = c/t$) generates a thermodynamic floor ($T = \hbar a_h / 2\pi k_B c$). Substituting the boundary limit $a_h$: 
\begin{align}
    T_{cosmic}(t) &= \frac{\hbar (c/t)}{2\pi k_B c}  \\
    T_{cosmic}(t) &= \frac{\hbar}{2\pi k_B t} \tag{\ref{eq:tcosmic}}
\end{align}

\textbf{2. Particle Entropy Quantum ($S_p$):} \\
The topological entropy ($S_p$) of a localized cavity evaluates via the Euler thermodynamic relation: 
\begin{equation}
    S_p = \frac{(\rho + P)V}{T} 
\end{equation}

Because a radiation-dominated acoustic cavity evaluates with an internal pressure $P = \rho/3$, substituting $P$ into the equation evaluates the numerator ($\rho + \rho/3 = 4\rho/3$). 
\begin{equation}
    S_p = \frac{\left(\frac{4\rho}{3}\right)V}{T} = \frac{4 \rho V}{3 T} 
\end{equation}

The internal baseline energy evaluates as the product of density and volume ($\rho V = E_{vac} = \hbar/t$). Substituting this energy and the Unruh temperature floor ($T = \hbar / 2\pi k_B t$) into the entropy equation yields: 
\begin{align}
    S_p &= \frac{\frac{4}{3} \left(\frac{\hbar}{t}\right)}{\left(\frac{\hbar}{2\pi k_B t}\right)}  \\
    S_p &= \frac{4}{3} \left( \frac{\hbar}{t} \right) \left( \frac{2\pi k_B t}{\hbar} \right) 
\end{align}

The parameters $\hbar$ and $t$ cancel: 
\begin{align}
    S_p &= \frac{4}{3} \left( 2\pi k_B \right)  \\
    S_p &= \frac{8\pi}{3} k_B \tag{\ref{eq:sp}}
\end{align}

\begin{tcolorbox}[recallbox, title={\ref{box:thermo_vacuum}. Thermodynamic Vacuum \& Entropy}]
\begin{align}
    a_h(t) &= \frac{c}{t} \implies T_{cosmic}(t) = \frac{\hbar}{2\pi k_B t} \tag{\ref{eq:tcosmic}} \\ 
    S_p &= \frac{(\rho + P)V}{T} \implies S_p = \frac{8\pi}{3} k_B \tag{\ref{eq:sp}}
\end{align}
\end{tcolorbox}



% ----------------------------------------------------------------------
\section{The Unified Equation of State (Vacuum Energy \& Topological Pressure)}
\label{sec:deriv_unified_eos}

The baseline energy of the continuum and the phase pressure of a particle evaluate as symmetric manifestations of a single equation of state. 

\textbf{1. Fundamental Vacuum Energy \& Null Action:} \\
For an acoustic standing wave bounded by the scale radius of the universe ($R_h = c \cdot t$), the wave vector evaluates as the inverse of this boundary ($k = 1/R_h$). The baseline energy ($E_{vac}$) of this boundary mode evaluates via the acoustic limit ($E = pc$, where $p = \hbar k$): 
\begin{align}
    E_{vac}(t) &= \left( \frac{\hbar}{R_h} \right) c  \\
    E_{vac}(t) &= \frac{\hbar c}{c \cdot t} = \frac{\hbar}{t} \tag{\ref{eq:vacuum_null_action}}
\end{align}

To evaluate the strain of this expanding boundary, the continuum Lagrangian ($\mathcal{L}_{vac}$) evaluates via the Legendre transform of the energy ($pv - E$). Because the boundary propagates at $c$, we substitute $v = c$ and $E_{vac} = pc$: 
\begin{align}
    \mathcal{L}_{vac}(c) &= p(c) - (pc)  \\
    \mathcal{L}_{vac}(c) &= 0 
\end{align}
Because the action evaluates to zero, the boundary expands as a null manifold, accumulating zero strain.

\textbf{2. The Topological Equation of State:} \\
When the continuum cavitates to form a particle, its thermodynamic state is governed by the Euler relation ($S_p = (\rho + P)V / T$). 

Because the trapped standing wave evaluates with an internal radiation pressure $P_{int} = \rho/3$, substituting $P_{int}$ into the equation evaluates the numerator ($\rho + \rho/3 = 4\rho/3$). 

Substituting this into the Euler relation yields $S_p = 4PV/T$. As established via the scaling of the continuum, the entropy of a spatial bit evaluates to $S_p = \frac{8\pi}{3} k_B$. Substituting this constant isolates the state equation: 
\begin{align}
    PV &= \frac{1}{4} \left( \frac{8\pi}{3} k_B \right) T  \\
    PV &= \frac{2\pi}{3} k_B T \tag{\ref{eq:unified_topological_eos}}
\end{align}
The molar quantity ($n$) utilized in classical thermodynamics evaluates for a particle as $2\pi/3$, dictated by the $2\pi$ phase wrap bounding the 3D spherical volume. 

\textbf{3. Rest Mass Equivalence \& Bag Pressure:} \\
For a radiation-dominated acoustic cavity, internal thermal energy evaluates as $U = \rho V = 3PV$. Substituting the derived Equation of State:
\begin{align}
    U &= 3 \left( \frac{2\pi}{3} k_B T \right)  \\
    U &= 2\pi k_B T 
\end{align}
Substituting the temperature of a boundary ($T = \frac{mc^2}{2\pi k_B}$) isolates the rest-mass energy ($U_{topo} = mc^2$). 

Isolating the radiation pressure from this energy equation ($P_{int} = \frac{m c^2}{3 V}$) and substituting the Euclidean bounding volume ($V = \frac{4}{3}\pi (\hbar/mc)^3$) yields: 
\begin{align}
    P_{int} &= \frac{m^4 c^5}{4\pi \hbar^3}  \\
    P_{int} &= 2\pi P_p \tag{\ref{eq:unified_pressure_ratio}}
\end{align}
The internal pressure of the particle evaluates as the Bag Pressure ($P_p$) multiplied by one full phase wrap ($2\pi$). 

\textbf{4. Scale Invariance (The Cosmic Loop):} \\
To prove the scale invariance of the metric, we apply the particle energy equation ($U = 2\pi k_B T$) to the cosmos. Substituting the temperature floor of the expanding universe ($T = \frac{\hbar}{2\pi k_B t}$) yields:
\begin{align}
    U_{vac} &= 2\pi k_B \left( \frac{\hbar}{2\pi k_B t} \right)  \\
    U_{vac} &= \frac{\hbar}{t} \tag{\ref{eq:vacuum_null_action}}
\end{align}
The metric's equation of state recovers the Fundamental Vacuum Energy derived from the horizon boundary in Step 1. 

\begin{tcolorbox}[recallbox, title={\ref{box:unified_eos}. The Unified Equation of State}]
\begin{align}
    E_{vac}(t) &= \frac{\hbar}{t} \implies \mathcal{L}_{vac}(c) = 0 \tag{\ref{eq:vacuum_null_action}} \\
    PV &= \frac{2\pi}{3} k_B T \implies U = 2\pi k_B T \tag{\ref{eq:unified_topological_eos}} \\
    U_{topo} &= mc^2 \implies P_{int} = 2\pi P_p \tag{\ref{eq:unified_pressure_ratio}}
\end{align}
\end{tcolorbox}




% ----------------------------------------------------------------------
\section{The Continuous Rheological Gradient}
\label{sec:deriv_rheological_gradient}

The mechanics of the continuum resolve via a rheological gradient governed by localized velocity ($v(r)$). 

\textbf{1. The Elastic Relation:} \\
The Poisson ratio ($\nu$) evaluates the fraction of transverse expansion to longitudinal extension. It is defined by the Bulk Modulus ($K$) and the Shear Modulus ($\mu_s$):
\begin{equation}
    \nu(r) = \frac{3K(r) - 2\mu_s}{2(3K(r) + \mu_s)} 
\end{equation}

\textbf{2. The Macroscopic Vacuum (Transverse Compliance):} \\
In the unperturbed vacuum ($r \to \infty$), velocity evaluates as negligible ($v \to 0$). The viscoelastic modulus ratio ($\alpha \approx 1/137$) evaluates as the fraction of shear capacity to bulk compressive capacity ($\alpha = \mu_s / K$). 

Substituting the bulk modulus ($K = \mu_s / \alpha$) into the relation derives the baseline mechanical behavior of the metric: 
\begin{align}
    \nu_{macro} &= \frac{3(\mu_s / \alpha) - 2\mu_s}{2(3(\mu_s / \alpha) + \mu_s)} \nonumber \\
    \nu_{macro} &= \frac{3/\alpha - 2}{6/\alpha + 2} \approx 0.496 \tag{\ref{eq:poisson_macro}} 
\end{align}
The vacuum evaluates as a nearly incompressible continuum ($\nu \approx 0.5$). It resists volumetric compression while maintaining compliance for transverse shear waves (electromagnetism and gravitation).

\textbf{3. The Compton Boundary (Phase Stiffening):} \\
Standard mass evaluates as a standing wave bounded by the Compton wavelength ($R_h = \lambda_c$). 

At the Compton boundary, internal circulation ($v$) is elevated, causing the Bulk Modulus to kinematically stiffen ($K \to \infty$). However, because $v \ll c$ at the macroscopic boundary, the continuum does not choke. 
\begin{equation}
    0.496 \le \nu(\lambda_c) < 0.5 \tag{\ref{eq:poisson_compton}} 
\end{equation}
This positive Poisson ratio provides structural compliance. The boundary yields to transverse strain without structural fracture.

\textbf{4. The Core (The Incompressible Anchor):} \\
As circulation approaches the cavitation boundary ($r \to r_c$), the velocity strikes the acoustic phase limit ($v \to c$). 

Because velocity cannot exceed $c$, resistance to volumetric compression diverges to infinity ($\lim_{v \to c} K = \infty$). Dividing the Poisson limit by $K$ isolates the state: 
\begin{align}
    \nu_{core} &= \lim_{K \to \infty} \frac{3 - 2\frac{\mu_s}{K}}{2(3 + \frac{\mu_s}{K})} \nonumber \\
    \nu_{core} &= \frac{3}{6} = 0.5 \tag{\ref{eq:poisson_local}}
\end{align}
At $\nu = 0.5$, the core becomes incompressible ($\nabla \cdot \mathbf{u} = 0$). This rigid limit acts as the mechanical anchor, preventing the boundary from collapsing under ambient pressure.

\vspace{0.2 cm}

\textbf{5. The Fine-Structure Compliance Margin ($\alpha/2$):} \\
Because the viscoelastic modulus ratio ($\alpha$) governs the relationship between the shear modulus and bulk modulus ($K = \mu_s / \alpha$), the volumetric compliance is tied to the electromagnetic transverse limit. 

Subtracting the baseline Poisson ratio from an incompressible boundary ($0.5$) isolates the exact volumetric compliance margin of the unperturbed coordinate space:
\begin{align}
    \Delta \nu &= 0.5 - \nu_{macro} \nonumber \\
    \Delta \nu &= \frac{1}{2} - \frac{3 - 2\alpha}{6 + 2\alpha} \nonumber \\
    \Delta \nu &= \frac{3\alpha}{6 + 2\alpha} \approx \frac{\alpha}{2}
\end{align}
The metric deviates from absolute rigidity ($\nu=0.5$) by approximately half the fine-structure constant. The metric's mechanical capacity to undergo volumetric compression (Gravity) is bounded by its capacity to sustain transverse shear (Electromagnetism).

\begin{tcolorbox}[recallbox, title={\ref{box:rheological_gradient}. The Rheological Gradient}]
\begin{align}
    \nu_{macro} &\approx 0.496 \tag{\ref{eq:poisson_macro}} \\
    \nu(\lambda_c) &\in [0.496, 0.5) \tag{\ref{eq:poisson_compton}} \\
    \nu_{core} &= 0.5 \implies \nabla \cdot \mathbf{u} = 0 \tag{\ref{eq:poisson_local}}
\end{align}
\end{tcolorbox}







% ----------------------------------------------------------------------
\section{The Fine-Structure Constant \& Transverse Shear Velocity}
\label{sec:deriv_fine_structure}

In continuum mechanics, the structural limits of a medium partition into a bulk compressive capacity ($P_c$) and a transverse shear capacity ($\Pi_{shear}$). The Fine-Structure Constant ($\alpha$) equates to the ratio of these orthogonal thermodynamic bounds: 
\begin{equation}
    \alpha = \frac{\Pi_{shear}}{P_{c}} \tag{\ref{eq:alpha_ratio}}
\end{equation}

\textbf{1. Kinematic Capacity Limits:} \\
The thermodynamic capacities ($P_c$ , $\Pi_{shear}$) scale against the kinematic wave limits of the baseline metric density ($\rho_\tau$). Longitudinal compression propagates up to the acoustic phase limit ($c$), structuring the bulk capacity: 
\begin{equation}
    P_{c} = \frac{1}{2} \rho_\tau c^2 \tag{\ref{eq:longitudinal_capacity}}
\end{equation}

Transverse tearing operates up to the bare shear wave velocity ($v_s$), setting the maximum transverse capacity:
\begin{equation}
    \Pi_{shear} = \frac{1}{2} \rho_\tau v_s^2 \label{eq:transverse_capacity}
\end{equation}

\textbf{2. The Transverse Velocity Limit ($v_s$):} \\
Substituting the kinematic capacity limits into the modulus ratio reduces to a velocity ratio: 
\begin{align}
    \alpha &= \frac{\frac{1}{2} \rho_\tau v_s^2}{\frac{1}{2} \rho_\tau c^2} \nonumber \\
    \alpha &= \frac{v_s^2}{c^2} \label{eq:alpha_velocity_ratio}
\end{align}

Rearranging the velocity ratio calculates the phase limit of bare transverse slip:
\begin{equation}
    v_s = c \sqrt{\alpha} \tag{\ref{eq:bare_shear_velocity}}
\end{equation}

While coupled electromagnetic waves propagate at $c$, this limit ($v_s \approx 0.085 c$) restricts uncoupled geometric shear. It predicts the maximum stable transverse orbital velocity of a macroscopic mass (e.g., the pericenter of a stellar orbit around a gravitational singularity). A topological boundary subjected to a transverse slip exceeding this limit fractures the continuum, identifying the geometric threshold for macroscopic tidal disruption.


\begin{tcolorbox}[recallbox, title={\ref{box:fine_structure}. The Fine-Structure Constant \& Shear Velocity}]
\begin{align}
    &\quad \alpha = \frac{\Pi_{shear}}{P_{c}} \tag{\ref{eq:alpha_ratio}} \\
    &\quad P_{c} = \frac{1}{2} \rho_\tau c^2, \quad \Pi_{shear} = \frac{1}{2} \rho_\tau v_s^2 \tag{\ref{eq:longitudinal_capacity}, \ref{eq:transverse_capacity}} \\
    &\quad \alpha = \frac{v_s^2}{c^2} \implies v_s = c \sqrt{\alpha} \tag{\ref{eq:bare_shear_velocity}}
\end{align}
\end{tcolorbox}









% ----------------------------------------------------------------------
\section{The Cosmological Deborah Number}
\label{sec:deriv_unified_deborah}

The rheological phase response and the relative strengths of the fundamental forces evaluate via a single ratio dictated by interaction energy. 

\textbf{1. The Predictive Equation (The Transverse Yield Limit):} \\
The Deborah Number ($De$) emerges as the ratio of the intrinsic structural relaxation time ($t_r$) to the applied strain timescale. Inverting to the frequency domain, it translates to the applied angular frequency ($\omega$) multiplied by the baseline relaxation time ($De_{base} = \omega t_r$). 

Intrinsic relaxation time ($t_r$) resolves via the yield limit. The continuum possesses a longitudinal spatial yield limit ($R$), beyond which it cannot sustain phase. Because stable boundaries conform via their Compton limit ($R = \hbar/mc$), this spatial yield limit defines the baseline mass ($\bar{\lambda}_e = \hbar/m_e c$). 

The baseline relaxation time equates to this yield limit divided by the phase velocity ($c$): 
\begin{equation}
    t_r = \frac{R_e}{c} = \frac{\hbar}{m_e c^2} 
\end{equation}

Because 3D manifolds are anisotropic, they absorb longitudinal compression but fracture under transverse shear. To calculate the strain capable of tearing the metric, the applied interaction energy ($E = \hbar\omega$) must be filtered by the transverse compliance fraction ($\alpha$). The active transverse frequency acts as $\omega_{\perp} = \alpha E / \hbar$. 

Substituting these parameters into the Deborah Number equation yields: 
\begin{equation}
    De = \omega_{\perp} t_r = \left( \frac{\alpha E}{\hbar} \right) \left( \frac{\hbar}{m_e c^2} \right) 
\end{equation}

The reduced Planck constant ($\hbar$) cancels. The Deborah Number scales as the active transverse interaction energy against the baseline yield: 
\begin{equation}
    De = \alpha \left( \frac{E}{m_e c^2} \right) \tag{\ref{eq:unified_deborah}}
\end{equation}

\textbf{2. Complex Dynamic Modulus ($G^*$):} \\
The frequency-dependent response is determined by applying the Laplace Transform to the Maxwell viscoelastic equation ($s \varepsilon = s \sigma / \mu_s + \sigma / \eta$). Isolating the ratio of stress to strain constructs the Complex Dynamic Modulus ($G^*$).  

Substituting the harmonic oscillation limit ($s = i\omega$) and the Deborah Number produces the real elastic Storage Modulus ($G'$) and the imaginary viscous Loss Modulus ($G''$): 
\begin{equation}
    G^*(E) = \mu_s \left[ \frac{De^2}{1 + De^2} \right] + i \mu_s \left[ \frac{De}{1 + De^2} \right] \tag{\ref{eq:deborah_modulus}}
\end{equation}

\textbf{3. The Viscoelastic Phase Angle (Angular Distribution):} \\
The phase angle ($\delta$) measures the partition between the dissipated energy (Loss Modulus, $G''$) and the stored elastic energy (Storage Modulus, $G'$). The loss tangent equates to their quotient: 
\begin{equation}
    \tan \delta = \frac{G''}{G'} = \frac{\frac{De}{1+De^2}}{\frac{De^2}{1+De^2}} = \frac{1}{De}
\end{equation}

The compliance fraction ($\alpha_s = g^2/4\pi$) acts as the viscous response ($G''$). The transverse coupling fraction ($\alpha$) acts as the elastic response ($G'$).  

The angular distribution between the strong force and the electromagnetic force projects as the phase angle. Equating the loss tangent isolates the structural ratio: 
\begin{equation}
    \frac{\alpha_s}{\alpha} = \frac{1}{De} \implies De = \frac{\alpha}{\alpha_s} \tag{\ref{eq:de_ratio}}
\end{equation}

Substituting the equation ($De = \alpha E / m_e c^2$) into the inverted ratio provides the dynamic coupling limit:
\begin{align}
    \frac{\alpha_s}{\alpha} &= \frac{m_e c^2}{\alpha E} \\
    \alpha_s &= \frac{m_e c^2}{E} \tag{\ref{eq:angular_distribution}} 
\end{align}
The transverse coupling fraction ($\alpha$) cancels from the equivalence. The scalar admittance of the continuum ($\alpha_s$) evaluates as the ratio of the baseline rest mass to the applied interaction energy. 

\vspace{0.2 cm}

\textbf{4. Confinement \& The Continuum Beta Function ($\beta$):} \\ 
At the boundary limit ($E = m_e c^2$), the ratio evaluates to $\alpha_s = 1$. The strong coupling parameter isolates to its bare scalar limit ($g^2 \approx 4\pi$), defining Color Confinement.  

The scale variance of this interaction evaluates as the logarithmic derivative of the coupling with respect to the applied interaction energy ($E$). This defines the Continuum Beta Function ($\beta(\alpha_s)$): 
\begin{equation}
    \beta(\alpha_s) = \frac{\partial \alpha_s}{\partial \ln E} = E \frac{\partial \alpha_s}{\partial E}
\end{equation}

Substituting the derived compliance state ($\alpha_s = m_e c^2 / E$), the linear differential limit evaluates to $\partial \alpha_s / \partial E = - m_e c^2 / E^2$:
\begin{align}
    \beta(\alpha_s) &= E \left( - \frac{m_e c^2}{E^2} \right) \\ 
    \beta(\alpha_s) &= - \frac{m_e c^2}{E}
\end{align}

Because the quotient isolates the baseline scalar fraction ($\alpha_s = m_e c^2 / E$), the beta function converges to a closed-form limit: 
\begin{equation}
    \beta(\alpha_s) = - \alpha_s \tag{\ref{eq:beta_function}}
\end{equation}

This negative differential limit structures the mechanics of Asymptotic Freedom. As interaction energy scales toward infinity ($E \to \infty$), the Deborah Number diverges ($De \gg 1$), and the phase angle collapses ($\tan \delta \to 0$). The scalar response vanishes ($\alpha_s \to 0$), freezing the metric into a solid. The apparent strengths of the fundamental forces evaluate as this frequency-dependent phase angle. 

\vspace{0.2 cm}

\textbf{5. The Glass Transition and Topological Fracture:} \\
The interaction energy where the continuum behaves like a solid is defined by the condition ($De \ge 1$).  
\begin{align}
    1 & \le \alpha \left( \frac{E_{crit}}{m_e c^2} \right) \\
    E_{crit} &\ge \frac{m_e c^2}{\alpha}
\end{align}

At this threshold, the metric loses compliance and locks into an amorphous solid state. Because the non-linear interaction applies unbounded stress to this rigid geometry, it cannot dissipate the energy via flow. The stress exceeds the tensile yield strength, causing viscoelastic fracture. 

Evaluating this limit using the electron mass ($0.511 \text{ MeV}$), the yield threshold bounds to $E_{crit} \approx 70.02 \text{ MeV}$. Because hadronization nucleates boundaries in symmetric pairs, the total cavitation energy requires:
\begin{equation}
    2 E_{crit} \ge 140.04 \text{ MeV} \equiv m_{\pi^{\pm}} \tag{\ref{eq:de_glass}}
\end{equation}
This identifies the rest mass of the charged pion ($m_{\pi^{\pm}} \approx 139.57 \text{ MeV}$). 

The spatial coordinate of this transition evaluates via the kinematic interaction length ($E = \hbar c / L$). Substituting the critical energy evaluates the transverse spatial fracture limit ($L_{crit}$): 
\begin{align}
    L_{crit} &= \frac{\hbar c}{E_{crit}} = \frac{\hbar c}{m_e c^2 / \alpha} \\
    L_{crit} &= \alpha \frac{\hbar}{m_e c} \equiv r_e \approx 2.81 \text{ fm} \tag{\ref{eq:de_range}}
\end{align}
The Classical Electron Radius ($r_e$) identifies the absolute spatial range of the strong interaction. Beyond this coordinate, the solid metric relaxes into the yielding continuum. 

\vspace{0.2 cm}

\textbf{6. The Leptonic-Hadronic Yield Threshold ($\gamma_{yield}$):} \\ 
The Deborah Number establishes the baseline phase state of elementary particles relative to the continuum limit ($m_e$). For an electron at rest, $De = \alpha \ll 1$, classifying it as compliant (Lepton). For a proton at rest, $De = \alpha (m_p/m_e) \approx 13.4 \gg 1$, classifying it as a rigid (Hadron). 

Relativistic translation for an electron evaluates as a coordinate transformation ($\alpha_s = 1/\gamma$). However, upon collision, this compliance collapses into structural strain. Substituting the Lorentz translation state of an electron into the dynamic ratio ($De = \alpha / \alpha_s$) evaluates its kinematic interaction limit ($De = \alpha \gamma$). 

An electron collision shatters the metric into hadronic geometries (pions) when its kinetic state hits the glass transition ($De = 1$): 
\begin{align}
    1 &= \alpha \gamma_{yield} \\
    \gamma_{yield} &= \frac{1}{\alpha} \approx 137 \tag{\ref{eq:gamma_yield}}
\end{align}

To accelerate an electron to this fracture limit, the applied kinetic energy evaluates as $E_k = (\gamma - 1) m_e c^2$. Dividing by the elementary charge ($e$) isolates the continuous acceleration potential ($V_{yield}$) required to shatter the metric:
\begin{align}
    V_{yield} &= \frac{\left(\frac{1}{\alpha} - 1\right) m_e c^2}{e} \\ 
    V_{yield} &\approx 69.51 \text{ MV} 
\end{align}
Electrons accelerated through a potential difference of $\approx 69.51 \text{ MV}$ cross the yield threshold upon collision, ejecting their kinetic energy as hadronic mass.

\vspace{0.2 cm}

\textbf{7. The Atomic Phase Limit ($Z \approx 137$):} \\ 
The internal circulation velocity of an atomic orbital scales with the central geometric charge ($v = Z \alpha c$). The structural limit of the atom evaluates against the kinematic yield threshold ($De \to 1$). 

When the central charge forces the orbital circulation to the phase limit ($v \to c$), the interaction strikes the yield boundary: 
\begin{align}
    Z \alpha c &= c \\
    Z &= \frac{1}{\alpha} \approx 137 \tag{\ref{eq:atomic_phase_limit}}
\end{align}
At $Z \approx 137$, internal circulation strikes the phase limit. The topological boundary crushes to the $r_e$ spatial floor, triggering structural shattering. This establishes a physical reason why the periodic table terminates before element $137$. 

\vspace{0.2 cm}

\textbf{8. Electro-Strong Phase Symmetry ($De = 1$):} \\ 
At the glass transition limit ($De = 1$), substituting the unity parameter into the ratio evaluates the phase symmetry limit: 
\begin{align}
    1 &= \frac{\alpha}{\alpha_s} \\
    \alpha_s &= \alpha \tag{\ref{eq:electro_strong_symmetry}}
\end{align}
Substituting the geometric definition of the scalar compliance fraction ($\alpha_s = g^2 / 4\pi$) isolates the bare scalar charge ($g^2$) at the yield threshold: 
\begin{align}
    \frac{g^2}{4\pi} &= \alpha \\
    g^2 &= 4\pi\alpha \tag{\ref{eq:strong_charge_symmetry}}
\end{align}
Calculating the phase angle leaves $\tan \delta = 1 \implies \delta = \pi/4$. Evaluating the Complex Dynamic Modulus (Equation \ref{eq:deborah_modulus}) at $De = 1$ partitions the stored and dissipated capacities into equal halves ($G'/\mu_s = G''/\mu_s = 1/2$). Substituting this $1/2$ capacity partition into the metric compliance ($\alpha_{yield} = \sqrt{1 - 1/2}$) isolates the structural yield compliance ($\alpha_{yield} = 1/\sqrt{2}$) bounding observable mass at fracture ($M = \alpha_{yield} M_{bare}$). % :::LLM_EDIT::: Derived alpha_yield = 1/sqrt(2) from De=1 modulus partition.

\vspace{0.2 cm}

\textbf{9. The Macroscopic Drag Limit (Grand Unification):} \\ 
Evaluating the Deborah Number at the kinematic minimum computes the ratio of elastic to viscous resistance. For a translating mass, the maximum observation time corresponds to the proper age of the universe ($t_{obs} = R_{h,0} / c$).
\begin{equation}
    De_{min} = \frac{t_r}{t_{obs}} = \frac{\alpha \hbar / m_e c^2}{R_{h,0} / c} = \alpha \left( \frac{\hbar}{m_e c R_{h,0}} \right)
\end{equation}

As derived in Section 4, the boundary acceleration translates to the scale radius ($\bar{a}_0 = c^2 / 2\pi R_{h,0} \implies R_{h,0} = c^2 / 2\pi \bar{a}_0$). Substituting this definition into the drag limit produces the ratio:
\begin{align}
    De_{min} &= \alpha \left( \frac{\hbar}{m_e c \left( \frac{c^2}{2\pi \bar{a}_0} \right)} \right) \\
    De_{min} &= \alpha \left( \frac{2\pi \hbar \bar{a}_0}{m_e c^3} \right) \label{eq:alpha_a0_bridge}
\end{align}

Separating the Unruh thermal energy of the cosmos ($E_{unruh} = \hbar \bar{a}_0 / 2\pi c$) and the rest energy of the electron ($E_{rest} = m_e c^2$):
\begin{equation}
    De_{min} = \alpha (2\pi)^2 \left( \frac{E_{unruh}}{E_{rest}} \right) \tag{\ref{eq:de_gravity}}
\end{equation}
The baseline rheological state of the universe ($De_{min}$) dictates the Viscoelastic Modulus Ratio ($\alpha$) scaled by the ratio of the Cosmic Unruh Energy ($\bar{a}_0$) to the Rest Mass ($m_e$). The galactic Bingham plug flow boundary ($\bar{a}_0$) is locked to the continuum yield limit ($\alpha$).

Substituting empirical values for the present epoch ($R_{h,0} \approx 4224 \text{ Mpc}$) yields the numerical hierarchy magnitude:
\begin{equation}
    De_{min} \approx 10^{-41} 
\end{equation}

\vspace{0.2 cm}

\textbf{10. The Kinematic Mass Limit (White Dwarf Collapse):} \\ 
In a stellar core, inward gravitational strain per nucleon evaluates as $E_G = G M_{bare} m_n / R$. The outward kinematic confinement strain evaluates as $E_e = \hbar c / d$, where $d$ represents the inter-particle spacing. The macroscopic radius scales by the nucleon count ($N = M_{bare}/m_n$) as $R = d (M_{bare}/m_n)^{1/3}$. % :::LLM_EDIT::: Labeled unthrottled mass capacity M_bare.

Equating the opposing strains and substituting the macroscopic radius yields:
\begin{equation}
    \frac{G M_{bare} m_n}{d \left( \frac{M_{bare}}{m_n} \right)^{1/3}} = \frac{\hbar c}{d} % :::LLM_EDIT::: Expressed strain balance in M_bare.
\end{equation}
The spatial coordinate ($d$) cancels, rendering the equilibrium state independent of physical radius:
\begin{equation}
    \frac{G M_{bare} m_n}{\left( \frac{M_{bare}}{m_n} \right)^{1/3}} = \hbar c % :::LLM_EDIT::: Preserved M_bare after d cancels.
\end{equation}
Isolating the bare mass ($M_{bare}$), substituting the Planck Mass ($M_P^2 = \hbar c / G$), and scaling by the yield compliance ($\alpha_{yield} = 1/\sqrt{2}$, Step 8) resolves the collapse limit: % :::LLM_EDIT::: Applied alpha_yield = 1/sqrt(2) compliance bound.
\begin{align}
    G M_{bare}^{2/3} m_n^{4/3} &= \hbar c \implies M_{bare} = \frac{M_P^3}{m_n^2} \\ % :::LLM_EDIT::: Isolated bare capacity M_bare.
    M &= \alpha_{yield} M_{bare} = \frac{M_P^3}{\sqrt{2}\, m_n^2} \approx 1.31 M_\odot \tag{\ref{eq:chandra_mass_limit}} % :::LLM_EDIT::: Projected observable collapse mass by 1/sqrt(2).
\end{align}
This identifies the Chandrasekhar mass limit. Mass exceeding this limit induces unhalted structural collapse, crushing the internal inter-particle spacing down to the solid metric yield limit ($L_{crit}$), initiating neutron degeneracy.

\vspace{0.2 cm}

\textbf{11. The Solid Mass Limit (Metric Fracture):} \\ 
When gravitational collapse restricts the continuum to its transverse relaxation limit ($L_{crit}$), the metric locks into a solid state. This rigidity evaluates as the ultimate structural floor, superseding classical Neutron Degeneracy Pressure. The core maintains a uniform density dictated by this confinement limit ($\rho_{nuc} = 3 m_n / 4\pi L_{crit}^3$). 

The core collapses into an event horizon when this density is forced across the glass boundary ($R = 2GM_{bare}/c^2$). Substituting this boundary radius into the spherical mass equation ($M_{bare} = \frac{4}{3}\pi R^3 \rho_{nuc}$) yields the bare cavitation capacity: % :::LLM_EDIT::: Labeled bare cavitation mass M_bare.
\begin{equation}
    M_{bare} = \frac{4}{3}\pi \left( \frac{2GM_{bare}}{c^2} \right)^3 \left( \frac{3 m_n}{4\pi L_{crit}^3} \right) % :::LLM_EDIT::: Substituted M_bare into spherical mass.
\end{equation}
Isolating $M_{bare}$, substituting the Planck Mass ($M_P^2 = \hbar c / G$), and scaling by the yield compliance ($\alpha_{yield} = 1/\sqrt{2}$, Step 8) resolves the astral limit: % :::LLM_EDIT::: Applied alpha_yield = 1/sqrt(2) compliance bound.
\begin{align}
    M_{bare}^2 &= \frac{c^6 L_{crit}^3}{8 G^3 m_n} = \frac{\alpha^3 \hbar^3 c^3}{8 G^3 m_n m_e^3} \\ % :::LLM_EDIT::: Solved bare squared capacity M_bare^2.
    M &= \alpha_{yield} M_{bare} = \frac{M_P^3}{4} \sqrt{ \frac{\alpha^3}{m_n m_e^3} } \approx 22.7 M_\odot \tag{\ref{eq:astral_mass_limit}} % :::LLM_EDIT::: Combined 1/sqrt(2) and 1/sqrt(8) into 1/4.
\end{align}
This identifies the structural ceiling of the metric, acting as the upper bound for the Tolman-Oppenheimer-Volkoff (TOV) mass limit. An astral core exceeding $\sim 22.7 M_\odot$ fractures the solid metric and collapses into a black hole. % :::LLM_EDIT::: Updated numerical bound to 22.7 M_sun.

\vspace{0.2 cm}

\textbf{12. The Glass Shell ($r_{glass}$):} \\ 
For a gravity well, compliance evaluates via the scalar potential ($\alpha_s^2 = 1 - \frac{2GM}{r c^2}$). The metric freezes at the glass transition ($De = 1 \implies \alpha_s = \alpha$). 

Substituting this limit isolates the spatial radius:
\begin{align}
    \alpha^2 &= 1 - \frac{2GM}{r_{glass} c^2} \\ 
    r_{glass} &= \frac{2GM}{c^2(1 - \alpha^2)}
\end{align}
Because $\alpha^2 \ll 1$, applying the binomial expansion isolates the boundary:
\begin{equation}
    r_{glass} \approx r_s (1 + \alpha^2) \tag{\ref{eq:r_glass}}
\end{equation}
The geometry transitions into an amorphous solid at a shell (scaled by $\alpha^2 \approx 5.3 \times 10^{-5}$) outside the event horizon ($r_s$). 
Incident mass translates into this rigid phase boundary, dictating viscoelastic fracture prior to coordinate crossing. Furthermore, the transverse structural lock-up halts volumetric collapse, mechanically prohibiting a Euclidean singularity ($r \to 0$).

\vspace{0.2 cm}

\textbf{13. Terminal Evaporation (The Bohr Radius Limit):} \\
The spatial thickness of the solid deformation zone ($\delta$) evaluates as the geometric gap between the glass shell and the event horizon ($R_s$):
\begin{align}
    \delta &= r_{glass} - R_s \\
    \delta &= \frac{R_s}{1 - \alpha^2} - R_s = R_s \left( \frac{\alpha^2}{1 - \alpha^2} \right)
\end{align}
The structural width of this zone ($\delta$) is bounded by the transverse relaxation limit ($r_e$). As a black hole evaporates, the horizon shrinks, narrowing the shell. The terminal limit evaluates when the shell thickness equals the fracture limit ($\delta = r_e$). 

Substituting the geometric identity ($r_e = \alpha^2 a_0$):
\begin{align}
    R_{s,crit} \left( \frac{\alpha^2}{1 - \alpha^2} \right) &= \alpha^2 a_0 \\
    R_{s,crit} &= a_0(1 - \alpha^2) \tag{\ref{eq:terminal_radius}}
\end{align}
At this coordinate, the black hole horizon shrinks to the Bohr Radius ($a_0$). The deformation zone reduces to the classical metric limit, triggering detonation.

\begin{tcolorbox}[recallbox, title={\ref{box:unified_deborah}. The Unified Cosmological Deborah Number}]
\begin{align}
    De &= \alpha \left( \frac{E}{m_e c^2} \right) \tag{\ref{eq:unified_deborah}} \\
    G^*(E) &= \mu_s \left[ \frac{De^2}{1 + De^2} \right] + i \mu_s \left[ \frac{De}{1 + De^2} \right] \tag{\ref{eq:deborah_modulus}} \\
    De &= \frac{\alpha}{\alpha_s} \tag{\ref{eq:de_ratio}} \\
    \alpha_s &= \frac{m_e c^2}{E} \tag{\ref{eq:angular_distribution}} \\
    \beta(\alpha_s) &= - \alpha_s \tag{\ref{eq:beta_function}} \\
    De \to 1 &\implies 2 E_{crit} = 140.04 \text{ MeV} \equiv m_{\pi^{\pm}} \tag{\ref{eq:de_glass}} \\
    L_{crit} &= r_e \approx 2.81 \text{ fm} \tag{\ref{eq:de_range}} \\ % :::LLM_EDIT::: Moved L_crit up to correctly match Step 5 position.
    \gamma_{yield} &= \frac{1}{\alpha} \approx 137 \tag{\ref{eq:gamma_yield}} \\
    Z &= \frac{1}{\alpha} \approx 137 \tag{\ref{eq:atomic_phase_limit}} \\
    De = 1 &\implies \alpha_s = \alpha \tag{\ref{eq:electro_strong_symmetry}} \\
    g^2 &= 4\pi\alpha \tag{\ref{eq:strong_charge_symmetry}} \\ % :::LLM_EDIT::: Added missing tag for the strong charge derivation.
    De_{min} &= \alpha (2\pi)^2 \left( \frac{E_{unruh}}{E_{rest}} \right) \tag{\ref{eq:de_gravity}} \\
    M &= \frac{M_P^3}{\sqrt{2}\, m_n^2} \approx 1.31 M_\odot \tag{\ref{eq:chandra_mass_limit}} \\ % :::LLM_EDIT::: Updated recall box with 1/sqrt(2) factor.
    M &= \frac{M_P^3}{4} \sqrt{ \frac{\alpha^3}{m_n m_e^3} } \approx 22.7 M_\odot \tag{\ref{eq:astral_mass_limit}} \\ % :::LLM_EDIT::: Updated recall box with 1/4 factor.
    r_{glass} &\approx r_s (1 + \alpha^2) \tag{\ref{eq:r_glass}} \\
    R_{s,crit} &= a_0(1 - \alpha^2) \tag{\ref{eq:terminal_radius}}
\end{align}
\end{tcolorbox}




% ----------------------------------------------------------------------
\section{Viscoelastic Phase Stiffening \& Time Dilation}
\label{sec:deriv_viscoelastic_stiffening}

The rheological stiffening of the continuum, the retardation of proper time, and the kinematic depletion of scalar admittance evaluate as conserved dynamic limits. 

\textbf{1. The Compliance Limit ($\alpha_s$):} \\
The total scalar admittance ($\alpha_s$) evaluates as the square root of the residual static capacity ratio ($P_{static} / P_c$). Because the continuous metric conserves total yield capacity via Bernoulli's principle ($P_c = P_{static} + P_{dyn} + P_{shear}$), substituting the kinematic definitions isolates the compliance limit: 
\begin{align}
    \alpha_s &= \sqrt{\frac{P_{static}}{P_c}} \\ 
    \alpha_s &= \sqrt{\frac{P_c - (P_{dyn} + P_{shear})}{P_c}} \\ 
    \alpha_s^2 &= 1 - \frac{P_{dyn} + P_{shear}}{P_c} \tag{\ref{eq:admittance_pressure_ratio}}
\end{align}
The compliance of the metric is depleted by both macroscopic kinematic velocity ($P_{dyn}$) and the distortional strain of localized rotational frequencies ($P_{shear}$). 

\textbf{2. The Yukawa-Bernoulli Identity:} \\
The total dynamic and shear pressure of the circulating metric evaluates as the interaction strength, equivalent to the squared Yukawa coupling constant ($g^2 \propto P_{dyn} + P_{shear}$). The acoustic yield pressure equates to the theoretical maximum coupling limit ($g_{max}^2 \propto P_c$). 

Substituting this normalization into the compliance equation derives the coupling identity:
\begin{equation}
    \alpha_s^2 = 1 - \frac{g^2}{g_{max}^2} \tag{\ref{eq:yukawa_identity}}
\end{equation}
The strong coupling constant ($\alpha_s$) evaluates as the residual structural compliance of the metric. 

\textbf{3. Relativistic Packing \& The Lorentz Pressure Limit:} \\
As a topological boundary translates through the continuous metric, macroscopic kinematic velocity forces the surrounding metric to undergo spatial compression. The effective local density of the $\boldsymbol{\tau}$-metric ($\rho_{eff}$) evaluates via the Lorentz factor ($\gamma$). 

Because $\gamma = 1/\alpha_s$, the geometric packing derives from the static pressure ratio:
\begin{align}
    \rho_{eff} &= \rho_{\tau} \gamma \tag{\ref{eq:relativistic_packing_admittance}} \\
    \gamma &= \sqrt{\frac{P_c}{P_c - (P_{dyn} + P_{shear})}} \tag{\ref{eq:lorentz_pressure}}
\end{align}
Time dilation evaluates as the physical degradation of available static capacity. As internal spin and macroscopic velocity consume the structural capacity, acoustic waves require more phase cycles to transit the coordinate space.

\textbf{4. The Kinematic Phase Lock (Thermodynamic Substitution):} \\
By the conservation of continuous phase capacity, the acoustic limit ($c$) dictates the geometric tradeoff between longitudinal translation and internal topological vorticity. The total kinetic energy of the coordinate cannot exceed the critical pressure ($P_c$).

Applying the admittance identity ($\gamma = 1/\alpha_s$) defines the Lorentz factor as the ratio of absolute capacity to residual capacity: 
\begin{equation}
    \gamma = \sqrt{\frac{P_c}{P_{static}}} \tag{\ref{eq:admittance_spin_identity}}
\end{equation}
Macroscopic strain evaluates as the continuous depletion of the internal static pressure required to sustain localized rotational phase.

\textbf{5. The Kinematic Yield Limit (Choked Flow):} \\ .
As the macroscopic kinematic drive accelerates toward the acoustic limit ($v_{\parallel} \to c$), the dynamic and shear pressures strike the yield ceiling ($(P_{dyn} + P_{shear}) \to P_c$). The normalized coupling maximizes ($g^2/g_{max}^2 \to 1$).

This saturation collapses the scalar admittance: 
\begin{align}
    \lim_{(P_{dyn} + P_{shear}) \to P_c} \alpha_s^2 &= 1 - 1 = 0 \\
    \lim_{v_{\parallel} \to c} \alpha_s &= 0 \implies \nabla \cdot \mathbf{v} = 0 \tag{\ref{eq:cavitation_collapse}}
\end{align}
This evaluates as choked flow. The spatial divergence of the metric evaluates to zero ($\nabla \cdot \mathbf{v} = 0$), forcing the coordinate into an incompressible solid state.

\vspace{0.2 cm}

\textbf{6. Acoustic Covariance (Phase Stiffening):} \\
Because the metric operates as an acoustic medium, the phase velocity of transverse shear waves is governed by the Newton-Laplace relation ($c = \sqrt{K_\tau/\rho_\tau}$). 

As the local effective density compacts under kinematic load ($\rho_{eff} = \rho_{\tau} / \alpha_s$), a proportional Phase Stiffening conserves the invariant wave limit. The effective bulk modulus ($K_{eff}$) must increase by the identical compliance fraction relative to the baseline stiffness ($K_\tau = \rho_{\tau} c^2$):
\begin{equation}
    K_{eff} = \frac{K_\tau}{\alpha_s} \label{eq:effective_bulk_modulus}
\end{equation}

Substituting the effective density and effective stiffness into the phase velocity equation cancels the compliance fractions:
\begin{align}
    c_{obs} &= \sqrt{\frac{K_{eff}}{\rho_{eff}}} \nonumber \\
    c_{obs} &= \sqrt{\frac{K_\tau / \alpha_s}{\rho_{\tau} / \alpha_s}} \nonumber \\
    c_{obs} &= \sqrt{\frac{K_\tau}{\rho_{\tau}}} \equiv c \label{eq:acoustic_covariance_proof}
\end{align}
The speed of light remains invariant across all reference frames because the structural stiffness of the metric increases in direct proportion to its relativistic density.

\begin{tcolorbox}[recallbox, title={\ref{box:viscoelastic_stiffening}. Viscoelastic Phase Stiffening \& Time Dilation}]
\begin{align}
    \alpha_s^2 &= 1 - \frac{P_{dyn} + P_{shear}}{P_c} \tag{\ref{eq:admittance_pressure_ratio}} \\
    \alpha_s^2 &= 1 - \frac{g^2}{g_{max}^2} \tag{\ref{eq:yukawa_identity}} \\
    \rho_{eff} &= \rho_{\tau} \gamma \tag{\ref{eq:relativistic_packing_admittance}} \\
    \gamma &= \sqrt{\frac{P_c}{P_{static}}} \tag{\ref{eq:lorentz_pressure}} \\
    \lim_{(P_{dyn}+P_{shear}) \to P_c} \alpha_s &= 0 \implies \nabla \cdot \mathbf{v} = 0 \tag{\ref{eq:cavitation_collapse}} 
\end{align}
\end{tcolorbox}


% ----------------------------------------------------------------------

\section{Kinematic Vacuum Impedance ($Z_{\tau}$)}
\label{sec:deriv_universal_impedance}

Inertial mass, optical density, quantum electrical resistance, and quantized spin evaluate as the directional limits of the continuum's resistance. Kinematic Vacuum Impedance ($Z_\tau$) represents the resistance to kinematic displacement. 

\vspace{0.2 cm}


\textbf{1. The Transverse Limit ($Y_{\tau(\perp)}$):} \\
Because mobility operates as the inverse of resistance, kinematic admittance ($Y$) is defined as the reciprocal of impedance:
\begin{equation}
    Y \equiv \frac{1}{Z}
\end{equation}

Admittance evaluates as the ratio of kinematic response to structural resistance. The kinematic response evaluates as the phase torque ($h$). Resistance evaluates as the squared metric charge ($e_\tau^2$), operating as transverse shear displacement.

The baseline transverse kinematic admittance ($Y_{\tau(\perp)}$) evaluates as the quotient: 
\begin{equation}
    Y_{\tau(\perp)} = \frac{h}{e_\tau^2} 
\end{equation}

Because electrical resistance maps to mechanical admittance scaled by a dimensional constant ($k_z$, derived in Section \ref{sec:deriv_em_lc_isomorphism}), the quantum of electrical resistance ($R_K$) evaluates as:
\begin{equation}
    R_K \equiv k_z Y_{\tau(\perp)} = k_z \left(\frac{h}{e_\tau^2}\right) \tag{\ref{eq:von_klitzing_impedance}}
\end{equation}
Since the constant is defined as the squared ratio of metric charge to SI charge ($k_z \equiv e_\tau^2 / e^2$), the parameters cancel, recovering the von Klitzing constant ($h/e^2$).


\vspace{0.2 cm}

\textbf{2. The Longitudinal Inertial Limit ($Z_{\tau(\parallel)}$):} \\
A transverse acoustic wave translating through a periodic topological boundary experiences continuous kinematic impedance. The group velocity ($v_g$) of the wave evaluates as the spatial derivative of the angular frequency ($\omega$) with respect to the wave vector ($k$): 
\begin{equation}
    v_g = \frac{d\omega}{dk} 
\end{equation}

Substituting the phase energy relation ($E = \hbar \omega \implies \omega = E/\hbar$):
\begin{equation}
    v_g = \frac{1}{\hbar} \frac{dE}{dk} 
\end{equation}

The effective acceleration ($a$) evaluates as the time derivative of this group velocity ($a = \frac{dv_g}{dt}$). Applying the chain rule transitions the derivative from temporal to spatial coordinates:
\begin{equation}
    a = \left( \frac{dv_g}{dk} \right) \frac{dk}{dt} 
\end{equation}

Substituting the spatial derivative of the group velocity ($\frac{dv_g}{dk} = \frac{1}{\hbar} \frac{d^2 E}{dk^2}$):
\begin{equation}
    a = \left( \frac{1}{\hbar} \frac{d^2 E}{dk^2} \right) \frac{dk}{dt} 
\end{equation}

External force ($F$) applied to the metric evaluates as the temporal rate of change of the structural phase momentum ($p = \hbar k$):
\begin{equation}
    F = \frac{d(\hbar k)}{dt} = \hbar \frac{dk}{dt} \implies \frac{dk}{dt} = \frac{F}{\hbar} 
\end{equation}

Substituting this momentum identity into the acceleration equation:
\begin{equation}
    a = F \left( \frac{1}{\hbar^2} \frac{d^2 E}{dk^2} \right) 
\end{equation}

Newtonian kinematics equates force to mass and acceleration ($F = m^* a \implies a = F/m^*$). Setting the kinematic equations as equal, isolates the apparent mass:
\begin{equation}
    \frac{1}{m^*} = \frac{1}{\hbar^2} \frac{d^2 E}{dk^2} 
\end{equation}

Inverting both sides extracts the longitudinal directional impedance:
\begin{equation}
    Z_{\tau(\parallel)} \equiv m^* = \hbar^2 \left( \frac{d^2 E}{dk^2} \right)^{-1} \tag{\ref{eq:effective_mass_impedance}}
\end{equation}
A particle possesses apparent inertial mass ($m^*$) because the boundary resists longitudinal displacement. 

\vspace{0.2 cm}

\textbf{3. The Transverse Refractive Limit ($Z_{\tau(\perp)}$):} \\
The refractive index ($n$) evaluates the kinematic impedance of the static pressure field ($P_{static}$), dictating the path of an acoustic shear wave. Evaluating the null transit path ($ds^2 = 0$) through the compliance limits:
\begin{align}
    -\left(\frac{P_{static}}{P_c}\right) c^2 dt^2 + \left(\frac{P_c}{P_{static}}\right) dr^2 &= 0 \\
    \left(\frac{P_c}{P_{static}}\right) dr^2 &= \left(\frac{P_{static}}{P_c}\right) c^2 dt^2 \\
    \frac{dr^2}{dt^2} &= c^2 \left(\frac{P_{static}}{P_c}\right)^2 
\end{align}

Taking the square root yields the local wave velocity ($c'$):
\begin{equation}
    c' = \frac{dr}{dt} = c \left(\frac{P_{static}}{P_c}\right) 
\end{equation}

The refractive impedance ($n(r)$) to transverse waves evaluates as the ratio of the acoustic limit ($c$) to the local velocity ($c'$):
\begin{align}
    Z_{\tau(\perp)} \propto n(r) &= \frac{c}{c'} \\ % :::LLM_EDIT::: Removed \nonumber per Protocol Rule 1.
    n(r) &= \frac{P_c}{P_{static}} \tag{\ref{eq:refractive_impedance}}
\end{align}

The static pressure deficit surrounding a centralized mass $M$ evaluates via the Bernoulli consumption of the static capacity, which equates to the classical gravitational limit: 
\begin{equation}
    \frac{P_{static}}{P_c} \approx 1 - \frac{2GM}{rc^2}
\end{equation}

Substituting this ratio into the refractive index yields the spatial curvature:
\begin{equation}
    n(r) \approx \left( 1 - \frac{2GM}{rc^2} \right)^{-1} 
\end{equation}

Applying the binomial expansion $(1-x)^{-1} \approx 1+x$:
\begin{equation}
    n(r) \approx 1 + \frac{2GM}{c^2 r} 
\end{equation}


The energy ($E(r)$) of a transverse wave refracting through this pressure gradient evaluates against its unperturbed effective mass ($m = E_0/c^2$) as $E(r) = \frac{mc^2}{n(r)}$. Applying the inverse binomial expansion $n(r)^{-1} \approx 1 - \frac{2GM}{c^2 r}$:
\begin{equation}
    E(r) \approx mc^2 \left( 1 - \frac{2GM}{c^2 r} \right) = mc^2 - \frac{2GMm}{r} 
\end{equation}

The effective force exerted upon the propagating wave evaluates as the negative spatial gradient of this energy state ($\mathbf{F} = -\nabla E$):
\begin{align}
    \mathbf{F}_g &= -\frac{d}{dr} \left( mc^2 - \frac{2GMm}{r} \right) \mathbf{\hat{r}} \\ 
    \mathbf{F}_g &= -\frac{2GMm}{r^2} \mathbf{\hat{r}} \label{eq:relativistic_limit}
\end{align}
Apparent gravitational trapping, effective rest mass, and electrical resistance evaluate as isomorphic directional impedances.

\vspace{0.2 cm}

\textbf{4. The Rotational Phase Limit ($\hbar$ \& Gyroscopic Inertia):} \\
Rotational kinematic impedance ($Z_{\tau(\circlearrowleft)}$) evaluates the physical resistance to topological vorticity ($\nabla \times \mathbf{v}_{\psi}$). In classical rotational mechanics, impedance evaluates as the ratio of applied rotational energy (torque, $E$) to the resulting angular phase velocity ($\omega$).
\begin{equation}
    Z_{\tau(\circlearrowleft)} = \frac{E}{\omega}
\end{equation}

Substituting the phase relation ($E = \hbar\omega$) isolates the baseline rotational resistance: 
\begin{equation}
    Z_{\tau(\circlearrowleft)} = \frac{\hbar\omega}{\omega} = \hbar \tag{\ref{eq:planck_impedance}}
\end{equation}
Planck's reduced constant ($\hbar$) evaluates as Rotational Impedance. The equation $E = \hbar\omega$ evaluates as the classical definition of mechanical impedance ($Z = E/\omega$). 

Furthermore, as the internal circulation of a topological boundary approaches the acoustic limit ($|\nabla \times \mathbf{v}_{\psi}| \to c$), the dynamic and shear pressures exhaust the thermodynamic budget ($P_{static} \to 0$). Because kinematic impedance evaluates as the inverse of the available static pressure capacity, the rotational impedance diverges to infinity: 
\begin{equation}
    \lim_{P_{static} \to 0} Z_{\tau(\circlearrowleft)} \to \infty \tag{\ref{eq:infinite_rotational_impedance}}
\end{equation}
This infinite rotational impedance evaluates as the quantization of intrinsic Spin (e.g., $\hbar/2$). Because the metric presents infinite resistance to additional vorticity at the cavitation boundary, particles evaluate as rigid gyroscopes. They cannot be spun up or slowed down through classical interaction, evaluating phenomenologically as Chiral Symmetry Breaking and Gyroscopic Stiffness.

In natural units, evaluating $\hbar = 1$ normalizes the baseline rotational impedance, against which all topological spin is quantized.
\vspace{0.2 cm}

\textbf{5. Total Kinematic Vacuum Admittance (Parallel Impedance):} \\
Orthogonal flow states driven by a shared pressure budget evaluate as parallel pathways. Because the metric possesses three orthogonal degrees of freedom (longitudinal translation, transverse shear, and rotational vorticity), the total admittance ($Y_\tau$) evaluates as the sum of the directional admittances ($Y_i = 1/Z_i$): 
\begin{equation}
    \frac{1}{Z_\tau} = \frac{1}{Z_{\tau(\parallel)}} + \frac{1}{Z_{\tau(\perp)}} + \frac{1}{Z_{\tau(\circlearrowleft)}} \label{eq:parallel_impedance}
\end{equation}
The total capacity of the vacuum to sustain kinematic displacement evaluates as the parallel sum of its orthogonal phases. As any kinematic state approaches structural yield ($Z_i \to \infty$), its dimensional phase capacity collapses ($1/Z_i \to 0$), choking the total kinematic admittance of the coordinate space.

\begin{tcolorbox}[recallbox, title={\ref{box:universal_impedance}. Universal Kinematic Vacuum Impedance ($Z_{\tau}$)}]
\begin{align}
    R_K &= k_z Y_{\tau(\perp)} = k_z \left(\frac{h}{e_\tau^2}\right) \tag{\ref{eq:von_klitzing_impedance}} \\
    Z_{\tau(\parallel)} &= \hbar^2 \left( \frac{d^2 E}{dk^2} \right)^{-1} \equiv m^* \tag{\ref{eq:effective_mass_impedance}} \\
    Z_{\tau(\perp)} &\propto n(r) = \frac{P_c}{P_{static}} \approx \left( 1 - \frac{2GM}{rc^2} \right)^{-1} \tag{\ref{eq:refractive_impedance}} \\
    Z_{\tau(\circlearrowleft)} &= \frac{E}{\omega} \equiv \hbar \tag{\ref{eq:planck_impedance}} \\
    \lim_{P_{static} \to 0} Z_{\tau(\circlearrowleft)} &\to \infty \implies \text{Spin Quantization } (\hbar/2) \tag{\ref{eq:infinite_rotational_impedance}}
\end{align}
\end{tcolorbox}

\vspace{0.2 cm}

\part*{PHASE III: TOPOLOGICAL MASS GENERATION}

\vspace{0.2 cm}

% ----------------------------------------------------------------------
\section{The Kinematic Origin of Inertia (Hydrodynamic Rest Mass)}
\label{sec:deriv_hydrodynamic_mass}

The resistance of a boundary to spatial acceleration (inertia) evaluates as metric drag.  
\textbf{1. The Massless Void ($m_{intrinsic}$):} \\
A topological boundary evaluates as a cavitated discontinuity generated by the circulating metric. Because the coordinate volume ($V_{r_c}$) is empty ($\rho_{void} = 0$), the intrinsic mass of the boundary evaluates to zero: 
\begin{equation}
    m_{intrinsic} = 0 \label{eq:intrinsic_mass}
\end{equation}
A void possesses no internal mechanism to resist spatial acceleration.

\textbf{2. Hydrodynamic Added Mass ($m_{added}$):} \\
In continuum mechanics, when a cavity accelerates through a medium of uniform baseline density ($\rho_{\tau}$), the applied force evaluates against the inertia of the surrounding displacement. The physical entrainment evaluates as volumetric added mass ($V_{added}$): 
\begin{equation}
    m_{added} = \rho_{\tau} V_{added} \tag{\ref{eq:added_mass}}
\end{equation}

\textbf{3. Effective Buoyancy Volume ($V_{eff}$):} \\
The rest mass of the boundary ($m_\tau$) evaluates as the sum of its intrinsic mass and the hydrodynamic added mass ($m_\tau = m_{intrinsic} + m_{added}$). Substituting the massless void limit zeroes the intrinsic term. 

Because the boundary evaluates as a high-frequency acoustic stator ($\Omega$), the rotational frequency generates transverse kinematic drag. The displacement of the metric required to satisfy this static pressure deficit evaluates as the effective buoyancy volume ($V_{eff}$). Inertial rest mass evaluates as the displacement of this volume: 
\begin{equation}
    m_\tau = \rho_{\tau} V_{eff} \label{eq:kinematic_mass}
\end{equation}

\textbf{4. Kinematic Strain Equivalence ($V_\tau$):} \\
The Archimedean buoyancy volume ($V_{eff}$) operates as a fraction of the total volumetric strain ($V_\tau$).  

Because the coordinate boundary acts as a specific geometric circulation, the metric displacement scales via its topological form factor ($C_m$), yielding $V_{eff} = C_m V_\tau$. Substituting this into the rest mass isolates the continuum definition of inertia: 
\begin{equation}
    m_\tau = C_m  \rho_\tau V_\tau \tag{\ref{eq:strain_volume_equivalence}}
\end{equation}

\begin{tcolorbox}[recallbox, title={\ref{box:hydrodynamic_mass}. The Kinematic Origin of Inertia}]
\begin{align}
    m_{added} &= \rho_{\tau} V_{added} \tag{\ref{eq:added_mass}} \\
    m_\tau &= \rho_{\tau} V_{eff} \tag{\ref{eq:kinematic_mass}} \\
    m_\tau &= C_m \rho_\tau V_\tau \tag{\ref{eq:strain_volume_equivalence}}
\end{align}
\end{tcolorbox}


% ----------------------------------------------------------------------
\section{Navier-Stokes Regularization (Topological Pinch-Off)}
\label{sec:deriv_ns_reg}

The finite-time singularity of the 3D Navier-Stokes equations becomes regularized when accounting for the acoustic limit ($c$). The singularity is prevented by a viscosity threshold. 

\textbf{1. The Kinematic Boundary Condition:} \\
In a turbulent cascade, enstrophic divergence dictates that the rotational velocity ($v$) of a shrinking eddy ($r \to 0$) increases continuously. The metric imposes a kinematic boundary condition: internal circulation cannot exceed the acoustic phase limit ($v \le c$). 

When the velocity reaches the limit ($v = c$), the radius of the vortex locks to its Compton radius ($L = \hbar / m c$). 

\textbf{2. Dynamic Viscosity ($\mu$):} \\
Dynamic viscosity ($\mu$) defines as the ratio of Shear Stress ($\tau_{shear}$) to the Velocity Gradient ($\nabla v$).

The maximum shear stress the continuum can support before failure maps to the Bernoulli dynamic pressure limit of the background ($\tau_{shear} = \frac{1}{2} \rho_{\tau} c^2$). 
The velocity gradient of the vortex acts as the kinematic limit ($c$) divided by the constrained boundary radius ($\nabla v = c / L$). Substituting the Compton length ($L = \hbar/mc$) yields:
\begin{equation}
    \nabla v = \frac{c}{\left(\frac{\hbar}{mc}\right)} = \frac{mc^2}{\hbar} 
\end{equation}

Dividing the shear stress by the velocity gradient isolates the dynamic viscosity of the metric:
\begin{align}
    \mu &= \frac{\tau_{shear}}{\nabla v}  \\
    \mu &= \frac{\frac{1}{2} \rho_{\tau} c^2}{\left( \frac{mc^2}{\hbar} \right)}  \\
    \mu &= \frac{1}{2} \frac{\rho_{\tau} \hbar}{m} \tag{\ref{eq:vacuum_viscosity}}
\end{align}

\textbf{3. Prandtl Mixing Length Equivalence:} \\
Because the Compton radius establishes the boundary of the acoustic cavity ($L = \hbar/mc$), rearranging the spatial relation yields $\hbar/m = cL$. Substituting this identity back into the viscosity isolates the origin of mass: 
\begin{equation}
    \mu = \frac{1}{2} \rho_{\tau} c L \tag{\ref{eq:prandtl_equivalence}}
\end{equation}

The metric models turbulent eddy viscosity via the kinematic mixing length ($\mu_t \sim \rho v L$) becoming a shear-thickening medium. As the internal circulation accelerates toward the kinematic limit ($c$) at the boundary length ($L$), the viscosity diverges to a maximum. 

\textbf{4. The Critical Reynolds Number ($Re_c$):} \\
The Critical Reynolds Number ($Re_c$) governing topological pinch-off is defined by the ratio of inertial forces to viscous forces:
\begin{equation}
    Re_c = \frac{\rho_{\tau} v L}{\mu} 
\end{equation}

Substituting the kinematic limit ($v=c$) and the Prandtl viscosity ($\mu = \frac{1}{2} \rho_{\tau} c L$) resolves the transition state:
\begin{align}
    Re_c &= \frac{\rho_{\tau} c L}{\frac{1}{2} \rho_{\tau} c L}  \\
    Re_c &= 2 \tag{\ref{eq:critical_reynolds}}
\end{align}

As a vortex approaches the $v=c$ limit, the Reynolds number converges to $2$. Because $Re \le 10$ defines creeping flow, viscous dampening prevents turbulent cascades. 

The finite-time singularity becomes bounded. The shear-thickening viscosity induces the cascade to pinch off into a stable, quantized particle. 

\vspace{0.2 cm}

\textbf{5. The Dynamic Elasticity Number ($El$) \& Topological Fracture:} \\
To establish the stability of the metric across kinematic states, the elastic stress ($De$) scales against the inertial flow ($Re$) via the Elasticity Number ($El = De/Re$). 

To construct a dynamic equation of state, both dimensionless numbers expand into their base variables. By substituting the dynamic viscosity ($\mu = \frac{1}{2} \rho_\tau c L$) into the Reynolds equation, the circulation acts as a function of velocity ($v$): 
\begin{align}
    Re &= \frac{\rho_\tau v L}{\frac{1}{2} \rho_\tau c L} \\ 
    Re &= \frac{2v}{c}
\end{align}
The Reynolds Number operates as twice the Mach fraction. 

Substituting the Reynolds Number and the Transverse Cosmological Deborah Number ($De = \alpha (E \cdot l_p) / \hbar c$) into the Elasticity quotient isolates the unified state equation:
\begin{align}
    El &= \frac{\left( \frac{\alpha (E \cdot l_p)}{\hbar c} \right)}{\left( \frac{2v}{c} \right)} \\ 
    El &= \frac{\alpha (E \cdot l_p)}{2 \hbar v} \tag{\ref{eq:dynamic_elasticity}}
\end{align}

This fraction dictates that elastic dominance is proportional to the applied transverse interaction energy ($E$) and inversely proportional to the flow velocity ($v$). 

As the velocity reaches the acoustic limit ($v \to c$), the denominator locks at $2\hbar c$. Because the numerator ($E \cdot l_p$) carries the mechanical dimensions of an applied Bending Moment ($N \cdot m^2$) and the denominator scales the baseline Flexural Rigidity of the vacuum ($\hbar c$), the Elasticity Number directly translates to the Euler-Bernoulli beam fraction: 
\begin{equation}
    El = \frac{\text{Applied Transverse Bending Moment}}{\text{Metric Flexural Rigidity}} \nonumber
\end{equation}

Because inertial flow is constrained by the phase limit ($c$), any accumulation of applied energy ($E$) drives this spatial bending ratio to diverge ($El \to \infty$). The metric cannot dissipate the applied torque via localized fluid flow. The elastic tension exceeds the yield limit, fracturing into a cavitation. The cascade pinches off, establishing the discrete kinematic volume of a quantized particle.

\begin{tcolorbox}[recallbox, title={\ref{box:ns_reg}. Navier-Stokes Regularization}]
\begin{align}
    \mu &= \frac{1}{2} \frac{\rho_{\tau} \hbar}{m} \tag{\ref{eq:vacuum_viscosity}} \\
    \mu &= \frac{1}{2} \rho_{\tau} c L \tag{\ref{eq:prandtl_equivalence}} \\
    Re_c &= \frac{\rho_{\tau} c L}{\mu} = 2 \tag{\ref{eq:critical_reynolds}} \\
    El &= \frac{\alpha (E \cdot l_p)}{2 \hbar v} \tag{\ref{eq:dynamic_elasticity}}
\end{align}
\end{tcolorbox}

% ----------------------------------------------------------------------
\section{Origin of Inertia (The Pressure Gradient \& Strain Compactness)}
\label{sec:deriv_fluid_inertia}

Apparent mass and kinematic inertia resolve as the scalar pressure deficit generated by the circulation traversing a rheological gradient. 

\textbf{1. Relativistic Strain Compactness (The Density-Pressure Identity):} \\
Because the metric operates with a volumetric density ($\rho_{\tau}$), accelerating the circulation velocity ($v_\phi$) subjects the boundary to kinematic packing. 

Substituting the Bernoulli compliance limit ($\alpha_s = \sqrt{P_{static}/P_c}$) into the Lorentz packing factor ($\rho_{eff} = \rho_{\tau} / \alpha_s$) isolates the continuum equation of state governing metric density: 
\begin{equation}
    \rho_{eff}(r) = \rho_{\tau} \left( \frac{P_{static}(r)}{P_c} \right)^{-1/2} = \rho_{\tau} \left( 1 - \frac{P_{dyn}(r) + P_{shear}(r)}{P_c} \right)^{-1/2} \tag{\ref{eq:fluid_inertia_gradient}}
\end{equation}
Effective density dictates the inverse of available static pressure. 

Substituting the kinematic definitions ($P_{dyn} = \frac{1}{2}\rho_{\tau} v^2$ and $P_c = \frac{1}{2}\rho_{\tau} c^2$) into an un-sheared boundary ($P_{shear} = 0$) recovers the Lorentz packing factor ($\gamma$): 
\begin{align}
    \rho_{eff}(v) &= \rho_{\tau} \left( 1 - \frac{\frac{1}{2}\rho_{\tau} v^2}{\frac{1}{2}\rho_{\tau} c^2} \right)^{-1/2} \\
    \rho_{eff}(v) &= \rho_{\tau} \left( 1 - \frac{v^2}{c^2} \right)^{-1/2} = \rho_{\tau} \gamma \label{eq:density_lorentz_recovery}
\end{align}
As the combined dynamic load ($P_{dyn}$) and distortional shear ($P_{shear}$) consumes the metric capacity ($P_c$), effective density compacts ($\rho_{eff}$). 


\textbf{2. The Isotropic Tension Identity ($\mathbb{T}_{local}$):} \\
Because the metric functions as an elastic continuum under geometric tautness, its baseline static pressure equates to an isotropic volumetric tension ($\mathbb{T}_c = -P_c$). Dimensionally, this tension translates to a continuous Cauchy stress ($\text{N/m}^2$), equivalent to localized energy density ($\text{J/m}^3$). 

Applying the Bernoulli Conservation Law ($P_{static} = P_c - (P_{dyn} + P_{shear})$) translates the pressure deficit into geometric tension:
\begin{align}
    \mathbb{T}_{local} &= - P_{static} \\ 
    \mathbb{T}_{local} &= - (P_c - (P_{dyn} + P_{shear})) \\ 
    \mathbb{T}_{local} &= -P_c + P_{dyn} + P_{shear} \\ 
    \mathbb{T}_{local} &= \mathbb{T}_c + P_{dyn} + P_{shear} \tag{\ref{eq:tension_identity}}
\end{align}
The internal circulation of the boundary creates a dynamic load. To maintain equilibrium, the spatial metric tightens. The total static tension of the particle ($\mathbb{T}_{local}$) acts as the baseline tension of the metric ($\mathbb{T}_c$) plus the dynamic load ($P_{dyn}$). The coordinate space pulls inward, causing the surrounding metric to compact into higher effective density. 


\textbf{3. The Geometric Equation of State (Pure Tension):} \\
To isolate the density of the metric, the dynamic load ($P_{dyn} + P_{shear} = \mathbb{T}_{local} - \mathbb{T}_c$) is substituted out of the Density-Pressure Identity. Because baseline tension is negative pressure ($\mathbb{T}_c = -P_c$), the substitution resolves as:
\begin{align}
    \rho_{eff}(r) &= \rho_{\tau} \left( 1 - \frac{\mathbb{T}_{local}(r) - \mathbb{T}_c}{-\mathbb{T}_c} \right)^{-1/2} \nonumber \\ 
    \rho_{eff}(r) &= \rho_{\tau} \left( 1 + \left( \frac{\mathbb{T}_{local}(r)}{\mathbb{T}_c} - 1 \right) \right)^{-1/2} \nonumber \\ 
    \rho_{eff}(r) &= \rho_{\tau} \left( \frac{\mathbb{T}_{local}(r)}{\mathbb{T}_c} \right)^{-1/2} \tag{\ref{eq:tension_density}}
\end{align}
The effective density resolves strictly as the inverse-root ratio of local tension to baseline tension.

\textbf{4. The Compliant Halo ($r \gg r_c$):} \\
While the core boundary acts as an acoustic horizon ($r_c$), the kinematic halo surrounding the defect functions as compliant. At macroscopic radii ($r \gg r_c$), circulation remains sub-luminal ($v_\phi \ll c$), and the dynamic load consumes a minority fraction of the yield capacity ($(P_{dyn} + P_{shear}) \ll P_c$). 

Expanding the root of the Density-Pressure equation via the binomial approximation ($(1-x)^{-1/2} \approx 1 + x/2$) reveals the origin of observable mass across this envelope: 
\begin{equation}
    \rho_{eff}(r) \approx \rho_{\tau} + \rho_{\tau} \left( \frac{P_{dyn}(r) + P_{shear}(r)}{2P_c} \right) \tag{\ref{eq:compliant_inertia}}
\end{equation}
The observable mass of a particle resolves as the baseline density ($\rho_{\tau}$) plus the excess generated by dynamic tension compressing the surrounding coordinate space. 

\vspace{0.2 cm}

% ----------------------------------------------------------------------
\textbf{5. The Cavitation Yield Limit ($r_c$):} \\ 
Treating the core as a dimensionless point ($r \to 0$) forces circulation to infinity ($v_\phi \to \infty$), requiring a dynamic tension exceeding metric capacity ($(P_{dyn} + P_{shear}) > P_c$). Substituting this into the Density-Pressure Identity yields an imaginary state:
\begin{equation}
    \rho_{eff}(0) = \rho_{\tau} (1 - \infty)^{-1/2} \implies \mathbb{I} 
\end{equation}

Because the continuum mandates real scalar limits ($\mathbb{R}$), internal circulation is bounded by the acoustic limit ($v_\phi \to c$). The dynamic load equals the total metric capacity ($(P_{dyn} + P_{shear}) \to P_c$), and the localized tension drops to zero ($\mathbb{T}_{local} \to 0$). The effective density diverges to a real infinity: 
\begin{equation}
    \lim_{\mathbb{T}_{local} \to 0} \rho_{eff}(r_c) \to \infty \label{eq:infinite_core_inertia}
\end{equation}
The spatial coordinate contracts under tension until it yields. The metric cavitates at the finite spatial boundary of the Compton radius ($r_c$), stabilizing the defect before crossing into an imaginary state. The minimum of this cavitation boundary isolates the fundamental relaxation length of the continuum ($r_c \ge l_p$).

\begin{tcolorbox}[recallbox, title={\ref{box:fluid_inertia}. Origin of Inertia (Strain Compactness \& Tension)}]
\begin{align}
    \rho_{eff}(r) &= \rho_{\tau} \left( 1 - \frac{P_{dyn}(r) + P_{shear}(r)}{P_c} \right)^{-1/2} \tag{\ref{eq:fluid_inertia_gradient}} \\
    \rho_{eff}(v) &= \rho_{\tau} \left( 1 - \frac{v^2}{c^2} \right)^{-1/2} = \rho_{\tau} \gamma \tag{\ref{eq:density_lorentz_recovery}} \\
    \mathbb{T}_{local} &= \mathbb{T}_c + P_{dyn} + P_{shear} \tag{\ref{eq:tension_identity}} \\
    \rho_{eff}(r) &= \rho_{\tau} \left( \frac{\mathbb{T}_{local}(r)}{\mathbb{T}_c} \right)^{-1/2} \tag{\ref{eq:tension_density}} \\
    r \gg r_c &\implies \rho_{eff}(r) \approx \rho_{\tau} + \rho_{\tau} \left( \frac{P_{dyn}(r) + P_{shear}(r)}{2P_c} \right) \tag{\ref{eq:compliant_inertia}} \\
    \lim_{\mathbb{T}_{local} \to 0} \rho_{eff}(r_c) &\to \infty \tag{\ref{eq:infinite_core_inertia}}
\end{align}
\end{tcolorbox}

% ----------------------------------------------------------------------
\section{Volumetric Approximations}
\label{sec:deriv_volumetric_inverse_cube}

Because matter evaluates as a cavitated boundary, its physical bounds reconcile the Euclidean footprint of the coordinate void ($V_{r_c}$) with the Riemannian metric strain displaced by its total kinematic volume ($V_\tau$). 

\vspace{0.2 cm}

\textbf{1. The Kinematic Strain Volume ($V_\tau$):} \\ 
Equating the derived origin of inertia ($m_\tau = C_m \rho_{\tau} V_\tau$) to the Compton boundary ($m_\tau = \hbar / r_c c$) fixes the effective strain volume ($V_\tau$): 
\begin{align}
    V_\tau &= \frac{m_\tau}{C_m\rho_\tau} = \frac{1}{C_m\rho_\tau} \left( \frac{\hbar}{r_c c} \right) \nonumber \\
    V_\tau &= \frac{2\pi l_p^4}{C_m r_c} \tag{\ref{eq:strain_volume_rc}}
\end{align}

\textbf{2. The Cavitation Volume ($V_{r_c}$):} \\
The coordinate volume ($V_{r_c}$) of a topological vortex is modeled as an isotropic 3D Euclidean bounding sphere parameterized by the Compton radius ($r_c = \hbar/m_\tau c$): 
\begin{equation}
    V_{r_c} = \frac{4}{3}\pi r_c^3 \tag{\ref{eq:coordinate_volume}}
\end{equation}

\textbf{3. The Inverse-Cube Substitution:} \\
To determine the geometric relationship between the bounding sphere ($V_{r_c}$) and the internal kinematic strain, the derived strain volume ($V_\tau$) substitutes into the product:
\begin{align}
    V_{r_c} \cdot V_\tau^3 &= \left( \frac{4}{3}\pi r_c^3 \right) \left( \frac{2\pi l_p^4}{C_m r_c} \right)^3 \nonumber \\
    V_{r_c} \cdot V_\tau^3 &= \left( \frac{4}{3}\pi r_c^3 \right) \left( \frac{8\pi^3 l_p^{12}}{C_m^3 r_c^3} \right) \nonumber \\
    V_{r_c} \cdot V_\tau^3 &= \frac{32\pi^4}{3C_m^3} l_p^{12} \tag{\ref{eq:volumetric_inverse_cube}} 
\end{align}

\textbf{4. Volumetric Equilibrium:} \\
Because Planck's reduced constant ($\hbar$), the baseline metric density ($\rho_{\tau}$), and the kinematic acoustic limit ($c$) project as metric invariants, the right side of the equation resolves to a constant constrained by the topological form factor ($C_m$). 

Inertia ($m_\tau$) cancels from the boundary condition. The coordinate volume ($V_{r_c}$) resolves proportionally to the inverse-cube of its total strain ($V_\tau$). As the radius of a boundary compresses, the total strain displaced into the surrounding metric dilates by a cubic ratio. 

\begin{tcolorbox}[recallbox, title={\ref{box:volumetric_inverse_cube}. The Volumetric Inverse-Cube Identity}]
\begin{align}
    V_\tau &= \frac{2\pi l_p^4}{C_m r_c} \tag{\ref{eq:strain_volume_rc}} \\ % :::LLM_EDIT::: Added V_tau derivation to recall box.
    V_{r_c} &= \frac{4}{3}\pi r_c^3 \tag{\ref{eq:coordinate_volume}} \\
    V_{r_c} \cdot V_\tau^3 &= \frac{32\pi^4}{3C_m^3} l_p^{12} \tag{\ref{eq:volumetric_inverse_cube}} 
\end{align}
\end{tcolorbox}

% ----------------------------------------------------------------------
\section{The Compliance-Modified Navier-Stokes Equation}
\label{sec:deriv_admittance_ns}

To determine the acceleration of a topological boundary, the momentum of the metric scales against the compliance and the modulus ratio ($\alpha$). % :::LLM_EDIT::: Cured "evaluate" repetition (determine, scales).
\vspace{0.2 cm}

\textbf{1. The Nested Compliance ($\alpha_g, \alpha_a, \alpha_s$):} \\
While the unperturbed metric is bounded by the critical pressure ($P_c$), a frame operating within a gravitational field is bounded by the ambient environmental pressure ($P_{ambient} = P_c - \Delta P_{macro}$). Because the capacity is pre-strained by the environment, $P_{ambient}$ serves as the available budget for kinematics ($P_{ambient} = P_{static} + P_{dyn} + P_{shear}$). 

This nested budget defines the compliance ratios. The Total Scalar Admittance ($\alpha_s$) resolves as the multiplicative product of all sequential geometric depletion factors: % :::LLM_EDIT::: Cured "evaluate" repetition (resolves).
\begin{align}
    \alpha_g &= \sqrt{\frac{P_{ambient}}{P_c}} \quad \text{(Gravitational Compliance)} \\
    \alpha_a &= \sqrt{\frac{P_{static}}{P_{ambient}}} \quad \text{(Ambient Kinematic Compliance)} \\
    \alpha_s &= \sqrt{\frac{P_{static}}{P_c}} = \prod_i \alpha_i = \alpha_g \alpha_a \dots \quad \text{(Total Scalar Admittance)}
\end{align}

The acceleration of a coordinate volume scales against its effective density ($\rho_{eff}$). Because the baseline density operates as an invariant property ($\rho_a = \rho_\tau$), the effective density of the boundary scales against the Total Scalar Admittance ($\alpha_s$): % :::LLM_EDIT::: Cured "evaluate" repetition (scales, operates).
\begin{equation}
    \rho_{eff} = \frac{\rho_a}{\alpha_s} = \frac{\rho_{\tau}}{\alpha_s} \tag{\ref{eq:ns_density_substitution}}
\end{equation}
The Navier-Stokes equation governs mechanics in the frame, scaling effective inertial density via this total product ratio. % :::LLM_EDIT::: Cured "evaluate" repetition (governs).

\textbf{2. The Reynolds Limit (Hadronization \& Decay):} \\
The Reynolds number ($Re$) defines the transition from laminar flow to turbulence via the ratio of inertial density to internal viscous shear ($Re = \rho v L / \mu$). % :::LLM_EDIT::: Cured "evaluate" repetition (defines).

To calculate the Reynolds Number ($Re_q$) of a boundary, substitute the effective density ($\rho_{eff} = \rho_\tau/\alpha_s$), the kinematic circulation limit of the core ($v = c_a$), and the dynamic viscosity of the continuum derived via the Prandtl boundary limit ($\mu = \frac{1}{2} \rho_\tau c_a L$): % :::LLM_EDIT::: Cured "evaluate" repetition (calculate).
\begin{align}
    Re_q &= \frac{\rho_{eff} c_a L}{\mu} \\ 
    Re_q &= \frac{\left( \frac{\rho_\tau}{\alpha_s} \right) c_a L}{\frac{1}{2} \rho_\tau c_a L} 
\end{align}

The ambient continuum constants ($\rho_\tau, c_a, L$) cancel from the numerator and denominator, isolating the transition limit of the metric:
\begin{equation}
    Re_q = \frac{2}{\alpha_s} \tag{\ref{eq:quantum_reynolds}} 
\end{equation}
At rest, the compliance of a defect reduces to its static rest state ($\alpha_{s, rest} = \alpha_g \sqrt{1 - P_{shear}/P_{ambient}}$). Because a particle possesses rest energy ($P_{shear} > 0$), its rest compliance remains less than $1$. Therefore, the ground state resolves as $Re_{q, rest} = 2 / \alpha_{s, rest}$. % :::LLM_EDIT::: Cured "evaluate" repetition (reduces, remains, resolves).

As dynamic pressure increases, the scalar admittance drops ($\alpha_s \to 0$), and the Reynolds number diverges ($Re_q \to \infty$). The continuum transitions into turbulence. To relieve strain, the metric divides its volume, nucleating additional boundaries (hadronization). 

Because heavier particles possess a larger internal shear stress ($P_{shear}$), they maintain a smaller resting compliance ($\alpha_{s,rest}$). Heavier particles maintain a higher resting Reynolds number, positioning them closer to the turbulent phase transition. % :::LLM_EDIT::: Cured "evaluate" repetition (maintain).

\vspace{0.2 cm}

\textbf{3. The Equivalence Principle \& Multi-Scale Phase Separation:} \\
The Navier-Stokes equation expresses equilibrium as Force Densities ($\text{N/m}^3$). There is a dimensional distinction between Body Forces and Surface Forces. % :::LLM_EDIT::: Cured "evaluate" repetition (expresses).
\begin{itemize}[nosep]
    \item \textbf{Body Forces ($\mathbf{F}_{body}$):} Gravity ($-\nabla \Phi_{env}$) acts upon the mass of the continuum. Because its kinematic unit equates to acceleration ($\text{m/s}^2$), it must be multiplied by the effective mass density ($\rho_{eff}$) to resolve dimensionally as Force Density ($\text{kg/m}^3 \times \text{m/s}^2 = \text{N/m}^3$). % :::LLM_EDIT::: Cured "evaluate" repetition (equates to, resolve).
    \item \textbf{Surface Forces ($\mathbf{F}_{surface}$):} Mechanical compression, planar shear, and internal viscosity act upon the boundary. Their spatial gradients act as Force Densities ($\text{N/m}^3$). % :::LLM_EDIT::: Cured "evaluate" repetition (act).
\end{itemize}

Expanding the Navier-Stokes momentum equation using these dimensions and substituting the effective density identity ($\rho_{eff} = \rho_\tau/\alpha_s$) yields: % :::LLM_EDIT::: Cured "evaluate" repetition (Expanding).
\begin{align}
    \rho_{eff} \frac{D\mathbf{v}}{Dt} &= -\rho_{eff} \nabla \Phi_{env} + \Sigma \mathbf{F}_{surface} \\ 
    \left( \frac{\rho_\tau}{\alpha_s} \right) \frac{D\mathbf{v}}{Dt} &= -\left( \frac{\rho_\tau}{\alpha_s} \right) \nabla \Phi_{env} + \Sigma \mathbf{F}_{surface} 
\end{align}

In a homogeneous continuum, convective acceleration houses the dynamic pressure gradient, preventing the metric from exerting a bow-shock against itself. However, because a defect operates as a phase boundary propagating through the background metric, the mechanics enforce a multi-scale separation. The formulation operates as an action-reaction pair. The Eulerian material derivative ($\frac{D\mathbf{v}}{Dt}$) represents the phase acceleration of the boundary (the Action). The surface forces on the right side of the equation represent the impedance of the surrounding metric pushing back against the translating phase (the Reaction). % :::LLM_EDIT::: Cured "evaluate" repetition (operates, enforce, represents).

To isolate the acceleration ($\frac{D\mathbf{v}}{Dt}$), the equation is multiplied by the inverse of the effective density ($\alpha_s / \rho_\tau$). The mass term cancels from the gravitational body force, transferring the compliance ($\alpha_s$) to the surface forces: 
\begin{equation}
    \frac{D\mathbf{v}}{Dt} = -\nabla \Phi_{env} + \frac{\alpha_s}{\rho_\tau} \Sigma \mathbf{F}_{surface} \tag{\ref{eq:admittance_ns}}
\end{equation}

Gravity acts independently of kinematic strain. Because gravity is the pressure gradient of the background metric, inertial density cancels with gravitational density. This derives the Weak Equivalence Principle from continuum mechanics; acceleration due to gravity applies independently of internal mass or boundary velocity, while mechanical propulsion ($\mathbf{F}_{surface}$) is constrained by the compliance ($\alpha_s$). % :::LLM_EDIT::: Cured "evaluate" repetition (acts, applies).

\textbf{4. Dimensional Phase-Space Mapping (The Couplings):} \\ 
Surface forces separate into normal pressure and deviatoric shear stress. To map these mechanics to the continuum limits, the surface forces are throttled by their respective coupling constraints: 

\begin{itemize}[leftmargin=*]
    \item \textbf{Total Scalar Admittance ($\alpha_s$):} Because all surface forces act upon the boundary ($\rho_{eff}$), they are throttled by the total scalar admittance ($\alpha_s$). 
    \item \textbf{Normal Volumetric Stress:} Longitudinal translation generates a dynamic pressure gradient ($-\nabla P_{dyn}$), acting as the bow-shock against the translating boundary. This acts longitudinally and carries no additional coupling. 
    \item \textbf{Transverse Shear \& Rotational Viscosity:} To prevent the modulus ratio from throttling longitudinal friction, the stress tensor is restricted to the transverse shear tensor ($\mathbb{T}_{\perp}$). Transverse shear ($\nabla \cdot \mathbb{T}_{\perp}$) and rotational viscous friction ($\mu \nabla^2 \mathbf{v}$) are coupled to the metric via the modulus ratio ($\alpha$). Because the core operates as an incompressible circulation ($\nabla \cdot \mathbf{v} = 0$), the viscous vector identity ($\nabla^2 \mathbf{v} = -\nabla \times \boldsymbol{\Omega}$) defines rotational friction as the resistance against the rotational phase ($\boldsymbol{\Omega} = \nabla\lambda \times \nabla\beta_\Omega$). 
\end{itemize}

Expanding the surface forces with these nested throttles yields: % :::LLM_EDIT::: Cured "evaluate" repetition (yields).
\begin{equation}
    \frac{D\mathbf{v}}{Dt} = -\nabla \Phi_{env} + \frac{\alpha_s}{\rho_\tau} \Big[ -\nabla P_{dyn} + \alpha \left( \nabla \cdot \mathbb{T}_{\perp} - \mu \nabla \times (\nabla\lambda \times \nabla\beta_\Omega) \right) \Big] \tag{\ref{eq:propulsion_master}} 
\end{equation}

\textbf{5. The Equation of Motion:} \\ 
To normalize the equation into kinematic gradients ($\text{m/s}^2$), the background density ($\rho_\tau$) is substituted into the stress operators. 

By the Pressure-Curvature Identity, the gravitational potential equates to the capacity of the environment ($\Phi_{env} = \mathcal{P}_{ambient}$). The body force acting on the particle resolves as the negative gradient ($-\nabla \mathcal{P}_{ambient}$). % :::LLM_EDIT::: Cured "evaluate" repetition (resolves).

\begin{itemize}[leftmargin=*]
    \item \textbf{Kinematic Pressure ($\mathcal{P}$):} Dynamic pressure resolves as $P_{dyn} = \rho_\tau \mathcal{P}_{dyn}$. The density cancels: $-\nabla P_{dyn} / \rho_\tau = -\nabla \mathcal{P}_{dyn}$. % :::LLM_EDIT::: Cured "evaluate" repetition (resolves).
    \item \textbf{Kinematic Transverse Shear ($\mathbb{T}_{\perp, kin}$):} The metric shear modulus equates to $\mu_s = \rho_\tau c_a^2$. The deviatoric stress isolates the density ($\mathbb{T}_{\perp} = \rho_\tau \mathbb{T}_{\perp, kin}$), which cancels: $\nabla \cdot \mathbb{T}_{\perp} / \rho_\tau = \nabla \cdot \mathbb{T}_{\perp, kin}$. % :::LLM_EDIT::: Cured "evaluate" repetition (equates to).
    \item \textbf{Kinematic Viscosity ($\nu$):} Dynamic viscosity derives from the Prandtl threshold ($\mu = \rho_\tau \frac{1}{2}c_a L$). Substituting this defines the mass-independent kinematic viscosity ($\nu = \frac{1}{2}c_a L$). The density cancels: $-\mu \nabla \times \boldsymbol{\Omega} / \rho_\tau = -\nu \nabla \times \boldsymbol{\Omega}$. % :::LLM_EDIT::: Cured "evaluate" repetition (derives from).
\end{itemize}

Applying these substitutions yields the Equation of Motion: 
\begin{equation}
    \frac{D\mathbf{v}}{Dt} = -\nabla \mathcal{P}_{ambient} + \alpha_s \Big[ -\nabla \mathcal{P}_{dyn} + \alpha(\nabla \cdot \mathbb{T}_{\perp, kin} - \nu \nabla \times (\nabla\lambda \times \nabla\beta_\Omega)) \Big] \tag{\ref{eq:kinematic_ns}}
\end{equation}

Applying the Bernoulli Conservation Law details the symmetry of acceleration. Taking the spatial derivative of the environmental budget ($P_{ambient} = P_{static} + P_{dyn} + P_{shear}$) yields the gradient equality: 
\begin{equation}
    \nabla \mathcal{P}_{ambient} = \nabla \mathcal{P}_{static} + \nabla(\mathcal{P}_{dyn} + \mathcal{P}_{shear}) \tag{\ref{eq:bernoulli_gradient}}
\end{equation}

Gravity and propulsion operate as mechanisms governed by the same scalar field. The continuum distinguishes them by interaction: the pressure gradient ($-\nabla \mathcal{P}_{ambient}$) applies independently of boundary strain, accelerating the metric into the low-pressure well. Propulsion and electromagnetism remain constrained by the total scalar admittance ($\alpha_s$) and coupling constant ($\alpha$). % :::LLM_EDIT::: Cured "evaluate" repetition (operate, applies).

\begin{tcolorbox}[recallbox, title={\ref{box:admittance_ns}. Compliance-Modified Navier-Stokes \hfill\makebox[0pt][r]{}}]
\begin{align}
    \rho_a &= \rho_\tau \implies \rho_{eff} = \frac{\rho_\tau}{\alpha_s} \tag{\ref{eq:ns_density_substitution}} \\ 
    Re_q &= \frac{2}{\alpha_s} \implies Re_{q,rest} = \frac{2}{\alpha_{s,rest}} \tag{\ref{eq:quantum_reynolds}} \\ 
    \frac{D\mathbf{v}}{Dt} &= -\nabla \Phi_{env} + \frac{\alpha_s}{\rho_\tau} \Sigma \mathbf{F}_{surface} \tag{\ref{eq:admittance_ns}} \\ 
    \frac{D\mathbf{v}}{Dt} &= -\nabla \Phi_{env} + \frac{\alpha_s}{\rho_\tau} \Big[ -\nabla P_{dyn} + \alpha(\nabla \cdot \mathbb{T}_{\perp} - \mu \nabla \times \boldsymbol{\Omega}) \Big] \tag{\ref{eq:propulsion_master}} \\ 
    \nabla \mathcal{P}_{ambient} &= \nabla \mathcal{P}_{static} + \nabla(\mathcal{P}_{dyn} + \mathcal{P}_{shear}) \tag{\ref{eq:bernoulli_gradient}} \\
    \frac{D\mathbf{v}}{Dt} &= -\nabla \mathcal{P}_{ambient} + \alpha_s \Big[ -\nabla \mathcal{P}_{dyn} + \alpha(\nabla \cdot \mathbb{T}_{\perp, kin} - \nu \nabla \times (\nabla\lambda \times \nabla\beta_\Omega)) \Big] \tag{\ref{eq:kinematic_ns}} 
\end{align}
\end{tcolorbox}

% ----------------------------------------------------------------------
\section{The Localized Metric (Topological Strain \& The Irrotational Vortex)}
\label{sec:deriv_localized_metric_strain}

The kinematic boundary conditions of a topological defect resolve as the spatial pressure inversion of the macroscopic expanding cosmos. 

\textbf{1. The Cavitation Boundary (Yield \& Displacement):} \\
For the macroscopic continuum, the scalar potential of expansion drops to zero at the origin and maximizes to the acoustic limit as $r \to \infty$. For an acoustic cavity, this boundary condition inverts. 

The volumetric displacement of the saturated core generates an outward restorative force ($F$) scaling as the inverse square of the spatial coordinate ($F \propto 1/r^2$). Integrating this 1D force gradient across the spatial coordinate yields the potential energy field ($E \propto 1/r$): % :::LLM_EDIT::: Replaced banned word "void" with "saturated core" to accurately reflect the internal phase limit.
\begin{equation}
    E(r) = \int F(r) \, dr \propto \int \frac{1}{r^2} \, dr \propto \frac{1}{r}
\end{equation}

Dividing this energy field by the constant volumetric displacement of the core extracts the 3D scalar potential pressure field ($P_\Phi \propto 1/r$). Because the internal confinement pressure cannot exceed the static yield pressure of the metric ($P_c = \frac{1}{2}\rho_{\tau} c^2$), this defines a physical cavitation radius ($r_c$) where the inverse scalar potential strikes the acoustic limit: 
\begin{equation}
    P_\Phi(r) = P_c \left( \frac{r_c}{r} \right) \implies P_\Phi(r_c) = P_c \tag{\ref{eq:pressure_cavitation_boundary}}
\end{equation}

\textbf{2. Optical Dilatation \& Refractive Strain:} \\
Because the background metric acts as a Gradient-Index (GRIN) refractive medium, the scalar admittance field ($\alpha_s(r)$) dictates the optical index ($n(r)$). While the spatial gradient of this index ($\nabla n$) governs kinematic trajectory and geometric lensing, the magnitude dictates the longitudinal spatial strain. Admittance represents the remaining capacity of the metric ($\alpha_s(r) = \sqrt{1 - P_\Phi(r) / P_c}$). 

Substituting the inverse scalar potential ($P_\Phi(r) = P_c(r_c/r)$) derives the spatial admittance field:
\begin{equation}
    \alpha_s(r) = \sqrt{1 - \frac{r_c}{r}}
\end{equation}

In continuum elastodynamics, the longitudinal geometric strain ($\varepsilon$) equates to the refractive index ($n$). Because strain expands the optical path length, it resolves as the inverse of the scalar admittance ($n(r) = 1/\alpha_s(r)$). 
\begin{equation}
    \varepsilon(r) = n(r) = \frac{1}{\alpha_s(r)} = \left( 1 - \frac{r_c}{r} \right)^{-1/2}
\end{equation}

By the geometric equivalence between refractive strain and spatial curvature, this optical dilatation dictates the macroscopic gravitational redshift ($1+z = n$). This unifies the continuum pressure gradient with the relativistic redshift identity:
\begin{equation}
    1+z(r) = \left( 1 - \frac{r_c}{r} \right)^{-1/2} \tag{\ref{eq:strain_redshift_identity}}
\end{equation}

\textbf{3. The Monopole Hydrodynamic Horizon ($R_\tau$):} \\
When integrating spatial strain, the potential requires a physical boundary limit. In matter clusters, this strain superimposes with the gravitational well ($P_{ambient}$) at short ranges. However, it eventually balances with the cosmic boundary. 

Deriving the background spatial threshold ($R_\tau$) where the inverse potential ($P_c(r_c/R_\tau)$) equilibrates with the quadratic cosmic scalar potential ($P_c(R_\tau/R_{h,0})^2$) yields: 
\begin{align}
    \frac{r_c}{R_\tau} &= \frac{R_\tau^2}{R_{h,0}^2} \nonumber \\
    R_\tau &= \left( r_c R_{h,0}^2 \right)^{1/3} \tag{\ref{eq:hydrodynamic_horizon_equivalence}}
\end{align}

\vspace{0.2 cm}

\textbf{4. The Irrotational Vortex (Transverse Kinematics \& The $2\pi$ Phase Constraint):} \\
To satisfy Eulerian continuity, the topological defect conserves angular momentum via Kelvin's Circulation Theorem ($\Gamma = \oint v_\phi \cdot dl$). For the circulation to remain constant ($\Gamma = 2\pi r v_\phi$), the transverse velocity field decays as the inverse of the radius. 

To maintain phase lock at the cavitation boundary ($v_\phi(r_c) = c$), the velocity field resolves as an irrotational vortex. This represents a relaxed equilibrium state bounded by the Israel-Stewart wave operator: 
\begin{equation}
    v_\phi(r) = c \left( \frac{r_c}{r} \right) \tag{\ref{eq:irrotational_velocity_field}}
\end{equation}

Substituting this velocity field into the circulation integral extracts the phase constraint of the topological state:
\begin{equation}
    \Gamma = 2\pi r \left( c \frac{r_c}{r} \right) = 2\pi r_c c \label{eq:kelvin_phase_constraint}
\end{equation}

This $1/r$ transverse rotation conserves topological spin and generates a $1/r^3$ transverse dynamic pressure gradient ($\nabla P_{dyn} \propto \nabla v_\phi^2$). This isolates the magnetic dipole moment, orthogonally decoupling quantum kinematics from $1/r^2$ longitudinal macroscopic gravitational strain. 

\vspace{0.2 cm}

\textbf{5. The Dipole Kinematic Horizon ($R_\omega$):} \\
The spatial boundary of the transverse spin footprint forms where the angular frequency ($\omega_{local}$) equilibrates with the temporal phase frequency ($\omega_0 = 1/2\pi t$). 

The angular frequency isolates as the velocity over the radial coordinate ($\omega_{local} = v_\phi / r$). Substituting the irrotational vortex field yields $\omega_{local}(r) = c(r_c/r^2)$. Equilibrating this frequency against the macroscopic phase establishes the kinematic boundary of the dipole field ($R_\omega$): 
\begin{equation}
    \frac{c \cdot r_c}{R_\omega^2} = \frac{1}{2\pi t}
\end{equation}

Because the macroscopic expansion velocity equates to the causal horizon over proper time ($c = R_{h,0}/t$), substituting $1/t = c/R_{h,0}$ coordinates the velocities:
\begin{equation}
    \frac{c \cdot r_c}{R_\omega^2} = \frac{c}{2\pi R_{h,0}}
\end{equation}

The acoustic phase limits ($c$) cancel. Rearranging isolates the Dipole Kinematic Horizon:
\begin{align}
    \frac{r_c}{R_\omega^2} &= \frac{1}{2\pi R_{h,0}} \nonumber \\
    R_\omega^2 &= 2\pi r_c R_{h,0} \nonumber \\
    R_\omega &= \sqrt{2\pi r_c R_{h,0}} \tag{\ref{eq:dipole_horizon_equivalence}}
\end{align}
Because the square-root scaling of the dipole horizon is less than the cube-root scaling of the monopole horizon, the $1/r^3$ transverse rotation field terminates prior to the $1/r^2$ longitudinal gravitational field. 

\vspace{0.2 cm}

\textbf{6. Strain Superposition ($R_{\tau(body)}$):} \\
$R_{\tau}$ marks the transition boundary where a boson's circulatory strain depletion hands over to the cosmic background. When fermions aggregate ($M_{body} = \sum m_i$), the individual scalar potentials superimpose. 

The summation across the aggregated mass generates the Newtonian gravitational potential ($\Phi = G M_{body} / r$). 

\vspace{0.2 cm}

For a single particle, $R_\tau$ scales based on the Compton radius ($R_\tau = (r_c R_{h,0}^2)^{1/3}$). Because $r_c \propto m^{-1}$, this horizon shrinks as mass increases ($R_\tau \propto m^{-1/3}$). 

The phase boundaries ($r_c$) do not superimpose. Rather, the metric evaluates the aggregate body via the potential ($\Phi = G M_{body} / r$). 

The boundary of influence for this mass ($R_{\tau(body)}$) evaluates at the kinematic limit of its bound state. In an expanding universe, this horizon occurs where the inward escape velocity of the mass ($v_{esc}^2 = 2GM/r$) is overcome by the outward expansion velocity of the metric. Substituting the low-range Taylor expansion ($r \ll R_{h,0}$) of the continuous velocity field (Equation \ref{eq:v0}), the outward velocity evaluates as $v_0(r) \approx c \frac{r}{R_{h,0}}$. Equating the specific kinetic energies ($\frac{1}{2}v_{esc}^2 = \frac{1}{2}v_0^2$) yields the aggregated horizon:
\begin{align}
    \frac{G M_{body}}{R_{\tau(body)}} &= \frac{1}{2}c^2 \left(\frac{R_{\tau(body)}}{R_{h,0}}\right)^2 \nonumber \\
    \frac{2 G M_{body}}{c^2 R_{\tau(body)}} &= \frac{R_{\tau(body)}^2}{R_{h,0}^2} \tag{\ref{eq:macroscopic_equilibration}}
\end{align}

Grouping the constant mass terms isolates the equivalent hydrodynamic core ($R_s = 2GM_{body}/c^2$) as the natural geometric anchor. Substituting this core simplifies the boundary scaling:
\begin{align}
    \frac{R_s}{R_{\tau(body)}} &= \frac{R_{\tau(body)}^2}{R_{h,0}^2} \nonumber \\
    R_{\tau(body)} &= \left( R_s R_{h,0}^2 \right)^{1/3} \tag{\ref{eq:body_halo_scaling}}
\end{align}

Because the aggregate core scales linearly with mass ($R_s \propto M_{body}$), the halo expands proportionally to the cube root of the total mass ($R_{\tau(body)} \propto M_{body}^{1/3}$). 

\vspace{0.2 cm}

This geometric inversion identifies the physical boundary where a single coherent defect would become its own event horizon. Because the quantum boundary shrinks ($r_c \propto M^{-1}$) and the gravitational boundary grows ($R_s \propto M$), they geometrically intersect ($r_c = R_s$). Equating these limits isolates the crossover mass: 
\begin{align}
    \frac{\hbar}{M_{body} c} &= \frac{2 G M_{body}}{c^2} \nonumber \\
    M_{body}^2 &= \frac{\hbar c}{2 G} \nonumber \\
    M_{body} &= \frac{1}{\sqrt{2}} \sqrt{\frac{\hbar c}{G}} = \frac{M_P}{\sqrt{2}} \tag{\ref{eq:planck_intersection}}
\end{align}
Evaluating this threshold yields a crossover mass of approximately $15.4 \text{ } \mu\text{g}$ ($M_P/\sqrt{2}$). 

This intersection reveals the floor of the continuum. The crossover occurs at the geometric yield radius ($r_c = R_s = \sqrt{2} l_p$). This threshold defines the upper mass limit of a particle. If a coherent elementary defect exceeded this limit, its gravitational core ($R_s$) would eclipse its internal quantum wavelength ($r_c$). Because phase cannot propagate within a cavitation, the metric cannot sustain mass beyond this threshold. To accumulate mass above this limit, matter must aggregate as distinct decoupled particles. 

Because the outer boundaries scale via the inner core, substituting this intersection dictates that $R_{\tau(body)}$ equals the quantum boundary at the crossover ($R_{\tau(body)} = R_\tau$). The spatial boundary of influence forms a V-shaped metric inversion. Below the Planck limit ($M_{body} \ll M_P/\sqrt{2}$), topological phase confines the metric, shrinking the halo as elementary mass increases ($R_\tau \propto m^{-1/3}$). Above the Planck limit, the pressure gradient expands the halo as aggregated mass increases ($R_{\tau(body)} \propto M_{body}^{1/3}$).


\vspace{0.2 cm}

\textbf{7. Topological Magnification \& The Mass Identity:} \\
Because the core scales inversely with rest mass ($r_c \propto m^{-1}$), the boundary of the surrounding halo ($R_\tau \propto m^{-1/3}$) contracts at a slower fractional rate. The Conformal Magnification ($Mag$) evaluates this ratio. Substituting the horizon equivalence isolates the variable:
\begin{align}
    Mag &= \frac{R_\tau}{r_c} \nonumber \\
    Mag &= \frac{(r_c R_{h,0}^2)^{1/3}}{r_c} \nonumber \\
    Mag &= \left( \frac{R_{h,0}}{r_c} \right)^{2/3} \tag{\ref{eq:conformal_mag_substituted}}
\end{align}
Because $r_c \propto m^{-1}$, the Conformal Magnification dictates the proportionality of inertia ($Mag \propto m^{2/3}$). The magnitude of observable mass evaluates proportionally to this leverage ($m \propto Mag^{3/2}$). 

\vspace{0.2 cm}

\textbf{8. Total Internal Reflection (The Effective Optical Cross-Section):} \\
To determine the stability of the topological state, the kinematic limits can be solved via geometric optics. In a spherically symmetric refractive medium, Snell's Law applies via Bouguer's Formula ($r \cdot n(r) \cdot \sin \theta = \text{const}$). 

Total Internal Reflection occurs where the radial wave trajectory cannot propagate ($\theta = \pi/2$), where the radial derivative of the optical potential equals zero ($d/dr [r \cdot n(r)] = 0$). Substituting the inverse-root refractive strain index isolates the geometric trap: 
\begin{align}
    \frac{d}{dr} \left[ r \left( 1 - \frac{r_c}{r} \right)^{-1/2} \right] &= 0 \nonumber \\
    1 - \frac{3r_c}{2r} &= 0 \nonumber \\
    r_{crit} &= 1.5 r_c \tag{\ref{eq:critical_optical_radius}}
\end{align}

This critical boundary ($r_{crit} = 1.5 r_c$) is the optical envelope of the particle. The metric entraps its kinematic energy as a total internal reflection boundary, establishing an effective cross-section for incident scattering. 

\vspace{0.2 cm}

\textbf{9. Deep Inelastic Scattering (Fractional Sub-Structure):} \\
The resolution of fractional sub-structure within the particle evaluates as the geometric consequence of four-momentum transfer ($Q^2$) collapsing the optical Fresnel zone.

The physical wavelength of an acoustic probe evaluates via its interaction energy ($\lambda = hc / \sqrt{Q^2 c^2}$). At low interaction momenta, the incident wavelength is larger than the total internal reflection boundary ($\lambda \ge 1.5 r_c$). The probe diffracts against the optical envelope, yielding a scattering cross-section.

As interaction momentum increases ($Q \propto 1/\lambda$), the wavelength compresses.  

When the wavelength drops below the optical boundary ($\lambda \ll 1.5 r_c$), the incident wave decouples from the refractive halo. It penetrates the envelope, evaluating the internal strain gradients of the core. Because a composite boundary (e.g., a Baryonic Trefoil) requires internal geometric intersections to maintain Eulerian continuity, its internal strain maximizes at these physical crossings.

In standard topology, the minimum number of these physical intersections is dictated by the Crossing Number ($c(K)$). To satisfy kinematic equilibrium at this limit, the macroscopic interaction transitions from an optical envelope to the sum of the scattering cross-sections ($\sigma_i$) of these crossing nodes:
\begin{equation}
    \sigma_{total}(Q^2 \gg m^2 c^2) = \sum_{i=1}^{c(K)} \sigma_{i} \tag{\ref{eq:deep_inelastic}} 
\end{equation}
Deep inelastic scattering reveals a fractional sub-structure without violating topological confinement. The wave collapses, permitting the incident momentum to bypass the optical trap and scatter against the crossing nodes comprising the core.

\begin{tcolorbox}[recallbox, title={\ref{box:localized_metric_strain}. The Localized Metric}]
\begin{align}
    P_\Phi(r_c) &= P_c \tag{\ref{eq:pressure_cavitation_boundary}} \\
    \varepsilon(r) &= \frac{1}{\alpha_s(r)} = \left( 1 - \frac{r_c}{r} \right)^{-1/2} \tag{\ref{eq:strain_redshift_identity}} \\
    R_\tau &= \left( r_c R_{h,0}^2 \right)^{1/3} \tag{\ref{eq:hydrodynamic_horizon_equivalence}} \\
    v_\phi(r) &= c \left( \frac{r_c}{r} \right) \tag{\ref{eq:irrotational_velocity_field}} \\
    \Gamma &= 2\pi r_c c \tag{\ref{eq:kelvin_phase_constraint}} \\
    R_\omega &= \sqrt{2\pi r_c R_{h,0}} \tag{\ref{eq:dipole_horizon_equivalence}} \\
    \frac{2 G M_{body}}{c^2 R_{\tau(body)}} &= \frac{R_{\tau(body)}^2}{R_{h,0}^2} \tag{\ref{eq:macroscopic_equilibration}} \\
    R_{\tau(body)} &= \left( R_s R_{h,0}^2 \right)^{1/3} \tag{\ref{eq:body_halo_scaling}} \\
    M_{body} &= \frac{M_P}{\sqrt{2}} \tag{\ref{eq:planck_intersection}} \\
    Mag &= \left( \frac{R_{h,0}}{r_c} \right)^{2/3} \tag{\ref{eq:conformal_mag_substituted}} \\
    r_{crit} &= 1.5 r_c \tag{\ref{eq:critical_optical_radius}} \\
    \sigma_{total}(Q^2 \gg m^2 c^2) &= \sum_{i=1}^{c(K)} \sigma_i \tag{\ref{eq:deep_inelastic}} 
\end{align}
\end{tcolorbox}






% ----------------------------------------------------------------------
\section{The Unified Topological Engine}
\label{sec:deriv_unified_engine}

Matter operates as a topological engine. To satisfy Eulerian continuity ($\nabla \cdot \mathbf{v} = 0$), the boundary acts as a stator, driving a circulation loop that generates scalar pressure deficits and rotational impedance. 

\textbf{1. The Dipole (Leptons) \& Clebsch Vorticity:} \\
By the Poincaré-Hopf theorem, a continuous tangent vector field on a spherical boundary requires at least two points where the rotational vector evaluates to zero ($\mathbf{v}_{\circlearrowleft} = 0$). These poles define the axis of rotation. 

To bypass the irrotational constraints of a smooth velocity potential, the circulation of the surrounding halo is parameterized by Clebsch variables ($\mathbf{v}_{\circlearrowleft} = \lambda \nabla \beta$). As the metric traverses the polar core, it curls orthogonally to establish the equatorial perimeter, governed by the Compton radius ($r_c = \hbar/mc$).

To maintain zero volumetric divergence, the continuum evaluates topologically as a torus of internal vorticity: 
\begin{equation}
    \nabla \cdot (\lambda \nabla \beta) = 0 \tag{\ref{eq:engine_transduction}}
\end{equation}
The lepton evaluates as this unknot. The Bernoulli deficit of the closed loop establishes rest mass, the phase limit ($c$) establishes the boundary, and the antisymmetric Clebsch gradients establish the electromagnetic gauge field. The orientation of the polar axis establishes the spin state.

\textbf{2. The Composite Trefoil (Baryons):} \\
To satisfy volumetric continuity for fractional lobes, three orthogonal tubes interlock. They weave around a central void to form a closed Trefoil knot. 

The topology of the knot dictates a dual-phase kinematic state. Deep within the core, the inward-bound tubes converge and bundle. Their vectors align, propagating as parallel and laminar ($\nabla \times \mathbf{v}_{core} = 0$). This parallel alignment eliminates relative shear between the streams, deriving frictionless internal circulation (Asymptotic Freedom). 

Upon exiting the core, the tubes splay outward into orthogonal spatial planes to form the three macroscopic lobes. At the perimeter, these orthogonal paths intersect to close the knot. The steep-angle collisions at these crossing nodes induce transverse shear and turbulence ($\nabla \times \mathbf{v}_{surface} \neq 0$), weaving the boundary (Topological Confinement).

\textbf{3. Hydrostatic Supremacy (The Legendre Phase Transformation):} \\
To prevent the macroscopic static pressure from crushing the void, the circulating engine generates an outward restorative dynamic pressure ($P_{out}$). To evaluate the work performed across the boundary without encountering a Cartesian radial singularity ($r \to 0$), the power density Lagrangian ($\mathcal{L}$) undergoes a Legendre transformation. 

The independent Cartesian coordinate ($\gamma_z$) swaps for the conjugate rotational phase angle ($\theta$), projecting the energy into the Hamiltonian ($\mathcal{H}$):
\begin{equation}
    \mathcal{H}(\theta, \tau_{\theta z}) = \theta \tau_{\theta z} - \mathcal{L}(\gamma_z) \label{eq:power_legendre}
\end{equation}
This transformation dictates that the 1D inward crush ($P_{in}$) maps against an angular differential. Because the boundary evaluates as a closed circulation loop, the total restorative dynamic pressure evaluates as the boundary integral across a full phase cycle: 
\begin{align}
    P_{out} &= \oint_{\partial S} P_{in} \, d\theta  \nonumber \\
    P_{out} &= P_{in} \int_{0}^{2\pi} d\theta \nonumber \\
    P_{out} &= P_{in} \Big[ \theta \Big]_{0}^{2\pi} \nonumber \\
    P_{out} &= 2\pi P_{in} \nonumber
\end{align}
The phase loop amplifies the dynamic pressure by $2\pi$. This internal restoring force balances the vacuum pressure, guaranteeing boundary stability:
\begin{equation}
    \frac{P_{out}}{P_{in}} = 2\pi \tag{\ref{eq:engine_supremacy}}
\end{equation}

\textbf{4. The Planck Redline ($r_c \to \sqrt{2} l_p$):} \\ % :::LLM_EDIT::: Aligned redline radius with Eq. planck_intersection.
The metric circulates at the acoustic limit ($v_{core} = c$) at the Compton boundary ($r_c$).

Because the Compton radius scales inversely with mass ($r_c \propto 1/m$), as mass increases, the core compresses. As established in Equation \ref{eq:planck_intersection}, the mass lockout occurs when the cavitation boundary intersects the gravitational yield radius ($r_c = R_s = \sqrt{2} l_p$), defining the compliance-bounded Planck Mass ($M_P / \sqrt{2}$) as the maximum displacement of a single core: % :::LLM_EDIT::: Aligned maximum single-core mass with M_P/sqrt(2).
\begin{equation}
    \lim_{r_c \to \sqrt{2} l_p} m \implies m_{max} = \frac{M_P}{\sqrt{2}} \tag{\ref{eq:engine_redline}} % :::LLM_EDIT::: Updated redline limit to M_P/sqrt(2).
\end{equation}

\textbf{5. Hydrostatic Instability ($m^4$ Decay Ratio):} \\
To satisfy mass conservation, the core volume ($V_{r_c}$) must balance the volume of the circulating halo ($V_\tau$). 

By the Kinematic Origin of Inertia, observable mass evaluates as hydrodynamic added mass ($m_\tau = C_{m,i} \rho_{\tau} V_\tau$). Rearranging isolates the strain volume ($V_\tau = m_\tau / (C_{m,i} \rho_{\tau})$). 

The boundary evaluates as the volume of Compton radius ($r_c = \hbar / m_\tau c$):
\begin{equation}
    V_{r_c} = \frac{4}{3}\pi \left(\frac{\hbar}{m_\tau c}\right)^3 = \frac{4\pi\hbar^3}{3m_\tau^3 c^3} 
\end{equation}

The Instability Ratio ($I$) of the defect evaluates as the geometric quotient of the required circulation footprint to the available core. Substituting the defined volumes evaluates the instability scaling limit:
\begin{align}
    I &= \frac{V_\tau}{V_{r_c}} \nonumber \\
    I &= \frac{\left(\frac{m_\tau}{C_{m,i} \rho_{\tau}}\right)}{\left(\frac{4\pi\hbar^3}{3m_\tau^3 c^3}\right)} \nonumber \\ 
    I &= \left(\frac{3c^3}{4\pi C_{m,i} \rho_{\tau}\hbar^3}\right) m_\tau^4 \tag{\ref{eq:engine_decay}} 
\end{align}

As mass increases, the volume required to house the circulating strain grows in direct proportion to the mass ($V_\tau \propto m_\tau$), while the core shrinks inversely to the cube ($V_{r_c} \propto 1/m_\tau^3$). Consequently, the instability ratio diverges to the fourth power of the mass ($m^4$). 

At extreme mass scales, the dynamic pressure of the disjointed circulation halo exceeds the restorative yield strength of the metric. The stator fractures into lower-mass configurations. This yield defines the physical mechanism of particle decay. 

\begin{tcolorbox}[recallbox, title={\ref{box:unified_engine}. The Unified Topological Engine}]
\begin{align}
    \nabla \cdot (\lambda \nabla \beta) &= 0 \tag{\ref{eq:engine_transduction}} \\
    \mathcal{H}(\theta, \tau_{\theta z}) &= \theta \tau_{\theta z} - \mathcal{L}(\gamma_z) \tag{\ref{eq:power_legendre}} \\
    \frac{P_{out}}{P_{in}} &= 2\pi \tag{\ref{eq:engine_supremacy}} \\
    \lim_{r_c \to \sqrt{2} l_p} v_{core} &\to c \implies m_{max} = \frac{M_P}{\sqrt{2}} \tag{\ref{eq:engine_redline}} \\ % :::LLM_EDIT::: Updated recall box to M_P/sqrt(2).
    \frac{V_\tau}{V_{r_c}} &= \left(\frac{3c^3}{4\pi C_{m,i} \rho_{\tau}\hbar^3}\right) m_\tau^4 \tag{\ref{eq:engine_decay}} 
\end{align}
\end{tcolorbox}


% ----------------------------------------------------------------------
\section{Confinement Pressure \& Work}
\label{sec:deriv_bag_pressure}

Because the topological boundary circulates at the acoustic limit ($v \to c$), it generates a localized Bernoulli dynamic pressure deficit. To balance this thermodynamic deficit and geometrically secure the acoustic cavity, the continuum exerts an inward hydrostatic centripetal proper acceleration evaluating as $a_{loc} = c/2\pi t_{loc}$. Because proper time evaluates as $t_{loc} = r_c/c$, the boundary acceleration isolates to:
\begin{align}
    a_{loc} &= \frac{c}{2\pi (r_c/c)} \nonumber \\
    a_{loc} &= \frac{c^2}{2\pi r_c}
\end{align}

The rest mass boundary evaluates as $r_c = \hbar/mc$. Substituting this dimension isolates the restorative Force ($F = m a_{loc}$) and surface Area ($A = 4\pi r_c^2$) constraints:
\begin{align}
    F &= m \left( \frac{c^2}{2\pi (\hbar/mc)} \right) \nonumber \\
    F &= m \left( \frac{m c^3}{2\pi \hbar} \right) = \frac{m^2 c^3}{2\pi \hbar} \\[8pt]
    A &= 4\pi \left( \frac{\hbar}{mc} \right)^2 \nonumber \\
    A &= \frac{4\pi \hbar^2}{m^2 c^2}
\end{align}

Dividing the restorative Force by the spherical surface Area isolates the Confinement Pressure (Bag Pressure, $P_p$) counteracting the Bernoulli deficit:
\begin{align}
    P_p &= \frac{F}{A} \nonumber \\
    P_p &= \frac{\left( \frac{m^2 c^3}{2\pi \hbar} \right)}{\left( \frac{4\pi \hbar^2}{m^2 c^2} \right)} \nonumber \\
    P_p &= \left( \frac{m^2 c^3}{2\pi \hbar} \right) \left( \frac{m^2 c^2}{4\pi \hbar^2} \right) \nonumber \\
    P_p &= \frac{m^4 c^5}{8\pi^2 \hbar^3} \label{eq:pp_derived}
\end{align}

Multiplying this Confinement Pressure by the physical Volume of the acoustic cavity ($V = \frac{4}{3}\pi r_c^3$) isolates the Volumetric Work ($W_v$) required to maintain topological confinement:
\begin{align}
    W_v &= P_p \cdot V \nonumber \\
    W_v &= \left( \frac{m^4 c^5}{8\pi^2 \hbar^3} \right) \left( \frac{4}{3}\pi r_c^3 \right) \nonumber
\end{align}
Substituting the boundary limit $r_c^3 = \hbar^3 / m^3 c^3$:
\begin{align}
    W_v &= \left( \frac{m^4 c^5}{8\pi^2 \hbar^3} \right) \left( \frac{4\pi \hbar^3}{3 m^3 c^3} \right) \nonumber \\
    W_v &= \frac{4\pi m^4 c^5 \hbar^3}{24\pi^2 m^3 c^3 \hbar^3} \nonumber \\
    W_v &= \frac{4 m c^2}{24 \pi} \nonumber \\
    W_v &= \frac{m c^2}{6\pi} \label{eq:wv_derived}
\end{align}

\begin{tcolorbox}[recallbox, title={\ref{box:bag_pressure}. Confinement Pressure \& Work}]
\begin{align}
    P_p &= \frac{m^4 c^5}{8\pi^2 \hbar^3} \tag{\ref{eq:pp}} \\
    W_v &= \frac{m c^2}{6\pi} \tag{\ref{eq:wv_derived}} 
\end{align}
\end{tcolorbox}

% ----------------------------------------------------------------------
% ----------------------------------------------------------------------
\section{Kinematic Circulation \& The Origin of Spin} 
\label{sec:deriv_quantum_spin}

In continuum mechanics, rotational strain resolves as kinematic vorticity ($\boldsymbol{\omega} = \nabla \times \mathbf{u} \neq 0$). Matter operates as a stable defect within the tensioned metric, sustaining continuous Clebsch circulation ($\lambda \nabla \beta$) to satisfy equilibrium. 

\textbf{1. Angular Momentum \& The Impedance Limit ($\hbar$):} \\
A vortex defines a state of internal vorticity ($\nabla \times (\lambda \nabla \beta)$). Total angular momentum ($\mathbf{S}$) isolates as the volume integral of the momentum density crossed with the radial boundary: 
\begin{equation}
    \mathbf{S} = \iiint_V \rho_{\tau} (\mathbf{r} \times (\lambda \nabla \beta)) \, dV 
\end{equation}
Circulation strikes the acoustic limit ($v = c$) at the cavitation boundary ($r = r_c$). Integrating the momentum density for the displaced mass ($m$) yields the scalar angular momentum ($S = mcr_c$). 

Because angular momentum cannot exceed the Rotational Impedance of the vacuum ($Z_{\tau(\circlearrowleft)} \equiv \hbar$), the boundary restricts to the inverse phase radius ($r_c = \hbar / mc$). 
\begin{align}
    S &= m c \left( \frac{\hbar}{mc} \right) \nonumber \\
    S &= \hbar \tag{\ref{eq:transverse_shear_spin}}
\end{align}

\textbf{2. The Baseline Circulation Limit:} \\
By Kelvin's Circulation Theorem, the total circulation ($\Gamma$) of a closed path resolves as the line integral of velocity around the contour ($C$): 
\begin{equation}
    \Gamma = \oint_C (\lambda \nabla \beta) \cdot d\mathbf{l} 
\end{equation}
For an acoustic cavity, the boundary velocity locks to the phase limit ($c$). The line integral simplifies to this limit multiplied by the circumference ($C = 2\pi r_c$): 
\begin{align}
    \Gamma &= c \oint_C d\mathbf{l}  \\
    \Gamma &= c (2\pi r_c) 
\end{align}
Substituting the cavitation radius ($r_c = \hbar / mc$): 
\begin{align}
    \Gamma &= c (2\pi) \left( \frac{\hbar}{m c} \right)  \\
    \Gamma &= \frac{2\pi \hbar}{m} 
\end{align}
Because the reduced Planck constant expands as $\hbar = h / 2\pi$, substituting this isolates the circulation:
\begin{equation}
    \Gamma = \frac{h}{m} \tag{\ref{eq:fluid_circulation}}
\end{equation}

\textbf{3. The Spinor Topology ($4\pi$ Metric Untangling):} \\
To translate circulation into quantized angular momentum ($S$), the circulation must distribute over the radian slice of the boundary. 

If the core existed in a discrete void, it would resolve via a standard $2\pi$ rotation. However, the continuous metric connects the core to the macroscopic environment via the hydrostatic pressure gradient. To rotate the stator without fracturing the surrounding continuum, the phase geometry must operate as a Dirac Spinor. 

The metric requires two full rotational phase wraps ($4\pi$ radians) to untangle the surrounding grid and return the boundary to its identity state. The intrinsic Angular Momentum ($S_{fermion}$) evaluates as the mass multiplied by the Clebsch circulation over this $4\pi$ slice: 
\begin{align}
    S_{fermion} &= m \left( \frac{\Gamma}{4\pi} \right)  \\
    S_{fermion} &= \frac{m}{4\pi} \left( \frac{h}{m} \right)  \\
    S_{fermion} &= \frac{h}{4\pi} 
\end{align}
Substituting $h = 2\pi\hbar$ isolates the derivation of the Fermionic spin anomaly:
\begin{align}
    S_{fermion} &= \frac{2\pi\hbar}{4\pi}  \\
    S_{fermion} &= \frac{\hbar}{2} \tag{\ref{eq:spinor_topology}}
\end{align}
Fermionic half-integer spin emerges not as an intrinsic quantum property, but as the continuous requirement to double-wrap the stator to maintain equilibrium with the background metric. 

\begin{tcolorbox}[recallbox, title={\ref{box:quantum_spin}. Kinematic Circulation \& The Origin of Spin}] % :::LLM_EDIT::: Replaced "Fluid" with "Kinematic".
\begin{align}
    \mathbf{S} &= \iiint_V \rho_{\tau} (\mathbf{r} \times (\lambda \nabla \beta)) \, dV \implies S = \hbar \tag{\ref{eq:transverse_shear_spin}} \\ % :::LLM_EDIT::: Upgraded \rho to \rho_{\tau} to match derivation.
    \Gamma &= \oint_C (\lambda \nabla \beta) \cdot d\mathbf{l} = \frac{h}{m} \tag{\ref{eq:fluid_circulation}} \\
    S_{fermion} &= m \left( \frac{\Gamma}{4\pi} \right) = \frac{\hbar}{2} \tag{\ref{eq:spinor_topology}}
\end{align}
\end{tcolorbox}

% ----------------------------------------------------------------------
\section{Kinetic-Potential Strain Conservation}
\label{sec:deriv_volumetric_energy}

Energy resolves as the displaced volume of the continuous metric. Potential and kinetic energies emerge as geometric partitions of the total scalar admittance ($\alpha_s$). 

\textbf{1. Kinetic Strain (The Dynamic Bow Shock):} \\
A rest boundary displaces a kinematic strain volume ($V_\tau$). Under spatial translation, the boundary compresses the surrounding metric. 

The kinematic drag factor ($\gamma$) equates to the inverse of the total scalar admittance ($\alpha_s$):
\begin{equation}
    \gamma = \frac{1}{\alpha_s} \tag{\ref{eq:strain_lorentz}}
\end{equation}

The physical volume of kinematic translation ($V_{kin}$) isolates as the effective dilated volume ($\gamma V_\tau$) minus the baseline rest volume ($V_\tau$). Substituting the admittance ratio expands the displaced volume as a function of metric compression:
\begin{align}
    V_{kin} &= \gamma V_\tau - V_\tau \nonumber \\
    V_{kin} &= \left(\frac{1}{\alpha_s}\right) V_\tau - V_\tau \nonumber \\
    V_{kin} &= \left(\frac{1}{\alpha_s} - 1\right) V_\tau \tag{\ref{eq:kinetic_volume}}
\end{align}
Multiplying by the kinematic energy density of the added mass ($C_{m,i}\rho_{\tau} c^2$) yields the kinetic energy limit ($E_k = (\gamma - 1)m_\tau c^2$). 

\textbf{2. Potential Energy (Static Elastic Deficit):} \\
Translating a boundary against a pressure gradient ($\nabla \Phi$) performs thermodynamic work against the continuous restorative force of the metric ($\mathbf{F}$). The physical volume of metric displaced into this static elastic strain ($V_{pot}$) derives via the energetic work integral, scaled by the species-specific added mass energy density ($C_{m,i}\rho_{\tau} c^2$): 
\begin{equation}
    V_{pot} = \frac{1}{C_{m,i}\rho_{\tau} c^2} \int_{x_1}^{x_2} \mathbf{F} \cdot d\mathbf{x} \tag{\ref{eq:potential_volume}}
\end{equation}
Potential energy resolves as displaced volume in the elastic tension of the environment.

\textbf{3. Isochoric Transmutation (Conservation of Energy):} \\
By Eulerian continuity, the total displaced volume ($V_{total}$) of a closed topological system equates to:
\begin{equation}
    V_{total} = V_\tau + V_{kin} + V_{pot} = \text{Constant} \nonumber
\end{equation}

Taking the proper time derivative of the volume isolates the volumetric phase exchange. Because the intrinsic added mass of the primary defect remains invariant ($\dot{V}_\tau = 0$), the derivative expands as:
\begin{align}
    \frac{d}{dt}(V_{total}) &= \frac{d}{dt}(V_\tau + V_{kin} + V_{pot}) \nonumber \\
    0 &= \dot{V}_\tau + \dot{V}_{kin} + \dot{V}_{pot} \nonumber \\
    0 &= 0 + \dot{V}_{kin} + \dot{V}_{pot} \nonumber \\
    \dot{V}_{kin} &= -\dot{V}_{pot} \tag{\ref{eq:isochoric_transmutation}}
\end{align}

When a boundary decelerates against a gravitational gradient, the metric compression of its dynamic bow shock relaxes ($\dot{V}_{kin} < 0$). This kinematic compliance shifts into the surrounding metric, increasing the static elastic strain of the gravitational well ($\dot{V}_{pot} > 0$).

\textbf{The Topographical Shift:} \\
Conservation of energy manifests as the continuous isochoric volumetric exchange between dynamic boundary strain and static elastic tension. In a gravitational orbit, the metric compression shifts from the longitudinal vector of the falling mass ($\mathbf{v}_{translation}$) into the topographical static pressure gradient of the planet ($\nabla \mathcal{P}_{ambient}$).

\begin{tcolorbox}[recallbox, title={\ref{box:volumetric_energy}. Kinetic-Potential Strain Conservation}]
\begin{align}
    \gamma &= \frac{1}{\alpha_s} \tag{\ref{eq:strain_lorentz}} \\
    V_{kin} &= \left( \frac{1}{\alpha_s} - 1 \right) V_\tau \tag{\ref{eq:kinetic_volume}} \\
    V_{pot} &= \frac{1}{C_{m,i}\rho_{\tau} c^2} \int_{x_1}^{x_2} \mathbf{F} \cdot d\mathbf{x} \tag{\ref{eq:potential_volume}} \\
    \dot{V}_{kin} &= -\dot{V}_{pot} \tag{\ref{eq:isochoric_transmutation}}
\end{align}
\end{tcolorbox}



\newpage

\part*{PHASE IV: WAVE KINEMATICS \& ELECTROMAGNETISM}

\vspace{0.2 cm}


% ----------------------------------------------------------------------
\section{The Unified Lagrangian \& The Origin of Action} 
\label{sec:deriv_unified_lagrangian}

The Lagrangian density ($\mathcal{L}$) resolves dimensionally as Energy per Volume ($\text{J/m}^3$), equating to mechanical Pressure ($\text{N/m}^2$). 

Physical Action ($\mathcal{S} = \int \mathcal{L} \, d^4x$) integrates this scalar pressure across spacetime. The Principle of Least Action ($\delta \mathcal{S} = 0$) dictates that the metric deforms geometrically to minimize the pressure gradient. 

\vspace{0.2 cm}

\textbf{1. Field Decomposition \& Exterior Calculus:} \\ 
Because the continuum operates as a viscoelastic medium, it possesses both a bulk modulus and a shear modulus. By the Helmholtz decomposition, longitudinal fields ($\nabla h$) govern compression of the bulk medium, while transverse and rotational fields ($\nabla \times \mathbf{A}, \nabla \times \boldsymbol{\Omega}$) isolate divergence-free shear. 

\textbf{Geometric Equivalence: }To map these 3D vector operators to a covariant manifold, the continuum translates into exterior calculus. The longitudinal gradient ($\nabla h$) maps to the exact 1-form ($dh$). The transverse curl ($\nabla \times \mathbf{A}$) translates to the Hodge dual of the exterior derivative ($\star d\mathbf{A}$). For internal Clebsch vorticity, the exterior derivative of the scalar phases constructs the invariant 2-form ($\boldsymbol{\Omega} = d\lambda_\Omega \wedge d\beta_\Omega$). As derived previously, the orthogonal capacities of these fields are structurally bounded by the invariant Fine-Structure modulus ratio ($\alpha$). 

\vspace{0.2 cm}

\textbf{2. The Local Kinematic Capacities \& Scalar Compliance ($\alpha_k$):} \\
To establish the boundary of motion, the maximum achievable dynamic pressure is bounded by the ambient capacity of the localized gravity well: 
\begin{equation}
    P_{ambient} = \frac{1}{2}\rho_a c_a^2 \tag{\ref{eq:kinematic_capacities}}
\end{equation}
Because the modulus ratio ($\alpha$) is invariant across the continuum, the transverse kinematic capacity scales as $\alpha P_{ambient}$. 

As an applied thermodynamic load ($P_{total}$) consumes this ambient background, the remaining structural availability of the metric is tracked by a dimensionless scalar, the Kinematic Compliance ($\alpha_k$):
\begin{equation}
    \alpha_k = \sqrt{1 - \frac{P_{total}}{P_{ambient}}} \label{eq:kinematic_compliance}
\end{equation}
This scalar maps as the covariant continuum equivalent of the inverse Lorentz factor ($1/\gamma$), bounding kinematic action by scaling from $1$ (an unperturbed ambient vacuum) to $0$ (capacity saturation).

\vspace{0.2 cm}

\textbf{3. The Covariant Particle Lagrangian ($\mathcal{L}_{particle}$):} \\
The covariant action of a defect maps to its internal boundary stress evolving over the compliant geometry of the environment. The macroscopic field Lagrangian isolates as the invariant rest energy density of the particle ($P_{shear}$) multiplied by this remaining scalar compliance ($\alpha_k$): 
\begin{equation}
    \mathcal{L}_{particle} = - P_{shear} \alpha_k \label{eq:base_lagrangian}
\end{equation}

The total applied load expands into the dynamic kinetic fields ($P_{dyn}$) and the rest mass stress ($P_{shear}$). Substituting the orthogonal Helmholtz fields against their respective limits expands the covariant Lagrangian: 
\begin{equation}
    \mathcal{L}_{particle} = - P_{shear} \sqrt{1 - \frac{\frac{1}{2}\rho_a (\nabla h)^2}{P_{ambient}} - \frac{\frac{1}{2}\rho_a (\nabla \times \mathbf{A})^2}{\alpha P_{ambient}} - \frac{\frac{1}{2}\rho_a (\nabla \times (\lambda \nabla \beta_\Omega))^2}{\alpha P_{ambient}} - \frac{P_{shear}}{P_{ambient}}} \nonumber
\end{equation}
\begin{equation}
    \mathcal{L}_{particle} = - P_{shear} \sqrt{1 - \frac{\frac{1}{2}\rho_a \Big[ (\nabla h)^2 + \frac{1}{\alpha}(\nabla \times \mathbf{A})^2 + \frac{1}{\alpha}(\nabla \times (\lambda \nabla \beta_\Omega))^2 \Big] + P_{shear}}{P_{ambient}}} \tag{\ref{eq:dynamic_lagrangian}}
\end{equation}

\textbf{4. The Parabolic Limit \& The Binding Energy Offset ($\mathcal{K} - \mathcal{P}_{net}$):} \\
For the Euler-Lagrange equations to extract equations of motion, the action surface must possess a parabolic minimum (a kinetic well). 

To map the continuum to mechanics, the unified Lagrangian is approximated at the low-velocity limit ($P_{dyn,eq} \ll P_{ambient}$). The parabolic kinetic well isolates as the local osculating curve of the radical, extracted via the Maclaurin series expansion ($\sqrt{1-x} \approx 1 - x/2$). 

Substituting $x = (P_{dyn,eq} + P_{shear}) / P_{ambient}$ expands the action surface into three discrete geometric components:
\begin{align}
    \mathcal{L}_{particle} &\approx - P_{shear} \left( 1 - \frac{P_{dyn,eq} + P_{shear}}{2 P_{ambient}} \right) \nonumber \\
    \mathcal{L}_{particle} &\approx - P_{shear} + \frac{1}{2} P_{shear} \left( \frac{P_{dyn,eq}}{P_{ambient}} \right) + \frac{P_{shear}^2}{2 P_{ambient}} \nonumber
\end{align}

The elements are grouped to define the Net Potential ($\mathcal{P}_{net}$). The expansion term ($\frac{P_{shear}^2}{2 P_{ambient}}$) isolates the permanent zero-point shift induced by the particle consuming a fraction of the ambient capacity (gravitational binding energy): 
\begin{equation}
    \mathcal{P}_{net} = P_{shear} \left( 1 - \frac{P_{shear}}{2 P_{ambient}} \right) \tag{\ref{eq:net_potential}}
\end{equation}

Expanding the middle term and substituting $P_{ambient} = \frac{1}{2}\rho_a c_a^2$ allows the baseline density to cancel, isolating the spatiotemporal ratio: 
\begin{align}
    \mathcal{K}_{macro} &= \frac{1}{2} P_{shear} \left( \frac{\frac{1}{2}\rho_a v_{eq}^2}{\frac{1}{2}\rho_a c_a^2} \right) \nonumber \\
    \mathcal{K}_{macro} &= \frac{1}{2} \left( \frac{P_{shear}}{c_a^2} \right) v_{eq}^2 = \frac{1}{2}\rho_0 v_{eq}^2 \nonumber
\end{align}
The scalar compliance resolves to the action well:
\begin{equation}
    \mathcal{L}_{particle} \approx - \mathcal{P}_{net} + \mathcal{K}_{macro} \tag{\ref{eq:classical_lagrangian_limit}}
\end{equation}
The open parabola ($\frac{1}{2}mv^2$) is generated by truncating the infinite series, discarding the higher-order geometric constraints. As dynamic pressure approaches the kinematic capacity, the higher-order spatial terms warp the parabola into a vertical asymptote, enforcing the limit of wave propagation.

\vspace{0.2 cm}

\textbf{5. Topological Confinement (The Cavitation Limit):} \\
The Deborah number ($De$) defines the phase transition between compliance and rigidity. As dynamic pressure increases, the structural relaxation time of the metric diverges ($De \to \infty$). 

Where the applied kinetic load saturates the remaining ambient capacity ($P_{dyn} \to P_{ambient} - P_{shear}$), the compliance radical drops to zero ($\alpha_k \to 0$). Because rest mass consumes a permanent fraction of the coordinate space, the maximum achievable transverse velocity is bounded by the defect:
\begin{equation}
    v_{\perp, max} = c_a \sqrt{\alpha \left( 1 - \frac{P_{shear}}{P_{ambient}} \right)} \tag{\ref{eq:max_transverse_velocity}}
\end{equation}

At this limit, the metric yields to cavitation. The continuum behaves as a solid, confining the particle from further acceleration:
\begin{equation}
    \lim_{v_\perp \to v_{\perp, max}} \alpha_k = 0 \implies De \to \infty \tag{\ref{eq:deborah_confinement}}
\end{equation}

\vspace{0.2 cm}

\textbf{6. The Kinematic Field Lagrangian ($\mathcal{L}_{kin}$):} \\
To isolate the action of the phase space, the ambient density ($\rho_a$) factors out of the Lagrangian. 

Dividing the Lagrangian by $\rho_a$ cancels the mass dimension, yielding the Lagrangian in spatiotemporal dimensions ($\text{m}^2/\text{s}^2$). The ambient pressure translates into ambient kinematic pressure ($\mathcal{P}_{ambient} = P_{ambient}/\rho_a$), and the rest mass shear translates to kinematic shear ($\mathcal{P}_{shear} = P_{shear}/\rho_a$):
\begin{equation}
    \mathcal{L}_{kin} = - \mathcal{P}_{shear} \sqrt{1 - \frac{\frac{1}{2} \Big[ (\nabla h)^2 + \frac{1}{\alpha}(\nabla \times \mathbf{A})^2 + \frac{1}{\alpha}(\nabla \times (\lambda \nabla \beta_\Omega))^2 \Big] + \mathcal{P}_{shear}}{\mathcal{P}_{ambient}}} \label{eq:kinematic_lagrangian}
\end{equation}

\textbf{7. The Phase-Space Lagrangian:} \\
To map the intensity of the boundary ($v_{tot}$) across its available degrees of freedom, the internal state is defined by its modal mixing angles ($\theta, \phi$). 

Applying the modal geometry partitions the intensity into its respective limits: 
\begin{align}
    v_{\parallel}^2 &= v_{tot}^2 \cos^2(\theta) \nonumber \\
    v_{\perp}^2 &= v_{tot}^2 \sin^2(\theta)\cos^2(\phi) \nonumber \\
    v_{\circlearrowleft}^2 &= v_{tot}^2 \sin^2(\theta)\sin^2(\phi) \nonumber
\end{align}

Substituting these angular components scaled by the modulus ratio ($\alpha$) bounds the phase space:
\begin{equation}
    \mathcal{L}_{kin}(v_{tot}, \theta, \phi) = - \mathcal{P}_{shear} \sqrt{1 - \frac{\frac{1}{2} v_{tot}^2 \Big[ \cos^2(\theta) + \frac{1}{\alpha}\sin^2(\theta) \big( \cos^2(\phi) + \sin^2(\phi) \big) \Big] + \mathcal{P}_{shear}}{\mathcal{P}_{ambient}}} \nonumber
\end{equation}

Because the azimuthal modes satisfy the Pythagorean identity ($\cos^2(\phi) + \sin^2(\phi) = 1$), the coordinate $\phi$ acts as a cyclic (ignorable) coordinate. The action state reduces to: 
\begin{equation}
    \mathcal{L}_{kin}(v_{tot}, \theta) = - \mathcal{P}_{shear} \sqrt{1 - \frac{\frac{1}{2} v_{tot}^2 \Big[ \cos^2(\theta) + \frac{1}{\alpha}\sin^2(\theta) \Big] + \mathcal{P}_{shear}}{\mathcal{P}_{ambient}}} \tag{\ref{eq:angular_lagrangian}}
\end{equation}
The absence of $\phi$ demonstrates that the metric models transverse wave propagation and internal Clebsch vorticity as thermodynamically degenerate. The action remains identical regardless of the energy distribution between these shear states, establishing degeneracy within the phase space. 

\vspace{0.2 cm}

\textbf{8. The Universal Baseline \& Gravitational Compliance ($\alpha_g$):} \\
The ambient capacity ($P_{ambient}$) acts as a scalar gradient embedded within a geometric boundary. The maximum capacity of an unperturbed metric is determined by the critical yield pressure ($P_c = \frac{1}{2}\rho_\tau c^2$). 

An accumulation of defects (a gravitational mass) imposes a pressure deficit ($\Delta P_{macro}$) upon the surrounding vacuum, depleting the baseline geometry. The background gradient translates into fractional Gravitational Compliance ($\alpha_g$):
\begin{align}
    \alpha_g &= \sqrt{\frac{P_c - \Delta P_{macro}}{P_c}} = \sqrt{\frac{P_{ambient}}{P_c}} \nonumber \\
    P_{ambient} &= P_c \alpha_g^2 \tag{\ref{eq:gravitational_compliance}}
\end{align}

Substituting the gravitational compliance back into the unexpanded covariant particle Lagrangian resolves the multi-scale action equation:
\begin{equation}
    \mathcal{L}_{particle} = - P_{shear} \sqrt{1 - \frac{P_{dyn} + P_{shear}}{P_c \alpha_g^2}} \tag{\ref{eq:nested_lagrangian}}
\end{equation}
This nests the action inside the gravitational geometry. As a defect propagates deeper into a gravity well, the gravitational compliance ($\alpha_g$) drops. This shrinks the denominator ($P_c \alpha_g^2$), causing the particle to reach its internal cavitation limit at a lower kinetic energy, enforcing a scaled coordinate velocity for the local speed of light. 

\begin{tcolorbox}[recallbox, title={\ref{box:unified_lagrangian}. The Unified Kinematic Lagrangian}]
\begin{align}
    \mathcal{L}_{particle} &= - P_{shear} \sqrt{1 - \frac{\frac{1}{2}\rho_a \Big[ (\nabla h)^2 + \frac{1}{\alpha}(\nabla \times \mathbf{A})^2 + \frac{1}{\alpha}(\nabla \times (\lambda \nabla \beta_\Omega))^2 \Big] + P_{shear}}{P_{ambient}}} \tag{\ref{eq:dynamic_lagrangian}} \\
    \mathcal{P}_{net} &= P_{shear} \left( 1 - \frac{P_{shear}}{2 P_{ambient}} \right) \tag{\ref{eq:net_potential}} \\
    \mathcal{L}_{particle} &\approx - \mathcal{P}_{net} + \frac{1}{2}\rho_0 v_{eq}^2 \equiv \mathcal{K} - \mathcal{P}_{net} \tag{\ref{eq:classical_lagrangian_limit}} \\
    v_{\perp, max} &= c_a \sqrt{\alpha \left( 1 - \frac{P_{shear}}{P_{ambient}} \right)} \tag{\ref{eq:max_transverse_velocity}} \\
    \lim_{v_\perp \to v_{\perp, max}} \alpha_k &= 0 \implies De \to \infty \tag{\ref{eq:deborah_confinement}} \\
    \mathcal{L}_{kin}(v_{tot}, \theta) &= - \mathcal{P}_{shear} \sqrt{1 - \frac{\frac{1}{2} v_{tot}^2 \Big[ \cos^2(\theta) + \frac{1}{\alpha}\sin^2(\theta) \Big] + \mathcal{P}_{shear}}{\mathcal{P}_{ambient}}} \tag{\ref{eq:angular_lagrangian}} \\
    P_{ambient} &= P_c \alpha_g^2 \tag{\ref{eq:gravitational_compliance}} \\
    \mathcal{L}_{particle} &= - P_{shear} \sqrt{1 - \frac{P_{dyn} + P_{shear}}{P_c \alpha_g^2}} \tag{\ref{eq:nested_lagrangian}}
\end{align}
\end{tcolorbox}


% ----------------------------------------------------------------------
\section{The Legendre Hamiltonian Translation}
\label{sec:deriv_legendre_hamiltonian}

The dynamics governed by the Unified Lagrangian dictate kinematic action. Extracting observables requires translating the system from velocity phase-space into momentum phase-space via the Legendre transformation. 

\textbf{1. Topological Conjugate Momentum ($\mathbf{p}_{macro}$):} \\
To isolate the energy state of the boundary, the covariant Lagrangian density integrates the geometry of General Relativity ($\alpha_g$) into the scalar compliance ($\alpha_k$).
\begin{equation}
    \mathcal{L}_{particle} = - P_{shear} \alpha_k = - P_{shear} \sqrt{1 - \frac{P_{dyn} + P_{shear}}{P_c \alpha_g^2}} \nonumber
\end{equation}

The momentum density ($\mathbf{p}_{macro}$) emerges as the partial derivative of this Lagrangian with respect to the coordinate velocity ($\mathbf{v}_{macro}$). Because $P_{dyn} = \frac{1}{2}\rho_a \mathbf{v}_{macro}^2$, executing the derivative via the chain rule isolates the mass: 
\begin{align}
    \mathbf{p}_{macro} &= \frac{\partial \mathcal{L}_{particle}}{\partial \mathbf{v}_{macro}} \nonumber \\
    \mathbf{p}_{macro} &= -P_{shear} \left[ \frac{1}{2\alpha_k} \right] \frac{\partial}{\partial \mathbf{v}_{macro}} \left( 1 - \frac{\frac{1}{2}\rho_a \mathbf{v}_{macro}^2 + P_{shear}}{P_c \alpha_g^2} \right) \nonumber \\
    \mathbf{p}_{macro} &= P_{shear} \left[ \frac{1}{2\alpha_k} \right] \left( \frac{\rho_a \mathbf{v}_{macro}}{P_c \alpha_g^2} \right) \nonumber
\end{align}
Substituting the capacity limit ($P_c \alpha_g^2 = P_{ambient} = \frac{1}{2} \rho_a c_a^2$) isolates the mass density ratio ($\frac{\rho_a}{P_{ambient}} = \frac{2}{c_a^2}$). The geometric fractions cancel:
\begin{equation}
    \mathbf{p}_{macro} = P_{shear} \left[ \frac{1}{2\alpha_k} \right] \left( \frac{2 \mathbf{v}_{macro}}{c_a^2} \right) = \frac{P_{shear}}{c_a^2} \frac{\mathbf{v}_{macro}}{\alpha_k} \nonumber
\end{equation}
To isolate the mass density ($\rho_0 \equiv P_{shear}/c^2$), the wave speed is expanded via the gravitational compliance ($c_a^2 = c^2 \alpha_g^2$):
\begin{equation}
    \mathbf{p}_{macro} = \frac{P_{shear}}{c^2 \alpha_g^2} \frac{\mathbf{v}_{macro}}{\alpha_k} \equiv \frac{\rho_0}{\alpha_g^2} \frac{\mathbf{v}_{macro}}{\alpha_k} \tag{\ref{eq:conjugate_momentum}}
\end{equation}
The Legendre transform dictates that the effective coordinate inertia of a mass propagating through a gravity well increases by the inverse square of the gravitational compliance ($1/\alpha_g^2$).

\textbf{2. The Covariant Energy Function ($\mathcal{H}_{particle}$):} \\
The total energy density ($\mathcal{H}_{particle}$) is defined by the Legendre transform of the Lagrangian. Because this scalar state is expressed in terms of velocity ($\mathbf{v}_{macro}$) rather than momentum, it operates in the tangent bundle (velocity phase-space) as the Covariant Energy Function, isomorphic to the Hamiltonian. 
\begin{align}
    \mathcal{H}_{particle} &= \mathbf{p}_{macro} \cdot \mathbf{v}_{macro} - \mathcal{L}_{particle} \nonumber \\
    \mathcal{H}_{particle} &= \left( \frac{P_{shear}}{c_a^2} \frac{\mathbf{v}_{macro}^2}{\alpha_k} \right) - (-P_{shear} \alpha_k) \nonumber
\end{align}
Substituting the kinematic pressure ratio ($\frac{\mathbf{v}_{macro}^2}{c_a^2} = \frac{P_{dyn}}{P_c \alpha_g^2}$) translates the momentum term into geometry:
\begin{equation}
    \mathcal{H}_{particle} = \frac{P_{shear} P_{dyn}}{P_c \alpha_g^2 \alpha_k} + P_{shear} \alpha_k \nonumber
\end{equation}
Factoring the compliance ($\alpha_k$) and substituting its equivalent limit ($\alpha_k^2 = 1 - \frac{P_{dyn} + P_{shear}}{P_c \alpha_g^2}$) resolves the covariant energy state:
\begin{align}
    \mathcal{H}_{particle} &= \frac{P_{shear} P_{dyn} + P_{shear} P_c \alpha_g^2 \left( 1 - \frac{P_{dyn} + P_{shear}}{P_c \alpha_g^2} \right)}{P_c \alpha_g^2 \alpha_k} \nonumber \\
    \mathcal{H}_{particle} &= \frac{P_{shear} P_{dyn} + P_{shear} P_c \alpha_g^2 - P_{shear} P_{dyn} - P_{shear}^2}{P_c \alpha_g^2 \alpha_k} \nonumber \\
    \mathcal{H}_{particle} &= \frac{P_{shear} (P_c \alpha_g^2 - P_{shear})}{P_c \alpha_g^2 \alpha_k} \nonumber \\
    \mathcal{H}_{particle} &= \frac{P_{shear} \left(1 - \frac{P_{shear}}{P_c \alpha_g^2}\right)}{\alpha_k} \nonumber
\end{align}
The kinematic pressure ($P_{dyn}$) cancels from the numerator. The remaining terms isolate a structural constant of the coordinate.

\textbf{3. The Gravitational Rest State ($\mathcal{H}_{rest}$):} \\
The bracketed term represents the square of the metric compliance when the localized particle is at macroscopic rest ($\alpha_{rest}^2$). Substituting this constant condenses the Hamiltonian into the unified energy state: 
\begin{equation}
    \mathcal{H}_{particle} = P_{shear} \frac{\alpha_{rest}^2}{\alpha_k} \tag{\ref{eq:covariant_hamiltonian}}
\end{equation}
When the particle is measured at rest ($v_{macro} = 0 \implies \alpha_k = \alpha_{rest}$), the Hamiltonian resolves the observable rest energy density: 
\begin{equation}
    \mathcal{H}_{rest} = P_{shear} \alpha_{rest} \tag{\ref{eq:hamiltonian_rest_mass}}
\end{equation}
This functions as gravitational binding energy. Because the ambient capacity in a gravity well is lower than the unperturbed vacuum, the rest compliance ($\alpha_{rest}$) drops below 1. The observable rest energy of a particle decreases in a gravitational field. 

\textbf{4. The Lorentz Factor ($\gamma$) \& Dispersion Relation:} \\
The compliance denominator ($\alpha_k$) maps to the relativistic Lorentz factor. Factoring the rest compliance ($\alpha_{rest}$) from the compliance radical yields: 
\begin{equation}
    \alpha_k = \sqrt{\alpha_{rest}^2 - \frac{v_{macro}^2}{c_a^2}} = \alpha_{rest} \sqrt{1 - \frac{v_{macro}^2}{c_a^2 \alpha_{rest}^2}} \nonumber
\end{equation}
Substituting the cavitation limit ($v_{max} = c_a \alpha_{rest}$) establishes the dynamic compliance ratio:
\begin{equation}
    \alpha_k = \alpha_{rest} \sqrt{1 - \frac{v_{macro}^2}{v_{max}^2}} \tag{\ref{eq:lorentz_compliance}}
\end{equation}
Substituting this compliance back into the energy state resolves the dynamic energy state equation:
\begin{equation}
    \mathcal{H}_{particle} = \frac{P_{shear} \alpha_{rest}}{\sqrt{1 - v_{macro}^2/v_{max}^2}} = \frac{\mathcal{H}_{rest}}{\sqrt{1 - v_{macro}^2/v_{max}^2}} \tag{\ref{eq:relativistic_energy}}
\end{equation}
Mapping this state against the conjugate momentum resolves the invariant dispersion relation. Squaring the energy state and subtracting the squared momentum scaled by the cavitation limit isolates the square of the rest state: 
\begin{equation}
    \mathcal{H}_{particle}^2 - (p_{macro} v_{max})^2 = \mathcal{H}_{rest}^2 \tag{\ref{eq:dispersion_relation}}
\end{equation}

\textbf{5. Topological Mass Extraction (The Invariant Bridge):} \\
While the energy state ($\mathcal{H}_{rest}$) scales against the gravity well, the mass of a defect is anchored to the unperturbed baseline vacuum. 

To extract the bare energy ($E_{bare}$), the flat-space stress ($P_{shear}$) is integrated across the volume of the quantum state ($V_{\psi}$):
\begin{equation}
    E_{bare} = \int P_{shear} \, dV = P_{shear} V_{\psi} \nonumber
\end{equation}
To map this quantum state, the bare rest energy of the defect equates to the volumetric displacement of the baseline metric. The displacement scales against the yield pressure ($P_c = \frac{1}{2}\rho_\tau c^2$) over the geometric volume ($V_\tau$): 
\begin{equation}
    P_{shear} V_{\psi} \equiv P_c V_{\tau} = \frac{1}{2}\rho_\tau c^2 V_{\tau} \tag{\ref{eq:rest_energy_integral}}
\end{equation}
Dividing by the square of the acoustic limit translates the displacement into the rest mass ($m_{bare}$). The mass isolates as the added mass: 
\begin{equation}
    m_{bare} = \frac{\frac{1}{2}\rho_\tau c^2 V_{\tau}}{c^2} = \frac{1}{2}\rho_\tau V_{\tau} \nonumber
\end{equation}
Because the dynamic pressure yields a coefficient of $1/2$, the geometry of an defect structurally reduces to the added-mass coefficient of a perfect sphere ($C_m = 1/2$). 

The observable physical mass ($m_{obs}$) isolates as the integral of the rest state ($\mathcal{H}_{rest}$), scaling the bare mass by the static compliance of the local coordinate: 
\begin{equation}
    m_{obs} = m_{bare} \alpha_{rest} \tag{\ref{eq:topological_mass}}
\end{equation}
Rest mass emerges as the added mass of a spherical displacement.

\textbf{6. The Kinematic Field State ($\mathcal{H}_{kin}$):} \\
To isolate the phase space, the ambient density ($\rho_a$) factors out of the dynamic Covariant Energy Function. Dividing by $\rho_a$ cancels the mass dimension, translating the dynamic pressures into kinematic pressures ($\mathcal{P}_{shear} = P_{shear}/\rho_a$) in spatiotemporal dimensions ($\text{m}^2/\text{s}^2$). The baseline capacity translates to critical kinematic pressure ($\mathcal{P}_c = P_c/\rho_a$). 

Substituting the Helmholtz fields into the local compliance radical ($\alpha_k$) expands the denominator:
\begin{equation}
    \mathcal{H}_{kin} = \frac{\mathcal{P}_{shear} \alpha_{rest}^2}{\sqrt{1 - \frac{\frac{1}{2} \Big[ (\nabla h)^2 + \frac{1}{\alpha}(\nabla \times \mathbf{A})^2 + \frac{1}{\alpha}(\nabla \times (\lambda \nabla \beta_\Omega))^2 \Big] + \mathcal{P}_{shear}}{\mathcal{P}_c \alpha_g^2}}} \tag{\ref{eq:kinematic_hamiltonian}}
\end{equation}

\textbf{7. The Reduced Phase-Space State:} \\
To map the kinematic intensity of the boundary ($v_{tot}$) across its internal degrees of freedom, the geometry is partitioned by its modal mixing angles ($\theta, \phi$). 

Substituting the sum of these angular components scaled by the fractional modulus ratio ($\alpha$) bounds the internal phase space:
\begin{equation}
    \mathcal{H}_{kin}(v_{tot}, \theta, \phi) = \frac{\mathcal{P}_{shear} \alpha_{rest}^2}{\sqrt{1 - \frac{\frac{1}{2} v_{tot}^2 \Big[ \cos^2(\theta) + \frac{1}{\alpha}\sin^2(\theta)\cos^2(\phi) + \frac{1}{\alpha}\sin^2(\theta)\sin^2(\phi) \Big] + \mathcal{P}_{shear}}{\mathcal{P}_c \alpha_g^2}}} \nonumber
\end{equation}
Because the azimuthal modes satisfy the Pythagorean identity ($\cos^2(\phi) + \sin^2(\phi) = 1$), the coordinate $\phi$ acts as a cyclic (ignorable) coordinate. The energy state reduces to: 
\begin{equation}
    \mathcal{H}_{kin}(v_{tot}, \theta) = \frac{\mathcal{P}_{shear} \alpha_{rest}^2}{\sqrt{1 - \frac{\frac{1}{2} v_{tot}^2 \Big[ \cos^2(\theta) + \frac{1}{\alpha}\sin^2(\theta) \Big] + \mathcal{P}_{shear}}{\mathcal{P}_c \alpha_g^2}}} \tag{\ref{eq:angular_hamiltonian}}
\end{equation}
The absence of $\phi$ demonstrates that the metric treats transverse wave propagation and vorticity as degenerate. The energy state remains identical regardless of the internal energy distribution between these shear states, maintaining kinematic degeneracy within the phase space. 

\begin{tcolorbox}[recallbox, title={\ref{box:legendre_hamiltonian}. The Legendre Hamiltonian Translation}]
\begin{align}
    \mathbf{p}_{macro} &= \frac{P_{shear}}{c^2 \alpha_g^2} \frac{\mathbf{v}_{macro}}{\alpha_k} \equiv \frac{\rho_0}{\alpha_g^2} \frac{\mathbf{v}_{macro}}{\alpha_k} \tag{\ref{eq:conjugate_momentum}} \\
    \mathcal{H}_{particle} &= P_{shear} \frac{\alpha_{rest}^2}{\alpha_k} \tag{\ref{eq:covariant_hamiltonian}} \\
    \alpha_{rest}^2 &= 1 - \frac{P_{shear}}{P_c \alpha_g^2} \implies \mathcal{H}_{rest} = P_{shear} \alpha_{rest} \tag{\ref{eq:hamiltonian_rest_mass}} \\
    \alpha_k &= \alpha_{rest} \sqrt{1 - \frac{v_{macro}^2}{v_{max}^2}} \tag{\ref{eq:lorentz_compliance}} \\
    \mathcal{H}_{particle} &= \frac{\mathcal{H}_{rest}}{\sqrt{1 - v_{macro}^2/v_{max}^2}} \tag{\ref{eq:relativistic_energy}} \\
    \mathcal{H}_{particle}^2 &- (p_{macro} v_{max})^2 = \mathcal{H}_{rest}^2 \tag{\ref{eq:dispersion_relation}} \\
    E_{bare} &= P_{shear} V_{\psi} \equiv P_c V_{\tau} = \frac{1}{2}\rho_\tau c^2 V_{\tau} \tag{\ref{eq:rest_energy_integral}} \\
    m_{bare} &= \frac{1}{2}\rho_\tau V_{\tau}, \quad m_{obs} = m_{bare} \alpha_{rest} \tag{\ref{eq:topological_mass}} \\
    \mathcal{H}_{kin} &= \mathcal{P}_{shear} \alpha_{rest}^2 \Bigg[ 1 - \frac{1}{\mathcal{P}_c \alpha_g^2} \Bigg( \frac{1}{2} \Big[ (\nabla h)^2  + \frac{1}{\alpha}(\nabla \times \mathbf{A})^2 + \frac{1}{\alpha}(\nabla \times (\lambda \nabla \beta_\Omega))^2 \Big] + \mathcal{P}_{shear} \Bigg) \Bigg]^{-1/2} \tag{\ref{eq:kinematic_hamiltonian}} \\
    \mathcal{H}_{kin}(v_{tot}, \theta) &= \mathcal{P}_{shear} \alpha_{rest}^2 \Bigg[ 1 - \frac{1}{\mathcal{P}_c \alpha_g^2} \Bigg( \frac{1}{2} v_{tot}^2 \Big[ \cos^2(\theta) + \frac{1}{\alpha}\sin^2(\theta) \Big] + \mathcal{P}_{shear} \Bigg) \Bigg]^{-1/2} \tag{\ref{eq:angular_hamiltonian}}
\end{align}
\end{tcolorbox}







% ----------------------------------------------------------------------
\section{Topological Vorticity Equilibrium (The Dirac Isomorphism)}
\label{sec:deriv_dirac_isomorphism}

Fermion kinematics resolve as the mechanics of a gyroscopically stabilized acoustic vortex. The abstract operators of the Dirac equation map to geometric identities.

\textbf{1. The Spinor (Cartan Tetrad):} \\
The state vector ($\psi$) defines the orthogonal geometric orientation of the vortex, defined by the Cartan moving frame ($\mathbf{\hat{e}}_i$): 
\begin{equation}
    \psi \equiv \mathbf{\hat{e}}_i \tag{\ref{eq:spinor_identity}}
\end{equation}
The Clifford algebra of the $\gamma$-matrices acts as the rotation generators coupling internal vorticity to the Eulerian grid. 

\textbf{2. The Kinematic Operator (The Clebsch Bridge):} \\
The kinetic operator ($i\gamma^\mu \partial_\mu$) maps to the curl ($\nabla \times \mathbf{v}_\psi$) of the spinor. This resolves as the Clebsch vorticity of the baseline metric: 
\begin{equation}
    i\gamma^\mu \partial_\mu \equiv \nabla \times \mathbf{v}_\psi \equiv \nabla \times (\lambda \nabla \beta_\Omega) \tag{\ref{eq:gamma_curl}}
\end{equation}

\textbf{3. The Quantum Rest Mass Identity ($P_{shear}$):} \\
To transition from a discrete energy to a continuous field interaction, the rest mass translates to the distortional shear stress ($P_{shear}$). By the Planck-Einstein relation, the total rotational energy equates to $E = \hbar \omega$. 

Because pressure is isomorphic to energy density, the distortional shear stress acts as the energy scaled by the core volume ($V_{\psi}$): 
\begin{equation}
    P_{shear} = \frac{\hbar \omega}{V_{\psi}} \nonumber
\end{equation}
Because the field density equates to $\bar{\psi} \psi = 1 / V_{\psi}$, the mass Lagrangian resolves as the negative product of the energy and the field distribution: 
\begin{align}
    \mathcal{L}_{mass} &= - (\hbar \omega) (\bar{\psi} \psi) \nonumber \\
    \mathcal{L}_{mass} &= - \frac{\hbar \omega}{V_{\psi}} \nonumber \\
    \mathcal{L}_{mass} &= - P_{shear} \tag{\ref{eq:Quantum_mass_lagrangian}}
\end{align}
This defines the mass coupling ($m\bar{\psi}\psi$) as the shear pressure of the vortex. 

\textbf{4. The Vorticity Equilibrium Lagrangian:} \\
The interaction Lagrangian balances the kinetic and potential energy densities of the topological state: 
\begin{equation}
    \mathcal{L}_{Dirac} = \bar{\psi} (i\gamma^\mu \partial_\mu) \psi - m\bar{\psi}\psi \nonumber
\end{equation}
Substituting the vorticity and the shear pressure, maps the operator to the state equation, evaluated against ambient density ($\rho_a$):
\begin{equation}
    \mathcal{L}_{Dirac} \equiv \frac{1}{2}\rho_a |\nabla \times (\lambda \nabla \beta_\Omega)|^2 - P_{shear} \tag{\ref{eq:dirac_fluid_lagrangian}}
\end{equation}
Setting the Lagrangian to zero dictates that the dynamic pressure of the vortex ($\frac{1}{2}\rho_a |\nabla \times (\lambda \nabla \beta_\Omega)|^2$) balances the distortional shear stress ($P_{shear}$) required to sustain it. The vortex circulates at equilibrium. 

\textbf{5. Antimatter (Chiral Inversion \& Annihilation):} \\
The adjoint ($\bar{\psi}$) acts as the chiral inverse of the circulation. A boundary and its conjugate resolve as $W_+ = +1$ and $W_- = -1$ circulations. Their spatial superposition dictates the interference: 
\begin{align}
    \bar{\psi}\psi &\equiv (W_+ + W_-) \nonumber \\
    \bar{\psi}\psi &= 0 \tag{\ref{eq:chiral_annihilate}}
\end{align}
This zero-sum interference reduces localized circulation to zero ($\mathbf{v}_{tot} = 0$). By Bernoulli's principle, the collapse of internal dynamic pressure restores the ambient pressure ($P_{ambient}$). The surrounding metric collapses acoustic cavity, resulting in structural annihilation. 

\begin{tcolorbox}[recallbox, title={\ref{box:dirac_isomorphism}. Topological Vorticity Equilibrium (The Dirac Isomorphism)}]
\begin{align}
    \psi &\equiv \mathbf{\hat{e}}_i \tag{\ref{eq:spinor_identity}} \\
    i\gamma^\mu \partial_\mu &\equiv \nabla \times \mathbf{v}_\psi \equiv \nabla \times (\lambda \nabla \beta_\Omega) \tag{\ref{eq:gamma_curl}} \\
    \bar{\psi} \psi &= \frac{1}{V_{\psi}} \implies \mathcal{L}_{mass} = - (\hbar \omega) \bar{\psi}\psi = - P_{shear} \tag{\ref{eq:Quantum_mass_lagrangian}} \\
    \mathcal{L}_{Dirac} &\equiv \frac{1}{2}\rho_a |\nabla \times (\lambda \nabla \beta_\Omega)|^2 - P_{shear} = 0 \tag{\ref{eq:dirac_fluid_lagrangian}} \\
    \bar{\psi}\psi &\equiv (W_+ + W_-) = 0 \tag{\ref{eq:chiral_annihilate}} 
\end{align}
\end{tcolorbox}


% ----------------------------------------------------------------------
\section{Continuous Wave Mechanics \& Topological Phase States}
\label{sec:deriv_wave_mechanics}

The Lagrangian maps the wave spectrum, discrete particle identities, and nuclear decay as geometric constraints.

\textbf{1. The Continuous Wave Solutions:} \\
The Helmholtz decomposition partitions the motion of the metric into a real-valued mechanical wave spectrum. Dynamic propagation resolves as the phase capacities ($c_a$ and $c_a\sqrt{\alpha}$) scaled by the phase angles ($\theta, \phi$), the temporal propagation phase ($\mathbf{k} \cdot \mathbf{r} - \omega t$), and their respective orthogonal unit vectors ($\hat{\mathbf{k}}, \hat{\mathbf{n}}, \hat{\mathbf{m}}$). 

The three orthogonal wave equations resolve without imaginary phase states: 
\begin{align}
    \mathbf{v}_{\parallel}(r,t) &= \left[ c_a \cos(\theta) \cos(\mathbf{k} \cdot \mathbf{r} - \omega t) \right] \hat{\mathbf{k}} \tag{\ref{eq:wave_p}} 
    \\
    \mathbf{v}_{\perp}(r,t) &= \left[ c_a \sqrt{\alpha} \sin(\theta)\cos(\phi) \cos(\mathbf{k} \cdot \mathbf{r} - \omega t) \right] \hat{\mathbf{n}} \tag{\ref{eq:wave_s}} 
    \\
    \mathbf{v}_{\circlearrowleft}(r,t) &= \left[ c_a \sqrt{\alpha} \sin(\theta)\sin(\phi) \cos(\mathbf{k} \cdot \mathbf{r} - \omega t) \right] \hat{\mathbf{m}} \tag{\ref{eq:wave_rot}} 
\end{align}

These velocity fields map the physical spectrum: 
\begin{itemize}
    \item \textbf{Longitudinal ($\mathbf{v}_{\parallel}$):} Operates as propagating P-waves. Because the longitudinal field governs metric compression, substituting the ambient velocity required to balance a dynamic pressure deficit ($v_{\parallel} = \sqrt{2GM/r}$) isolates the gravitational potential: 
    \begin{equation}
        \Phi(r,t) = - \frac{(\mathbf{v}_{\parallel})^2}{2} \tag{\ref{eq:p_wave_potential}} 
    \end{equation}
    Dynamic fluctuations in this compressive field propagate as Gravitational Radiation.
    \item \textbf{Transverse ($\mathbf{v}_{\perp}$):} Operates as volume-preserving S-waves (Electromagnetic Radiation). 
    \item \textbf{Rotational ($\mathbf{v}_{\circlearrowleft} = \lambda \nabla \beta_\Omega$):} Operates as Clebsch vortices (Topological Rest Mass). 
\end{itemize}

\textbf{2. Kinematics \& The Lorentz Angle ($\theta$):} \\
The Divergence Angle ($\theta$) partitions the kinetic capacity between longitudinal translation ($\nabla h$) and topological structure. 

For a boundary translating through the metric, the velocity ($v_{macro} \equiv |\mathbf{v}_{\parallel}|$) equates to the longitudinal cosine projection of the mass-adjusted cavitation limit ($v_{max} = c_a \alpha_{rest}$): 
\begin{align}
    v_{macro} &= v_{max} \cos(\theta) \nonumber \\
    \frac{v_{macro}^2}{v_{max}^2} &= \cos^2(\theta) \nonumber
\end{align}
Applying the Pythagorean identity ($\sin^2\theta = 1 - \cos^2\theta$) isolates the sine projection:
\begin{align}
    \sin(\theta) &= \sqrt{1 - \cos^2(\theta)} \nonumber \\
    \sin(\theta) &= \sqrt{1 - \frac{v_{macro}^2}{v_{max}^2}} \tag{\ref{eq:lorentz_angle}} % :::LLM_EDIT::: Swapped label for tag ref to inherit from cheatbox.
\end{align}
The sine projection maps to the dynamic compliance ratio ($\alpha_k / \alpha_{rest}$), defining the inverse Lorentz factor ($1/\gamma$). Relativistic kinematics emerge as a geometric record of the boundary closing its Divergence Angle proportional to velocity through the capacity. 

\vspace{0.2 cm}

\textbf{3. Topological Charge \& The Phase Map ($\phi$):} \\
While the vacuum phase space exhibits continuous gauge symmetry across the azimuthal angle ($\phi$), a stable particle operates as a closed boundary. Satisfying internal phase interference, closed metric boundaries establish quantized harmonic fractions (standing waves). 

Mapped at the transverse equator ($\theta = \pi/2$), topological identity is governed by the Curl Angle ($\phi$), partitioning the shear capacity between rotational mass and transverse electromagnetism. Conventional electric charge ($Q$) acts as the geometric fraction of kinematic energy extruded as transverse shear, signed by the chiral handedness of the rotational flux: 
\begin{equation}
    Q = \text{sgn}(\phi) \cos^2(\phi) \tag{\ref{eq:charge_projection}} 
\end{equation}

Table \ref{tab:wave_phases} maps the composite states as discrete angular constraints of the phase space. While the phase angles ($\theta, \phi$) dictate the shape and charge coupling of the boundary, the dynamic mass hierarchy scales via the resonant frequency harmonic ($\omega_n$) of the internal circulation. The chiral phase ($+\phi$ vs $-\phi$) governs the direction of rotational flux, dictating the conventional charge sign differentiating matter from antimatter. 

\begin{table}[h]
\centering
\begin{tabular}{l c c c c l}
\hline
\textbf{Particle State} & \textbf{Harmonic ($\omega_n$)} & \textbf{Divergence ($\theta$)} & \textbf{Curl ($\phi$)} & \textbf{Charge ($Q$)} & \textbf{Mechanical State} \\
\hline
\multicolumn{6}{c}{\textbf{Closed Topological Boundaries (Fermions)}} \\
\hline
Electron ($e^-$) & Base ($\omega_0$) & Dipole & $-\pi$ & -1 & Transverse Vortex \\
Positron ($e^+$) & Base ($\omega_0$) & Dipole & $+\pi$ & +1 & Transverse Anti-Vortex \\
Up Quark ($u$) & Base ($\omega_0$) & Dipole & $+\arcsin(\sqrt{1/3})$ & +2/3 & Primary Triad Lobe \\
Down Quark ($d$) & Base ($\omega_0$) & Dipole & $-\arcsin(\sqrt{2/3})$ & -1/3 & Recessive Triad Lobe \\
Top Quark ($t$) & 3rd Gen ($\omega_2$) & Dipole & $+\arcsin(\sqrt{1/3})$ & +2/3 & Extreme Triad Lobe \\
\hline
\multicolumn{6}{c}{\textbf{Composite Topological Boundaries (Baryons)}} \\
\hline
Proton ($p^+$) & Composite & Dipole & Trefoil ($uud$) & +1 & Stable Composite Knot \\
Neutron ($n^0$) & Composite & Dipole & Trefoil ($udd$) & 0 & Unstable Composite Knot \\
\hline
\multicolumn{6}{c}{\textbf{Open Propagating Waves (Bosons \& Uncontained Strain)}} \\
\hline
Photon ($\gamma$) & Spectrum ($\omega$) & $\pi/2$ & $0$ & $0$ (Carrier) & Transverse S-Wave \\
Neutrino ($\nu$) & Spectrum ($\omega$) & $\to \pi$ & $+\pi/2$ & 0 & Left-Handed Slip-Knot \\
Antineutrino ($\bar{\nu}$) & Spectrum ($\omega$) & $\to 0$ & $-\pi/2$ & 0 & Right-Handed Slip-Knot \\
\hline
\end{tabular}
\caption{Phase Angle Mapping of Fundamental and Composite Metric States}
\label{tab:wave_phases}
\end{table}

\textbf{4. The Dynamic Neutrino Limit:} \\
A neutrino operates as longitudinal strain ($\phi = \pi/2$). Its cosine projection reduces to zero, yielding zero electromagnetic charge. Lacking a rigid volume, it propagates approaching the macroscopic kinematic limit ($v_{macro} \to v_{max}$).

To possess a non-zero mass defect, the Divergence Angle approaches the longitudinal axis ($\theta \to \pi$ for left-handed, $\theta \to 0$ for right-handed). This limit maintains a trace rotational shear ($P_{shear}$) while eliminating the transverse cross-section. Substituting these limits into the Kinematic Lagrangian ($\mathcal{L}_{kin}$) preserves a trace phase-space action: 
\begin{equation}
    \mathcal{L}_{kin}(\to \pi, \pi/2) \approx 0 \quad \text{and} \quad \mathcal{L}_{kin}(\to 0, -\pi/2) \approx 0 \nonumber
\end{equation}
The neutrino resolves as a longitudinal wave-packet propagating without the transverse shear required to couple with electromagnetism. 
\vspace{0.2 cm}

\textbf{5. Antimatter \& Annihilation:} \\ 
Antimatter acts as the geometric inverse of a Clebsch vortex. Where a matter vortex circulates as $\nabla \times (+\lambda \nabla \beta_\Omega)$, an anti-vortex reverses phase as $\nabla \times (-\lambda \nabla \beta_\Omega)$. Because the Lagrangian operates on the square of the internal velocities, the chiral phase inversion neutralizes: 
\begin{equation}
    (\nabla \times (-\lambda \nabla \beta_\Omega))^2 = (\nabla \times (+\lambda \nabla \beta_\Omega))^2 \nonumber
\end{equation}
Antimatter projects identical mass and metric curvature. 

Annihilation resolves as destructive interference. When a vortex ($+\lambda$) and an anti-vortex ($-\lambda$) overlap, their geometric boundaries mechanically sum to zero ($\sum \lambda \nabla \beta_\Omega = 0$). 

This eliminates the shear of the boundaries ($P_{shear} \to 0$). Without internal shear to consume the ambient capacity, the Net Potential ($\mathcal{P}_{net}$) vanishes. The scalar compliance radical returns to the unperturbed baseline ($\alpha_k \to \alpha_g$), eliminating the cavitation limit. The boundary yields, and the dynamic pressure radiates outward as compliant transverse S-waves (gamma radiation).
\begin{equation}
    \sum (\lambda \nabla \beta_\Omega) = 0 \implies P_{shear} \to 0 \implies \text{Boundary Yield} \tag{\ref{eq:annihilation_limit}} 
\end{equation}

\textbf{6. Mechanical Yield \& Beta Decay ($\beta^-$):} \\
Nuclear decay maps as metric yield. A neutron equates to an isolated trefoil triad ($u, d, d$). Outside the compressive pressure of a nucleus, the triad operates outside equilibrium. The heavy down lobe yields, relaxing into the lower-strain phase of an up quark ($u$). 

The excess transverse phase ($\Delta Q = -1$) and kinetic energy sheared during this phase transition extrude from the triad, stabilizing into an independent transverse vortex (the electron) and a longitudinal wake (the antineutrino):
\begin{equation}
    \mathcal{L}_{kin, n}(udd) \to \mathcal{L}_{kin, p}(uud) + \mathcal{L}_{kin, e}(\phi=-\pi) + \mathcal{L}_{kin, \bar{\nu}}(\theta \to 0, \phi = -\pi/2) \tag{\ref{eq:neutron_rupture}}
\end{equation}
Standard Model decay channels record the continuous angular metric volume resolving to structural equilibrium. 

\begin{tcolorbox}[recallbox, title={\ref{box:wave_mechanics}. Continuous Wave Mechanics \& Topological Phase States}] 
\begin{align}
    \mathbf{v}_{\parallel}(r,t) &= \left[ c_a \cos(\theta) \cos(\mathbf{k} \cdot \mathbf{r} - \omega t) \right] \hat{\mathbf{k}} \tag{\ref{eq:wave_p}} \\
    \mathbf{v}_{\perp}(r,t) &= \left[ c_a \sqrt{\alpha} \sin(\theta)\cos(\phi) \cos(\mathbf{k} \cdot \mathbf{r} - \omega t) \right] \hat{\mathbf{n}} \tag{\ref{eq:wave_s}} \\
    \mathbf{v}_{\circlearrowleft}(r,t) &= \left[ c_a \sqrt{\alpha} \sin(\theta)\sin(\phi) \cos(\mathbf{k} \cdot \mathbf{r} - \omega t) \right] \hat{\mathbf{m}} \tag{\ref{eq:wave_rot}} \\
    \Phi(r,t) &= - \frac{(\mathbf{v}_{\parallel})^2}{2} \tag{\ref{eq:p_wave_potential}} \\
    \sin(\theta) &= \sqrt{1 - \frac{v_{macro}^2}{v_{max}^2}} \equiv \frac{\alpha_k}{\alpha_{rest}} \equiv \frac{1}{\gamma} \tag{\ref{eq:lorentz_angle}} \\
    Q &= \text{sgn}(\phi) \cos^2(\phi) \tag{\ref{eq:charge_projection}} \\
    \sum (\lambda \nabla \beta_\Omega) = 0 &\implies P_{shear} \to 0 \implies \text{Boundary Yield} \tag{\ref{eq:annihilation_limit}} \\
    \mathcal{L}_{kin, n}(udd) &\to \mathcal{L}_{kin, p}(uud) + \mathcal{L}_{kin, e}(\phi=-\pi) + \mathcal{L}_{kin, \bar{\nu}}(\theta \to 0, \phi = -\pi/2) \tag{\ref{eq:neutron_rupture}}
\end{align}
\end{tcolorbox}







% ----------------------------------------------------------------------
\section{EM Kinematics \& The LC-Acoustic Isomorphism}
\label{sec:deriv_em_lc_isomorphism}

The phase velocity of a transverse shear wave propagating through the metric resolves equivalently via the Newton-Laplace acoustic equation ($c^2 = K_\tau/\rho_{\tau}$) and the Maxwell electromagnetic limit ($c^2 = 1/\epsilon_{\tau}\mu_{\tau}$). 

\vspace{0.2 cm}

\textbf{1. Continuum Vorticity \& Kinematic Acceleration:} \\
The transverse velocity of the metric $\mathbf{v}_{\perp}$ and the scalar kinematic pressure deficit $\mathcal{P}_{dyn} \equiv P_{dyn}/\rho_\tau$ directly derive the electromagnetic field. The spatial vorticity ($\boldsymbol{\Omega}$) acts as the curl of the transverse velocity ($\boldsymbol{\Omega} = \nabla \times \mathbf{v}_{\perp}$). The divergence of a curl resolves to zero:
\begin{equation}
    \nabla \cdot \boldsymbol{\Omega} = \nabla \cdot (\nabla \times \mathbf{v}_{\perp}) = 0 \tag{\ref{eq:maxwell_gauss_faraday}a}
\end{equation}

The transverse acceleration ($\mathbf{a}_\perp$) resolves as the spatial gradient of the kinematic pressure combined with the temporal rate of spatial sliding ($\mathbf{a}_\perp = -\nabla \mathcal{P}_{dyn} - \partial \mathbf{v}_{\perp}/\partial t$). 

Taking the curl of the acceleration cancels the scalar gradient ($\nabla \times \nabla \mathcal{P}_{dyn} = 0$). Commuting the independent spatial and temporal derivative operators upon the remaining velocity field isolates the induction state:
\begin{align}
    \nabla \times \mathbf{a}_\perp &= \nabla \times \left(-\nabla \mathcal{P}_{dyn} - \frac{\partial \mathbf{v}_{\perp}}{\partial t}\right) \nonumber \\
    \nabla \times \mathbf{a}_\perp &= -\frac{\partial}{\partial t}(\nabla \times \mathbf{v}_{\perp}) = -\frac{\partial \boldsymbol{\Omega}}{\partial t} \tag{\ref{eq:maxwell_gauss_faraday}b}
\end{align}

\textit{Electromagnetic Equivalence:} The classical electromagnetic vector potential ($\mathbf{A}_\perp$) acts as the transverse continuum velocity ($\mathbf{v}_\perp$), and the electric scalar potential ($\Phi$) acts as the kinematic pressure deficit ($\mathcal{P}_{dyn}$). Consequently, the magnetic field ($\mathbf{B}$) operates as spatial vorticity ($\boldsymbol{\Omega}$), and the electric field ($\mathbf{E}$) operates as acceleration ($\mathbf{a}_\perp$).

\vspace{0.2 cm}

\textbf{2. Dimensional Unit Normalization:} \\
By mapping the classical fields to metric kinematics, their fundamental dimensions translate. 

Spatial vorticity evaluates dimensionally as velocity divided by distance, yielding angular frequency:
\begin{equation}
    \boldsymbol{\Omega} = \nabla \times \mathbf{v}_{\perp} \implies \left[\frac{1}{\text{m}}\right] \left[\frac{\text{m}}{\text{s}}\right] = \left[\frac{1}{\text{s}}\right] \label{eq:b_frequency}
\end{equation}
Magnetism governs the rotational frequency of the metric ($\mathbf{B} \equiv \boldsymbol{\Omega}$). 

Transverse acceleration evaluates as velocity divided by time:
\begin{equation}
    \mathbf{a}_\perp = -\nabla \mathcal{P}_{dyn} - \frac{\partial \mathbf{v}_{\perp}}{\partial t} \implies \left[\frac{\text{m}}{\text{s}^2}\right] \label{eq:e_acceleration}
\end{equation}
The electric field governs the transverse kinematic acceleration ($\mathbf{E} \equiv \mathbf{a}_\perp$). 

\vspace{0.2 cm}

\textbf{3. Mass Equivalent Strain (Gauss's Law):} \\
The spatial divergence of local acceleration ($\nabla \cdot \mathbf{a}_\perp$) resolves as the restorative frequency of the metric ($[1/\text{s}^2]$). Multiplying this divergence by the 1D Linear Density ($\epsilon_{\tau}$, $[\text{kg/m}]$) isolates mechanical pressure ($[\text{N/m}^2]$) or energy density ($[\text{J/m}^3]$). 

By mass-energy equivalence ($E = mc^2$), this transverse energy density ($\sigma_\perp$) resolves as an effective mass density ($\rho_e \equiv \rho_{mass}$, $[\text{kg/m}^3]$) multiplied by the square of the phase limit:
\begin{equation}
    \sigma_\perp = \rho_e c^2 \nonumber
\end{equation}
Charge density ($\rho_e$) translates to the equivalent density of metric strain. Equating the geometric surface integral of the kinematic acceleration to this strain density yields the metric state:
\begin{align}
    \int_V \rho_e c^2 \, dV &= \int_S \epsilon_{\tau} \mathbf{a}_\perp \cdot d\mathbf{A} \nonumber \\
    \epsilon_\tau (\nabla \cdot \mathbf{a}_\perp) &= \rho_e c^2 \nonumber \\
    \nabla \cdot \mathbf{a}_\perp &= \frac{\rho_e c^2}{\epsilon_{\tau}} \tag{\ref{eq:maxwell_gaussE}}
\end{align}
Dividing the mass-energy density ($[\text{kg}/(\text{m}\cdot\text{s}^2)]$) by the 1D linear density ($[\text{kg/m}]$) yields the acceleration divergence limit ($[1/\text{s}^2]$).

\vspace{0.2 cm}

\textit{Electromagnetic Equivalence:} Substituting the classical field aliases ($\mathbf{E} \equiv \mathbf{a}_\perp$) recovers Gauss's Law for electricity ($\nabla \cdot \mathbf{E} = \rho_e c^2 / \epsilon_\tau$).

\vspace{0.2 cm}

\textbf{4. The 1D vs. 3D Dimensional Mapping:} \\
Under the Electro-Mechanical Mobility mapping, the geometry underlying 1D permittivity maps to the 1D Linear Density ($\epsilon_{\tau}$) of a continuous topological flux tube. Metric permeability maps to the 1D Metric Inverse Tension ($\mu_{\tau}$) of that topological tether. 

To transition from the 1D tether to the 3D continuum ($\rho_{\tau}$), the 1D linear density divides by the minimal transverse Planck shear area ($A_{shear, p}$). Substituting the tension limit ($1/\mu_{\tau} = F_p$) isolates the baseline 3D elastic bulk modulus ($K_\tau$):
\begin{align}
    \rho_{\tau} &= \frac{\epsilon_{\tau}}{A_{shear, p}} \nonumber \\
    K_\tau &= \frac{1}{\mu_{\tau} A_{shear, p}} = \frac{F_p}{A_{shear, p}} \tag{\ref{eq:lc_constants}}
\end{align}

\textbf{5. The Isomorphic Circuit Elements:} \\
\textit{Capacitance ($C \equiv m_\tau$):} Observable rest mass translates to 3D hydrodynamic added mass ($m_\tau = \frac{1}{2} \rho_{\tau} V_\tau$). Because 3D volume divided by 2D area reduces to 1D length ($V_\tau/A_{shear, p} = L_\tau$), observable inertia resolves as 1D capacitance:
\begin{equation}
    m_\tau = \frac{1}{2} \epsilon_{\tau} L_\tau \tag{\ref{eq:lc_mass}}
\end{equation}

\textit{Transverse Velocity ($v_{\perp}$):} Calculating the energy of a capacitor ($E_{cap}$) and substituting the isomorphic metric ($C \to m_\tau$) maps to classical kinetic energy:
\begin{equation}
    E_{cap} = E_{kin} = \frac{1}{2} m_\tau v_{\perp}^2 \nonumber
\end{equation}
Electrical potential maps to the transverse kinematic velocity ($v_{\perp}$) of the metric. Substituting this identity into the Bernoulli equation dictates that kinetic energy generates dynamic pressure:
\begin{equation}
    P_{dyn} = \frac{1}{2}\rho_{\tau} v_{\perp}^2 \tag{\ref{eq:velocity_pressure}}
\end{equation}

\textit{Inductance ($L_H \equiv 1/k$):} Electrical current maps to Macroscopic Mechanical Force ($I \to F$), and inductance maps to the Mechanical Compliance ($L_H \to 1/k$). Because Hookean elastic force dictates $F = kx$, dynamic inductive energy resolves as static elastic potential energy:
\begin{equation}
    E_{ind} = \frac{1}{2} L_H I^2 = \frac{1}{2} \left(\frac{1}{k}\right) (kx)^2 = \frac{1}{2} k x^2 \tag{\ref{eq:lc_inductance}}
\end{equation}

\textbf{6. The 1D Kinematic Projection (Ampère-Maxwell):} \\
Applying the 3D Navier-Cauchy equation for transverse waves yields:
\begin{equation}
    \rho_{\tau} \frac{\partial^2 \mathbf{v}_{\perp}}{\partial t^2} + K_\tau \nabla \times (\nabla \times \mathbf{v}_{\perp}) = \mathbf{J}_{3D} \nonumber
\end{equation}
Substituting the kinematic operators ($\frac{\partial \mathbf{v}_{\perp}}{\partial t} = -\mathbf{a}_\perp$ and $\nabla \times \mathbf{v}_{\perp} = \boldsymbol{\Omega}$) and the dimensional 1D/3D geometric projections ($\rho_{\tau} = \epsilon_{\tau}/A_{shear, p}$ and $K_\tau = 1/(\mu_{\tau} A_{shear, p})$):
\begin{align}
    \left( \frac{1}{\mu_{\tau} A_{shear, p}} \right) \nabla \times \boldsymbol{\Omega} &= \mathbf{J}_{3D} + \left( \frac{\epsilon_{\tau}}{A_{shear, p}} \right) \frac{\partial \mathbf{a}_\perp}{\partial t} \nonumber \\
    \nabla \times \boldsymbol{\Omega} &= \mu_{\tau} (A_{shear, p} \mathbf{J}_{3D}) + \mu_{\tau} \epsilon_{\tau} \frac{\partial \mathbf{a}_\perp}{\partial t} \nonumber
\end{align}
Because the interaction operates as the 1D topological projection of the 3D continuum, the 1D current density ($\mathbf{J}_{1D}$) resolves as the 3D macroscopic current ($\mathbf{J}_{3D}$) projected across the transverse yield area ($\mathbf{J}_{1D} = A_{shear, p} \mathbf{J}_{3D}$). Substituting this current isolates the relation:
\begin{equation}
    \nabla \times \boldsymbol{\Omega} = \mu_{\tau} \mathbf{J}_{1D} + \mu_{\tau} \epsilon_{\tau} \frac{\partial \mathbf{a}_\perp}{\partial t} \tag{\ref{eq:maxwell_curlB}}
\end{equation}

\textit{Electromagnetic Equivalence:} Substituting the classical fields ($\mathbf{B} \equiv \boldsymbol{\Omega}$, $\mathbf{E} \equiv \mathbf{a}_\perp$) and current ($\mathbf{J}_{EM} \equiv \mathbf{J}_{1D}$) recovers the Ampère-Maxwell law.

\vspace{0.2 cm}

\textbf{7. Kinematic Energy Flux (The Poynting Vector):} \\
The directional transport of energy represents the propagation of transverse kinetic energy through the metric. This defines the Kinematic Flux vector ($\mathbf{S}_{kin}$): 
\begin{equation}
    \mathbf{S}_{kin} = \frac{1}{\mu_{\tau}} (\mathbf{a}_\perp \times \boldsymbol{\Omega}) \tag{\ref{eq:poynting_flux}}
\end{equation}

This flux isolates as the spatial gradient of asymmetric dynamic pressure (the Acoustic Magnus Effect). 

When a uniform transverse velocity ($\mathbf{v}_\perp$) superimposes across a field of spatial vorticity ($\mathbf{v}_\circlearrowleft$), the total dynamic pressure ($\mathcal{P}_{dyn}$) expands via the square of the combined velocity fields:
\begin{align}
    \mathcal{P}_{dyn} &= \frac{1}{2}(\mathbf{v}_\perp + \mathbf{v}_\circlearrowleft)^2 \nonumber \\
    \mathcal{P}_{dyn} &= \frac{1}{2}v_\perp^2 + \frac{1}{2}v_\circlearrowleft^2 + (\mathbf{v}_\perp \cdot \mathbf{v}_\circlearrowleft) \label{eq:asymmetric_pressure_expansion}
\end{align}

The cross-term $(\mathbf{v}_\perp \cdot \mathbf{v}_\circlearrowleft)$ generates a high-pressure zone on one side of the rotational axis and a low-pressure zone on the opposing side causing a restorative force driven by the negative spatial gradient of this asymmetric pressure. 

Applying the vector calculus identity for the gradient of a dot product isolates the geometric cross-product:
\begin{align}
    -\nabla (\mathbf{v}_\perp \cdot \mathbf{v}_\circlearrowleft) &= - \mathbf{v}_\perp \times (\nabla \times \mathbf{v}_\circlearrowleft) \nonumber \\
    -\nabla (\mathbf{v}_\perp \cdot \mathbf{v}_\circlearrowleft) &= - \mathbf{v}_\perp \times \boldsymbol{\Omega}
\end{align}

Taking the temporal rate of change of this spatial pressure gradient translates the displacement into kinematic flux, converting the slip velocity ($\mathbf{v}_\perp$) into kinematic acceleration ($\mathbf{a}_\perp$):
\begin{equation}
    \partial_t \Big( -\nabla (\mathbf{v}_\perp \cdot \mathbf{v}_\circlearrowleft) \Big) \propto \mathbf{a}_\perp \times \boldsymbol{\Omega}
\end{equation}

In exterior calculus, this pressure squeeze structures the wedge product of a 1-form acceleration and a 2-form vorticity ($\mathbf{a}_\perp \wedge \boldsymbol{\Omega}$), constructing a 3-form volumetric flux. The Hodge dual of this flux outputs the 1-form longitudinal propagation vector. Electromagnetic energy flux acts as the acoustic Magnus effect of the metric, driving kinetic energy forward independent of mass transport. 
\vspace{0.2 cm}

\textit{Electromagnetic Equivalence:} Substituting the classical fields recovers the Poynting vector ($\mathbf{S} = \frac{1}{\mu_\tau}(\mathbf{E} \times \mathbf{B})$).

\vspace{0.2 cm}

\textbf{8. The Dual Transverse Geometries (Linear vs. Rotational):} \\
Transverse metric kinematics bifurcate into two independent spatial deformations within the 2D plane orthogonal to the axis of longitudinal propagation. 

The kinematic pressure deficit ($\mathcal{P}_{dyn}$) dictates the planar sliding of the metric, defining the linear transverse shear velocity ($\mathcal{V} \equiv v_{\perp}$). Circulation dictates the axial twisting of the metric, defining the rotational vorticity ($\boldsymbol{\Omega} \equiv \nabla \times \mathbf{v}_{\perp}$). 
\begin{equation}
    \mathcal{V} \equiv v_{\perp} \text{ (Linear)} \quad \text{\&} \quad \boldsymbol{\Omega} \equiv \nabla \times \mathbf{v}_{\perp} \text{ (Rotational)} \tag{\ref{eq:dual_transverse}}
\end{equation}

Translational sliding (Voltage) and axial twisting (Magnetism) operate as mechanically distinct shear states. A continuum can sustain sliding without twisting, or twisting without sliding. However, because both deformations occupy the same 2D transverse geometry, their simultaneous operation forces the velocity fields to overlap. This generates the asymmetric pressure dot-product $(\mathbf{v}_\perp \cdot \mathbf{v}_\circlearrowleft)$ derived in Equation \ref{eq:asymmetric_pressure_expansion}, driving the geometry into longitudinal propagation.

\vspace{0.2 cm}

\textbf{9. The Dimensional Calibration of Vacuum Permittivity ($\epsilon_0$):} \\
Because the metric defines charge as shear displacement ($e_{\tau}$), and SI units parameterize charge discretely (Coulombs, $e$), empirical vacuum permittivity ($\epsilon_0$) bridges via the dimensional ratio of these definitions. 

The transverse coupling identity defines charge as $e_{\tau}^2 = 4\pi \epsilon_{\tau} \hbar c \alpha$. The parameterized SI equivalent defines charge as $e^2 = 4\pi \epsilon_0 \hbar c \alpha$. Dividing the continuous derivation by the empirical parameterization eliminates the macroscopic wave limit and the geometric coupling constant:
\begin{align}
    \frac{\epsilon_{\tau}}{\epsilon_0} &= \frac{e_{\tau}^2}{e^2} \nonumber \\
    \epsilon_0 &= \epsilon_{\tau} \left( \frac{e}{e_{\tau}} \right)^2 \tag{\ref{eq:epsilon_0_derivation}}
\end{align}
Vacuum permittivity ($\epsilon_0$) maps to the 1D linear density ($\epsilon_{\tau}$), dimensionally scaled by the square of the ratio between the SI charge unit and continuum momentum. It does not represent an independent physical property of the vacuum, but acts as the dimensional conversion factor bridging electrical units and continuum mechanics. 

Defining this conversion scalar as $k_z = \epsilon_{\tau}/\epsilon_0$, the capacitive energy equivalence ($\frac{1}{2} C \mathcal{V}^2 = \frac{1}{2} m_\tau v_{\perp}^2$) maps scalar potential ($\mathcal{V}$) to transverse velocity:
\begin{equation}
    \mathcal{V} = \sqrt{k_z} \cdot v_{\perp} \approx (1.037 \times 10^{19}) \cdot v_{\perp} 
\end{equation}
The maximum scalar potential of the metric resolves when the transverse velocity strikes the phase bottleneck ($v_{\perp} = c$):
\begin{equation}
    \mathcal{V}_{max} = \sqrt{k_z} \cdot c \approx 3.11 \times 10^{27} \text{ Volts} \label{eq:universal_voltage_b}
\end{equation}

\textbf{10. The Spin Tensor (Faraday Matrix):} \\
To parameterize the deformation of the metric, the velocity gradient tensor ($\partial^\mu u^\nu$) separates volumetric expansion (gravity) from volume-preserving rotational shear (the kinematic spin tensor, $\Omega^{\mu\nu}$). 

The matrix separates the dual transverse geometries into dimensional planes based on their orthogonal shear ($\Omega^{\mu\nu} = \frac{1}{2}(\partial^\mu u^\nu - \partial^\nu u^\mu)$): 
\begin{equation}
    \Omega^{\mu\nu} = \begin{bmatrix} 
    0 & -a_{\perp,x}/c & -a_{\perp,y}/c & -a_{\perp,z}/c \\ 
    a_{\perp,x}/c & 0 & -\Omega_z & \Omega_y \\ 
    a_{\perp,y}/c & \Omega_z & 0 & -\Omega_x \\ 
    a_{\perp,z}/c & -\Omega_y & \Omega_x & 0 
    \end{bmatrix} \tag{\ref{eq:rosetta_faraday}}
\end{equation}

Because transverse acceleration ($\mathbf{a}_\perp$) evaluates as ($\text{m/s}^2$) and is divided by the phase limit $c$ ($\text{m/s}$), every element in the matrix resolves to the dimension of Frequency ($[1/\text{s}]$). The spin tensor operates as the angular velocity matrix. 

Because the deformations are transverse (volume-conserving), normal stress components vanish, driving the principal diagonal to zero. The remaining elements isolate the two behaviors: 

\textit{Space-Time Shear (Linear / Potential):} The intersection of the temporal and spatial dimensions isolates the temporal rate of spatial sliding ($\partial_t u_x$). This maps the transverse acceleration ($\mathbf{a}_\perp$) to the space-time components. 

\textit{Space-Space Shear (Rotational / Vorticity):} The intersection of spatial dimensions isolates the classical 2D curl ($\partial_x u_y - \partial_y u_x$). This maps the vorticity ($\boldsymbol{\Omega}$) to the space-space components. 

\textit{Electromagnetic Equivalence:} The classical electromagnetic Faraday tensor ($F^{\mu\nu}$) resolves as the spin tensor ($F^{\mu\nu} \equiv \Omega^{\mu\nu}$). Space-time shear maps the electric field ($\mathbf{E} \equiv \mathbf{a}_\perp$), and space-space shear maps the magnetic field ($\mathbf{B} \equiv \boldsymbol{\Omega}$). The electromagnetic field resolves as the spin of the metric.

\begin{tcolorbox}[recallbox, title={\ref{box:em_lc_isomorphism}. Continuum Kinematics \& The LC-Acoustic Isomorphism}]
\begin{align}
    \nabla \cdot \boldsymbol{\Omega} = 0 \quad &\text{\&} \quad \nabla \times \mathbf{a}_\perp = -\frac{\partial \boldsymbol{\Omega}}{\partial t} \tag{\ref{eq:maxwell_gauss_faraday}} \\
    \boldsymbol{\Omega} \equiv \text{Frequency } [1/\text{s}], \quad &\text{\&} \quad \mathbf{a}_\perp \equiv \text{Acceleration } [\text{m/s}^2] \tag{\ref{eq:b_frequency}, \ref{eq:e_acceleration}} \\
    \nabla \cdot \mathbf{a}_\perp &= \frac{\rho_e c^2}{\epsilon_{\tau}} \tag{\ref{eq:maxwell_gaussE}} \\
    \rho_{\tau} = \frac{\epsilon_{\tau}}{A_{shear, p}}, \quad &K_\tau = \frac{F_p}{A_{shear, p}} \tag{\ref{eq:lc_constants}} \\
    m_\tau &= \frac{1}{2} \epsilon_{\tau} L_\tau \tag{\ref{eq:lc_mass}} \\
    E_{cap} = E_{kin} &= \frac{1}{2} m_{\tau} v_{\perp}^2 \implies P_{dyn} = \frac{1}{2}\rho_{\tau} v_{\perp}^2 \tag{\ref{eq:velocity_pressure}} \\
    E_{ind} = \frac{1}{2} L_H I^2 &= \frac{1}{2} k x^2 \tag{\ref{eq:lc_inductance}} \\
    \nabla \times \boldsymbol{\Omega} &= \mu_{\tau} \mathbf{J}_{1D} + \mu_{\tau} \epsilon_{\tau} \frac{\partial \mathbf{a}_\perp}{\partial t} \tag{\ref{eq:maxwell_curlB}} \\
    \mathbf{S}_{kin} &= \frac{1}{\mu_{\tau}} (\mathbf{a}_\perp \times \boldsymbol{\Omega}) \tag{\ref{eq:poynting_flux}} \\
    \mathcal{V} \equiv v_{\perp} \text{ (Linear)} \quad &\text{\&} \quad \boldsymbol{\Omega} \equiv \nabla \times \mathbf{v}_{\perp} \text{ (Rotational)} \tag{\ref{eq:dual_transverse}} \\
    \epsilon_0 &= \epsilon_{\tau} \left( \frac{e}{e_{\tau}} \right)^2 \tag{\ref{eq:epsilon_0_derivation}} \\
    \mathcal{V}_{max} &= \sqrt{k_z} \cdot c \approx 3.11 \times 10^{27} \text{ V} \tag{\ref{eq:universal_voltage_b}}
\end{align}
\end{tcolorbox}

% ----------------------------------------------------------------------
\section{1D Fractional Phase Flux (Dimensional Projections)}
\label{sec:deriv_tensor_projections}

Fractional phase states emerge as tensor artifacts of projecting an isotropic 3D circulation onto a 1D observation axis. 

\vspace{0.2 cm}

\textbf{1. The Isotropic 3D Circulation ($SU(3)$ Traceless Strain):} \\
As established, the $SU(3)$ traceless deviatoric strain ($\text{Tr}(\varepsilon) = 0$) acts as an incompressible, volume-preserving defect. Gyroscopic stability distributes the phase flux ($\Phi_{\circlearrowleft}$) symmetrically across all $3$ orthogonal spatial dimensions. 

\vspace{0.2 cm}

\textbf{2. The Longitudinal Dimensional Projection ($\parallel$):} \\
When a linear probe impacts this 3D isotropic vortex, the momentum transfer measures the flux along a 1D Cartesian vector (the longitudinal observation axis). 

The projection of a uniform 3D spherical distribution onto a 1D longitudinal axis integrates the squared directional cosine over the full 3D solid angle ($\Omega$). The polar integration computes the 1D longitudinal flux fraction ($\Phi_{\parallel}$): 
\begin{align}
    \Phi_{\parallel} &= \Phi_{\circlearrowleft} \left( \frac{1}{4\pi} \int_{0}^{2\pi} d\phi \int_{0}^{\pi} \cos^2 \theta \sin \theta \, d\theta \right) \nonumber \\
    \Phi_{\parallel} &= \frac{1}{3} \Phi_{\circlearrowleft} \tag{\ref{eq:tensor_long}}
\end{align}
Because the longitudinal projection measures $1$ of the $3$ orthogonal spatial dimensions, it equals one-third of the total 3D circulation. 

\textbf{3. The Transverse Dimensional Projection ($\perp$):} \\
The remainder of the isotropic 3D flux projects onto the coupled transverse plane orthogonal to the observation axis. This projection applies the squared sine over the solid angle resulting in the transverse flux fraction ($\Phi_{\perp}$): 
\begin{align}
    \Phi_{\perp} &= \Phi_{\circlearrowleft} \left( \frac{1}{4\pi} \int_{0}^{2\pi} d\phi \int_{0}^{\pi} \sin^2 \theta \sin \theta \, d\theta \right) \nonumber \\
    \Phi_{\perp} &= \frac{2}{3} \Phi_{\circlearrowleft} \tag{\ref{eq:tensor_trans}}
\end{align}
Because the transverse projection measures $2$ of the $3$ orthogonal dimensions, it represents two-thirds of the total 3D circulation. 

\vspace{0.2 cm}

\textbf{4. Flux Conservation:} \\
Because the sum of the orthogonal projections equals unity, the total flux is conserved during the 1D measurement: 
\begin{align}
    \Phi_{\parallel} + \Phi_{\perp} &= \frac{1}{3}\Phi_{\circlearrowleft} + \frac{2}{3}\Phi_{\circlearrowleft} \nonumber \\
    \Phi_{\parallel} + \Phi_{\perp} &= \Phi_{\circlearrowleft} \tag{\ref{eq:projection_conservation}}
\end{align}
The fractional phase fluxes ($1/3, 2/3$) describe the orthogonal spatial cross-sections of a single 3D volume, independent of sub-components. 

\vspace{0.2 cm}

\textbf{5. The Subatomic Fractional States:} \\
The linear phase flux projecting across a 1D longitudinal observation axis ($\Phi_0$) acts as the $1/3$ dimensional projection: 
\begin{equation}
    \Phi_0 = \frac{1}{3} \Phi_{\circlearrowleft} \tag{\ref{eq:1d_phase_flux}}
\end{equation}

Because charge operates as a volumetric phase flux, subatomic fractional charge states function as the tensor projections of a single isotropic defect. This derivation removes fractional charge as a particulate property, redefining the subatomic sector as the spatial projection of 3D Euclidean space. 

\begin{tcolorbox}[recallbox, title={\ref{box:tensor_projections}. 1D Fractional Phase Flux (Dimensional Projections)}]
\begin{align}
    \Phi_{\parallel} &= \frac{1}{3} \Phi_{\circlearrowleft} \tag{\ref{eq:tensor_long}} \\
    \Phi_{\perp} &= \frac{2}{3} \Phi_{\circlearrowleft} \tag{\ref{eq:tensor_trans}} \\
    \Phi_{\parallel} + \Phi_{\perp} &= \Phi_{\circlearrowleft} \tag{\ref{eq:projection_conservation}} \\
    \Phi_0 &= \frac{1}{3} \Phi_{\circlearrowleft} \tag{\ref{eq:1d_phase_flux}}
\end{align}
\end{tcolorbox}

% ----------------------------------------------------------------------
\section{Continuum Constants, Scale Invariance, and UV Regularization}
\label{sec:deriv_continuum_constants}

The integration of momentum states resolves to a maximum, bounded by the geometric yield limits. This yield limit is not a static constant. The background tension operates as a sliding scale, where the local yield strength is dictated by the volume of the region under stress.

\vspace{0.2 cm}

\textbf{1. The Planck Baseline (Density \& 1D Stiffness)}\\
The boundary is analyzed at the minimum geometric limit: the Planck scale. 

To derive the baseline metric density ($\rho_\tau$), the background stress tensor equates to the isotropic tension of the manifold. By divergence equilibrium, the total outward spatial force exerted by the vacuum stress ($P_c$) across the minimum resolvable coordinate boundary (a sphere of Planck radius, $4\pi l_p^2$) equates to the maximum tension limit ($c^4/G$): 
\begin{align}
    P_c (4\pi l_p^2) &= \frac{c^4}{G} \nonumber \\
    P_c &= \frac{c^4}{4\pi G l_p^2} = \frac{c^7}{4\pi \hbar G^2}
\end{align}

Substituting this stiffness into the dynamic pressure equivalence ($P_c = \frac{1}{2}\rho_\tau c^2$) extracts the baseline metric density:
\begin{equation}
    \rho_\tau = \frac{2 P_c}{c^2} = \frac{1}{2\pi} \left( \frac{c^5}{\hbar G^2} \right) \tag{\ref{eq:rho_tau_corrected}}
\end{equation}

When the maximum tension ($F_p = c^4/G$) is applied across the minimal 1D tether (the Planck length, $l_p$), the baseline 1D transverse stiffness ($\kappa_\tau$) of the untensioned vacuum resolves to:
\begin{equation}
    \kappa_\tau = \frac{F_p}{l_p} = \frac{c^4}{G l_p} \tag{\ref{eq:kappa_tau_corrected}}
\end{equation}
This value represents the maximum linear tension limit of the continuum before the 1D metric fails.

\vspace{0.2 cm}

\textbf{2. The Microscopic Limit (Shear Area, Pressure \& 1D Inertia)}\\
When the 1D metric column undergoes uniaxial tension, structural brittle fracture does not resolve along a perpendicular plane. By the Tresca Yield Criterion, the material yields along the plane of maximum transverse shear, orienting at $45^\circ$ to the longitudinal axis.

For the minimum Planck boundary, the cross-sectional area is a Planck circle ($\pi l_p^2$). The $45^\circ$ shear stretches this boundary into an ellipse, establishing the minimum Planck Shear Area ($A_{shear, p}$):
\begin{equation}
    A_{shear, p} = \sqrt{2}\pi l_p^2 \tag{\ref{eq:planck_shear_area}}
\end{equation}

Distributing the Planck Force across this minimal shear area establishes the maximum structural pressure (the Planck Shear, $\tau_p$) the continuum can sustain:
\begin{equation}
    \tau_p = \frac{F_p}{\sqrt{2}\pi l_p^2} = \frac{c^4}{\sqrt{2}\pi G l_p^2} \label{eq:planck_shear_pressure}
\end{equation}

The inertial limit of the 1D topological tether (the 1D Linear Density, $\epsilon_\tau$) resolves as the 3D metric density ($\rho_\tau$) projected across this minimal transverse shear area:
\begin{align}
    \epsilon_\tau &= \rho_\tau A_{shear, p} \nonumber \\
    \epsilon_\tau &= \left( \frac{1}{2\pi} \frac{c^5}{\hbar G^2} \right) \left( \sqrt{2}\pi \frac{\hbar G}{c^3} \right) \nonumber \\
    \epsilon_\tau &= \frac{1}{\sqrt{2}} \frac{c^2}{G} \tag{\ref{eq:epsilon_tau_density}} 
\end{align}
This dimensional projection establishes the equivalent density of the 1D metric strain, isolating the inertia required to propagate a transverse wave.

\vspace{0.2 cm}

\textbf{3. The Macroscopic Limit (The Fermi Constant as a Shear Area)}\\
A stress event with a larger cross-sectional volume possesses more geometric leverage. As the fracture area increases, the local pressure required to shear the metric proportionally dilutes.

The weak interaction is defined by the Fermi Constant ($G_F$). Dimensionally, $G_F$ operates as an area (inverse energy squared). Treating $G_F$ as the shear area of the weak interaction ($A_F = G_F$) reveals that the weak interaction is the Planck Force distributed over a diluted fracture plane.

The equivalent shear pressure of the weak interaction ($\tau_F$) resolves as the Planck Force diluted by this macroscopic area:
\begin{equation}
    \tau_F = \frac{F_p}{A_F} \label{eq:fermi_shear_pressure}
\end{equation}
Because $A_F$ is orders of magnitude larger than $A_{shear, p}$, the metric tension drops identically, triggering cavitation at the laboratory-observed weak energy limits. 

\vspace{0.2 cm}

\textbf{4. The Scale-Invariant Yield Limit (Generalization)}\\
This geometric dilution establishes a scale-invariant law: the yield strength of the vacuum is a variable dependent on the area undergoing shear. The generalized shear pressure ($\tau$) for any cross-sectional fracture area ($A$) is dictated by:
\begin{equation}
    \tau(A) = \frac{F_p}{A} \label{eq:generalized_shear_pressure}
\end{equation}

Because energy scales inversely with the square root of physical area ($E \propto 1/\sqrt{A}$), the generalized energy required to fracture the metric ($E_{yield}$) at any physical scale resolves as the Planck Energy ($E_p$) scaled by the geometric ratio of the fracture:
\begin{equation}
    E_{yield}(A) = E_p \sqrt{\frac{A_{shear, p}}{A_{macroscopic}}} \label{eq:generalized_yield_energy}
\end{equation}
This establishes a sliding scale of vacuum expectation values. The background metric relaxes and dilutes its structural resistance in proportion to the geometric footprint of the event.

\vspace{0.2 cm}


\textbf{5. The Geometric Floor (Intrinsic Finitude)}\\
As distance scales contract, the interaction area tears at the Planck Shear Area ($A_{shear, p} = \sqrt{2}\pi l_p^2$). 

Evaluating the shear pressure at this boundary establishes the maximum tension limit ($\tau_{max}$):
\begin{equation}
    \tau_{max} = \frac{F_p}{A_{shear, p}} = \frac{c^4}{\sqrt{2}\pi G l_p^2} = \tau_p \label{eq:geometric_floor_tau}
\end{equation}

Because the fracture area cannot contract below this threshold, shear tension does not exceed the Planck Shear ($\tau_p$), binding the interaction energy to the finite Planck limit ($E_p$).

\vspace{0.2 cm}

\textbf{6. Analytical Demonstration: Deriving the Vacuum Expectation Value (VEV)}\\
To prove this volumetric scaling, the geometric yield equation is applied to the weak interaction fracture area ($G_F \approx 1.166 \times 10^{-5} \text{ GeV}^{-2}$). Setting $A_{macroscopic} = G_F$ isolates the yield energy:
\begin{equation}
    E_{yield}(G_F) = E_p \sqrt{\frac{\sqrt{2}\pi l_p^2}{G_F}} = E_p l_p \sqrt{\frac{\sqrt{2}\pi}{G_F}} \nonumber
\end{equation}
Because length and energy scale inversely, their product constitutes the geometric invariant ($E_p l_p = \hbar c$). In natural units ($\hbar=c=1$), the Planck constants cancel, leaving the ratio:
\begin{equation}
    E_{yield}(G_F) = \sqrt{\frac{\sqrt{2}\pi}{G_F}} \nonumber
\end{equation}
The Fermi constant defines the Vacuum Expectation Value ($v \approx 246 \text{ GeV}$) as $G_F = 1/(\sqrt{2}v^2)$. Substituting this definition back into the continuum equation extracts the energy scale of the weak force fracture:
\begin{align}
    E_{yield}(G_F) &= \sqrt{\sqrt{2}\pi (\sqrt{2}v^2)} \nonumber \\
    E_{yield}(G_F) &= \sqrt{2\pi v^2} \nonumber \\
    E_{yield}(G_F) &= v\sqrt{2\pi} \nonumber
\end{align}
The continuum yield equation outputs the Standard Model VEV ($v$), scaled by $\sqrt{2\pi}$. Geometrically, the surface area of a sphere is $4\pi r^2$, and a hemisphere is $2\pi r^2$. Taking the square root of this 2D area projection to isolate the 1D linear phase limit yields the $\sqrt{2\pi}$ multiplier.

Normalizing this 2D boundary reveals that the 1D yield strength of the vacuum at the Fermi scale equates to the VEV:
\begin{equation}
    v = \frac{E_{yield}(G_F)}{\sqrt{2\pi}} \tag{\ref{eq:vev_derivation}}
\end{equation}
This confirms that the weak interaction is the geometric dilution of the Planck Force stretching across a boundary, rendering the Ginzburg-Landau thermodynamic potential redundant.

\vspace{0.2 cm}

\textbf{7. Topological Mass Density (The $m^4$ Strain Compounding)}\\
When the metric yields, the fracture defines mass. This interaction energy derives from the baseline density ($\rho_\tau$) acting across dual, inversely proportional geometric limits: the effective volumetric strain ($V_\tau$) and the physical coordinate boundary footprint ($V_{eff}$).

The observable mass of a boundary acts as the hydrodynamic added mass ($m = \frac{1}{2} \rho_{\tau} V_{\tau}$). Multiplying by the phase limit isolates the total localized rest energy:
\begin{equation}
    E_\tau = \frac{1}{2} \rho_{\tau} V_{\tau} c^2 \label{eq:hydro_energy}
\end{equation}

Because internal circulation strikes the acoustic phase limit at the cavitation boundary ($v \to c$), the metric reaches incompressibility ($\nu_{core} = 0.5 \implies \nabla \cdot \mathbf{u} = 0$). The boundary acts as an inner geometric limit, where strain is excluded from the core and displaced into the surrounding continuum.

The phase distribution of this energy projects across the geometry of this boundary. Depending on the topological spin state, this footprint may resolve as a 3D spherical exclusion volume or a highly polarized 2D transverse shear plane (a dipole). We define the effective spatial footprint ($V_{eff}$) using the Topological Form Factor ($C_m$), a dimensionless coefficient scaling the cubed Compton radius ($V_{eff} = C_m r_c^3$). The geometric field density ($\bar{\psi}\psi$) equates to the inverse of this footprint:
\begin{equation}
    \bar{\psi} \psi = \frac{1}{V_{eff}} = \frac{1}{C_m r_c^3} \label{eq:field_inverse_volume}
\end{equation}

To transition from total energy to a continuous field interaction, the rest energy is distributed across the spatial boundary. The mass Lagrangian equates to the negative product of the total kinematic energy and the coordinate field distribution:
\begin{equation}
    \mathcal{L}_{mass} = - E_\tau (\bar{\psi} \psi) = - \frac{1}{2} \rho_{\tau} c^2 \left( \frac{V_{\tau}}{C_m r_c^3} \right) \nonumber
\end{equation}

Because the hydrodynamic added mass defines the strain volume ($V_\tau = 2m/\rho_\tau$), substituting this equivalence cancels the baseline density of the vacuum ($\rho_\tau$) entirely from the interaction:
\begin{equation}
    \mathcal{L}_{mass} = - \frac{1}{2} \rho_{\tau} c^2 \left( \frac{ \frac{2m}{\rho_\tau} }{C_m r_c^3} \right) = - \frac{mc^2}{C_m r_c^3}
\end{equation}
The Lagrangian density structurally equates to the rest energy of the defect ($mc^2$) confined within its physical geometric void.

Substituting the physical length of the Compton radius ($r_c = \hbar/mc$) extracts the universal scaling limit:
\begin{equation}
    \mathcal{L}_{mass} = - \frac{mc^2}{C_m \left(\frac{\hbar}{mc}\right)^3} = - \frac{c^5}{C_m \hbar^3} m^4 \tag{\ref{eq:absolute_mass_lagrangian}}
\end{equation}

Evaluating this boundary density for the proton mass ($m_p$) computes a mechanical pressure on the order of $\sim 10^{38} \text{ Pa}$. While resolving the specific geometric coefficient ($C_m$) requires exact topological modeling of the species-specific dipole shear plane, the Lagrangian density compounds to the fourth power ($\mathcal{L}_{mass} \propto m^4$). Because the effective strain volume is proportional to mass ($V_\tau \propto m$) and the coordinate footprint is inversely proportional to the cube of the mass ($V_{eff} \propto 1/m^3$), the resulting energy density compounds. Heavy particles possess a high strain displacement confined around a compressed inner boundary, dynamically generating the exponentially higher $m^4$ Lagrangian energy densities of Quantum Field Theory from continuum yield limits.

\vspace{0.2 cm}

\textbf{8. The Spin-Zero Geometric Constraint}\\
Because the core generated by the fracture resolves as an irrotational volumetric expansion ($\mathbf{u} = -\nabla \Phi$), the vorticity ($\boldsymbol{\Omega}$) equals zero:
\begin{equation}
    \boldsymbol{\Omega} = \nabla \times (-\nabla \Phi) = 0 \nonumber
\end{equation}
Angular momentum (Spin, $\mathbf{S}$) is the volume integral of vorticity. Substituting this irrotational constraint cancels the vector cross-product:
\begin{align}
    \mathbf{S} &\propto \iiint_V \rho_{\tau} (\mathbf{r} \times \boldsymbol{\Omega}) \, dV \nonumber \\
    \mathbf{S} &= 0 \label{eq:spin_zero_proof}
\end{align}
The particle generated by the metric yield (the Higgs equivalent) possesses a spin of zero. It is a kinematic mandate for a longitudinal P-wave cavitation, driven by the yield limit of the background geometry.

\begin{tcolorbox}[recallbox, title={\ref{box:continuum_constants}. Constants \& Scale Invariance}]
\begin{align}
    \rho_{\tau} &= \frac{1}{2\pi} \frac{c^5}{\hbar G^2} \approx 8.21 \times 10^{95} \text{ kg/m}^3 \tag{\ref{eq:rho_tau_corrected}} \\
    \kappa_{\tau} &= \frac{c^4}{G l_p} \approx 7.49 \times 10^{78} \text{ N/m} \tag{\ref{eq:kappa_tau_corrected}} \\
    A_{shear, p} &= \sqrt{2}\pi l_p^2 \approx 1.16 \times 10^{-69} \text{ m}^2 \tag{\ref{eq:planck_shear_area}} \\
    \epsilon_{\tau} &= \frac{1}{\sqrt{2}} \frac{c^2}{G} \approx 9.52 \times 10^{26} \text{ kg/m} \tag{\ref{eq:epsilon_tau_density}} \\
    \tau(A) &= \frac{F_p}{A} \tag{\ref{eq:generalized_shear_pressure}} \\
    E_{yield}(G_F) &= E_p \sqrt{\frac{\sqrt{2}\pi l_p^2}{G_F}} \implies v = \frac{E_{yield}(G_F)}{\sqrt{2\pi}} \approx 246 \text{ GeV} \tag{\ref{eq:vev_derivation}} \\
    \mathcal{L}_{mass} &= - \frac{c^5}{C_m \hbar^3} m^4 \implies \mathcal{L}_{mass}(m_p) \sim 10^{38} \text{ Pa} \tag{\ref{eq:absolute_mass_lagrangian}} \\
    \mathbf{S} &= 0 \tag{\ref{eq:spin_zero_proof}}
\end{align}
\end{tcolorbox}

% ----------------------------------------------------------------------
\section{The Viscoelastic Modulus Ratio ($\alpha$) \& Topological Impedance}
\label{sec:deriv_fine_structure2}

The metric operates as a viscoelastic continuum governed by kinematic capacities. The continuum supports both longitudinal compression and transverse shear. The Fine-Structure Constant ($\alpha \approx 1/137$) functions as the Viscoelastic Modulus Ratio, defined as the fractional capacity of the metric to sustain transverse shear relative to its longitudinal bulk capacity ($\Pi_{shear} = \alpha P_{ambient}$). 

Because the metric's shear capacity represents a boundary condition, the transfer of energy between the vorticity ($\boldsymbol{\Omega} = \nabla\lambda \times \nabla\beta_\Omega$) of a boundary and the propagating transverse shear ($\mathbf{v}_{\perp}$) of the background metric is heavily throttled.

\vspace{0.2 cm}

\textbf{1. Transverse Shear \& Kinematic Throttling:} \\
Because the boundary acts as an imperfect acoustic container, the spatial rotation of the phase ($\nabla \times (\lambda \nabla \beta_\Omega)$) strains the surrounding metric. The baseline restorative force of rotation aligns with the limit of the quantum harmonic oscillator ($E = \hbar c / r \implies F = -dE/dr$):
\begin{equation}
    F_{\circlearrowleft}(r) = \frac{\hbar c}{r^2} \tag{\ref{eq:internal_rotational_force}}
\end{equation}

A fraction of this rotational energy undergoes a secondary transduction, shearing into the surrounding metric to manifest as transverse strain. This leakage is throttled by the Viscoelastic Modulus Ratio ($\alpha$). 

The electromagnetic force ($F_{EM}$) emerges as the rotational force ($F_{\circlearrowleft}$) scaled by this transverse limit:
\begin{align}
    F_{EM}(r) &= \alpha \cdot F_{\circlearrowleft}(r) \nonumber \\
    F_{EM}(r) &= \alpha \frac{\hbar c}{r^2} \tag{\ref{eq:electromagnetic_leakage}}
\end{align}

Opposing phase orientations induces compound velocity between boundaries, dropping the static pressure below the ambient field to create Bernoulli attraction. Identical phase orientations induce stagnation at the interaction plane; this stall retains the static pressure at the ambient baseline. Because the surrounding perturbed metric remains in continuous motion (retaining a lower static pressure), this stalled pressure boundary drives the defects apart.

\vspace{0.2 cm}


\textbf{2. Macroscopic vs. Topological Impedance:} \\
The resistance of the metric to this energy transfer resolves as the ratio of two boundaries. The acoustic transverse impedance of the unperturbed metric ($Z_{macro}$) equates to the product of the 1D transverse density ($\epsilon_\tau$) and the coordinate phase limit ($c$):
\begin{equation}
    Z_{macro} = \epsilon_{\tau} c \tag{\ref{eq:macroscopic_impedance}}
\end{equation}

This mechanical baseline ($\sim 2.86 \times 10^{35} \text{ kg/s}$) is the physical origin of the characteristic impedance of free space ($Z_0 \approx 376.7 \text{ } \Omega$). Under the Electro-Mechanical Mobility mapping (derived in Section \ref{sec:deriv_em_lc_isomorphism}), converting this mechanical limit to SI electrical units requires dividing by the dimensional ratio of the metric density to the SI permittivity ($\epsilon_\tau / \epsilon_0$). 

Because mechanical impedance maps to electrical admittance ($Y_0$), dividing the baseline by this dimensional ratio yields the admittance of free space:
\begin{equation}
    Y_0 = \frac{Z_{macro}}{(\epsilon_\tau / \epsilon_0)} = \frac{\epsilon_\tau c}{\epsilon_\tau / \epsilon_0} = \epsilon_0 c
\end{equation}
Inverting this admittance isolates the classical characteristic impedance:
\begin{equation}
    Z_0 = \frac{1}{Y_0} = \frac{1}{\epsilon_0 c} \approx 376.7 \text{ } \Omega
\end{equation}

\vspace{0.2 cm}


\textbf{3. Topological Impedance ($Z_{topo}$) \& The von Klitzing Constant:} \\
Because the electromagnetic field operates as the transverse leakage of the core's circulation, the mechanical impedance of the boundary ($Z_{topo}$) is coupled to the impedance of the vacuum ($Z_{macro}$). 

The Viscoelastic Modulus Ratio ($\alpha$) dictates this transverse leakage fraction. Because the boundary constitutes a 2D closed manifold (a sphere) embedded within the 3D metric, the ratio of its impedance to the macroscopic impedance is governed by this fraction, scaled by $1/2$ to account for the geometric bisection of the shear plane:
\begin{equation}
    \alpha = \frac{1}{2} \left( \frac{Z_{topo}}{Z_{macro}} \right) \tag{\ref{eq:tensor_coupling_derivation}}
\end{equation}

Rearranging the modulus ratio isolates the mechanical resistance of the boundary ($Z_{topo}$):
\begin{align}
    Z_{topo} &= 2 \alpha Z_{macro} \nonumber \\
    Z_{topo} &= 2 \alpha \epsilon_{\tau} c \tag{\ref{eq:topological_impedance}}
\end{align}
Because $\alpha$ is a small fraction ($\approx 1/137$), the boundary possesses significantly lower mechanical resistance than the vacuum ($Z_{topo} \approx 4.16 \times 10^{33} \text{ kg/s}$). 

Under the Electro-Mechanical Mobility mapping (derived in Section \ref{sec:deriv_em_lc_isomorphism}), electrical impedance maps to admittance (the inverse of impedance). Converting impedance to SI electrical resistance requires the dimensional calibration scalar ($k_z = \epsilon_\tau / \epsilon_0 \approx 1.075 \times 10^{38} \text{ kg}^2\cdot\text{m}^2/\text{s}^2\cdot\text{C}^2$):
\begin{equation}
    R_K  = \frac{\epsilon_\tau / \epsilon_0}{2\alpha \epsilon_\tau c} = \frac{1}{2\alpha \epsilon_0 c} \approx 25,812 \text{ } \Omega \label{eq:von_klitzing_derived}
\end{equation}
This derives the von Klitzing constant from the transverse acoustic impedance of the metric, demonstrating that electrical resistance is the scaled mobility of a boundary.

\begin{tcolorbox}[recallbox, title={\ref{box:fine_structure2}. The Viscoelastic Modulus Ratio ($\alpha$) \& Topological Impedance}]
\begin{align}
    F_{\circlearrowleft}(r) &= \frac{\hbar c}{r^2} \tag{\ref{eq:internal_rotational_force}} \\
    F_{EM}(r) &= \alpha \frac{\hbar c}{r^2} \tag{\ref{eq:electromagnetic_leakage}} \\
    Z_{macro} &= \epsilon_{\tau} c \tag{\ref{eq:macroscopic_impedance}} \\
    \alpha &= \frac{1}{2} \left( \frac{Z_{topo}}{Z_{macro}} \right) \tag{\ref{eq:tensor_coupling_derivation}} \\
    Z_{topo} &= 2 \alpha \epsilon_{\tau} c \tag{\ref{eq:topological_impedance}} \\
    R_K &= \frac{\epsilon_\tau / \epsilon_0}{Z_{topo}} \approx 25,812 \text{ } \Omega \tag{\ref{eq:von_klitzing_derived}}
\end{align}
\end{tcolorbox}

% ----------------------------------------------------------------------
\section{The Kinematic Origin of Charge \& The Magnetic Anomaly}
\label{sec:deriv_kinematic_charge}

Electromagnetism resolves as shear.

\textbf{1. The Shear Transduction Limit:} \\
The kinetic energy of the strain circulation ($v_\phi$) radiating across the equatorial perimeter transduces into transverse shear. The conversion of kinetic rotation into an elastic shear wave resolves via energy equivalence:
\begin{equation}
    \frac{1}{2}\rho_{\tau} v_\phi^2 = \frac{1}{2}\epsilon_{\tau} (\nabla \times \mathbf{v}_{\perp})^2 \tag{\ref{eq:shear_transduction_limit}}
\end{equation}
The kinetic rotation of the density ($\rho_{\tau}$) equates to the transverse shear of the 1D linear density ($\epsilon_{\tau}$).

\textit{Electromagnetic Equivalence:} The transverse velocity ($\mathbf{v}_\perp$) maps to the classical magnetic vector potential ($\mathbf{A}_\perp$), mapping shear energy to magnetic field energy.

\vspace{0.2 cm}

\textbf{2. Continuum Charge ($e_{\tau}$):} \\
Metric charge ($e_{\tau}$) acts as the kinematic momentum displaced by the shear stress. It derives as the closed surface integral of the transverse spatial vorticity ($\nabla \times \mathbf{v}_{\perp}$) scaled by the 1D linear density ($\epsilon_{\tau}$):
\begin{equation}
    e_{\tau} = \int_S \epsilon_{\tau} (\nabla \times \mathbf{v}_{\perp}) \cdot d\mathbf{A} \tag{\ref{eq:charge_displacement}}
\end{equation}

The electrostatic force dictates the interaction of two displacements across the linear density of the metric ($F = e_{\tau}^2 / 4\pi \epsilon_{\tau} r^2$). Equating this force to the coupling limit ($F = \alpha \hbar c / r^2$), isolates the magnitude of the displaced momentum:
\begin{align}
    \frac{e_{\tau}^2}{4\pi \epsilon_{\tau} r^2} &= \alpha \frac{\hbar c}{r^2} \nonumber \\
    e_{\tau}^2 &= 4\pi \epsilon_{\tau} \hbar c \alpha \nonumber \\
    e_{\tau} &= \sqrt{4\pi \epsilon_{\tau} \hbar c \alpha} \tag{\ref{eq:absolute_charge_derived}}
\end{align}
Solving this limit via the transverse wave constants isolates the magnitude of charge: $e_{\tau} \approx 1.66 \text{ kg}\cdot\text{m/s}$.

\vspace{0.2 cm}

\textbf{3. The Elimination of the Coulomb (Charge as Momentum):} \\
By substituting the dimensions of the continuum variables ($[\epsilon_{\tau}] = \text{kg/m}$, $[\hbar] = \text{kg}\cdot\text{m}^2/\text{s}$, $[c] = \text{m/s}$) into the derivation, the Coulomb is eliminated. Charge reduces to kinematic momentum:
\begin{equation}
    [e_{\tau}] = \sqrt{ \left(\frac{\text{kg}}{\text{m}}\right) \left(\frac{\text{kg}\cdot\text{m}^2}{\text{s}}\right) \left(\frac{\text{m}}{\text{s}}\right) } = \text{kg}\cdot\text{m/s} \label{eq:charge_dimension}
\end{equation}
Substituting this momentum back into the electrostatic force equation ($F = e_{\tau}^2 / 4\pi \epsilon_{\tau} r^2$) correctly outputs Newtons ($\text{kg}\cdot\text{m/s}^2$). The directional density ($\epsilon_{\tau}$) and spherical geometry ($4\pi$) cancel symmetrically, reducing the interaction to the established transverse coupling force ($F = \alpha \hbar c / r^2$).

\vspace{0.2 cm}

\textbf{4. The Dirac Moment ($\mu_{B\tau}$):} \\
Mapping the momentum displacement ($e_{\tau}$) and the angular momentum limit ($\hbar$) against the boundary mass ($m$) defines the 2D baseline magnetic moment:
\begin{equation}
    \mu_{B\tau} = \frac{e_{\tau} \hbar}{2m} \tag{\ref{eq:bohr_magneton}}
\end{equation}

\vspace{0.2 cm}

\textbf{5. The Toroidal Phase Distribution \& True Magnetic Moment:} \\
Because the torus operates as a closed 3D manifold, the transverse coupling limit ($\alpha$) is continuously distributed across the entire rotational phase loop of the boundary. 

The geometric anomaly resolves as the transverse coupling distribution divided by the complete angular phase cycle ($2\pi$). Adding this distributed phase strain to the Dirac loop extracts the magnetic moment:
\begin{equation}
    a_e = \frac{\alpha}{2\pi} \implies \mu_{true} = \mu_{B\tau} \left( 1 + \frac{\alpha}{2\pi} \right) \tag{\ref{eq:true_magnetic_moment}}
\end{equation}

\begin{tcolorbox}[recallbox, title={\ref{box:kinematic_charge}. The Kinematic Origin of Charge \& The Magnetic Anomaly}]
\begin{align}
    \frac{1}{2}\rho_{\tau} v_\phi^2 &= \frac{1}{2}\epsilon_{\tau} (\nabla \times \mathbf{v}_{\perp})^2 \tag{\ref{eq:shear_transduction_limit}} \\
    e_{\tau} &= \int_S \epsilon_{\tau} (\nabla \times \mathbf{v}_{\perp}) \cdot d\mathbf{A} \tag{\ref{eq:charge_displacement}} \\
    e_{\tau} &= \sqrt{4\pi \epsilon_{\tau} \hbar c \alpha} \approx 1.66 \text{ kg}\cdot\text{m/s} \tag{\ref{eq:absolute_charge_derived}} \\
    \mu_{B\tau} &= \frac{e_{\tau} \hbar}{2m} \tag{\ref{eq:bohr_magneton}} \\
    a_e = \frac{\alpha}{2\pi} &\implies \mu_{true} = \mu_{B\tau} \left( 1 + \frac{\alpha}{2\pi} \right) \tag{\ref{eq:true_magnetic_moment}}
\end{align}
\end{tcolorbox}
% ----------------------------------------------------------------------
% ----------------------------------------------------------------------
\section{Coulomb's Law (Spherical Bernoulli \& Phase Flux)}
\label{sec:deriv_coulomb_bernoulli}

Electric charge defines the rotational chirality of an acoustic cavity. The electrostatic force mirrors the pressure gradient generated by the Bernoulli effect.

\vspace{0.2 cm}

\textbf{1. The Kinematic Velocity Field \& Volumetric Action:} \\
Because a defect operates as a steady-state flux (an engine), it processes volumetric circulation ($\Phi_{\circlearrowleft}$). To conserve Eulerian volumetric continuity, the defect cannot create or destroy metric volume; therefore, the net scalar sum of its kinematic action must resolve to zero. This circulation ($\Phi_{\circlearrowleft}$), the longitudinal radial intake/exhaust ($\Phi_{\parallel}$), and the transverse spill rate ($\Phi_{\perp}$) sum to zero.
\begin{equation}
    \Phi_{\circlearrowleft} + \Phi_{\parallel} + \Phi_{\perp} = 0 \label{eq:flux_decomposition}
\end{equation}
The longitudinal flux operates as a circulation channel traversing the central axis and wrapping the exterior geometry. When perfectly balanced, this closed intake/exhaust loop acts as the magnetic dipole. Any asymmetric net longitudinal flux acts as an action-reaction engine, manifesting as propulsion (inertia). 

The long-range electrostatic field emerges from the transverse spill rate ($\Phi_{\perp}$) shearing against the surrounding metric. To conserve the action as it propagates outward, the transverse circulation velocity ($v_\perp$) must decrease proportionally to the spherical surface area ($A = 4\pi r^2$): 
\begin{equation}
    v_{\perp}(r) = \frac{\Phi_{\perp}}{4\pi r^2} \label{eq:transverse_velocity_dilution}
\end{equation}

\textbf{2. Transverse Dynamic Pressure \& Binomial Expansion:} \\
When two transverse kinematic fields superimpose ($\mathbf{v}_{\perp} = \mathbf{v}_{\perp(1)} + \mathbf{v}_{\perp(2)}$), the baseline kinetic energy density (dynamic pressure) of the metric expands via the binomial theorem against the unperturbed fluid density ($\rho_{\tau}$):
\begin{equation}
    P_{dyn} = \left[ \frac{1}{2}\rho_{\tau} v_{\perp(1)}^2 \right] + \left[ \frac{1}{2}\rho_{\tau} v_{\perp(2)}^2 \right] + \left[ \rho_{\tau} (\mathbf{v}_{\perp(1)} \cdot \mathbf{v}_{\perp(2)}) \right]
\end{equation}
The independent squared terms represent the isolated self-energy (mass equivalent) of the individual topological defects. The geometric dot product of the cross-term dictates the active interaction energy (attraction/repulsion) circulating around adjacent boundaries.

\textit{Attraction (Opposite Winding Numbers):} If the defects possess opposite chiralities, their adjacent transverse circulations align ($\mathbf{v}_{\perp(1)} \cdot \mathbf{v}_{\perp(2)} > 0$). This alignment drops the static pressure ($P_{static} \downarrow$) between them:
\begin{equation}
    \mathbf{v}_{\perp(1)} \cdot \mathbf{v}_{\perp(2)} > 0 \implies P_{static} \downarrow \tag{\ref{eq:attraction}}
\end{equation}

\textit{Repulsion (Identical Winding Numbers):} If the defects possess identical chiralities, their adjacent transverse circulations oppose each other ($\mathbf{v}_{\perp(1)} \cdot \mathbf{v}_{\perp(2)} < 0$). This collision results in stagnation (the static pressure remains at the ambient baseline):
\begin{equation}
    \mathbf{v}_{\perp(1)} \cdot \mathbf{v}_{\perp(2)} < 0 \implies P_{static} \uparrow \tag{\ref{eq:repulsion}}
\end{equation}

\textbf{3. The Volume Integral \& Compliance-Modified Force:} \\
The isolated interaction pressure deficit ($\Delta P_{int}$) resolves from the transverse cross-term:
\begin{equation}
    \Delta P_{int}(r) = - \rho_{\tau} (v_{\perp(1)} v_{\perp(2)}) = - \frac{\rho_{\tau} \Phi_{\perp(1)} \Phi_{\perp(2)}}{16\pi^2 r^4}
\end{equation}

The total raw interaction potential energy ($U_{raw}$) acts as the volume integral ($\iiint \Delta P_{int} \, dV$) over the spherical metric from the boundary ($r$) to macroscopic infinity, scaled by a factor of $1/2$ to correct for mutual interaction double-counting:
\begin{align}
    U_{raw}(r) &= \frac{1}{2} \int_{r}^{\infty} \left( - \frac{\rho_{\tau} \Phi_{\perp(1)} \Phi_{\perp(2)}}{16\pi^2 r'^4} \right) 4\pi r'^2 dr' \nonumber \\
    U_{raw}(r) &= - \frac{\rho_{\tau} \Phi_{\perp(1)} \Phi_{\perp(2)}}{8\pi r}
\end{align}

Because the Bernoulli pressure gradient operates as a mechanical Surface Force, the modified Navier-Stokes equation mandates that its spatial application is throttled by the transverse coupling constant ($\alpha$). The restoring force ($\mathbf{F}_{restore}$) translates to the spatial gradient ($-\nabla U_{raw}$) multiplied by this transverse limit:
\begin{equation}
    \mathbf{F}_{restore} = \alpha (-\nabla U_{raw}) = - \frac{\alpha \rho_{\tau} \Phi_{\perp(1)} \Phi_{\perp(2)}}{8\pi r^2} \mathbf{\hat{r}} \tag{\ref{eq:bernoulli_force}}
\end{equation}

\textbf{4. Elimination of the Coupling Constant ($\alpha$):} \\
The continuum force mapping to the Coulomb limit resolves via the action:
\begin{equation}
    F_e = \frac{\alpha \hbar c}{r^2}
\end{equation}
Equating this force to the kinematic limit eliminates parameterized constants ($k_e, \epsilon_0, q$):
\begin{equation}
    \frac{\alpha \rho_{\tau} \Phi_{\perp(e)}^2}{8\pi r^2} = \frac{\alpha \hbar c}{r^2} \tag{\ref{eq:force_equivalence}}
\end{equation}
Because the fine-structure constant ($\alpha$) operates as the coupling limit, the factor cancels:
\begin{align}
    \Phi_{\perp(e)}^2 &= \frac{8\pi \hbar c}{\rho_{\tau}}
\end{align}

\textbf{5. The Continuous Volumetric Phase Flux ($\Phi_{\perp(e)}$):} \\
Taking the square root isolates the volumetric phase flux of the charge boundary:
\begin{equation}
    \Phi_{\perp(e)} = \sqrt{\frac{8\pi \hbar c}{\rho_{\tau}}} 
\end{equation}
Substituting the metric density ($\rho_{\tau} = \frac{1}{2\pi} \frac{c^5}{\hbar G^2}$) into the equivalence isolates the identity: 
\begin{align}
    \Phi_{\perp(e)} &= \sqrt{ \frac{8\pi \hbar c}{\frac{1}{2\pi} \frac{c^5}{\hbar G^2}} } \nonumber \\ 
    \Phi_{\perp(e)} &= 4\pi \frac{\hbar G}{c^2} \approx 9.84 \times 10^{-61} \text{ m}^3/\text{s} \tag{\ref{eq:charge_flux_reduced}} 
\end{align}
Electric charge acts as a continuous volumetric flow rate, dictated by density and action limits. The transverse coupling constant ($\alpha$) is eliminated from the geometry of the charge boundary, demonstrating that it operates as a surface force transmission limit.

\vspace{0.2 cm}

\textbf{6. The Kinematic Force Constant ($\mathcal{K}_\tau$):} \\
Extracting the constants governing the Compliance-Modified Bernoulli force isolates the Force Constant ($\mathcal{K}_\tau$):
\begin{equation}
    \mathcal{K}_\tau = \frac{\alpha \rho_{\tau} \Phi_{\perp(e)}^2}{8\pi}
\end{equation}
Substituting the identity ($\Phi_{\perp(e)}^2 = 8\pi \hbar c / \rho_{\tau}$) allows the density ($\rho_{\tau}$) and the spherical geometry ($8\pi$) to cancel symmetrically. The force constant reduces to the transverse coupling of the action limit:
\begin{equation}
    \mathcal{K}_\tau = \alpha \hbar c \tag{\ref{eq:force_constant_reduced}}
\end{equation}
Substituting the mechanical dimensions of the variables ($[\alpha] = 1$, $[\hbar] = \text{kg}\cdot\text{m}^2/\text{s}$, $[c] = \text{m/s}$) isolates the structural units:
\begin{align}
    [\mathcal{K}_\tau] &= (1) \cdot \left(\frac{\text{kg}\cdot\text{m}^2}{\text{s}}\right) \cdot \left(\frac{\text{m}}{\text{s}}\right) \nonumber \\
    [\mathcal{K}_\tau] &= \frac{\text{kg}\cdot\text{m}^3}{\text{s}^2} \equiv \text{N}\cdot\text{m}^2
\end{align}
When this constant is divided by the spatial variance ($r^2$, measured in $\text{m}^2$) in the Coulomb force equation, the square meters cancel, outputting Newtons of force ($F_e = \mathcal{K}_\tau / r^2$). 
This unifies Coulomb's constant and the elementary charge squared ($k_e e^2 \equiv \mathcal{K}_\tau$), demonstrating the interaction resolves entirely from the baseline angular momentum ($\hbar$), phase limit ($c$), and transverse coupling ($\alpha$).

\begin{tcolorbox}[recallbox, title={\ref{box:coulomb_bernoulli}. Coulomb's Law (Spherical Bernoulli \& Phase Flux)}]
\begin{align}
    \mathbf{v}_{\perp(1)} \cdot \mathbf{v}_{\perp(2)} > 0 &\implies P_{static} \downarrow \tag{\ref{eq:attraction}} \\
    \mathbf{v}_{\perp(1)} \cdot \mathbf{v}_{\perp(2)} < 0 &\implies P_{static} \uparrow \tag{\ref{eq:repulsion}} \\
    \mathbf{F}_{fluid} &= - \frac{\alpha \rho_{\tau} \Phi_{\perp(1)} \Phi_{\perp(2)}}{8\pi r^2} \mathbf{\hat{r}} \tag{\ref{eq:bernoulli_force}} \\
    \frac{\alpha \rho_{\tau} \Phi_{\perp(e)}^2}{8\pi r^2} &= \frac{\alpha \hbar c}{r^2} \tag{\ref{eq:force_equivalence}} \\
    \Phi_{\perp(e)} = 4\pi \frac{\hbar G}{c^2} &\approx 9.84 \times 10^{-61} \text{ m}^3/\text{s} \tag{\ref{eq:charge_flux_reduced}} \\ % :::LLM_EDIT::: Purged 2/3 artifact and updated numeric value
    \mathcal{K}_\tau = \alpha \hbar c &\approx 2.307 \times 10^{-28} \text{ N}\cdot\text{m}^2 \tag{\ref{eq:force_constant_reduced}}
\end{align}
\end{tcolorbox}

\newpage

% ======================================================================
% PHASE V: TENSOR TRANSLATION & FORCE UNIFICATION
% ======================================================================
\part*{PHASE V: TENSOR TRANSLATION \& FORCE UNIFICATION}






% ----------------------------------------------------------------------
\section{Tensor Evolution \& Geometric Bending (The Rosetta Stone)}
\label{sec:deriv_rosetta_stone}

Tensor curvature and geometric gauge evolution resolve as the kinematic velocity field, spatial pressure gradients, vorticity transport, and structural incompatibility of the continuous metric.

\textbf{1. The Metric Tensor ($[\boldsymbol{g}_{\mu\nu}]$) \& The Acoustic Horizon:} \\
The metric tensor acts as the Unruh-Visser Acoustic Metric. The curved spacetime of General Relativity translates to the degradation of Eulerian compliance ($\alpha_s$). 

The spacetime geometry projects as a $4 \times 4$ diagonal tensor. In classical relativity, the diagonal components ($g_{tt}, g_{rr}, g_{\theta\theta}, g_{\phi\phi}$) govern time dilation, radial stretching, and the spherical coordinate surface. 

These components map to functions of the Total Scalar Admittance ($\alpha_s$). Whether modeling a black hole (compliance dropping to zero at $r_s$) or the macroscopic cosmos (compliance dropping to zero at the causal boundary $R_h = c \cdot t$), the tensor dictates identical geometry:
\begin{equation}
    [\boldsymbol{g}_{\mu\nu}] = \begin{bmatrix}
    -\alpha_s^2 & 0 & 0 & 0 \\
    0 & \alpha_s^{-2} & 0 & 0 \\
    0 & 0 & r^2 & 0 \\
    0 & 0 & 0 & r^2 \sin^2(\theta)
    \end{bmatrix} \tag{\ref{eq:rosetta_metric_matrix}}
\end{equation}

Expanding the compliance limits ($\alpha_s^2 = 1 - P_{dyn}/P_c$), the invariant spacetime interval ($ds^2 = g_{\mu\nu} dx^\mu dx^\nu$) translates into the classical acoustic wave equation for a tensioned medium:
\begin{equation}
    ds^2 \equiv -\left( 1 - \frac{P_{dyn}}{P_c} \right) c^2 dt^2 + \left( 1 - \frac{P_{dyn}}{P_c} \right)^{-1} dr^2 + r^2 d\Omega^2 \label{eq:rosetta_acoustic_metric}
\end{equation}

\textbf{2. The Covariant Derivative (Material Evolution):} \\
The Covariant Derivative ($\nabla_\mu$) tracks the spatial evolution of a tensor field ($A^\nu$). It combines the partial derivative and the connection terms:
\begin{equation}
    \nabla_\mu A^\nu = \partial_\mu A^\nu + \Gamma^\nu_{\mu\lambda} A^\lambda 
\end{equation}
The transport of a property across the manifold acts as the Eulerian Material Derivative ($D/Dt$):
\begin{equation}
    \frac{D A_i}{Dt} = \frac{\partial A_i}{\partial t} + v_j \partial_j A_i 
\end{equation}
The Covariant Derivative maps to the Material Derivative. The tensor is sheared and rotated because it propagates along the velocity field.
\begin{equation}
    \nabla_\mu \equiv \frac{D}{Dt} \tag{\ref{eq:rosetta_covariant}}
\end{equation}

\textbf{3. The Christoffel Symbols (Acceleration \& Shear):} \\
The geometric connection bending the metric tensor projects as both the acceleration and the deformation rate of the continuum. 

The temporal connection driving gravitational acceleration operates as the spatial derivative of the metric ($\Gamma^i_{00} = -\frac{1}{2}\partial^i g_{00}$). Because the temporal metric resolves to the dimensionless pressure capacity ($g_{00} = -1 + P_{dyn}/P_c$), its spatial derivative isolates the Bernoulli pressure gradient: 
\begin{align}
    \Gamma^i_{00} &= -\frac{1}{2}\partial^i \left( -1 + \frac{P_{dyn}}{P_c} \right) \nonumber \\
    \Gamma^i_{00} &= -\frac{1}{2 P_c} \partial^i P_{dyn} \nonumber
\end{align}
Substituting the acoustic yield pressure ($P_c = \frac{1}{2}\rho_{\tau} c^2$) and recalling that the static pressure gradient balances the dynamic pressure ($\nabla^i P_{static} = -\partial^i P_{dyn}$) isolates the Navier-Stokes acceleration ($a^i = -c^2 \Gamma^i_{00}$):
\begin{align}
    -c^2 \Gamma^i_{00} &= \frac{c^2}{2 \left(\frac{1}{2}\rho_{\tau} c^2\right)} \partial^i P_{dyn} \nonumber \\
    -c^2 \Gamma^i_{00} &= \frac{1}{\rho_{\tau}} \partial^i P_{dyn} \nonumber \\
    c^2 \Gamma^i_{00} &\equiv -\frac{1}{\rho_{\tau}} \nabla^i P_{static} \tag{\ref{eq:rosetta_christoffel_pressure}}
\end{align}

Conversely, the spatial connections translate to the velocity gradient tensor. 
\begin{equation}
    \Gamma^k_{ij} \equiv \frac{1}{c} \partial_j v_i \tag{\ref{eq:rosetta_christoffel}}
\end{equation}
The Christoffel symbols map to the continuous pressure and velocity gradients.

\vspace{0.2 cm}

\textbf{4. Riemann Curvature (Saint-Venant Compatibility):} \\
The failure of a manifold to close a coordinate loop under parallel transport projects as spatial incompatibility. The displacement ($u_i$) generates a symmetric Cauchy strain tensor ($\varepsilon_{ij}$):
\begin{equation}
    \varepsilon_{ij} = \frac{1}{2}(\partial_j u_i + \partial_i u_j) 
\end{equation}
The metric tensor ($g_{ij}$) behaves as the flat Euclidean background ($\delta_{ij}$) deformed by this strain ($g_{ij} = \delta_{ij} + 2\varepsilon_{ij}$). 

Substituting the metric strain equivalence into the Christoffel symbols yields the connection as the spatial derivative of the strain:
\begin{equation}
    \Gamma_{ijk} = \partial_j \varepsilon_{ik} + \partial_i \varepsilon_{jk} - \partial_k \varepsilon_{ij} 
\end{equation}
Substituting this connection into the linear derivative terms of the Riemann tensor ($R_{ijkl} = \partial_i \Gamma_{jkl} - \partial_j \Gamma_{ikl}$) generates the second-order spatial derivatives. Subtracting the symmetric cross-terms yields:
\begin{equation}
    R_{ijkl} = \partial_j \partial_l \varepsilon_{ik} + \partial_i \partial_k \varepsilon_{jl} - \partial_j \partial_k \varepsilon_{il} - \partial_i \partial_l \varepsilon_{jk} \nonumber
\end{equation}

An arbitrary strain field ($\varepsilon_{ij}$) is constrained via the Saint-Venant Compatibility Condition ($d \wedge \boldsymbol{\varepsilon} \wedge d = 0$). The incompatibility tensor ($\text{Inc}(\boldsymbol{\varepsilon})$) generates the identical combination of second-order spatial derivatives. A cavitation maps as a Volterra dislocation. Because the strain cannot close without tearing, Riemann curvature conforms to strain incompatibility: 
\begin{equation}
    R_{ijkl} \equiv \text{Inc}(\boldsymbol{\varepsilon})_{ijkl} \tag{\ref{eq:rosetta_riemann}}
\end{equation}

\textbf{5. Ricci \& Weyl Curvature (Volumetric vs. Deviatoric Strain):} \\
In General Relativity, the Riemann curvature tensor decomposes into the Ricci tensor ($R_{\mu\nu}$) and the Weyl tensor ($C_{\mu\nu\rho\sigma}$). This translates to the elastodynamic decomposition of strain.

The Ricci tensor quantifies the change in spatial volume. It maps to the Isotropic Volumetric Strain (Dilatation), dictating the compactification of density: 
\begin{equation}
    R_{\mu\nu} \equiv \text{Inc}\left(\frac{1}{3}\text{Tr}(\varepsilon)\delta_{ij}\right) \implies \text{Volumetric Strain} \tag{\ref{eq:rosetta_ricci}}
\end{equation}
The Weyl tensor captures tidal curvature that conserves volume. It translates to Deviatoric Shear Strain. Gravitational waves behave as macroscopic transverse deviatoric shear waves: 
\begin{equation}
    C_{\mu\nu\rho\sigma} \equiv \text{Inc}(\varepsilon'_{ij}) \implies \text{Deviatoric Shear} \tag{\ref{eq:rosetta_weyl}}
\end{equation}

\textbf{6. The Stress-Energy \& Einstein Tensors (Mechanical Equilibrium):} \\
The source of spacetime curvature ($T_{\mu\nu}$) resolves as the Cauchy Stress Tensor ($\sigma_{\mu\nu}$) combined with the convective momentum flux of the continuous fluid. Because the Cauchy tensor maps normal stress to hydrostatic pressure ($\sigma_{ij} = -P \delta_{ij} + \tau_{ij}$), it defines the scalar Bernoulli pressure deficit, internal fluid shear, and kinematic momentum density at any localized coordinate.
\begin{equation}
    T_{\mu\nu} \equiv \sigma_{\mu\nu} + \rho u_\mu u_\nu \tag{\ref{eq:rosetta_stress_energy}}
\end{equation}

The Einstein Field Equation ($G_{\mu\nu} = \frac{8\pi G}{c^4} T_{\mu\nu}$) equates structural geometry ($G_{\mu\nu}$) to thermodynamic stress ($T_{\mu\nu}$). This models as Newton's Third Law of mechanical equilibrium applied to an elastic medium.  Internal geometric strain must equal the applied stress, scaled by the geometric tension limit ($c^4/8\pi G$):
\begin{equation}
    G_{\mu\nu} \equiv \text{Hookean Strain Balance of } T_{\mu\nu} \tag{\ref{eq:rosetta_einstein_tensor}}
\end{equation}

\textbf{7. The Faraday Tensor (Kinematic Vorticity):} \\
The electromagnetic gauge maps to the Vorticity Tensor ($\Omega_{\mu\nu}$). With the four-potential operating as the velocity ($A_\mu = c u_\mu$), the tensor isolates the antisymmetric covariant derivative of the velocity field. Because the spacetime gradient ($\partial_\mu$) has dimensions of inverse length ($[1/\text{m}]$) and the normalized 4-velocity ($u_\mu$) is dimensionless, their derivative evaluates as spatial curvature ($[1/\text{m}]$). Multiplying by the acoustic phase limit ($c$) converts this geometric curvature into physical rotational frequency ($[1/\text{s}]$):
\begin{equation}
    F_{\mu\nu} \equiv \Omega_{\mu\nu} = c (\partial_\mu u_\nu - \partial_\nu u_\mu) \tag{\ref{eq:rosetta_faraday}}
\end{equation}

\textbf{8. Yang-Mills Evolution (Convective Acceleration):} \\
Non-Abelian gauge theories (such as the $SU(3)$ strong force) incorporate a geometric self-interaction commutator ($[A_\mu, A_\nu]$). This gauge self-interaction behaves as the metric displacing its own momentum. This non-linear auto-displacement operates as the Convective Acceleration tensor ($\mathbf{a}_{conv} = (\mathbf{v} \cdot \nabla)\mathbf{v} = v_j \partial_j v_i$). The non-Abelian Yang-Mills equation maps to convective inertia: 
\begin{equation}
    \partial_\mu F^{\mu\nu} + [A_\mu, F^{\mu\nu}] \equiv (\mathbf{v} \cdot \nabla)\mathbf{v} \tag{\ref{eq:rosetta_yang_mills}} 
\end{equation}

\textbf{9. Spinors \& $SU(2)$ (Cartan Moving Frames):} \\
The intrinsic angular momentum of fermions resolves via the orthogonal rotation generators of 3D axes undergoing gyroscopic precession. The spatial axes of a cavity map to a Cartan moving frame (a tetrad, $\mathbf{e}^a$). The evolution of these spatial axes translates via the Cartan structure equations, upgrading the cross product to the exterior wedge product of the fluid's spin connection 1-form ($\boldsymbol{\omega}^a_{\phantom{a}b}$): 
\begin{equation}
    d\mathbf{e}^a + \boldsymbol{\omega}^a_{\phantom{a}b} \wedge \mathbf{e}^b = \mathbf{T}^a \tag{\ref{eq:rosetta_cartan}} 
\end{equation}
Because the axes are tethered to the 3D bulk, a $2\pi$ rotation geometrically twists the surrounding metric. Resolving this spatial torsion 2-form ($\mathbf{T}^a$) requires a $4\pi$ geometric phase rotation. The $SU(2)$ matrices operate as the orthogonal rotation generators of this constraint.

\vspace{0.2cm}

\textbf{10. The Unified Asymmetric Field Tensor ($\hat{U}_{\mu\nu}$) \& Clebsch Parameterization:} \\
A requirement for the emergence of the electromagnetic gauge field ($F_{\mu\nu}$) is the capacity for vorticity. A vacuum defined by a globally smooth velocity potential ($u_\mu = \partial_\mu \phi$) results in irrotational flow ($\partial_{[\nu} u_{\mu]} = 0$), eliminating the antisymmetric tensor. 

To satisfy the requirements of a rotational continuum, the macroscopic kinematic velocity field is parameterized utilizing Clebsch variables:
\begin{equation}
    u_\mu = \partial_\mu \phi + \lambda \partial_\mu \beta \label{eq:clebsch_velocity}
\end{equation}
where $\phi$ governs the irrotational background expansion, and the conjugate pair $(\lambda, \beta)$ parameterize vorticity.

Substituting this Clebsch-parameterized velocity field into the general asymmetric tensor ($\hat{U}_{\mu\nu} = \partial_\nu u_\mu$) partitions the unified state into its orthogonal geometric components. To preserve dimensional orthogonality across the continuum, the antisymmetric transverse vorticity rate ($[1/T]$) must be converted into a dimensionless spatial deformation. 

By geometrically parameterizing the Faraday rate tensor against the temporal relaxation limit, the Planck time ($t_p$), the dynamic rate is normalized into accumulated shear strain. The unification resolves as the geometric decomposition:
\begin{equation}
    \hat{U}_{\mu\nu} = \hat{U}_{(\mu\nu)} + \hat{U}_{[\mu\nu]} \equiv g_{\mu\nu} + t_p F_{\mu\nu} \tag{\ref{eq:unified_tensor_decomposition}}  
\end{equation}

To satisfy covariance, the metric is constructed via the continuous 4-velocity ($u_\mu$) and the spatial projection tensor ($h_{\mu\nu}$). The symmetric sector ($S_{\mu\nu}$) and the antisymmetric sector ($A_{\mu\nu}$) split as: % :::LLM_EDIT::: Inserted covariant tensor construction
\begin{equation}
    \hat{U}_{\mu\nu} = \underbrace{-\alpha_s^2 u_\mu u_\nu + \frac{1}{\alpha_s^2} h_{\mu\nu}}_{S_{\mu\nu}} + \underbrace{\alpha \left( \frac{u_\mu v_{\perp, \nu} - u_\nu v_{\perp, \mu}}{c} \right) + \alpha(t_p) \omega_{\mu\nu}}_{A_{\mu\nu}} \label{eq:covariant_tensor}
\end{equation}

Where the covariant tensor variables are:
    \begin{itemize}[noitemsep]
    \item $u_\mu$: \textbf{Continuum 4-Velocity.} The normalized, timelike vector field ($u_\mu u^\mu = -1$) representing the local rest frame of the coordinate metric. In this comoving frame, the spatial components evaluate to zero ($u^i = 0$) and the temporal component maps to $u^0 = 1$, isolating bulk geometric strain to the symmetric diagonals. 
    \item $h_{\mu\nu}$: \textbf{Spatial Projection Tensor.} The symmetric tensor ($h_{\mu\nu} = g_{\mu\nu} + u_\mu u_\nu$) that projects the metric onto the orthogonal spatial hypersurface, where $g_{\mu\nu}$ is the unperturbed Minkowski background signature $(-,+,+,+)$.
    \item $\alpha_s$: \textbf{Total Scalar Admittance.} The dimensionless phase capacity bounded by thermodynamic saturation, representing the scalar contraction of the symmetric sector. 
    \item $\alpha$: \textbf{Viscoelastic Modulus Ratio.} The dimensionless transverse coupling constant. Tracks the fractional capacity of the metric to sustain transverse shear relative to bulk compression. 
    \item $v_{\perp, \mu}$: \textbf{Transverse Slip Vector.} The orthogonal spatial velocity representing solenoidal slip, coupling to the time-space components via the antisymmetric wedge product.
    \item $\omega_{\mu\nu}$: \textbf{Vorticity Tensor.} The antisymmetric component representing the exterior derivative of the Clebsch rotational phase.
    \item $t_p$: \textbf{Structural Relaxation Time.} The temporal phase limit ($t_p = l_p/c$) that normalizes the vorticity into dimensionless strain.
\end{itemize}

Because the exterior derivative of the scalar gradient equates to zero ($d^2\phi = 0$), $\phi$ cancels from the antisymmetric decomposition. Multiplying by the acoustic phase limit ($c$) to conserve the frequency established in Equation \ref{eq:rosetta_faraday}, the Faraday tensor ($F$) translates to the exterior wedge product of the Clebsch variables ($F = c \, d\lambda \wedge d\beta$). The gauge field behaves as metric vorticity. % :::LLM_EDIT::: Upgraded Clebsch crossed-gradients to formal exterior wedge product

Evaluating this field within a local Cartesian coordinate chart $x^\mu = (ct, x, y, z)$, the temporal components of the normalized gauge field ($t_p E_i / c$) extract the ratio $c/t_p$. This ratio defines the Planck acceleration limit ($a_p = c/t_p$). Substituting this limit, the dimensionless matrix is represented as: 
\begin{equation}
    \hat{U}_{\mu\nu} = \begin{bmatrix} -\alpha_s^2 & -E_x/a_p & -E_y/a_p & -E_z/a_p \\ E_x/a_p & \alpha_s^{-2} & -t_p B_z & t_p B_y \\ E_y/a_p & t_p B_z & \alpha_s^{-2} & -t_p B_x \\ E_z/a_p & -t_p B_y & t_p B_x & \alpha_s^{-2} \end{bmatrix}
\end{equation}

This resolves the kinematic state. Because the Electric field ($E$) normalizes against the Planck acceleration ($a_p$), $E$ operates dimensionally and physically as the kinematic acceleration of the vacuum. Because the Magnetic field ($B$) normalizes against the Planck time ($t_p$), $B$ projects as spatial vorticity, or spin frequency.

Consequently, classical electromagnetic gauge constraints map to geometric limits of continuum mechanics. Faraday's Law of Induction ($d\mathbf{E} = -\partial \mathbf{B} / \partial t$), for example, acts as a kinematic identity: the exterior spatial derivative of the acceleration must equal the temporal rate of change of spin. % :::LLM_EDIT::: Upgraded 3D curl to exterior spatial derivative

Because both the geometric properties and the electromagnetic gauge derive from the same tensor, they are subject to non-linear strain coupling. The symmetric diagonals (Gravity) bound the kinematic capacity of the antisymmetric off-diagonals (Electromagnetism). As gravitational curvature increases, the scalar admittance ($\alpha_s \to 0$) is exhausted. 

The formal derivation of physical observables from this unified matrix is executed in the companion manuscript, \href{https://doi.org/10.5281/zenodo.21321264}{\textit{The Asymmetric Metric-Tensor and the Thermodynamic Equation of State}}.

\begin{tcolorbox}[recallbox, title={\ref{box:rosetta_stone}. Tensor Evolution \& Geometric Bending}]
\begin{align}
    [\boldsymbol{g}_{\mu\nu}] &= \text{diag}(-\alpha_s^2, \, \alpha_s^{-2}, \, r^2, \, r^2\sin^2\theta) \tag{\ref{eq:rosetta_metric_matrix}} \\
    ds^2 &\equiv -\left( 1 - \frac{P_{dyn}}{P_c} \right) c^2 dt^2 + \left( 1 - \frac{P_{dyn}}{P_c} \right)^{-1} dr^2 + r^2 d\Omega^2 \tag{\ref{eq:rosetta_acoustic_metric}} \\
    \nabla_\mu &\equiv \frac{D}{Dt} = \frac{\partial}{\partial t} + (\mathbf{v} \cdot \nabla) \tag{\ref{eq:rosetta_covariant}}  \\
    c^2 \Gamma^i_{00} &\equiv -\frac{1}{\rho_{\tau}} \nabla^i P_{static}, \quad \Gamma^k_{ij} \equiv \frac{1}{c} \partial_j v_i \tag{\ref{eq:rosetta_christoffel_pressure}} \\
    R_{ijkl} &\equiv \text{Inc}(\boldsymbol{\varepsilon})_{ijkl} \tag{\ref{eq:rosetta_riemann}} \\
    R_{\mu\nu} &\equiv \text{Volumetric Strain}, \quad C_{\mu\nu\rho\sigma} \equiv \text{Deviatoric Shear} \tag{\ref{eq:rosetta_ricci}} \\
    T_{\mu\nu} &\equiv \sigma_{\mu\nu} + \rho u_\mu u_\nu \tag{\ref{eq:rosetta_stress_energy}} \\
    G_{\mu\nu} &\equiv \text{Hookean Strain Balance of } T_{\mu\nu} \tag{\ref{eq:rosetta_einstein_tensor}} \\
    F_{\mu\nu} &\equiv \Omega_{\mu\nu} = c (\partial_\mu u_\nu - \partial_\nu u_\mu) \tag{\ref{eq:rosetta_faraday}}  \\
    \partial_\mu F^{\mu\nu} &+ [A_\mu, F^{\mu\nu}] \equiv (\mathbf{v} \cdot \nabla)\mathbf{v} \tag{\ref{eq:rosetta_yang_mills}}  \\
    d\mathbf{e}^a &+ \boldsymbol{\omega}^a_{\phantom{a}b} \wedge \mathbf{e}^b = \mathbf{T}^a \tag{\ref{eq:rosetta_cartan}} \\
    \hat{U}_{\mu\nu} &= \hat{U}_{(\mu\nu)} + \hat{U}_{[\mu\nu]} \equiv g_{\mu\nu} + t_p F_{\mu\nu} \tag{\ref{eq:unified_tensor_decomposition}} \end{align}
\end{tcolorbox}


% ----------------------------------------------------------------------
% ----------------------------------------------------------------------
\section{Continuous Pressure Dynamics (Grand Unification)}
\label{sec:deriv_pressure_gradient}

The four fundamental forces resolve as identical geometric projections of a single continuous pressure gradient. The metric gradient derives via three independent mathematical paths which converge on the same equation.

\vspace{0.2 cm}

\textbf{1. Path I: The Kinematic Tensor Proof (Isochoric Velocity Scaling):} \\
The continuous velocity field satisfies the trace-free isochoric constraint ($\nabla \cdot \mathbf{v} = 0$). Expanding the divergence operator in spherical coordinates $(r, \theta, \phi)$ cancels the transverse angular gradients ($\partial_\theta = 0$, $\partial_\phi = 0$): 
\begin{equation}
    \frac{1}{r^2} \frac{\partial}{\partial r}(r^2 v_r) = 0 \implies v_r(r) \equiv v(r) = \frac{C}{r^2} 
\end{equation}

Substituting this velocity into the Bernoulli capacity equation constructs the dynamic pressure ($P_{dyn} = \frac{1}{2}\rho_\tau v^2$):
\begin{equation}
    P_{dyn}(r) = \frac{1}{2}\rho_\tau \left( \frac{C}{r^2} \right)^2 = \frac{\rho_\tau C^2}{2 r^4}
\end{equation}

To define the integration constant ($C$), the geometric Action ($S$) of a Clebsch twist constitutes an invariant constant ($\kappa \equiv h$). Action defines as the structural energy ($E$) integrated over the period ($T = 2\pi r/c$):
\begin{align}
    E &= \int_r^\infty P_{dyn}(r') \cdot 4\pi (r')^2 \, dr' = \frac{2\pi \rho_\tau C^2}{r} \nonumber \\
    h &= E \cdot T = \left( \frac{2\pi \rho_\tau C^2}{r} \right) \left( \frac{2\pi r}{c} \right) = \frac{4\pi^2 \rho_\tau C^2}{c}
\end{align}

Isolating the constant ($C^2 = h c / 4\pi^2 \rho_\tau$) and substituting it back into the dynamic pressure equation evaluates the spatial confinement gradient:
\begin{align}
    P(r) &= \frac{\rho_\tau}{2 r^4} \left( \frac{h c}{4\pi^2 \rho_\tau} \right) = \frac{h c}{8\pi^2 r^4}  \\
    P(r) &= \frac{\hbar c}{4\pi r^4} \tag{\ref{eq:pressure_thermo}}
\end{align}

\textbf{2. Path II: The Thermodynamic Proof (Volumetric Strain Energy):} \\
Pressure ($P$) acts as volumetric strain energy density ($E/V$). For a continuous acoustic boundary, this internal state is governed by the radiation equation of state ($P = \frac{1}{3} \frac{E}{V}$). 

The topological rest energy ($E$) required to integrate continuous force across the boundary resolves via the limit ($E(r) = \hbar c / r$). The 3D Euclidean spherical volume defines as $V(r) = \frac{4}{3}\pi r^3$. 

Substituting the phase energy and the spherical volume into the acoustic equation of state isolates the continuous pressure gradient:
\begin{align}
    P(r) &= \frac{1}{3} \frac{E(r)}{V(r)}   \\
    P(r) &= \frac{1}{3} \frac{\left( \frac{\hbar c}{r} \right)}{\left( \frac{4}{3}\pi r^3 \right)} \nonumber \\
    P(r) &= \frac{\hbar c}{4\pi r^4} \tag{\ref{eq:pressure_thermo}}
\end{align}

\textbf{3. Path III: The Newtonian Limit ($F=ma$):} \\
The unified gradient emerges from classical mechanics when evaluated at the kinematic boundary limits. By Newton's second law ($F = ma$), force evaluates as mass accelerated across a given geometry. 

Mass ($m$) represents the volumetric strain of the boundary, defined by the inverse phase limit ($m = \hbar / c r$).  The maximum radial acceleration ($a_{max}$) the metric boundary can sustain at the acoustic phase limit ($v = c$) is the centripetal bound:
\begin{equation}
    a_{max} = \frac{v^2}{r} = \frac{c^2}{r}
\end{equation}

Substituting the volumetric strain ($m$) and the acceleration limit ($a_{max}$) directly into Newton's equation outputs the spatial force gradient:
\begin{align}
    F(r) &= m \cdot a_{max} \nonumber \\
    F(r) &= \left( \frac{\hbar}{c r} \right) \left( \frac{c^2}{r} \right) \nonumber \\
    F(r) &= \frac{\hbar c}{r^2} \tag{\ref{eq:pressure_gradient}}
\end{align}

Continuous pressure defines as force distributed over a boundary surface area ($A = 4\pi r^2$), projecting this Newtonian limit across the Euclidean sphere recovers the exact pressure gradient:
\begin{align}
    P(r) &= \frac{F(r)}{A} \nonumber \\
    P(r) &= \frac{\left( \frac{\hbar c}{r^2} \right)}{4\pi r^2} \nonumber \\
    P(r) &= \frac{\hbar c}{4\pi r^4} 
\end{align}
All three (tensor kinematics, thermodynamics, and classical Newtonian mechanics) converge on the identical pressure state.

\vspace{0.2 cm}

\textbf{4. Dimensional Translation Matrices and Curvature Scaling:} \\
The continuum evaluates temporal and spatial curvature via the action limit ($\hbar$).

Temporal curvature ($1/t$) scales to thermodynamic energy via the reduced action limit:
\begin{equation}
    E(t) = \hbar \left( \frac{1}{t} \right)
\end{equation}

Spatial curvature ($1/r$) scales to thermodynamic energy via the acoustic phase limit ($c$). Operating as the geometric conversion factor between spatial and temporal dimensions ($r = ct$), this dictates the spatial scaling matrix:
\begin{equation}
    E(r) = \hbar c \left( \frac{1}{r} \right) \tag{\ref{eq:energy_2pi}}
\end{equation}

\vspace{0.2 cm}



\textbf{5. Kinematic Derivatives and Geometric Transport:} \\
Evaluating the rate of change of these energy states derives the mechanical operators. The spatial derivative extracts continuous force ($F$):
\begin{align}
    F &= - \frac{dE}{dr} \nonumber \\
    F &= -\frac{d}{dr} \left( \frac{\hbar c}{r} \right) \nonumber \\
    F &= - \left( - \frac{\hbar c}{r^2} \right) \nonumber \\
    F &= \frac{\hbar c}{r^2} 
\end{align}

The temporal derivative extracts kinematic power ($P$):
\begin{equation}
    P = -\frac{d}{dt} \left( \frac{\hbar}{t} \right) = \frac{\hbar}{t^2}
\end{equation}

Because the spatial and temporal states are coupled by the acoustic limit ($dr = c \, dt$), applying the chain rule equates the derivatives:
\begin{align}
    \frac{dE}{dt} &= \frac{dE}{dr} \frac{dr}{dt} \nonumber \\
    P &= F c
\end{align}

Applying this derivative equivalence to the unified pressure gradient ($P_{dyn} = \hbar c / 4\pi r^4$) evaluates the transport. Substituting the phase limit ($r = ct$) into the differential operators yields:
\begin{align}
    \frac{d P_{dyn}}{dt} &= c \frac{d P_{dyn}}{dr} \nonumber \\
    \partial_t P_{dyn} - c \partial_r P_{dyn} &= 0
\end{align}
This recovers the one-way continuous wave equation. The unified pressure gradient does not represent the physical flow of the background metric; rather, it propagates as strain transmitted across the coordinate space at the acoustic limit ($c$).

\vspace{0.2 cm}

\textbf{6. Deterministic Phase Boundaries (The Uncertainty Limits):} \\
Isolating the temporal state yields the energy-time boundary constraint:
\begin{align}
    E &= \frac{\hbar}{t} \nonumber \\
    E \cdot t &= \hbar
\end{align}

For the spatial state, the kinematic energy of a propagating transverse wave equates to the product of its momentum and velocity ($E = pc$). Substituting this into the spatial curvature equation isolates the position-momentum boundary constraint:
\begin{align}
    p c &= \frac{\hbar c}{r} \nonumber \\
    p \cdot r &= \hbar
\end{align}

The Heisenberg Uncertainty limits ($\Delta E \Delta t \ge \hbar$ and $\Delta p \Delta r \ge \hbar$) resolve as geometric boundaries. The momentum of the transverse shear wave is locked to the curvature of its spatial radius. A boundary requires proportional internal momentum to sustain a spatial footprint against the metric.

\vspace{0.2 cm}
\textit{The Fourier Heat Transform \& Geometric Diffusion:}\\
The Fourier transform of the heat equation ($\partial_t U = D \nabla^2 U$) evaluates spatial diffusion. Transitioning from coordinate space ($x$) to spatial frequency space ($k$), the Laplacian ($\nabla^2$) maps to the squared wavenumber ($-k^2$), dictating that spatial strain dissipates at a rate proportional to the square of its curvature ($\partial_t \hat{U} \propto -D k^2$). For a defect defined by a circumference ($2\pi r$), the Fourier wavenumber equates to its spatial curvature ($k = 1/r$). 

If a defect is compressed into a dimensionless point ($r \to 0$), the curvature diverges ($k \to \infty$), causing the metric to instantly dissipate the strain back to the baseline ($\partial_t \hat{U} \to -\infty$). Therefore, a point-particle cannot mechanically exist within a continuous metric. To prevent the metric from diffusing the strain, the boundary resolves as a closed loop. Because momentum is the generator of spatial translation, sustaining a specific curvature ($1/r$) against the metric's heat diffusion requires a conserved magnitude of kinematic momentum ($p$). This constraint restricts the stable rest state of the boundary to $p = \hbar k = \hbar/r$. Rearranging this isolates the origin of the Heisenberg Uncertainty Principle ($p \cdot r = \hbar$). Quantum uncertainty acts as the phase limit of a continuous geometry, dictating that mass stabilizes into equilibrium when its rotational momentum equals the Fourier diffusion limit of its spatial curvature.

\vspace{0.2 cm}

\textbf{7. The Geometric Translation Limit ($E = mc^2$):} \\
The translation matrix resolves as the work stored within the deformation integrated across the metric ($Work = \int F \cdot dr$):
\begin{equation}
    \int_r^\infty \frac{\hbar c}{(r')^2} \, dr' = \frac{\hbar c}{r} \tag{\ref{eq:energy_2pi}}
\end{equation}

Equating this work to the rest energy of the defect ($mc^2$) isolates the reduced Compton radius ($r_c = \hbar/mc$). The baseline rest mass of a physical state ($E=mc^2$) operates identically as the continuous thermodynamic work the metric performs to sustain the boundary.

\vspace{0.2 cm}


\textbf{8. Orthogonal Curvature Superposition (The Relativistic Invariant):} \\
The metric supports two distinct boundary states.  Rest mass operates as locked internal structural curvature defined by the Compton radius ($1/r_c$), projecting a rest energy of $E_{rest} = \hbar c / r_c$. Kinematic momentum operates as translational curvature ($1/r_k$), projecting a kinematic energy of $E_{kin} = \hbar c / r_k$.

For a defect translating through the coordinate space, the metric sustains both states simultaneously. Because internal locked rotation and external linear translation project as orthogonal vectors, their total capacity ($E$) sums as a Pythagorean superposition of their independent curvatures:
\begin{align}
    E^2 &= E_{rest}^2 + E_{kin}^2 \nonumber \\
    E^2 &= \left( \frac{\hbar c}{r_c} \right)^2 + \left( \frac{\hbar c}{r_k} \right)^2 
\end{align}

Substituting the classical variables for rest mass ($mc^2 = \hbar c / r_c$) and momentum ($pc = \hbar c / r_k$) isolates the total energy of the boundary:
\begin{equation}
    E^2 = (mc^2)^2 + (pc)^2 
\end{equation}
The relativistic energy-momentum relation resolves not as a kinematic postulate, but as the superposition of orthogonal curvature within a continuous metric.

\vspace{0.2 cm}

\textbf{9. The Interaction Hierarchy (Force Projections):} \\
The continuous force equation unifies the four fundamental forces. The isolated interaction strengths resolve as the identical equation ($F(r)$) projected across four distinct spatial boundaries ($r_i$).

\textbf{I. Gravity (The Core Boundary):} \\
Applying the spatial gradient at the geometric minimum of the acoustic cavity ($l_p = \sqrt{\hbar G / c^3}$) extracts the tension limit:
\begin{equation}
    F_{grav} = F(l_p) = \frac{\hbar c}{(\sqrt{\hbar G / c^3})^2} = \frac{c^4}{G} \approx 1.21 \times 10^{44} \text{ N} \tag{\ref{eq:force_gravity}} 
\end{equation}

\textbf{II. The Weak Force (The Cavitation Yield):} \\
Calculating the continuous pressure across the Compton boundary of the W-Boson ($R_W = \hbar / m_W c$) specifies the volumetric rupture limit: 
\begin{equation}
    F_{weak} = F(R_W) = \frac{\hbar c}{(\hbar / m_W c)^2} = \frac{m_W^2 c^3}{\hbar} \approx 5.25 \times 10^9 \text{ N} \tag{\ref{eq:force_weak}} 
\end{equation}

\textbf{III. The Strong Force (The Acoustic Confinement):} \\
Projecting the spatial gradient at the viscoelastic relaxation envelope ($L_{crit} = \alpha \hbar / m_e c$) establishes the boundary: 
\begin{equation}
    F_{strong} = F(L_{crit}) = \frac{\hbar c}{(\alpha \hbar / m_e c)^2} = \frac{m_e^2 c^3}{\alpha^2 \hbar} \approx 3980 \text{ N} \tag{\ref{eq:force_strong}} 
\end{equation}

\textbf{IV. Electromagnetism (The Transverse Leakage):} \\
Applying the spatial gradient at the classical Bohr radius ($a_0 = \hbar / \alpha m_e c$) resolves the bare electromagnetic baseline tension:
\begin{equation}
    F_{EM_{base}} = F(a_0) = \frac{\hbar c}{(\hbar / \alpha m_e c)^2} = \frac{\alpha^2 m_e^2 c^3}{\hbar} \approx 1.13 \times 10^{-5} \text{ N} \tag{\ref{eq:force_em}}
\end{equation}

\textbf{10. The Transverse Phase Flux \& Coulomb Equivalence:} \\
Electromagnetism operates distinctly from gravity because it is governed by transverse shear strain (curl) rather than volumetric strain (divergence). A charge boundary actively twists the surrounding metric, creating an outward-radiating transverse spiral. 

\vspace{0.2 cm}
\textit{The Kinematic Decay of the Electric Field:} \\
Because this transverse shear wave propagates outward at the phase limit ($r = c t$), the angular frequency of the twist decays temporally as the spiral expands. The kinematic rotation of the field at any given moment is dictated by its curvature ($\omega = c/r$).

The quantum of action ($\hbar$) functions as the geometric translation, mapping curvature to capacity:
\begin{equation}
    E(t) = \hbar \omega(t) = \frac{\hbar}{t}
\end{equation}
Substituting the radial phase limit ($t = r/c$) into the energy state recovers the energy capacity of the boundary ($E = \hbar c / r$). Taking the derivative ($F = -dE/dr$) recovers the pressure gradient:
\begin{align}
    F(r) &= - \frac{dE}{dr} \nonumber \\
    F(r) &= - \frac{d}{dr} \left( \frac{\hbar c}{r} \right) \nonumber \\
    F(r) &= - \left( - \frac{\hbar c}{r^2} \right) \nonumber \\
    F(r) &= \frac{\hbar c}{r^2} \tag{\ref{eq:pressure_gradient}}
\end{align}
The force gradient applies equally to transverse shear strain as it does to volumetric strain.

\vspace{0.2 cm}
\textit{The Coulomb Leakage Limit ($\alpha$):} \\
While volumetric strain displaces the bulk capacity of the metric ($P_c$), transverse shear is restricted by the transverse shear capacity ($\Pi_{shear}$). 

The Fine-Structure Constant ($\alpha = \Pi_{shear}/P_c$) operates as the Viscoelastic Modulus Ratio of the continuum. 

It dictates the geometric fraction of metric tension that transfers laterally through the boundary without dissipating into continuous slip. 

The electrostatic force is therefore the continuous pressure gradient bounded by this transverse leakage identity ($\alpha$):
\begin{equation}
    F_{coulomb}(r) = \alpha \frac{\hbar c}{r^2}
\end{equation}

\vspace{0.2 cm}
\textit{The Macroscopic Equivalence:} \\
Projecting this transverse gradient at the boundary of the Bohr radius ($a_0 = \hbar / \alpha m_e c$) yields:
\begin{align}
    F_{coulomb}(a_0) &= \alpha \cdot F(a_0) \nonumber \\
    F_{coulomb}(a_0) &= \alpha \frac{\hbar c}{a_0^2} \approx 8.23 \times 10^{-8} \text{ N} \tag{\ref{eq:force_coulomb}}
\end{align}

Equating this derivation to the classical Coulomb measurement ($F_{coulomb} = e^2 / 4\pi \epsilon_0 a_0^2 \approx 8.23 \times 10^{-8} \text{ N}$) reveals the classical definition of the fine structure constant:
\begin{equation}
    \alpha \frac{\hbar c}{a_0^2} = \frac{e^2}{4\pi \epsilon_0 a_0^2} \nonumber
\end{equation}

\begin{equation}
    \alpha \hbar c = \frac{e^2}{4\pi \epsilon_0} \nonumber
\end{equation}

\begin{equation}
    \alpha = \frac{e^2}{4\pi \epsilon_0 \hbar c} \label{eq:alpha_empirical_recovery}
\end{equation}

\vspace{0.2 cm}
\textit{The Kinematic Phase Hierarchy (The Deborah Translation):} \\
The emergence of the coupling coefficient ($\alpha$) across these boundaries reveals that the fundamental interactions are not distinct forces, but rather the identical continuous force gradient evaluated across the three distinct rheological phase states. 

By translating the boundary radii into their corresponding energy states ($E = \hbar c / r$), the scaling action of the Fine-Structure Constant defines the phase transitions relative to the baseline rest mass ($m_e c^2$):
\begin{align}
    E_{strong}(r_e) &= \frac{m_e c^2}{\alpha} \quad &\text{(Solid Metric / Glass Transition)} \\
    E_{mass}(r_c) &= m_e c^2 \quad &\text{(Liquid Metric / Yield Boundary)} \\
    E_{EM}(a_0) &= \alpha m_e c^2 \quad &\text{(Compliant Metric / Transverse Field)}
\end{align}
$\alpha$ dictates the transition of the Deborah Number ($De$). At the inner bound ($r_e$), the applied interaction energy brings $De \to 1$, freezing the metric into an inelastic solid (The Strong Force). At the baseline yield limit ($r_c$), the metric acts as a fluid vortex (Rest Mass). Stepping outward to the relaxed boundary ($a_0$), the phase response drops $De \ll 1$, allowing the metric to freely transmit transverse shear strain without cavitation (The Coulomb Force).

\vspace{0.2 cm}
\textit{The Illusion of Relative Force Strengths ($\alpha^4$ vs $\alpha^5$):} \\
Relative forces are artifacts of the inverse-square law ($F \propto 1/r^2$) evaluated across phase boundaries. 

Because the spatial separation between the Solid metric limit ($r_e$) and the Compliant field limit ($a_0$) spans exactly two structural phase steps ($r_e/a_0 = \alpha^2$), evaluating the bare geometric gradient at these boundaries isolates this identity:
\begin{equation}
    \frac{F_{EM_{base}}}{F_{strong}} = \frac{F(a_0)}{F(r_e)} = \frac{\alpha^2}{\alpha^{-2}} = \alpha^4 \label{eq:force_em_ratio}
\end{equation}
This $\alpha^4$ identity demonstrates that the difference in interaction strength is the consequence of evaluating the inverse-square law at two boundaries separated by a factor of $\alpha^2$.

Because electromagnetism is further throttled by the transverse compliance limit ($\alpha$), the ratio of the observed macroscopic Coulomb force ($F_{coulomb}$) to the strong boundary isolates an additional factor of $\alpha$:
\begin{equation}
    \frac{F_{coulomb}}{F_{strong}} = \alpha^5 \label{eq:force_coulomb_ratio}
\end{equation}

The physical strength of electromagnetism relative to the strong force ($\alpha^5$) emerges from a geometric $\alpha^4$ coordinate separation multiplied by the $\alpha$ transverse compliance limit.





\vspace{0.2 cm}

\textbf{11. The Macroscopic Baseline (The Horizon Floor \& $\Lambda$):} \\
The unified spatial gradient ($F = \hbar c / r^2$) resolves at the macroscopic scale as the geometric dilution of metric tension. As the universe expands, this tension distributes across the causal horizon ($R_{h,0}$). 

\vspace{0.2 cm}
\textit{The Cosmological Constant ($\Lambda$) as Holographic Curvature:} \\
Because this baseline tension is defined by its dilution across the 2D spherical boundary area of the causal horizon ($4\pi R_{h,0}^2$), the corresponding curvature is dictated by the differential geometry of that boundary surface. 

In differential geometry, the Gaussian curvature of a 2-sphere evaluates as $1/r^2$. Its total intrinsic scalar curvature (the Ricci scalar) evaluates as exactly twice the Gaussian curvature ($2/r^2$). Thus, the Cosmological Constant ($\Lambda$) emerges as the curvature of the holographic boundary at radius $R_{h,0}$:
\begin{equation}
    \Lambda = \frac{2}{R_{h,0}^2} \label{eq:lambda_fluid_derivation}
\end{equation}

\vspace{0.2 cm}
\textit{Empirical Numerical Evaluation:} \\
Substituting the present-day Hubble radius ($R_{h,0} = c / H_0 \approx 1.37 \times 10^{26} \text{ m}$) into the geometric curvature equation calculates the local snapshot of the metric: 
\begin{align}
    \Lambda &= \frac{2}{(1.37 \times 10^{26} \text{ m})^2} \nonumber \\
    \Lambda &\approx 1.06 \times 10^{-52} \text{ m}^{-2} 
\end{align}
This derivation reproduces the Planck 2018 empirical measurement ($\Lambda_{obs} \approx 1.1 \times 10^{-52} \text{ m}^{-2}$) from continuous boundary geometry.

\vspace{0.2 cm}
\textit{Temporal Relaxation \& Cosmological Redshift ($z$):} \\
Because the causal horizon expands at the acoustic phase limit ($R_h = c t$), the baseline force floor operates as a time-dependent viscoelastic relaxation of the fluid metric:
\begin{equation}
    F_{cosmic}(t) = \frac{\hbar c}{(c t)^2} = \frac{\hbar}{c t^2}
\end{equation}
As the universe ages ($t \to \infty$), the geometric tension continuously dissipates. Cosmological redshift ($z$) is measured as the proportional expansion of the metric between the time a photon is emitted ($t_{emit}$) and the time it is observed ($t_{obs}$), defined kinematically as $1+z = R_{h,obs} / R_{h,emit}$. 

By comparing the tension of the vacuum at the moment of emission ($F_{emit}$) to the tension at observation ($F_{obs}$), the ratio isolates the source of redshift:
\begin{align}
    \frac{F_{emit}}{F_{obs}} &= \frac{\left( \frac{\hbar c}{R_{h,emit}^2} \right)}{\left( \frac{\hbar c}{R_{h,obs}^2} \right)} = \frac{R_{h,obs}^2}{R_{h,emit}^2} \nonumber \\
    \frac{F_{emit}}{F_{obs}} &= (1+z)^2 \implies 1+z = \frac{R_{h,obs}}{R_{h,emit}} \label{eq:redshift_tension}
\end{align}
The wavelength elongates as the continuum sheds its tension and reduces its minimum curvature.

\vspace{0.2 cm}

\textbf{12. Macroscopic Combinatorial Aggregates (Newtonian Gravity):} \\
For a macroscopic object, mass partitions as a combinatorial superposition of boundaries.

The Planck mass is the maximum mass used as a scale factor:
\begin{equation}
    M_P = \sqrt{\frac{\hbar c}{G}} 
\end{equation}

The total number of interacting boundaries ($N$) within a boundary ($M$) equates to the ratio of the aggregate mass to this boundary ($N = M / M_P$). The force ($F_{macro}$) between two aggregate masses ($M_1$ and $M_2$) separated by a coordinate radius ($r$) acts as the combinatorial product of their individual boundaries ($N_1 \cdot N_2$) applied to the unified continuum force:
\begin{align}
    F_{macro} &= (N_1 \cdot N_2) \cdot F(r) \nonumber \\
    F_{macro} &= \left( \frac{M_1}{M_P} \right) \left( \frac{M_2}{M_P} \right) \left( \frac{\hbar c}{r^2} \right)
\end{align}

Substituting the identity ($M_P^2 = \hbar c / G$) recovers the classical format:
\begin{align}
    F_{macro} &= \frac{M_1 M_2}{\left( \frac{\hbar c}{G} \right)} \left( \frac{\hbar c}{r^2} \right) \nonumber \\
    F_{macro} &= M_1 M_2 \left( \frac{G}{\hbar c} \right) \left( \frac{\hbar c}{r^2} \right) = \frac{G M_1 M_2}{r^2}
\end{align}


\begin{tcolorbox}[recallbox, title={\ref{box:pressure_gradient}. Continuous Pressure Dynamics}]
\begin{align}
    P(r) &= \frac{1}{3} \frac{E}{V} = \frac{\hbar c}{4\pi r^4} \tag{\ref{eq:pressure_thermo}} \\
    F(r) &= P(r) \cdot 4\pi r^2 = \frac{\hbar c}{r^2} \tag{\ref{eq:pressure_gradient}} \\
    E(r) &= \int F(r) \cdot dr = \frac{\hbar c}{r} \implies E_{rest} = mc^2 \tag{\ref{eq:energy_2pi}} \\
    F(l_p) &= \frac{c^4}{G} \tag{\ref{eq:force_gravity}} \\
    F(R_W) &= \frac{m_W^2 c^3}{\hbar} \tag{\ref{eq:force_weak}} \\
    F(L_{crit}) &= \frac{m_e^2 c^3}{\alpha^2 \hbar} \tag{\ref{eq:force_strong}} \\
    F_{EM_{base}}(a_0) &= \frac{\hbar c}{a_0^2} \implies \frac{F_{EM_{base}}}{F_{strong}} = \alpha^4 \tag{\ref{eq:force_em_ratio}} \\
    F_{coulomb}(a_0) &= \alpha \frac{\hbar c}{a_0^2} \implies \frac{F_{coulomb}}{F_{strong}} = \alpha^5 \tag{\ref{eq:force_coulomb_ratio}}
\end{align}
\end{tcolorbox}

% ----------------------------------------------------------------------
\section{Generalized Hooke's Law \& $SU(3)$ Symmetry}
\label{sec:deriv_hookes_su3}

A baryon evaluates as a 3D isochoric cavity. The 8 geometric degrees of freedom of a volume-preserving deformation tensor evaluate as the $SU(3)$ symmetry group of the Strong Interaction.

\vspace{0.2 cm}

\textbf{1. Tensor Decomposition (Volumetric vs. Deviatoric):} \\
By classical elastodynamics, the strain tensor ($\varepsilon_{ij}$) evaluates as the linear superposition of its isotropic expansion and its shape-distorting deviatoric shear ($\varepsilon'_{ij}$):
\begin{equation}
    \varepsilon_{ij} = \underbrace{\frac{1}{3}\text{Tr}(\varepsilon)\delta_{ij}}_{\text{Volumetric Strain}} + \underbrace{\varepsilon'_{ij}}_{\text{Deviatoric Strain}} 
\end{equation}
The corresponding Cauchy stress tensor ($\sigma_{ij}$) evaluates via Generalized Hooke's Law, partitioning identical geometric boundaries governed by the Bulk Modulus ($K$) and the Shear Modulus ($\mu_s$):
\begin{equation}
    \sigma_{ij} = K \text{Tr}(\varepsilon)\delta_{ij} + 2\mu_s \varepsilon'_{ij} 
\end{equation}

\textbf{2. The Incompressible Annihilation (Zero-Trace):} \\
A stable Baryon cavity circulates at the acoustic limit ($v \to c$). The core evaluates as incompressible ($\nabla \cdot \mathbf{u} = 0$). The volumetric trace evaluates to zero:
\begin{equation}
    \text{Tr}(\varepsilon) = \varepsilon_{11} + \varepsilon_{22} + \varepsilon_{33} = 0 \tag{\ref{eq:isochoric_trace}}
\end{equation}
Substituting this zero-divergence limit into the strain decomposition eliminates the volumetric component ($\varepsilon_{ij} = \varepsilon'_{ij}$). For an incompressible cavity, the total geometric strain resolves as traceless deviatoric shear.

\textbf{3. The Deviatoric Stress State (Hookean $SU(3)$):} \\
Substituting the zero-trace volumetric boundary condition into Generalized Hooke's Law isolates the structural restorative stress of the cavity:
\begin{align}
    \sigma_{ij} &= K(0)\delta_{ij} + 2\mu_s \varepsilon'_{ij}  \\
    \sigma_{ij} &= 2\mu_s \varepsilon'_{ij} \tag{\ref{eq:deviatoric_stress}}
\end{align}
Because the Deviatoric Strain Tensor ($\varepsilon'_{ij}$) evaluates as a $3 \times 3$ symmetric, traceless matrix, setting the trace to zero removes the volumetric degree of freedom from the standard $9$. The deformation is thus governed by $8$ independent generators.

\textbf{4. The Diagonal Generators (Longitudinal Strain):} \\
The continuous longitudinal strains evaluate as diagonal, traceless matrices. To form a complete mathematical basis, these generators evaluate as orthogonal and normalize to the trace inner product ($\text{Tr}(\lambda_a \lambda_b) = 2\delta_{ab}$).

\textit{Basis 1 ($\lambda_3$):} 
Evaluating a longitudinal stretch along the $x$-axis ($\varepsilon_{11} = a$) while conserving volume within the $x$-$y$ plane dictates $\varepsilon_{22} = -a$ and $\varepsilon_{33} = 0$. Applying the normalization constraint ($\text{Tr}(\lambda_3^2) = 2$) isolates $a = 1$:
\begin{equation}
    \lambda_3 = \begin{pmatrix} 1 & 0 & 0 \\ 0 & -1 & 0 \\ 0 & 0 & 0 \end{pmatrix} 
\end{equation}

\textit{Basis 2 ($\lambda_8$):} 
The second longitudinal state must conserve 3D volume ($\text{Tr}(\lambda_8) = 0$) and evaluate orthogonal to $\lambda_3$ ($\text{Tr}(\lambda_3 \lambda_8) = 0$). Let the matrix evaluate as $\text{diag}(x, y, z)$.
\begin{align}
    x + y + z &= 0 \nonumber \\
    (1)(x) + (-1)(y) + (0)(z) &= 0 \implies x = y \nonumber
\end{align}
Substituting $x = y$ into the isochoric constraint dictates $x + x + z = 0$, yielding $z = -2x$. Applying the normalization constraint ($\text{Tr}(\lambda_8^2) = 2$) isolates the geometric fraction ($x = 1/\sqrt{3}$):
\begin{equation}
    \lambda_8 = \frac{1}{\sqrt{3}} \begin{pmatrix} 1 & 0 & 0 \\ 0 & 1 & 0 \\ 0 & 0 & -2 \end{pmatrix} 
\end{equation}
The $1/\sqrt{3}$ fraction evaluates as the normalization requirement for a 3D orthogonal volume-conserving strain.

\textbf{5. The Off-Diagonal Generators (Shear \& Vorticity):} \\
The remaining $6$ geometric degrees of freedom evaluate as off-diagonal spatial transitions, partitioned into symmetric and antisymmetric components:

\textit{Symmetric Transverse Shear:} The real, symmetric matrices evaluate the $3$ degrees of isochoric planar sliding (e.g., $\lambda_1, \lambda_4, \lambda_6$).
\begin{equation}
    \lambda_1 = \begin{pmatrix} 0 & 1 & 0 \\ 1 & 0 & 0 \\ 0 & 0 & 0 \end{pmatrix} 
\end{equation}

\textit{Antisymmetric Vorticity:} The imaginary, antisymmetric matrices evaluate the $3$ degrees of rotation (e.g., $\lambda_2, \lambda_5, \lambda_7$). The imaginary unit ($i$) evaluates as the requirement for rotational transformations to evaluate as Hermitian observables ($\lambda^\dagger = \lambda$).
\begin{equation}
    \lambda_2 = \begin{pmatrix} 0 & -i & 0 \\ i & 0 & 0 \\ 0 & 0 & 0 \end{pmatrix} 
\end{equation}
The $8$ Gell-Mann matrices evaluate as the orthogonal basis vectors for this macroscopic deviatoric fluid strain tensor. 

\vspace{0.2 cm}

\textbf{6. Kinematic Non-Commutativity (Finite Strain):} \\
Infinitesimal wave strains superimpose linearly. However, extreme spatial deformation at the core evaluates via the non-linear (Almansi) Strain Tensor:
\begin{equation}
    \varepsilon_{ij} = \frac{1}{2} \left( \frac{\partial u_i}{\partial x_j} + \frac{\partial u_j}{\partial x_i} - \frac{\partial u_k}{\partial x_i} \frac{\partial u_k}{\partial x_j} \right) \tag{\ref{eq:su3_finite_strain}}
\end{equation}
Because the orthogonal displacement vectors ($u_k$) are multiplied against each other within the quadratic tensor product, the strain field couples with its own displacement gradients. 

Finite rotational shears do not commute. The operation of shearing the fluid followed by rotation yields a different state than rotating followed by shearing. This incompatibility dictates that the matrix multiplication of the continuous generators evaluates as non-commutative ($[\lambda_a, \lambda_b] \neq 0$), yielding the structural Lie bracket:
\begin{equation}
    [\lambda_a, \lambda_b] = 2i f_{abc} \lambda_c \tag{\ref{eq:non_abelian_commutator}}
\end{equation}
The non-Abelian self-interaction evaluates as the non-linear kinematic product of finite 3D Eulerian fluid strain.

\begin{tcolorbox}[recallbox, title={\ref{box:hookes_su3}. Generalized Hooke's Law \& $SU(3)$}]
\begin{align}
    \text{Tr}(\varepsilon) &= 0 \implies \varepsilon_{ij} = \varepsilon'_{ij} \tag{\ref{eq:isochoric_trace}} \\
    \sigma_{ij} &= 2\mu_s \varepsilon'_{ij} \implies \text{Deviatoric } SU(3) \text{ Basis} \tag{\ref{eq:deviatoric_stress}} \\
    \varepsilon_{ij} &= \frac{1}{2} \left( \frac{\partial u_i}{\partial x_j} + \frac{\partial u_j}{\partial x_i} - \frac{\partial u_k}{\partial x_i} \frac{\partial u_k}{\partial x_j} \right) \tag{\ref{eq:su3_finite_strain}} \\
    [\lambda_a, \lambda_b] &= 2i f_{abc} \lambda_c \tag{\ref{eq:non_abelian_commutator}}
\end{align}
\end{tcolorbox}




\end{document}
